\documentclass[11pt,reqno]{amsart}
\usepackage[T1]{fontenc}
\usepackage[utf8]{inputenc}
\usepackage{lmodern}
\usepackage[margin=1in]{geometry}
\usepackage{amssymb,amsmath,amsthm,amsfonts,mathtools}
\usepackage{mathrsfs}
\usepackage{xcolor}
\usepackage{graphicx}
\usepackage{tikz}
\usetikzlibrary{arrows.meta,backgrounds,calc,fit,positioning}

\newcommand{\proofdrawcolor}{black!60}
\newcommand{\prooffillcolor}{gray!7}
\newcommand{\proofpanelcolor}{black!28}
\newcommand{\proofpanelfill}{gray!2}

\usepackage{microtype}
\usepackage[shortlabels]{enumitem}
\usepackage{aliascnt}
\usepackage[colorlinks=true]{hyperref}
\usepackage[nameinlink,capitalise,noabbrev]{cleveref}

\allowdisplaybreaks
\setlist{listparindent=\parindent,parsep=0pt}
\setlist[itemize]{topsep=4pt,itemsep=2pt,parsep=1pt}
\setlist[enumerate]{topsep=4pt,itemsep=3pt,parsep=1pt}

\numberwithin{equation}{section}

\theoremstyle{plain}
\newtheorem{thm}{Theorem}[section]
\newaliascnt{prop}{thm}
\newtheorem{prop}[prop]{Proposition}
\aliascntresetthe{prop}
\newaliascnt{lemma}{thm}
\newtheorem{lemma}[lemma]{Lemma}
\aliascntresetthe{lemma}
\newaliascnt{cor}{thm}
\newtheorem{cor}[cor]{Corollary}
\aliascntresetthe{cor}
\theoremstyle{remark}
\newaliascnt{remark}{thm}
\newtheorem{remark}[remark]{Remark}
\aliascntresetthe{remark}

\DeclarePairedDelimiter{\pa}{\lparen}{\rparen}
\DeclarePairedDelimiter{\brac}{\lbrace}{\rbrace}
\DeclarePairedDelimiter{\abs}{\lvert}{\rvert}
\DeclarePairedDelimiter{\norm}{\lVert}{\rVert}
\DeclarePairedDelimiter{\ipp}{\langle}{\rangle}

\renewcommand{\d}{\mathsf{d}}
\newcommand{\md}{\mathrm{d}}
\newcommand{\R}{{\mathbb{R}}}
\newcommand{\N}{{\mathbb{N}}}
\newcommand{\E}{{\mathbb{E}}}
\newcommand{\Q}{{\mathbb{Q}}}
\newcommand{\Pb}{{\mathbb{P}}}
\newcommand{\K}{{\mathsf{K}}}
\newcommand{\g}{{\mathsf{g}}}
\newcommand{\tl}{\widetilde}
\newcommand{\Fr}{{\mathsf{F}}}
\newcommand{\Hr}{{\mathsf{H}}}
\newcommand{\Er}{{\mathsf{E}}}
\newcommand{\XN}{X_N}
\newcommand{\Xk}{X_k}
\newcommand{\be}{\beta}
\newcommand{\s}{\mathsf{s}}
\newcommand{\Tc}{T_\mu}
\newcommand{\Pp}{\mathcal{P}}
\newcommand{\Chi}{\chi^2}
\newcommand{\TV}{\operatorname{TV}}
\Crefname{thm}{Theorem}{Theorems}
\crefname{thm}{theorem}{theorems}
\Crefname{prop}{Proposition}{Propositions}
\crefname{prop}{proposition}{propositions}
\Crefname{lemma}{Lemma}{Lemmas}
\crefname{lemma}{lemma}{lemmas}
\Crefname{cor}{Corollary}{Corollaries}
\crefname{cor}{corollary}{corollaries}
\Crefname{remark}{Remark}{Remarks}
\crefname{remark}{remark}{remarks}

\DeclareMathOperator{\supp}{supp}
\DeclareMathOperator{\Var}{Var}
\DeclareMathOperator{\Tr}{Tr}
\newcommand{\detTwo}{\operatorname{det}_2}
\newcommand{\one}{\mathbf 1}

% Multi-stage proofs: use only when the logical decomposition materially helps.
\newcommand{\step}[1]{\par\medskip\noindent\textbf{#1}}
\newcommand{\parheading}[1]{\par\medskip\noindent\textbf{#1}\ }

\title[Sharp two-sided size of chaos estimates]{Sharp two-sided size of chaos estimates for diffusive gases in the Hilbert--Schmidt regime}

\author[M.G. Delgadino]{Matias G. Delgadino}
\address{Matias G. Delgadino, Department of Mathematics, The University of Texas at Austin, PMA 8.100, 2515 Speedway, Austin, TX 78712, USA}
\email{matias.delgadino@math.utexas.edu}
\urladdr{https://math.utexas.edu/directory/matias-delgadino}

\author[R. Gvalani]{Rishabh Gvalani}
\address{Rishabh Gvalani, School of Mathematics, University of Edinburgh, James Clerk Maxwell Building, Peter Guthrie Tait Road, Edinburgh EH9 3FD, United Kingdom}
\email{rishabh.gvalani@ed.ac.uk}
\urladdr{https://webhomes.maths.ed.ac.uk/~rgvalani/}

\author[M. Rosenzweig]{Matthew Rosenzweig}
\address{Matthew Rosenzweig, Department of Mathematical Sciences, Carnegie Mellon University, 7127 Wean Hall, 5000 Forbes Avenue, Pittsburgh, PA 15213, USA}
\email{mrosenz2@andrew.cmu.edu}
\urladdr{https://matthewrosenzweigwork-max.github.io/}

\thanks{M.G.D. was supported by NSF grant DMS-2205937. RG is partially funded by the Deutsche Forschungsgemeinschaft (DFG, German Research Foundation) - SPP 2410 Hyperbolic Balance Laws in Fluid Mechanics: Complexity, Scales, Randomness (CoScaRa). M.R. was supported by NSF grants DMS-2441170, DMS-2345533, and DMS-2342349.}

\hypersetup{
  colorlinks=true,
  linkcolor=blue!45!black,
  citecolor=magenta,
  urlcolor=blue!45!black,
  pdftitle={Sharp two-sided size of chaos estimates for diffusive gases in the Hilbert-Schmidt regime},
  pdfauthor={Matias G. Delgadino, Rishabh Gvalani, and Matthew Rosenzweig}
}

% Public working-draft presentation (2026-10-28).
\usepackage{draftwatermark}
\SetWatermarkText{DRAFT}
\SetWatermarkAngle{45}
\SetWatermarkScale{0.8}
\SetWatermarkColor[gray]{0.92}

\begin{document}

\raggedbottom

\begin{abstract}
We identify the complete first-order correction to independence for equilibrium mean-field Gibbs ensembles and quantify its approximation uniformly in the marginal size.  The correction is an explicit sum of singleton and Hoeffding-centered pair terms.  An abstract theorem for symmetric interactions with square-integrable centered kernels, suitable row and source moments, and an amplified centered-pair exponential bound gives algebraic approximation rates without a sign condition on the centered interaction operator.  For repulsive logarithmic and Riesz gases below the Hilbert--Schmidt threshold, these hypotheses hold at every fixed positive inverse temperature for bounded background densities, with a polynomial-moment condition in the logarithmic case and no additional spatial moment in the Riesz case.  The approximation determines the leading constants in chi-square divergence, relative entropy, and total variation for fixed marginals, and the normalized limits for every growing marginal containing a vanishing fraction of the particles.  Its proof combines quantitative Laplace comparison, correlated replicas, screened response, and a second-order cavity expansion, followed by orthogonal Hoeffding projection.  Under a separate, more flexible cavity criterion, we also prove matching order bounds: total variation is of order $k/N$, and relative entropy and chi-square divergence are of order $k^2/N^2$, uniformly for $2\leq k\leq N$; the upper bounds include $k=1$.  Under compatible moment and cavity assumptions, a subquadratic partition-function moment suffices for order-$N^{-1}$ convergence of every fixed marginal in $L^\tau$ and total variation, for $1<\tau<2$.  At the critical Riesz exponent $\s=\d/2$, logarithmically enhanced fixed-particle lower bounds and a sequential Hellinger argument obstruct the classical scales.  In the strictly supercritical range $\d/2<\s<\d$, depleted close pairs give the corresponding mesoscopic obstruction.  The quantitative exponents are sufficient and are not claimed optimal.
\end{abstract}

\maketitle




% STAGE2 MOVE: sec:introduction-main-results
\section{Introduction}
\label{sec:introduction-main-results}

% STAGE2 MOVE: A1-three-questions-opening
We study the leading dependence among labeled particles in fixed-temperature mean-field Gibbs ensembles.  What is the first correction to independence, and how accurately does it describe a marginal whose size grows with the number of particles?  Which leading constants does this correction determine in statistical divergences?  How does the singularity of the interaction limit these conclusions?  Our principal results identify and quantify the marginal correction under abstract analytic hypotheses, with applications to logarithmic and Riesz interactions.  A separate order theory and ultraviolet obstructions complement these asymptotic results.

% STAGE2 MOVE: A1-physical-setup-and-order
Let \(W:(\R^\d)^2\to\R\) be a symmetric two-body interaction potential and let \(V:\R^\d\to\R\) be a confining potential.  In the mean-field scaling, the Hamiltonian of an \(N\)-particle configuration \(\XN=(x_1,\ldots,x_N)\) is
\begin{align}\label{eq:intro-physical-Hamiltonian}
 \mathcal H_N(\XN)
 \coloneqq
 \frac{1}{2N}\sum_{1\le i\ne j\le N}W(x_i,x_j)
 +\sum_{i=1}^N V(x_i),
\end{align}
and the associated canonical Gibbs measure at inverse temperature \(\be>0\) is
\begin{align}\label{eq:intro-physical-Gibbs}
 \md\Pb_{N,\be}(\XN)
 =\frac{1}{Z_{N,\be}}e^{-\be\mathcal H_N(\XN)}\,\md\XN.
\end{align}
When the corresponding overdamped interacting diffusion is well posed, \eqref{eq:intro-physical-Gibbs} is its equilibrium law.  We work at fixed inverse temperature as \(N\to\infty\), which we call the \emph{diffusive temperature regime}.

An absolutely continuous thermal equilibrium measure \(\mu\) is characterized by the Euler--Lagrange relation
\begin{align}\label{eq:intro-mean-field-equilibrium}
 \md\mu(x)=Z_\mu^{-1}e^{-\be(V(x)+h_\mu(x))}\,\md x,
 \qquad
 h_\mu(x)\coloneqq\int_{\R^\d}W(x,y)\,\md\mu(y).
\end{align}
When this thermal equilibrium measure is unique, qualitative chaos asserts that, for every fixed \(k\), the labeled marginal \(\Pb_{N,\be}^{(k)}\) converges to \(\mu^{\otimes k}\).  Equivalently, the empirical measure converges in law to the deterministic measure \(\mu\); see Kac~\cite{Kac1956}, Sznitman~\cite{Sznitman1991}, and Hauray--Mischler~\cite{HaurayMischler2014}.  The quantitative question is finer: how much dependence remains among \(k\) labeled particles when \(k\) is fixed or grows with \(N\)?

The principal positive result is an explicit singleton-plus-pair approximation of the marginal density, with algebraic errors uniform in the number of observed particles.  It identifies the leading divergence constants for fixed marginals and, after normalization, for every growing marginal containing a vanishing fraction of the particles.  The underlying abstract theorem allows general symmetric interactions satisfying row, source, and amplified exponential-moment bounds, with no sign condition on the centered interaction operator; the logarithmic and Riesz statements are applications.  The precise results are given in \Cref{subsec:exact-fixed-marginal-profile}, and the abstract theorem in \Cref{thm:abstract-quantitative-marginal}.

Following Lacker~\cite{Lacker2022}, it is useful to distinguish this \emph{local} question from the \emph{global} problem of controlling the full law or the empirical measure.  The distinction is already visible in relative entropy.  A global estimate
\begin{align}\label{eq:intro-global-entropy-O1}
 \Hr(\Pb_{N,\be}\mid\mu^{\otimes N})=O(1)
\end{align}
combined with entropy subadditivity gives only
\begin{align}\label{eq:intro-subadditive-local-rate}
 \Hr(\Pb_{N,\be}^{(k)}\mid\mu^{\otimes k})=O\pa{\frac{k}{N}}.
\end{align}
For regular continuous Gibbs measures, Lacker showed that this loss is not intrinsic: under a perturbative force-level transport-information hypothesis, he proved a relative-entropy upper bound of quadratic order in \(k/N\) and exhibited a Gaussian model attaining that order, showing that the exponent cannot in general be improved~\cite{Lacker2022}.  Ren and Wang subsequently obtained the same entropy scale nonperturbatively in several convex regimes by means of a hierarchy of conditional entropies~\cite{RenWang2024}.  Singular logarithmic and Riesz interactions lie outside both theories.

The independent order theorem \Cref{thm:chi2-chaos-two-sided} and its model verification in \Cref{sec:verification-log-riesz} give the following sharp scales.  For every fixed $\be>0$, a bounded background density, and either $0<\s<\d/2$ or the logarithmic moment criterion of \Cref{lem:log-polynomial-moment-criterion}, there are positive comparison constants and $N_0$, independent of $k$, such that, for every $N\geq N_0$ and $2\leq k\leq N$,
\begin{align}\label{eq:intro-informal-scales}
 \TV\pa{\Pb_{N,\be}^{(k)},\mu^{\otimes k}}
 &\asymp \frac{k}{N},
 &
 \Hr\pa{\Pb_{N,\be}^{(k)}\mid\mu^{\otimes k}}
 &\asymp \frac{k^2}{N^2},
 &
 \Chi\pa{\Pb_{N,\be}^{(k)}\mid\mu^{\otimes k}}
 &\asymp \frac{k^2}{N^2}.
\end{align}
The corresponding upper bounds also hold for $k=1$.  The logarithmic order criterion is weaker than the hypothesis used for the quantitative marginal approximation.  These conclusions do not assert a positive singleton lower bound: for a uniform background and a translation-invariant interaction on the flat torus, the one-point marginal is exactly the background.


For a Riesz singularity of order $\s$, the centered pair interaction is locally square-integrable precisely when $2\s<\d$.  At and above this threshold, near-diagonal configurations produce lower bounds that rule out the classical fixed-particle chi-square and entropy scales and the uniform total-variation scale.  The latter obstruction occurs for a vanishing fraction of the particles.  These are one-sided conclusions: we do not determine matching ultraviolet upper bounds or limiting profiles.  The optimal convergence rates and extensions under weaker exponential-moment hypotheses remain open.


% STAGE2 MOVE: A6-historical-orientation-footnote
The pair-correlation and quadratic-information mechanisms underlying these scales have antecedents in the Bogolyubov hierarchy for connected correlations~\cite{Bogolyubov1962} and in classical multiparticle expansions of equilibrium entropy.\footnote{The Bogolyubov hierarchy predicts that the connected two-particle correction has order \(N^{-1}\), and more generally that the connected \((m+1)\)-particle correlation has order \(N^{-m}\); see~\cite{Bogolyubov1962}.  Duerinckx proved this scaling, in suitable negative Sobolev norms, for smooth classical mean-field dynamics~\cite{DuerinckxGlauber2021}.  In the entropy expansions of H.~S.~Green, Stratonovich, and Nettleton--Green~\cite{Green1952,Stratonovich1955,NettletonGreen1958}, the two-body contribution is governed by the relative-entropy integrand \(g\log g-g+1\) and hence begins quadratically at the uncorrelated state \(g=1\).  Taken together, these observations suggest the \(k^2/N^2\) local entropy scale, but they do not provide a nonasymptotic relative-entropy estimate for \(k\)-particle marginals.  The present paper shows that this \(k^2/N^2\) prediction fails for sufficiently singular interactions; the older statistical-mechanics expansions did not isolate the relevant regularity threshold.}

% STAGE2 MOVE: sec:main-results-proof-strategy
\section{Main results and proof strategy}
\label{sec:main-results-proof-strategy}

% STAGE2 MOVE: subsec:modulated-formulation
\subsection{Modulated Gibbs formulation and divergence conventions}
\label{subsec:modulated-formulation}

% STAGE2 MOVE: A1-modulated-setup-and-divergences
We now introduce the centered formulation used throughout the paper.  Let $\mu\in\Pp(\R^\d)$ be a Borel probability measure, and let $W$ be real, symmetric, measurable, and integrable with respect to $\mu^{\otimes2}$.  No absolute continuity or bounded-density assumption is needed for these definitions.  When a measure is absolutely continuous, we use the same symbol for the measure and its Lebesgue density.  The integrability assumption makes the potentials below finite almost everywhere; all identities are understood with respect to the appropriate product of $\mu$.  Set

\begin{align}\label{eq:h-I-tildeg-statement}
 h_\mu(x)&\coloneqq\int_{\R^\d}W(x,y)\,\md\mu(y),
 &
 I_\mu&\coloneqq\frac12\int_{(\R^\d)^2}W(x,y)\,\md\mu(x)\md\mu(y),
\end{align}
and define the Hoeffding centering of a two-body kernel by
\begin{align}\label{eq:tildeg-def-statement}
 \tl f(x,y)
 \coloneqq
 f(x,y)-\int f(x,z)\,\md\mu(z)-\int f(z,y)\,\md\mu(z)
 +\int f(z,w)\,\md\mu^{\otimes2}(z,w).
\end{align}
Thus,
\begin{align}\label{eq:centered-W-explicit}
 \tl W(x,y)=W(x,y)-h_\mu(x)-h_\mu(y)+2I_\mu,
 \qquad
 \int\tl W(x,y)\,\md\mu(y)=0.
\end{align}
The modulated energy is
\begin{align}\label{eq:FN-def-intro}
 \Fr_N(\XN,\mu)
 =\frac{1}{2N^2}\sum_{1\le i\ne j\le N}\tl W(x_i,x_j)
 -\frac{1}{N^2}\sum_{i=1}^N h_\mu(x_i)
 +\frac{I_\mu}{N},
\end{align}
and its partition function takes values in $(0,\infty]$; whenever it is finite, the associated Gibbs measure is defined by

\begin{align}\label{eq:K-Q-def-intro}
 \K_{N,\be}(\mu)
 &\coloneqq
 \int e^{-\be N\Fr_N(\XN,\mu)}\,\md\mu^{\otimes N}(\XN),\\
 \md\Q_{N,\be}(\mu)(\XN)
 &\coloneqq
 \frac{1}{\K_{N,\be}(\mu)}e^{-\be N\Fr_N(\XN,\mu)}\,\md\mu^{\otimes N}(\XN).
\end{align}
The hypotheses of each result below ensure finiteness of every partition function used to define a Gibbs measure in that result.  The exact splitting formula
\begin{align}\label{eq:intro-splitting-formula}
 \frac{1}{2N}\sum_{1\le i\ne j\le N}W(x_i,x_j)
 =N\Fr_N(\XN,\mu)+\sum_{i=1}^N h_\mu(x_i)-NI_\mu
\end{align}
shows that, when \(\mu\) satisfies \eqref{eq:intro-mean-field-equilibrium}, the modulated Gibbs measure \(\Q_{N,\be}(\mu)\) coincides with the physical Gibbs ensemble \(\Pb_{N,\be}\).  More generally, \(\Q_{N,\be}(\mu)\) is the natural Gibbs tilt of the prescribed product background \(\mu^{\otimes N}\).  Its normalizing constant \(\K_{N,\be}(\mu)\) is the exact global cost of this change of reference.  Uniform control of this cost bounds the full Gibbs density, but it does not by itself yield the sharp local gain in \eqref{eq:intro-informal-scales}.  For the splitting formula and the modulated-energy formalism in Coulomb and Riesz gases, see \cite[Chapter~5, Section~5.1]{SerfatyLN}.

For \(1\le k\le N\), the labeled \(k\)-point marginal is
\begin{align}\label{eq:labeled-k-marginal-def}
 \Q_{N,\be}^{(k)}(\mu)
 \coloneqq
 (\pi_{1,\ldots,k})_\#\Q_{N,\be}(\mu),
 \qquad
 \pi_{1,\ldots,k}(x_1,\ldots,x_N)\coloneqq(x_1,\ldots,x_k).
\end{align}
For \(\nu\ll\pi\), we respectively denote the chi-square divergence, relative entropy or Kullback--Leibler divergence, and total-variation distance by
\begin{align}\label{eq:divergence-definitions-intro}
 \Chi(\nu\mid\pi)
 &\coloneqq\int\pa{\frac{\md\nu}{\md\pi}-1}^2\,\md\pi,\\
 \Hr(\nu\mid\pi)
 &\coloneqq\int\frac{\md\nu}{\md\pi}\log\frac{\md\nu}{\md\pi}\,\md\pi,\\
 \TV(\nu,\pi)
 &\coloneqq\frac12\int\abs{\frac{\md\nu}{\md\pi}-1}\,\md\pi.
\end{align}

% STAGE2 MOVE: A1-density-definition-before-profile
We write
\begin{align}\label{eq:gibbs-marginal-density-ratio}
 f_{N,\be}^{(k)}
 \coloneqq
 \frac{\md\Q_{N,\be}^{(k)}(\mu)}{\md\mu^{\otimes k}},
\end{align}
and set $f_{N,\be}=f_{N,\be}^{(N)}$.


% STAGE2 MOVE: subsec:main-order-bounds
\subsection{Quantitative marginal approximation}
\label{subsec:exact-fixed-marginal-profile}

% STAGE2 MOVE: A1-quantitative-statement-transition
The next two theorems are applications of \Cref{thm:abstract-quantitative-marginal}, which treats general real symmetric interactions under analytic row, source, and exponential-moment assumptions.  No sign condition is imposed on the centered interaction operator.  The first theorem gives the complete marginal approximation and its divergence errors; the second identifies the normalized divergence limits for growing marginals.  Their proofs are common to the abstract class; \Cref{subsubsec:signed-model-verification} verifies the logarithmic and Riesz hypotheses.

For the model statements, use
\begin{align}\label{eq:intro-log-Riesz-kernel}
 \g_\s(z)
 =\begin{cases}
 -\log\abs{z}, & \s=0,\\[0.2em]
 \s^{-1}\abs{z}^{-\s}, & 0<\s<\d.
 \end{cases}
\end{align}
We introduce the common profiles and rate before stating the two theorems.

% STAGE2 MOVE: A1-integral-operator-definition
For a symmetric centered kernel $\tl W\in L^2(\mu^{\otimes2})$, we use the integral operator
\begin{align}\label{eq:T-def}
 \Tc f(x)\coloneqq \int_{\R^\d}\tl W(x,y)f(y)\,\md\mu(y).
\end{align}
Since \(\tl W\in L^2(\mu^{\otimes2})\), the operator \(\Tc\) is self-adjoint and Hilbert--Schmidt on \(L^2(\mu)\).  Canonicity gives \(\Tc\one=0\), but no sign condition is imposed on the remaining spectrum.

% STAGE2 MOVE: A1-model-resolvent-setup
Whenever $\mathrm{Id}+\be T_\mu$ is strictly positive and boundedly invertible, set


% STAGE2 MOVE: A1-resolvent-definition
\begin{align}\label{eq:weak-second-resolvent-def}
 R_{\be,\mu}\coloneqq(\mathrm{Id}+\be\Tc)^{-1}.
\end{align}

% STAGE2 MOVE: A1-pair-kernel-setup
The operator
\begin{equation}\label{eq:quantitative-screened-operator}
 B_{\be,\mu}\coloneqq R_{\be,\mu}-\mathrm{Id}=-\be T_\mu R_{\be,\mu}
\end{equation}
is Hilbert--Schmidt and has a real symmetric canonical kernel $b_{\be,\mu}$.  For the repulsive log/Riesz models, the required invertibility holds at every $\be\geq0$ by nonnegativity of $T_\mu$.  Under the abstract quantitative hypotheses in \Cref{thm:abstract-quantitative-marginal}, it follows instead from the amplified exponential bound.  Define the centered one-body function and kernel sections by
% STAGE2 MOVE: A1-common-profile-and-log-constants
\begin{align}
 \psi_\mu&\coloneqq h_\mu-2I_\mu,\label{eq:exact-psi-def}\\
 w_x(y)&\coloneqq\tl W(x,y).\label{eq:exact-row-def}
\end{align}
The row response and the singleton profile are
\begin{align}
 \omega_{\be,\mu}(x)
 &\coloneqq\be^2\ipp*{w_x,R_{\be,\mu}w_x}_{L^2(\mu)},
 \qquad
 \omega_{\be,\mu}^{\circ}
 \coloneqq\omega_{\be,\mu}-\int\omega_{\be,\mu}\,\md\mu,
 \label{eq:exact-omega-def}\\
 a_{\be,\mu}
 &\coloneqq R_{\be,\mu}\left(\be\psi_\mu+\frac12\omega_{\be,\mu}^{\circ}\right).
 \label{eq:exact-singleton-profile-def}
\end{align}
The estimates below show that these functions belong to $L^2(\mu)$ under either the model or the abstract quantitative hypotheses.
  Both $\psi_\mu$ and $\omega_{\be,\mu}^{\circ}$ are centered, and $R_{\be,\mu}$ preserves $L^2_0(\mu)$; hence, $a_{\be,\mu}$ is centered.  For each integer $k\geq1$, define
\begin{align}
 G_{k,\be,\mu}(x_1,\ldots,x_k)
 &\coloneqq\sum_{i=1}^k a_{\be,\mu}(x_i)
 +\sum_{1\leq i<j\leq k}b_{\be,\mu}(x_i,x_j).
 \label{eq:exact-fixed-profile-def}
\end{align}

Canonicity makes distinct singleton and pair terms orthogonal.  Consequently,
\begin{equation}\label{eq:exact-Gk-L2}
 \norm{G_{k,\be,\mu}}_{L^2(\mu^{\otimes k})}^2
 =k\norm{a_{\be,\mu}}_{L^2(\mu)}^2
 +\binom{k}{2}\norm{b_{\be,\mu}}_{L^2(\mu^{\otimes2})}^2.
\end{equation}
The kernel/operator correspondence gives $\norm{b_{\be,\mu}}_2=\norm{B_{\be,\mu}}_{\mathrm{HS}}$.

For the logarithmic model, we use the explicit cavity-envelope constants
\begin{align}
 C_{\mathrm{loc}}
 &\coloneqq\norm{\mu}_{L^\infty}
 \int_{\abs{z}<1}-\log\abs{z}\,\md z,
 \label{eq:exact-log-local-constant}\\
 \mathsf c_{\log,\be,\mu}
 &\coloneqq\be\left(3C_{\mathrm{loc}}+4\abs{I_\mu}+1\right).
 \label{eq:exact-log-cavity-constant}
\end{align}
This agrees with \eqref{eq:log-cavity-weight-constant} for the envelope constants specified in \Cref{lem:exact-log-analytic}.

In the Riesz case, fix
\begin{equation}\label{eq:quantitative-alpha-riesz}
 0<\alpha\leq1,\qquad 2+\alpha<\d/\s,
\end{equation}
and, in the logarithmic case, set $\alpha=1$.  Define the model rate
\begin{equation}\label{eq:quantitative-profile-rate}
 \mathfrak r_N\coloneqq
 \begin{cases}
 N^{-\alpha/8},&0<\s<\d/2,\\
 N^{-1/8},&\s=0.
 \end{cases}
\end{equation}
For example, $\alpha=\min\{1,(\d/\s-2)/2\}$ is admissible in the Riesz
case.  Constants in all the estimates below may depend on this fixed
choice of $\alpha$, on the model parameters, and on the stated moment
data, but not on $N$ or $k$.  The exponents and the logarithmic moment
threshold are sufficient and are not claimed optimal.
\begin{thm}[Quantitative exact marginal profile]
\label{thm:quantitative-exact-marginal-profile}
Assume one of the following model hypotheses:
\begin{enumerate}[(i)]
\item $\d\geq1$, $0<\s<\d/2$, $W(x,y)=1/(\s\abs{x-y}^{\s})$, $\mu$ is a probability measure on $\R^\d$ with bounded density, and $\be\geq0$;
\item $\d\geq1$, $\s=0$, $W(x,y)=-\log\abs{x-y}$, $\be>0$, $\mu$ is a probability measure on $\R^\d$ with bounded density, and, for some $q>0$,
\begin{align}
 &\int_{\R^\d}(e+\abs{x})^q\,\md\mu(x)<\infty,
 \label{eq:exact-log-polynomial-moment}\\
 &q>\max\left\{4\be,\,36\mathsf c_{\log,\be,\mu}+8\right\}.
 \label{eq:exact-log-moment-threshold}
\end{align}
\end{enumerate}
Then, for every $N\geq2$ and $1\leq k\leq N$,
\begin{equation}\label{eq:quantitative-natural-profile}
 \left\|\frac Nk(f_{N,\be}^{(k)}-1)-\frac1kG_{k,\be,\mu}\right\|_{L^2(\mu^{\otimes k})}
 \leq C\left(\mathfrak r_N+\sqrt{\frac{k}{N}}\right).
\end{equation}
Moreover,
\begin{align}
 \left|N^2\Chi\!\left(\Q_{N,\be}^{(k)}(\mu)\mid\mu^{\otimes k}\right)-\norm{G_{k,\be,\mu}}_2^2\right|
 &\leq C\left(k\mathfrak r_N+k^2\mathfrak r_N^2+\frac{k^3}{N}\right),
 \label{eq:quantitative-uniform-chi}\\
 \left|N^2\Hr\!\left(\Q_{N,\be}^{(k)}(\mu)\mid\mu^{\otimes k}\right)-\frac12\norm{G_{k,\be,\mu}}_2^2\right|
 &\leq C_\alpha\left(k\mathfrak r_N+k^2\mathfrak r_N^2+k^2(k/N)^\alpha\right),
 \label{eq:quantitative-uniform-entropy}\\
 \left|N\TV\!\left(\Q_{N,\be}^{(k)}(\mu),\mu^{\otimes k}\right)-\frac12\norm{G_{k,\be,\mu}}_1\right|
 &\leq C\left(k\mathfrak r_N+\frac{k^{3/2}}{\sqrt N}\right).
 \label{eq:quantitative-uniform-TV}
\end{align}
\end{thm}

For fixed $k$, the estimates identify the strong first-order marginal limit and the leading constants in all three divergences; the constants in the two quadratic divergences are explicit in \eqref{eq:exact-Gk-L2}.  The logarithmic moment threshold is sufficient and nonoptimal, while no additional moment assumption is imposed in the Riesz case.
  The logarithmic conclusion at $\be=0$ is immediate, since the Gibbs law is the product law and $a_{0,\mu}$, $b_{0,\mu}$, and $G_{k,0,\mu}$ vanish.

Let $(\lambda_j(B_{\be,\mu}))_j$ be the nonzero eigenvalues of $B_{\be,\mu}$, counted with multiplicity, and let $(Z_j)_j$ be independent standard real Gaussian variables.  Define the centered Gaussian quadratic form
\begin{equation}\label{eq:quantitative-centered-Gaussian-quadratic-form}
 \mathcal Q_{\be,\mu}\coloneqq\sum_j\lambda_j(B_{\be,\mu})(Z_j^2-1)
 \quad\text{as a limit in }L^2.
\end{equation}
This definition uses only the Hilbert--Schmidt property; $B_{\be,\mu}$ need not be trace class.  In particular,
\begin{equation}\label{eq:quantitative-Gaussian-variance}
 \E\left[\mathcal Q_{\be,\mu}^2\right]=2\norm{B_{\be,\mu}}_{\mathrm{HS}}^2.
\end{equation}

\begin{thm}[Quantitative divergence limits for growing marginals]
\label{thm:quantitative-growing-marginal-constants}
Under the same hypotheses, for every $N\geq2$ and $1\leq k\leq N$,
\begin{align}
 \left|\frac{N^2}{k^2}\Chi\!\left(\Q_{N,\be}^{(k)}(\mu)\mid\mu^{\otimes k}\right)-\frac12\norm{B_{\be,\mu}}_{\mathrm{HS}}^2\right|
 &\leq C\left(\mathfrak r_N^2+\frac1k+\frac{k}{N}\right),
 \label{eq:quantitative-growing-chi}\\
 \left|\frac{N^2}{k^2}\Hr\!\left(\Q_{N,\be}^{(k)}(\mu)\mid\mu^{\otimes k}\right)-\frac14\norm{B_{\be,\mu}}_{\mathrm{HS}}^2\right|
 &\leq C_\alpha\left(\mathfrak r_N^2+\frac1k+(k/N)^\alpha\right),
 \label{eq:quantitative-growing-entropy}\\
 \left|\frac Nk\TV\!\left(\Q_{N,\be}^{(k)}(\mu),\mu^{\otimes k}\right)-\frac14\E\left[\abs{\mathcal Q_{\be,\mu}}\right]\right|
 &\leq C_\alpha\left(\mathfrak r_N+\sqrt{\frac{k}{N}}+k^{-\alpha/(2(2+\alpha))}\right).
 \label{eq:quantitative-growing-TV}
\end{align}
Thus, all three normalized divergences converge to the displayed constants along every sequence $k_N\to\infty$ with $k_N=o(N)$.  In the logarithmic case, the last error in \eqref{eq:quantitative-growing-TV} is $k^{-1/6}$, and the right-hand sides of \eqref{eq:quantitative-growing-chi} and \eqref{eq:quantitative-growing-entropy} are both bounded by $C(\mathfrak r_N^2+k^{-1}+k/N)$.
\end{thm}

% STAGE2 MOVE: A1-natural-scale-interpretation
The normalized limits require only $k_N\to\infty$ and $k_N/N\to0$.
Requiring the undivided density error
$\norm{N(f_{N,\be}^{(k)}-1)-G_{k,\be,\mu}}_2$ and the divergence errors
in \eqref{eq:quantitative-uniform-chi}--\eqref{eq:quantitative-uniform-TV}
to vanish is stronger; \Cref{cor:quantitative-literal-growing-window}
gives sufficient windows.  On their common hypotheses, the quantitative
estimates recover the order bounds.  The separate criterion in \Cref{thm:chi2-chaos-two-sided} retains the weaker logarithmic hypothesis and its broader weighted cavity formulation; its model specializations are recorded in \Cref{sec:verification-log-riesz}.
% STAGE2 MOVE: subsec:main-ultraviolet-obstructions
\subsection{Ultraviolet obstructions}
\label{subsec:main-ultraviolet-obstructions}

% STAGE2 MOVE: A1-ultraviolet-statement-and-interpretation
The exponent \(\s=\d/2\) is the exact Hilbert--Schmidt threshold.  For bounded densities, the centered Riesz kernel belongs to \(L^2(\mu^{\otimes2})\) precisely when \(2\s<\d\).  The authors' companion paper~\cite[Proposition~5.10]{DGR26} also proves that the modulated partition function diverges for every fixed \(\be>0\) throughout \(\d/2\le\s<\d\).  Below the threshold, its boundedness is an input to our upper bounds.  At and above the threshold, however, partition-function divergence does not imply a local marginal lower bound: the results below are proved independently from the near-diagonal geometry of the Riesz interaction.  The two papers therefore identify global and local manifestations of the same ultraviolet obstruction, but neither supercritical statement is a formal consequence of the other.

The upper endpoint \(\s<\d\) in the next theorem is imposed by the modulated mean-field formulation, not by the existence of the bare finite-particle repulsive Gibbs law.  The latter may remain meaningful in the hypersingular regime \(\s\ge\d\), but the continuum potential \(h_\mu\) and energy \(I_\mu\) are then generally infinite for an absolutely continuous background.  Consequently, the Hoeffding centering and splitting formula used here are not defined without a different scaling and renormalization.

\begin{thm}[Ultraviolet obstructions at and above the Hilbert--Schmidt threshold]\label{thm:intro-ultraviolet-sharpness}
Let \(\be>0\), let \(\d/2\le\s<\d\), set \(W(x,y)=\g_\s(x-y)\), and let \(\mu\) be a bounded probability density.  Set
\begin{align}\label{eq:intro-uv-scale}
 \mathfrak a_N
 \coloneqq
 \begin{cases}
  N^{-2}\log N,&\s=\d/2,\\
  N^{-\d/\s},&\d/2<\s<\d.
 \end{cases}
\end{align}
At the critical exponent \(\s=\d/2\), assume additionally that \(\mu\ge a>0\) on some ball.
\begin{enumerate}[(i)]
\item For every fixed \(k\ge2\), there are positive constants \(c_\chi,c_H,c_{\mathrm{TV}}\) such that, for all sufficiently large \(N\),
\begin{subequations}\label{eq:intro-uv-fixed-k-lower}
\begin{align}
 \Chi\pa{\Q_{N,\be}^{(k)}(\mu)\mid\mu^{\otimes k}}
 &\ge c_\chi\mathfrak a_N,\label{eq:intro-uv-fixed-k-chi}\\
 \Hr\pa{\Q_{N,\be}^{(k)}(\mu)\mid\mu^{\otimes k}}
 &\ge c_H\mathfrak a_N,\label{eq:intro-uv-fixed-k-H}\\
 \TV\pa{\Q_{N,\be}^{(k)}(\mu),\mu^{\otimes k}}
 &\ge c_{\mathrm{TV}}\mathfrak a_N.\label{eq:intro-uv-fixed-k-TV}
\end{align}
\end{subequations}
In particular, the fixed-particle \(N^{-2}\) chi-square and entropy scales fail at and above \(\s=\d/2\).

\item There are \(c_0,c_1>0\) such that, for all sufficiently large \(N\),
\begin{align}\label{eq:intro-mesoscopic-TV-lower}
 \TV\pa{\Q_{N,\be}^{(k)}(\mu),\mu^{\otimes k}}
 \ge c_1k^2\mathfrak a_N
\end{align}
for every integer
\begin{align}\label{eq:intro-mesoscopic-k-window}
 2\le k\le c_0\mathfrak a_N^{-1/2}.
\end{align}
Consequently, for every sufficiently small fixed \(\kappa>0\), the sequence
\begin{align}\label{eq:intro-mesoscopic-kN}
 k_N\coloneqq\left\lfloor\kappa\mathfrak a_N^{-1/2}\right\rfloor
\end{align}
satisfies
\begin{align}
 \liminf_{N\to\infty}\TV\pa{\Q_{N,\be}^{(k_N)}(\mu),\mu^{\otimes k_N}}&>0,
 \label{eq:intro-mesoscopic-positive-liminf}\\
 \frac{k_N}{N}&\longrightarrow0.
 \label{eq:intro-mesoscopic-vanishing-fraction}
\end{align}
Thus, the uniform \(k/N\) total-variation bound fails throughout \(\d/2\le\s<\d\).  At criticality, \(\mathfrak a_N^{-1/2}=N/\sqrt{\log N}\); in the strictly supercritical range, \(\mathfrak a_N^{-1/2}=N^{\d/(2\s)}\).
\end{enumerate}
\end{thm}

By Pinsker's inequality and the elementary bound \(\Chi(\nu\mid\pi)\ge4\TV(\nu,\pi)^2\), part~\textup{(ii)} also implies that the corresponding \(k_N^2/N^2\) entropy and chi-square scales fail.  Indeed, \(k_N^2/N^2\asymp1/\log N\to0\) at \(\s=\d/2\), while \(k_N^2/N^2\asymp N^{\d/\s-2}\to0\) when \(\d/2<\s<\d\).  Thus, the Hilbert--Schmidt threshold is the exact boundary for the classical local scales in all three divergences: the fixed-particle chi-square and entropy scales fail at the threshold, and the uniform total-variation scale fails at and above it.

The lower comparison in \Cref{cor:f-divergence-consequences} shows, in particular, that the critical sequence \(k_N\asymp N/\sqrt{\log N}\) also rules out the classical uniform \((k/N)^\tau\) scale for every centered \(L^\tau\)-divergence, \(1\le\tau\le2\).

% STAGE2 MOVE: subsec:related-work
\subsection{Proof strategy and logical structure}
\label{subsec:global-proof-strategy}

The proof of \Cref{thm:abstract-quantitative-marginal} identifies the first two Hoeffding components of the Gibbs density and controls the remaining components together.  We first explain why this suffices for the marginal approximation, then describe how the two leading kernels are obtained.  \Cref{fig:intro-proof-architecture} distinguishes this argument from the complementary order lower bounds and ultraviolet obstructions.

\begin{figure}[!htbp]
\centering
\begin{tikzpicture}[
 font=\small,
 box/.style={draw=black!55,rounded corners=2pt,fill=gray!5,align=center,inner sep=5pt,text width=5.9cm},
 wide/.style={box,text width=13.2cm},
 result/.style={wide,fill=gray!12,line width=.7pt},
 arr/.style={-{Stealth[length=2mm]},semithick,draw=black!60}]
\node[box] (inputs) at (-3.55,0) {Common density moments, one-point bound, and Hoeffding projection};
\node[box] (smooth) at (3.55,0) {Smooth Gaussian comparison};
\node[box] (scalar) at (3.55,-1.15) {Scalar exponential truncation:\\quantitative Laplace estimate};
\node[box] (responses) at (3.55,-2.45) {Replicas: squared pair norm\\Response: inner product with $b_{\be,\mu}$};
\node[box] (pair) at (3.55,-3.8) {Extraction and polarization:\\strong pair approximation};
\node[box] (sing) at (-3.55,-3.8) {Second-order cavity expansion:\\singleton approximation};
\node[result] (marg) at (0,-5.2) {Orthogonal assembly and collective remainder bound:\\quantitative marginal approximation};
\node[result] (div) at (0,-6.8) {Chi-square: squared norm; entropy: profile-moment comparison\\Total variation: $L^1$ approximation and absolute first moment};
\draw[arr] (smooth)--(scalar); \draw[arr] (scalar)--(responses);
\draw[arr] (responses)--(pair);
\draw[arr] (inputs.east)--(0,0)--(0,-2.45)--(responses.west);
\draw[arr] (inputs)--(sing); \draw[arr] (pair)--(sing);
\draw[arr] (pair)--(marg); \draw[arr] (sing)--(marg);
\draw[arr] (inputs.west)--(-7.0,0)--(-7.0,-5.2)--(marg.west);
\draw[arr] (marg)--(div);
\draw[arr,dashed] (smooth.east)--(7.05,0)--(7.05,-6.8)--(div.east);
\node[font=\footnotesize,rotate=90,fill=white,inner sep=2pt] at (7.42,-3.4) {Absolute first-moment comparison};
\node[box] (order) at (-3.55,-8.3) {Order lower bounds:\\Jensen gap and pair-sum tests};
\node[box] (uv) at (3.55,-8.3) {Ultraviolet obstructions:\\critical annuli; depleted close pairs};
\end{tikzpicture}
\caption{The common estimates enter pair extraction, the cavity expansion, and marginal projection.  The dashed arrow is the absolute-first-moment application of the smooth Gaussian comparison, needed for growing-marginal total variation.  The lower boxes are not inputs to the quantitative theorem.}
\label{fig:intro-proof-architecture}
\end{figure}

\parheading{Why two components suffice.}
Marginalization is conditional expectation under the product law: it removes every Hoeffding component involving an unobserved coordinate.  By exchangeability, the components involving $r$ coordinates have equal squared norms.  The fraction of their total squared norm that survives in the $k$-particle marginal is therefore
\begin{equation}\label{eq:intro-hoeffding-marginal-ratio}
 \frac{\binom{k}{r}}{\binom{N}{r}}
 \leq\left(\frac{k}{N}\right)^r,
 \qquad 1\leq r\leq k\leq N.
\end{equation}
The density-moment estimate bounds the full $N$-particle chi-square divergence and hence the sum of these squared norms.  After the first two components are removed, \eqref{eq:intro-hoeffding-marginal-ratio} bounds the squared marginal remainder by $C(k/N)^3$; see \eqref{eq:exact-higher-tail-fixed}.  At the density normalization $N/k$, its norm is at most $C\sqrt{k/N}$.  Thus, quantitative identification of the first two kernels suffices throughout every sublinear marginal range.  Higher components are present, but no separate estimate for each degree and no hierarchy of coupled marginal or correlation equations is required.

\parheading{Common upper estimates.}
The same projection argument gives the marginal upper bounds once the first component is bounded at order $N^{-1}$.  For this estimate, we condition on the other $N-1$ particles and expand the normalized conditional exponential.  Writing $u_N=f_{N,\be}^{(1)}-1$, averaging over their actual Gibbs marginal law gives
\begin{align}\label{eq:intro-screened-equation}
 \pa{\mathrm{Id}+\be\frac{N-1}{N}T_\mu}u_N
 =\frac{\be}{N}(h_\mu-2I_\mu)+r_N,
 \qquad \norm{r_N}_{L^2(\mu)}\lesssim N^{-1}.
\end{align}
The order hypotheses supply a uniform inverse bound through the modulated partition moment at $2\be$; the quantitative hypotheses supply it through their amplified centered-pair exponential moment.  The two assumptions are compared after \Cref{thm:abstract-quantitative-marginal}.  In either case, the mean response of the reservoir remains on the left-hand side of \eqref{eq:intro-screened-equation}, rather than being bounded as a perturbative error.  The inverse bound yields $\norm{u_N}_2\leq C/N$.  Combining this with the full-density bound and projection proves the all-marginal upper estimates; see \Cref{sec:gibbs-identities-upper-bounds}.

\parheading{Strong pair approximation.}
For the two-variable Hoeffding kernel $v_N$ defined in \eqref{eq:weak-second-model-v-def}, we estimate both $\norm{Nv_N}_2^2$ and $\ipp{Nv_N,b_{\be,\mu}}$.  Their limits determine the squared error through
\begin{equation}\label{eq:intro-pair-polarization}
 \norm{Nv_N-b_{\be,\mu}}_2^2
 =\norm{Nv_N}_2^2+\norm{b_{\be,\mu}}_2^2
       -2\ipp{Nv_N,b_{\be,\mu}}.
\end{equation}
The two measurements play different roles: convergence against fixed tests would identify a candidate limit but would not exclude an additional squared norm.  Replicas determine the norm, and a perturbation in the direction of the candidate kernel determines the inner product.

Both computations use the quantitative Laplace estimate of \Cref{thm:quantitative-signed-Laplace}.  A one-coordinate Gaussian replacement, in the spirit of \cite[Section~2]{Chatterjee2006}, controls smooth tests of the canonical quadratic statistic.  We apply it to a truncated scalar exponential.  The amplified exponential bound controls the discarded part under the original product law and the limiting Gaussian law, so no exponential estimate at the intermediate replacement laws is needed.  The truncation is applied to the test function, not to obtain a quantitative approximation rate for the interaction kernel.  This gives the determinant comparison without a sign condition on the centered operator.

For the norm calculation, sample a configuration under $\mu^{\otimes N}$ and form a second configuration by independently keeping each coordinate with probability $t$ or resampling it from $\mu$.  The underlying identity is the usual noise-stability identity for product resampling; its finite-space formulation is given in \cite[Proposition~8.28 and Section~8.3]{ODonnell2014}.  Canonicity makes resampling act by $t^r$ on an order-$r$ Hoeffding component.  Consequently, the expected product of the two Gibbs density ratios is a polynomial whose degree-$r$ coefficient is the total squared norm of that level, including its subset multiplicity.  The configurations themselves have product marginals; they are not independent Gibbs samples.  The Gibbs exponential expresses this overlap as a paired partition function, to which the scalar Laplace estimate applies.  Extracting its quadratic coefficient gives the pair squared norm, with the full-density bound controlling the extraction error; see \Cref{subsubsec:exact-profile-strong-pair}.

For the inner product, perturb the centered interaction in the direction $b_{\be,\mu}$ while leaving the one-body source unchanged.  Differentiation of the logarithmic partition function at zero gives $-(N-1)\ipp{v_N,b_{\be,\mu}}/2$.  Quantitative control for both signs of the perturbation permits comparison of this derivative with that of the limiting determinant.  The resulting inner product and the replica norm estimate give strong $L^2(\mu^{\otimes2})$ convergence by \eqref{eq:intro-pair-polarization}; see \Cref{subsec:screened-response-proof,subsec:strong-pair-estimate}.

\parheading{The singleton correction.}
The first-order cavity bound leaves an error of the same order as $u_N$, so identifying the singleton profile requires the quadratic term of the normalized exponential.  Its averaged contribution is one half of the spatially centered cavity second moment $q_N$ defined in \eqref{eq:exact-qN-def}.  This is a second moment under the actual Gibbs reservoir law, not a variance: centering in the trial coordinate does not make the reservoir expectation vanish.  Separating equal and distinct reservoir labels expresses it through the row square and the two-variable Hoeffding kernel.  The strong pair estimate then identifies $\omega_{\be,\mu}$ as the $L^2(\mu)$ limit of $Nq_N$, with the quantitative bound \eqref{eq:signed-cavity-variance-limit}.  Applying the screened inverse in \eqref{eq:intro-screened-equation} identifies the profile $a_{\be,\mu}$ in \eqref{eq:exact-singleton-profile-def}.  The estimates take the Gibbs expectation before the spatial norm and combine row integrability with Hilbert-valued moment bounds; no fourth pointwise moment of the singular interaction is required.

\parheading{Marginal approximation and the three divergences.}
Orthogonality combines the two kernel errors with the collective remainder already described.  This is the density approximation in \Cref{thm:quantitative-exact-marginal-profile}; the proof is given in \Cref{subsubsec:exact-profile-assembly}.  For chi-square divergence, the squared norm is exact.  The leading term follows from the profile norm identity \eqref{eq:exact-Gk-L2}, while the sharper quantitative error uses the separately proved pair squared-norm estimate, not just the strong pair bound.

For relative entropy, a quadratic comparison produces one half of the chi-square leading term.  Its justification uses a moment above two only for the explicit profile $G_{k,\be,\mu}/k$, uniformly in $k$; it requires neither pointwise closeness of the density to one nor a higher moment of the density error.  The comparison preserves the squared approximation error.  This is the role of \Cref{quant-tools:lem-profilemoments,quant-tools:lem-entropy}.

For total variation, the $L^2$ density approximation gives an $L^1$ approximation.  When $k$ grows, the normalized singleton sum vanishes and the remaining constant is the absolute first moment of the canonical pair sum.  Its comparison with $\mathcal Q_{\be,\mu}/2$ uses the same smooth Gaussian-comparison lemma as the Laplace estimate, now applied to a smooth approximation of absolute value; see \Cref{quant-tools:prop-absolute}.  The limiting quadratic form need not be Gaussian.  \Cref{subsec:growing-divergence-proof} completes these three evaluations of the growing-marginal constants.

\parheading{The complementary order theory.}
The order theorem has its own analytic hypotheses, including a weighted lower-tail alternative.  Its lower bound follows from nonzero product variance, higher-moment control, and a strict Jensen gap for the scalar log-Laplace transform.  Differentiation gives a negative pair response of order $N^{-1}$.  Testing with the centered pair sum gives the chi-square lower bound; truncating that sum gives total variation, and Pinsker's inequality gives entropy.  Under the common hypotheses and a nontrivial interaction, the quantitative estimates and contraction under marginalization also give the order lower bounds for every $2\leq k\leq N$, once $N$ is sufficiently large independently of $k$.  We prove them directly in \Cref{sec:chi2-tv-lower-bound-proof} to establish \Cref{thm:chi2-chaos-two-sided} without the quantitative theorem's additional assumptions.

The subquadratic theorem instead pairs a full-density $L^\tau$ bound with canonical ordered $U$-statistics by duality.  A refined degree-two moment estimate preserves the order-$N^{-1}$ contribution to the norm; higher fixed orders are negligible at that scale.  This proves convergence of every fixed marginal under the separate hypotheses of \Cref{thm:fixed-marginal-Ltau}.

\parheading{Ultraviolet obstructions.}
Square-integrability of the centered pair kernel is the finite-variance condition behind the pair-statistic and Hilbert--Schmidt estimates above.  For a Riesz singularity, it holds precisely below $2\s=\d$.  It is not, by itself, a substitute for the positive results' other analytic assumptions.  At and above this threshold, the lower obstructions use instead the explicit Boltzmann factor of an observed pair.  For $\s>\d/2$, repulsion depletes sufficiently close pairs, and a union over observed pairs, with their overlaps controlled, gives the mesoscopic obstruction.

At $\s=\d/2$, the logarithmic enhancement accumulates across transition annuli.  A marked-reservoir representation gives conditional Hellinger-affinity deficits near suitable isolated observed particles.  Their expected number is bounded below by a fixed multiple of the number of observed labels, uniformly over the stated marginal window.  The exact affinity recursion accumulates these deficits to prove the critical lower bound; no high-probability assertion about that count is required.  \Cref{sec:uv-sharpness-counterexamples} gives the argument under the respective hypotheses of \Cref{thm:intro-ultraviolet-sharpness}.  The conclusions are one-sided obstructions and do not follow from divergence of the partition function.

\begin{remark}[Relation with Gaussian fluctuations]\label{rem:intro-gaussian-fluctuations}
Under the hypotheses of the companion's static linear-statistics theorem~\cite[Proposition~2.11]{DGR26}, the limiting covariance for bounded empirical tests is
\begin{align}\label{eq:intro-gaussian-fluctuation-covariance}
 \Sigma_\be(\phi,\psi)
 =\left\langle\bar\phi,(\mathrm{Id}+\be T_\mu)^{-1}\bar\psi\right\rangle_{L^2(\mu)},
 \qquad \bar\phi=\phi-\int\phi\,\md\mu.
\end{align}
This is the same screened resolvent as in the marginal correction.  Covariance convergence alone does not give a strong pair estimate or identify the singleton correction; neither is deduced from that fluctuation theorem in the present proof.
\end{remark}

\subsection{Relation to previous work}
\label{subsec:related-work}

Classical chaos results are primarily qualitative.  Large-deviation methods identify the mean-field equilibrium and yield convergence of the empirical measure, hence convergence of every fixed marginal.  Ben Arous and Zeitouni obtained an increasing-propagation-of-chaos result under uniqueness and nondegeneracy of the mean-field minimizer~\cite{BenArousZeitouni1999}.  At the level relevant here, such global arguments typically give an \(O(1)\) full relative entropy and therefore only the local estimate \eqref{eq:intro-subadditive-local-rate}.

\parheading{Local divergence estimates.}
For local divergence orders, Lacker's work~\cite{Lacker2022} is the principal direct predecessor.  Its central theorem introduces a nonasymptotic hierarchy relating the \(k\)- and \((k+1)\)-particle relative Fisher informations.  Under additional transport-information, Poincar\'e, or convexity hypotheses, this yields corresponding relative-entropy and squared-Wasserstein bounds.  To the best of our knowledge, this was the first general nonasymptotic theorem to prove an upper bound at the quadratic local relative-entropy scale for a broad class of continuous mean-field Gibbs measures.  The abstract assumptions are force-level: a square-integrability bound for the centered interaction force and a transport-information inequality with constant \(\gamma<\be^{-2}\).  Bounded- or Lipschitz-gradient conditions furnish sufficient regimes.  A Gaussian model attains the quadratic order in \(k/N\), showing that the exponent in the upper bound cannot in general be improved.  This is a model-based optimality statement, rather than a matching lower-bound theorem throughout a class of interactions.  Singular interactions are explicitly left outside the physical scope of the result.

Ren and Wang~\cite{RenWang2024} replace the hierarchy of marginal Fisher informations by a hierarchy of conditional entropies.  Comparing neighboring conditional laws produces the second difference of the marginal entropy sequence and introduces a nonlinear Fisher-information dissipation tied to the mean-field free-energy landscape.  This yields nonperturbative entropy upper bounds of order \(k^2/N^2\) under suitable nonlinear and conditional transport inequalities.  In their flat semiconvex setting, the interaction is decomposed as \(W=W_+-W_-\), with both \(W_+\) and \(W_-\) of positive type---that is, conditionally positive definite on zero-mass signed measures---and either the relevant interaction forces are bounded or \(W_\pm\) are twice continuously differentiable with globally bounded mixed Hessians.  In their displacement-convex setting, \(W(x,y)=W'(x-y)\), where \(W'\) is even, convex, twice continuously differentiable, and has globally bounded Hessian.  These hypotheses exclude the diagonal logarithmic and Riesz singularities considered here.

Related recent work of Bernou and Duerinckx~\cite{BernouDuerinckx2025} obtains quantitative correlation estimates, Gibbs-relaxation bounds, and cross estimates for Brownian particle systems with bounded interaction forces, initially in a weak-interaction regime and with partial extensions beyond it.  Duerinckx and Jabin~\cite{DuerinckxJabin2026} prove high-temperature partition-function estimates in a homogeneous setting and use the resulting weighted-\(L^p\) control of the Gibbs equilibrium as an input to the derivation of the linearized Vlasov dynamics near equilibrium.  These works provide one-sided quantitative control of Gibbs equilibria and their correlations, but do not give matching lower bounds for local marginal divergences.

\parheading{Marginal and correlation approximations.}
The terminology and systematic dynamical study of the size of chaos go back in particular to Paul, Pulvirenti, and Simonella~\cite{PaulPulvirentiSimonella2019}, who made the pair-counting mechanism explicit and proved an \(O(k^2/N)\) marginal estimate in their correlation-error framework.  Duerinckx subsequently established the Bogolyubov scaling \(G_N^{m+1}=O(N^{-m})\) for connected \((m+1)\)-particle correlations, in suitable negative Sobolev norms, for smooth classical mean-field dynamics by Glauber calculus~\cite{DuerinckxGlauber2021}.  Lacker developed the corresponding local entropy hierarchy for mean-field diffusions~\cite{Lacker2023}.  Subsequent sharp dynamical results include the uniform-in-time estimates of Lacker and Le Flem~\cite{LackerLeFlem2023}, the higher-order \(L^2\) expansion of Hess-Childs and Rowan~\cite{HessChildsRowan2025}, the coupled entropy--Fisher-information hierarchy of Grass, Guillin, and Poquet~\cite{GrassGuillinPoquet2025}, and the extension by Arnese and Lacker to coefficients depending nonlinearly on the measure argument, with applications to mean-field Langevin dynamics, control, and games~\cite{ArneseLacker2026}.  For singular diffusions, Jabin and Wang obtained quantitative propagation of chaos in global relative entropy for \(W^{-1,\infty}\) interaction forces~\cite{JabinWang2018}; see also Guillin, Le Bris, and Monmarch\'{e}~\cite{GuillinLeBrisMonmarche2025} for subsequent uniform-in-time improvements.  Building on this line, Songbo Wang proved sharp local entropy and weighted-\(L^2\) estimates in the same rough-force class~\cite{Wang2026SharpLocal}.

Several complementary approaches control singular dynamical marginals or correlations in norms other than statistical divergences.  Bresch, Jabin, and Soler~\cite{BreschJabinSoler2025} developed uniform-in-\(N\) weighted-\(L^p\) estimates for the hierarchy of marginals in Vlasov--Fokker--Planck dynamics and also recorded a simpler first-order counterpart, for which no velocity weight is required.  Bresch, Duerinckx, and Jabin~\cite{BreschDuerinckxJabin2024} introduced a dual-correlation method for first- and second-order systems with square-integrable interaction forces, allowing possibly vanishing temperature.  For Brownian systems with square-integrable singular interactions, Duerinckx and Jabin~\cite{DuerinckxJabin2025Correlations} obtained systematic correlation estimates and derived the Bogolyubov correction and a central limit theorem.  B\'{e}jar-L\'{o}pez, Blaustein, Jabin, and Soler~\cite{BejarLopezBlausteinJabinSoler2026} proved uniform-in-time and uniform-in-\(N\) \(L^2\) bounds for marginals in singular Liouville dynamics, including configurations far from chaos.  These results are complementary to, but distinct from, local estimates in statistical divergences.

Nonexchangeable systems form a separate direction.  Jabin, Poyato, and Soler~\cite{JabinPoyatoSoler2025} established mean-field limits for nonexchangeable systems on sparse, possibly directed graphs through an extended-graphon framework.  Lacker, Yeung, and Zhou~\cite{LackerYeungZhou2024} proved sharp relative-entropy estimates for the marginals associated with arbitrary particle subsets in heterogeneous diffusions, extending Lacker's local entropy hierarchy by means of first-passage-percolation estimates.  These works address different questions: the former develops a broad nonexchangeable mean-field framework, whereas the latter gives sharp quantitative control of subset marginals.

  A stationary result of Zhang gives sharp one-particle R\'enyi-divergence upper bounds under strong isoperimetry and weak interaction assumptions~\cite{Zhang2026Renyi}.  The results considered here instead concern singular equilibrium marginals of every size, together with their quantitative leading profiles and ultraviolet obstructions.

\parheading{Partition functions and the present argument.}
The partition-function input has a more specific recent lineage.  Delgadino and Gvalani~\cite[Theorem~2.5]{DelgadinoGvalani2025Log} proved an all-temperature, \(N\)-uniform bound for the repulsive logarithmic gas under bounded-background assumptions.  The companion paper~\cite{DGR26} extends this result to logarithmic and Riesz interactions throughout the Hilbert--Schmidt range and additionally identifies the Carleman--Fredholm determinant normalization and the static fluctuation law.

Concurrent work of Wang and Zhao~\cite[Assumption~A and Theorems~1.1 and~1.3]{WangZhao2026Partition} proves all-temperature uniform partition-function estimates for logarithmic, Riesz, and Bessel--Riesz kernels throughout \(0\le\s<\d/2\), together with the second-Wiener-chaos and Carleman--Fredholm determinant limit.  Their logarithmic hypothesis includes a stretched-exponential spatial moment, whereas the companion uniform-bound theorem used here assumes finite absolute logarithmic energy.  Their proof uses a connected-graph expansion and heat-kernel truncation, while \cite{DGR26} uses positive-definite truncations and a low/high-frequency decomposition.  The two near-simultaneous works overlap at the level of the companion partition-function branch; no priority claim between the two approaches is intended.

Our equilibrium argument works at the level of the interaction rather than its derivatives.  It combines the explicit Gibbs density with square-integrability of the Hoeffding-centered interaction, cavity estimates, and partition-function control of the screened response.  This permits the logarithmic and Riesz singularities in the Hilbert--Schmidt regime.  Beyond the upper bounds, the abstract order theorem gives matching local-divergence lower bounds for every nontrivial interaction satisfying its hypotheses.  The abstract quantitative theorem identifies the first-order marginal profile and its growing-marginal divergence constants under row, source, and amplified centered-pair exponential bounds, without a sign condition on the centered operator.  The logarithmic and Riesz models satisfy these assumptions at every fixed inverse temperature under the stated background hypotheses.
  The ultraviolet theorem supplies lower obstructions to extending the classical scales beyond this regularity range.

\subsection{Organization of the paper}
\label{subsec:organization}

\Cref{sec:analytic-hypotheses-constants} states the abstract quantitative theorem, followed by the separate order and subquadratic results.  \Cref{sec:gibbs-identities-upper-bounds} collects the exact Gibbs identities, density and screened one-point estimates, and orthogonal marginal projection.  \Cref{sec:second-hoeffding-component,sec:singleton-quantitative-limits} prove the strong pair and singleton approximations and the statistical-divergence conclusions.  \Cref{sec:verification-log-riesz} verifies the respective logarithmic and Riesz hypotheses.  \Cref{sec:subquadratic-marginals,sec:chi2-tv-lower-bound-proof} give the subquadratic argument and the matching order lower bounds.  \Cref{sec:uv-sharpness-counterexamples} proves the ultraviolet obstructions.  \Cref{app:weighted-log-cavity,app:dependence-constants,app:probability-estimates} contain the weighted cavity proof, the constant ledger, and the shared probability estimates.

\subsection{Statement on AI use}

This working manuscript was generated with the assistance of generative AI tools, including OpenAI's ChatGPT and Codex, for literature identification, the development and testing of mathematical arguments, exposition, and consistency checks. The authors have not yet reviewed and independently verified all AI-generated content. They are in the process of formalizing the results in Lean and revising the natural-language presentation. The manuscript is therefore being shared as a draft, and its statements, proofs, references, and exposition remain subject to correction. Any completed formal verification applies only to the statements and versions explicitly identified in the accompanying verification records.

% STAGE2 MOVE: A3-abstract-criterion-opening
\section{Abstract results}
\label{sec:analytic-hypotheses-constants}

We first state the abstract theorem giving the quantitative marginal and divergence estimates of \Cref{subsec:exact-fixed-marginal-profile}.  The order and subquadratic results that follow have separate moment and cavity hypotheses.  The comparison after the first theorem explains this distinction.  All model-specific verification is deferred to \Cref{sec:verification-log-riesz}.

\subsection{Quantitative marginal approximation}
\label{subsec:abstract-quantitative-marginal}

The assumptions in this subsection are independent of those in
\Cref{subsec:quantitative-hypotheses}.  Let $\mu$ be a Borel probability measure
on $\R^\d$, let $W\in L^1(\mu^{\otimes2})$ be real and symmetric, and use
its centering and one-body function from \eqref{eq:centered-W-explicit} and
\eqref{eq:exact-psi-def}.  Fix $\be>0$, $p>2$, and
\begin{equation}\label{eq:signed-vartheta}
 0<\vartheta<\min\{1,p-2\}.
\end{equation}
Assume that, for some finite nonnegative $C_-,C_+$,
\begin{align}
 &\tl W\geq-C_-,\qquad
 M_p\coloneqq\operatorname*{ess\,sup}_{x}
          \norm{\tl W(x,\cdot)}_{L^p(\mu)}<\infty,
 \label{eq:signed-kernel-assumptions}\\
 &\psi_\mu\in L^6(\mu),\qquad \psi_\mu\leq C_+.
 \label{eq:signed-source-assumptions}
\end{align}
All essential bounds are relative to the appropriate product of $\mu$.
For a canonical kernel $c$ and independent coordinates $x_1,\ldots,x_N$
with law $\mu$, write
\begin{equation}\label{eq:signed-ordered-statistic}
 \mathsf W_N(c)\coloneqq\frac1N\sum_{i\ne j}c(x_i,x_j).
\end{equation}
Set
\begin{equation}\label{eq:signed-temperature-amplification}
 r_*\coloneqq\max\left\{12,\frac{p}{p-2-\vartheta},
                         \frac{2(2+\vartheta)}{\vartheta}\right\},
 \qquad \Gamma\coloneqq\be r_*,
\end{equation}
and assume the centered-pair exponential bound
\begin{equation}\label{eq:signed-centered-pair-moment}
 \sup_{N\geq2}\E_{\mu^{\otimes N}}
       \left[e^{-\Gamma\mathsf W_N(\tl W)/2}\right]<\infty.
\end{equation}
This bound does not include the one-body term in the modulated Gibbs density.
Its implication for the density moments is recorded in
\Cref{subsec:signed-first-cavity-estimate}.  No sign condition is imposed on
$T_\mu$, and no nontriviality assumption is imposed on $\tl W$.

\begin{thm}[Abstract quantitative marginal approximation]
\label{thm:abstract-quantitative-marginal}
Assume \eqref{eq:signed-vartheta}--\eqref{eq:signed-source-assumptions} and
\eqref{eq:signed-centered-pair-moment}, with $r_*,\Gamma$ as in
\eqref{eq:signed-temperature-amplification}.  Fix $0<\alpha\leq1$ with
$2+\alpha\leq p$, and set
\begin{equation}\label{eq:signed-profile-rate}
 \mathfrak r_N\coloneqq N^{-\alpha/8}+N^{-\vartheta/2}.
\end{equation}
The profiles in \eqref{eq:weak-second-resolvent-def}--\eqref{eq:exact-fixed-profile-def}
are well defined.  With this rate and these profiles, the estimates
\eqref{eq:quantitative-natural-profile},
\eqref{eq:quantitative-uniform-chi}--\eqref{eq:quantitative-uniform-TV}, and
\eqref{eq:quantitative-growing-chi}--\eqref{eq:quantitative-growing-TV}
hold for every $N\geq2$ and $1\leq k\leq N$.
Their constants depend only on the fixed exponents and hypothesis data,
not on $N$ or $k$.  In particular, the normalized density error vanishes
uniformly on every sublinear marginal range, and all three normalized
divergences converge to the displayed constants whenever $k_N\to\infty$
and $k_N/N\to0$.  At $\be=0$, the conclusions hold trivially whenever the
modulated definitions are well defined.
\end{thm}

Both \Cref{thm:chi2-chaos-two-sided,thm:abstract-quantitative-marginal} allow
positive and negative spectrum of the centered interaction operator.
Their hypotheses serve different purposes.  The order theorem permits the
weighted logarithmic lower-tail condition and uses modulated partition
bounds at the exponents in \eqref{eq:pf-input}.  The quantitative theorem
instead assumes uniform row integrability, a uniformly lower-bounded
centered kernel, and the stronger centered-pair exponential bound
\eqref{eq:signed-centered-pair-moment}; its source may be unbounded below
but must belong to $L^6(\mu)$.  These assumptions give quantitative Laplace
estimates for the replica and response families, not merely an inverse
bound for the screened operator.  Neither theorem inherits the hypotheses
of the other.  The proof of \Cref{thm:abstract-quantitative-marginal} is
completed in \Cref{sec:singleton-quantitative-limits}, after the strong pair
estimate in \Cref{sec:second-hoeffding-component}.

\begin{remark}[Spectral stability and exponential control]
\label{rem:signed-spectral-not-sufficient}
The amplified exponential bound cannot be replaced by spectral stability.
Let $\epsilon=e^{-4}$ and let a centered function $\xi$ equal $1$ on a set
of $\mu$-measure $\epsilon$, and $-\epsilon/(1-\epsilon)$ otherwise.  For
$W=-\xi\otimes\xi$, one has $\tl W=W$, $\psi_\mu=0$, and the only nonzero
eigenvalue is $\lambda=-\epsilon/(1-\epsilon)$.  At $\be=10$,
$1+2\be\lambda>0$, but the event on which every coordinate lies in the
first set gives
\begin{equation}\label{eq:signed-spectral-counterexample}
 \E_{\mu^{\otimes N}}\left[e^{-\be\mathsf W_N(\tl W)/2}\right]
 \geq\epsilon^N e^{5(N-1)}=e^{N-5}.
\end{equation}
Conversely, the assumptions admit genuinely mixed-sign interactions.
For independent symmetric signs $\sigma,\tau$ under $\mu$, take
$W=a\sigma\otimes\sigma-b\tau\otimes\tau$, with $a,b>0$.  Its nonzero
eigenvalues are $a,-b$, and $\psi_\mu=0$.  Gaussian integration and
$\cosh u\leq e^{u^2/2}$ give, whenever $\Gamma b<1$,
\begin{equation}\label{eq:signed-mixed-example}
 \E_{\mu^{\otimes N}}\left[e^{-\Gamma\mathsf W_N(W)/2}\right]
 \leq e^{\Gamma(a-b)/2}\E_Z
       \left[\cosh\left(\sqrt{\Gamma b/N}\,Z\right)^N\right]
 \leq\frac{e^{\Gamma(a-b)/2}}{\sqrt{1-\Gamma b}},
\end{equation}
where $Z$ is a standard Gaussian.  The other assumptions are immediate,
and no restriction on $a$ is needed.
\end{remark}

% STAGE2 MOVE: sec:verification-log-riesz
\subsection{Hypotheses for the order bounds}
\label{subsec:quantitative-hypotheses}

The structural upper-bound assumption is square-integrability of the centered interaction,
\begin{align}\label{eq:structural-L2-intro}
 \tl W\in L^2(\mu^{\otimes2}).
\end{align}
No sign condition is imposed on the associated Hilbert--Schmidt operator.  The two-sided theorem below separately assumes that the centered interaction is nontrivial, because its lower bounds are generated by a nonzero pair fluctuation.

For \(\gamma\ge0\), define the normalized partition-function constant by
\begin{align}\label{eq:normalized-pf-constant}
 \mathfrak P_\gamma(\mu)
 \coloneqq e^{-\gamma I_\mu}\sup_{M\ge1}\K_{M,\gamma}(\mu).
\end{align}
Jensen's inequality gives \(\K_{M,\gamma}(\mu)\ge e^{\gamma I_\mu}\), and therefore \(1\le\mathfrak P_\gamma(\mu)<\infty\) whenever the corresponding partition-function bound holds.

\begin{lemma}[Interpolation of normalized partition functions]\label{lem:normalized-partition-interpolation}
Let \(\be\ge0\) and \(1\le s\le t<\infty\).  Then, for every \(N\ge1\),
\begin{align}\label{eq:partition-interpolation-finite-N}
 \K_{N,s\be}(\mu)
 \le \K_{N,t\be}(\mu)^{s/t}.
\end{align}
Consequently,
\begin{align}\label{eq:normalized-partition-interpolation}
 \mathfrak P_{s\be}(\mu)
 \le \mathfrak P_{t\be}(\mu)^{s/t},
 \qquad
 \mathfrak P_{s\be}(\mu)^{1/s}
 \le \mathfrak P_{t\be}(\mu)^{1/t}.
\end{align}
\end{lemma}

\begin{proof}
Apply H\"older's inequality to \(Y_N=e^{-\be N\Fr_N}\):
\begin{align}\label{eq:partition-interpolation-holder}
 \E_{\mu^{\otimes N}}\left[Y_N^s\right]
 \le \E_{\mu^{\otimes N}}\left[Y_N^t\right]^{s/t}.
\end{align}
This is \eqref{eq:partition-interpolation-finite-N}.  Multiplication by \(e^{-s\be I_\mu}\), followed by the supremum over \(N\), gives \eqref{eq:normalized-partition-interpolation}.
\end{proof}

Fix
\begin{align}\label{eq:p-choice-quantitative}
 p>2,
 \qquad
 r\coloneqq\frac{p}{p-2},
\end{align}
and assume
\begin{align}\label{eq:mu-controls}
 &\norm{h_\mu-2I_\mu}_{L^2(\mu)}
 +\norm{h_\mu-2I_\mu}_{L^p(\mu)}
 +\norm{\tl W}_{L^2(\mu^{\otimes2})}
 +\norm{\tl W}_{L^p(\mu^{\otimes2})}<\infty.
\end{align}
For the lower bounds only, we additionally impose the nontriviality condition
\begin{align}\label{eq:interaction-nontriviality}
 \norm{\tl W}_{L^2(\mu^{\otimes2})}>0.
\end{align}  For later use, we write
\begin{align}\label{eq:mixed-section-norm}
 \norm{\tl W}_{L_x^\infty L_y^p(\mu^{\otimes2})}
 \coloneqq
 \sup_{x\in\R^\d}\norm{\tl W(x,\cdot)}_{L^p(\mu)}.
\end{align}
The joint condition \(\tl W\in L^p(\mu^{\otimes2})\) and the mixed-norm condition \(\tl W\in L_x^\infty L_y^p(\mu^{\otimes2})\) below serve different purposes as assumptions: the former controls product-space moments, whereas the latter gives a fixed-\(x\) moment bound for the section \(y\mapsto\tl W(x,y)\).

For \(N\ge2\), write \(\xi_1,\ldots,\xi_N\) for the coordinate variables on \((\R^\d)^N\), set \(\hat\Xi_1=(\xi_2,\ldots,\xi_N)\), and define the centered one-particle cavity field
\begin{align}\label{eq:cavity-B-def}
 B_N(x;\hat\Xi_1)
 \coloneqq \frac{\be}{N}\pa{h_\mu(x)-2I_\mu}
 -\frac{\be}{N}\sum_{j=2}^N \tl W(x,\xi_j).
\end{align}
We impose one of the following two cavity conditions.

\parheading{Bounded lower-tail condition.}
Assume
\begin{align}\label{eq:cavity-tail-bounded}
 \norm{h_\mu-2I_\mu}_{L^\infty}
 +\norm{\tl W}_{L_x^\infty L_y^p(\mu^{\otimes2})}<\infty,
 \qquad
 \operatorname*{ess\,inf}_{(x,y)\in(\R^\d)^2}\tl W(x,y)>-\infty.
\end{align}
This condition gives a deterministic upper bound on \(B_N\), uniformly in \(N\) and in the configuration of the remaining \(N-1\) particles.

\parheading{Logarithmic lower-tail condition.}
Let \(L(x)\coloneqq\log(e+|x|)\).  Assume that there are finite constants \(a_0,a_1,a_h\) such that
\begin{align}\label{eq:cavity-tail-log}
 (\tl W)_-(x,y)&\le a_0+a_1\pa{L(x)+L(y)},
 &
 (h_\mu-2I_\mu)_+(x)&\le a_h\pa{1+L(x)}
\end{align}
for \(\mu^{\otimes2}\)-a.e. \((x,y)\) and \(\mu\)-a.e. \(x\), respectively.  Define
\begin{align}
 \mathcal L_N(\hat\Xi_1)&\coloneqq \frac1N\sum_{j=2}^N L(\xi_j),\label{eq:log-spatial-average}\\
 \mathsf c_{\log,\be,\mu}&\coloneqq \be(a_0+a_1+a_h),\label{eq:log-cavity-weight-constant}\\
 \Phi_N(x;\hat\Xi_1)&\coloneqq
 \mathsf c_{\log,\be,\mu}\pa{1+L(x)+\mathcal L_N(\hat\Xi_1)}.\label{eq:log-cavity-weight}
\end{align}
Then the positive part of the cavity field satisfies
\begin{align}\label{eq:B-log-upper}
 B_N(x;\hat\Xi_1)_+\le \Phi_N(x;\hat\Xi_1).
\end{align}
Under this condition, assume the finiteness of the following weighted cavity constants:
\begin{subequations}\label{eq:log-weighted-cavity-constants}
\begin{align}
 \mathsf E_{\log,\be,\mu}
 &\coloneqq
 \frac12\sup_{N\ge2}N
 \norm{\E_{\Q_{N,\be}(\mu)}\left[
 e^{\Phi_N(\cdot;\hat\Xi_1)}
 B_N(\cdot;\hat\Xi_1)^2\right]}_{L^2(\mu)}<\infty,\label{eq:log-weighted-cavity-E}\\
 \mathsf D_{\log,\be,\mu}
 &\coloneqq
 \frac12\sup_{N\ge2}N
 \E_{\Q_{N,\be}(\mu)}\left[
 \int e^{\Phi_N(x;\hat\Xi_1)}
 B_N(x;\hat\Xi_1)^2\,\md\mu(x)\right]<\infty,\label{eq:log-weighted-cavity-D}\\
 \mathsf V_{\log,\be,\mu}
 &\coloneqq
 \sup_{N\ge2}N
 \E_{\Q_{N,\be}(\mu)}\left[
 \left(\int e^{\Phi_N(x;\hat\Xi_1)}|B_N(x;\hat\Xi_1)|\,\md\mu(x)\right)^2\right]<\infty,\label{eq:log-weighted-cavity-V}\\
 \mathsf B_{\log,\be,\mu}
 &\coloneqq
 \sup_{N\ge2}N
 \E_{\Q_{N,\be}(\mu)}\left[
 \int B_N(x;\hat\Xi_1)^2\,\md\mu(x)\right]<\infty.\label{eq:log-weighted-cavity-B}
\end{align}
\end{subequations}
The bound \eqref{eq:B-log-upper} introduces polynomial weights into the nonlinear cavity estimate.  \Cref{lem:weighted-logarithmic-cavity} gives an explicit sufficient condition for the four estimates above.

\subsection{Abstract two-sided chaos theorem}
\label{subsec:abstract-two-sided-theorem}

With the structural, moment, cavity, and partition-function assumptions now fixed, we can state the abstract result.  The interaction remains arbitrary in this theorem; the logarithmic and Riesz kernels enter only through the verification in \Cref{sec:verification-log-riesz}.

\begin{thm}[Abstract two-sided chaos bounds]\label{thm:chi2-chaos-two-sided}
Fix \(\be\ge0\).  Assume \eqref{eq:structural-L2-intro}, \eqref{eq:mu-controls}, and \eqref{eq:interaction-nontriviality} for some \(p>2\).  Assume either the bounded lower-tail cavity condition \eqref{eq:cavity-tail-bounded}, or the logarithmic lower-tail condition \eqref{eq:cavity-tail-log} together with the weighted estimates \eqref{eq:log-weighted-cavity-constants}; in the latter case the weighted estimates may instead be verified through \Cref{lem:weighted-logarithmic-cavity}.  With \(r=p/(p-2)\), assume
\begin{align}\label{eq:pf-input}
 \sup_{N\ge1}\K_{N,r\be}(\mu)<\infty,
 \qquad
 \sup_{N\ge1}\K_{N,2\be}(\mu)<\infty.
\end{align}
Then there are finite constants \(\mathsf C_{\TV,\be,\mu}^{\mathrm{up}}\), \(\mathsf C_{\Hr,\be,\mu}^{\mathrm{up}}\), and \(\mathsf C_{\chi,\be,\mu}^{\mathrm{up}}\) such that, for every \(1\le k\le N\),
\begin{align}\label{eq:unified-upper-final}
 \TV\pa{\Q_{N,\be}^{(k)}(\mu),\mu^{\otimes k}}
 &\le \mathsf C_{\TV,\be,\mu}^{\mathrm{up}}\frac{k}{N},\\
 \Hr\pa{\Q_{N,\be}^{(k)}(\mu)\mid\mu^{\otimes k}}
 &\le \mathsf C_{\Hr,\be,\mu}^{\mathrm{up}}\frac{k^2}{N^2},\\
 \Chi\pa{\Q_{N,\be}^{(k)}(\mu)\mid\mu^{\otimes k}}
 &\le \mathsf C_{\chi,\be,\mu}^{\mathrm{up}}\frac{k^2}{N^2}.
\end{align}
If \(\be>0\), then there are \(\mathsf N_{\be,\mu}<\infty\) and positive constants \(\mathsf c_{\TV,\be,\mu}^{\mathrm{low}}\), \(\mathsf c_{\Hr,\be,\mu}^{\mathrm{low}}\), and \(\mathsf c_{\chi,\be,\mu}^{\mathrm{low}}\) such that, for every \(N\ge\mathsf N_{\be,\mu}\) and \(2\le k\le N\),
\begin{align}\label{eq:unified-lower-final}
 \TV\pa{\Q_{N,\be}^{(k)}(\mu),\mu^{\otimes k}}
 &\ge \mathsf c_{\TV,\be,\mu}^{\mathrm{low}}\frac{k}{N},\\
 \Hr\pa{\Q_{N,\be}^{(k)}(\mu)\mid\mu^{\otimes k}}
 &\ge \mathsf c_{\Hr,\be,\mu}^{\mathrm{low}}\frac{k^2}{N^2},\\
 \Chi\pa{\Q_{N,\be}^{(k)}(\mu)\mid\mu^{\otimes k}}
 &\ge \mathsf c_{\chi,\be,\mu}^{\mathrm{low}}\frac{k^2}{N^2}.
 \label{eq:unified-lower-chi}
\end{align}
\end{thm}

% STAGE2 MOVE: A1-general-divergence-formalism
These three quantities belong to the Csisz\'ar--Morimoto class of \(f\)-divergences~\cite{AliSilvey1966,Csiszar1967}.  For a convex function \(f:[0,\infty)\to(-\infty,\infty]\) with \(f(1)=0\), define
\begin{align}\label{eq:f-divergence-definition}
 \mathscr D_f(\nu\mid\pi)
 \coloneqq
 \int f\!\left(\frac{\md\nu}{\md\pi}\right)\,\md\pi,
 \qquad \nu\ll\pi.
\end{align}
Adding a multiple of \(t-1\) to \(f\) does not change \(\mathscr D_f\), so one may impose \(f'(1)=0\) when this derivative exists.  The three quantities above are generated by
\begin{align}\label{eq:principal-f-divergence-generators}
 f_{\mathrm{TV}}(t)=\frac12|t-1|,
 \qquad
 f_{\mathrm{KL}}(t)=t\log t-t+1,
 \qquad
 f_{\chi^2}(t)=(t-1)^2.
\end{align}
Near the product law,
\begin{align}
 f_{\mathrm{TV}}(1+u)&=\frac12|u|,
 \label{eq:f-TV-local-order}\\
 f_{\mathrm{KL}}(1+u)&=\frac12u^2+O(|u|^3),
 \qquad
 f_{\chi^2}(1+u)=u^2.
 \label{eq:f-quadratic-local-order}
\end{align}
Thus, total variation is first order in the density perturbation, whereas relative entropy and chi-square divergence are quadratic; this accounts for the two scales in \eqref{eq:intro-informal-scales}.

% STAGE2 MOVE: A5-divergence-corollary-and-subquadratic-theorem
\begin{cor}[Consequences for \(f\)-divergences]\label{cor:f-divergence-consequences}
Assume the hypotheses of \Cref{thm:chi2-chaos-two-sided}.
\begin{enumerate}[(i)]
\item For \(1\le\tau\le2\), define the centered \(L^\tau\)-divergence by
\begin{align}\label{eq:centered-Ltau-divergence}
 \mathscr D_\tau(\nu\mid\pi)
 \coloneqq
 \int\left|\frac{\md\nu}{\md\pi}-1\right|^\tau\,\md\pi,
 \qquad \nu\ll\pi.
\end{align}
Its endpoints are
\begin{align}\label{eq:centered-Ltau-endpoints}
 \mathscr D_1(\nu\mid\pi)=2\TV(\nu,\pi),
 \qquad
 \mathscr D_2(\nu\mid\pi)=\Chi(\nu\mid\pi).
\end{align}
For every \(1\le k\le N\),
\begin{align}\label{eq:centered-Ltau-upper}
 \mathscr D_\tau\!\left(\Q_{N,\be}^{(k)}(\mu)\mid\mu^{\otimes k}\right)
 \le
 \left(\mathsf C_{\chi,\be,\mu}^{\mathrm{up}}\right)^{\tau/2}
 \left(\frac{k}{N}\right)^\tau.
\end{align}
If \(\be>0\), then, for every \(N\ge\mathsf N_{\be,\mu}\) and \(2\le k\le N\),
\begin{align}\label{eq:centered-Ltau-lower}
 \mathscr D_\tau\!\left(\Q_{N,\be}^{(k)}(\mu)\mid\mu^{\otimes k}\right)
 \ge
 \left(2\mathsf c_{\TV,\be,\mu}^{\mathrm{low}}\right)^\tau
 \left(\frac{k}{N}\right)^\tau.
\end{align}

\item Let \(f:[0,\infty)\to(-\infty,\infty]\) be convex with \(f(1)=0\), and let \(\mathscr D_f\) be as in \eqref{eq:f-divergence-definition}.  Suppose that, for some \(\tau\in[1,2]\) and constants \(a_f,b_f>0\), the global comparisons
\begin{align}\label{eq:global-f-divergence-comparisons}
 a_f\TV(\nu,\pi)^\tau
 \le \mathscr D_f(\nu\mid\pi)
 \le b_f\Chi(\nu\mid\pi)^{\tau/2}
\end{align}
hold for every pair of probability measures \(\nu\ll\pi\).  Then, for every \(1\le k\le N\),
\begin{align}\label{eq:comparable-f-divergence-upper}
 \mathscr D_f\!\left(\Q_{N,\be}^{(k)}(\mu)\mid\mu^{\otimes k}\right)
 \le
 b_f\left(\mathsf C_{\chi,\be,\mu}^{\mathrm{up}}\right)^{\tau/2}
 \left(\frac{k}{N}\right)^\tau.
\end{align}
If \(\be>0\), then, for every \(N\ge\mathsf N_{\be,\mu}\) and \(2\le k\le N\),
\begin{align}\label{eq:comparable-f-divergence-lower}
 \mathscr D_f\!\left(\Q_{N,\be}^{(k)}(\mu)\mid\mu^{\otimes k}\right)
 \ge
 a_f\left(\mathsf c_{\TV,\be,\mu}^{\mathrm{low}}\right)^\tau
 \left(\frac{k}{N}\right)^\tau.
\end{align}
\end{enumerate}
Total variation is recovered with \(\tau=1\); relative entropy is recovered with \(\tau=2\), using Pinsker's inequality and \(\Hr\le\Chi\); and chi-square divergence is recovered with \(\tau=2\).
\end{cor}

\begin{proof}
Set \(g=\md\nu/\md\pi-1\).  Since \(\pi\) is a probability measure, for \(1\le\tau\le2\),
\begin{align}\label{eq:L1-Ltau-L2-comparison}
 \norm{g}_{L^1(\pi)}
 \le \norm{g}_{L^\tau(\pi)}
 \le \norm{g}_{L^2(\pi)}.
\end{align}
Consequently,
\begin{align}\label{eq:centered-Ltau-endpoint-comparison}
 \left(2\TV(\nu,\pi)\right)^\tau
 \le \mathscr D_\tau(\nu\mid\pi)
 \le \Chi(\nu\mid\pi)^{\tau/2}.
\end{align}
The estimates in part~\textup{(i)} follow from the total-variation lower bound and the chi-square upper bound in \Cref{thm:chi2-chaos-two-sided}.  Part~\textup{(ii)} follows in the same way from \eqref{eq:global-f-divergence-comparisons}.  The final identifications use \eqref{eq:centered-Ltau-endpoints} and the standard inequalities \(\Hr(\nu\mid\pi)\ge2\TV(\nu,\pi)^2\) and \(\Hr(\nu\mid\pi)\le\Chi(\nu\mid\pi)\).
\end{proof}

\begin{thm}[Fixed marginals under a subquadratic partition-function moment]
\label{thm:fixed-marginal-Ltau}
Fix \(\be>0\) and \(1<\tau<2\), and set
\begin{align}\label{eq:fixed-marginal-conjugate-exponents}
 q\coloneqq\frac{\tau}{\tau-1},
 \qquad
 p\ge2q,
 \qquad
 r\coloneqq\frac{p}{p-2}\le\tau.
\end{align}
Assume \eqref{eq:structural-L2-intro} and the finiteness conditions in \eqref{eq:mu-controls} with this exponent \(p\).  Assume either the bounded lower-tail condition \eqref{eq:cavity-tail-bounded}, or the logarithmic lower-tail condition \eqref{eq:cavity-tail-log} together with the weighted estimates \eqref{eq:log-weighted-cavity-constants}; in the latter case the weighted estimates may instead be verified through \Cref{lem:weighted-logarithmic-cavity}.  Finally, assume
\begin{align}\label{eq:fixed-marginal-partition-input}
 \sup_{N\ge1}\K_{N,\tau\be}(\mu)<\infty.
\end{align}
No sign condition is imposed on \(T_\mu\) or on its quadratic form, and the nontriviality condition \eqref{eq:interaction-nontriviality} is not assumed.

Let
\begin{align}\label{eq:fixed-marginal-Gibbs-constants}
 \mathsf M_{\tau,\be,\mu}
 &\coloneqq \mathfrak P_{\tau\be}(\mu)^{1/\tau},
 &
 \mathsf U_{\tau,\be,\mu}
 &\coloneqq
 \frac{\tau}{\tau-1}
 \left(
 \be\norm{h_\mu-2I_\mu}_{L^2(\mu)}
 +\mathsf R_{\mathrm{cav},\be,\mu}
 \right),
\end{align}
where \(\mathsf R_{\mathrm{cav},\be,\mu}\) is the cavity remainder constant in \eqref{eq:cavity-error-bound}, with the needed \(L^r\) density bound supplied by \Cref{lem:normalized-partition-interpolation}.  Then, for every \(N\ge2\) and \(2\le k\le N\),
\begin{align}\label{eq:fixed-marginal-Gibbs-explicit}
 \norm{f_{N,\be}^{(k)}-1}_{L^\tau(\mu^{\otimes k})}
 &\le \frac{\mathsf U_{\tau,\be,\mu}\sqrt{k}}{N}
 +(\mathsf M_{\tau,\be,\mu}+1)\mathsf B_q
 \left[
 \frac{\sqrt{\binom{k}{2}}}{N-1}
 +\frac{4\binom{k}{2}N^{1/q-1/2}}{N-1}
 \right]\notag\\
 &\quad +(\mathsf M_{\tau,\be,\mu}+1)
 \sum_{j=3}^k\binom{k}{j}
 \frac{2^j\mathsf A_{j,q}}{\sqrt{(N)_j}}.
\end{align}
Here, the empty sum is zero, and the constants \(\mathsf A_{j,q}\) and \(\mathsf B_q\) are those in \Cref{lem:fixed-order-canonical-U-moment,lem:refined-canonical-pair-moment}.  For \(k=1\),
\begin{align}\label{eq:fixed-marginal-Gibbs-one-point}
 \norm{f_{N,\be}^{(1)}-1}_{L^\tau(\mu)}
 \le \frac{\mathsf U_{\tau,\be,\mu}}{N}.
\end{align}
Consequently, for every fixed \(k\ge2\),
\begin{align}\label{eq:fixed-marginal-Gibbs-limsup}
 \limsup_{N\to\infty}
 N\norm{f_{N,\be}^{(k)}-1}_{L^\tau(\mu^{\otimes k})}
 \le
 \mathsf U_{\tau,\be,\mu}\sqrt{k}
 +(\mathsf M_{\tau,\be,\mu}+1)\mathsf B_q
 \sqrt{\binom{k}{2}}.
\end{align}
In particular, for every fixed \(k\ge1\),
\begin{align}\label{eq:fixed-marginal-Gibbs-rates}
 \norm{f_{N,\be}^{(k)}-1}_{L^\tau(\mu^{\otimes k})}
 &=O_{k,\tau,\be,\mu}(N^{-1}),\\
 \TV\pa{\Q_{N,\be}^{(k)}(\mu),\mu^{\otimes k}}
 &=O_{k,\tau,\be,\mu}(N^{-1}).
\end{align}
The corresponding centered divergence and relative entropy satisfy
\begin{align}\label{eq:fixed-marginal-Gibbs-divergences}
 \mathscr D_\tau\pa{\Q_{N,\be}^{(k)}(\mu)\mid\mu^{\otimes k}}
 &=O_{k,\tau,\be,\mu}(N^{-\tau}),\\
 \Hr\pa{\Q_{N,\be}^{(k)}(\mu)\mid\mu^{\otimes k}}
 &=O_{k,\tau,\be,\mu}(N^{-\tau}).
 \label{eq:fixed-marginal-Gibbs-entropy}
\end{align}
The entropy estimate in \eqref{eq:fixed-marginal-Gibbs-entropy} is only a nonsharp consequence.  No quantitative growing-\(k\) range is asserted.
\end{thm}

For the logarithmic and Riesz models below, whenever the remaining moment and cavity hypotheses of \Cref{thm:fixed-marginal-Ltau} hold for a given \(\tau\), its only external partition-function input is supplied by the companion uniform-bound theorem~\cite[Theorem~2.2\textup{(i)}]{DGR26} applied at the single inverse temperature \(\tau\be\).  Equivalently, this follows from the compact-temperature statement in \cite[Corollary~5.7]{DGR26}; in the logarithmic case, finite absolute logarithmic energy is the relevant additional hypothesis.

\section{Gibbs identities and marginal upper bounds}
\label{sec:gibbs-identities-upper-bounds}\label{sec:one-point-upper-estimate}

The exact identities below separate the observed particles from the reservoir.  Orthogonal projection then reduces the marginal upper bound to the first Hoeffding component and the full density.  We obtain the required one-point estimate from the conditional exponential and spectral stability.  The order and quantitative assumptions are invoked separately when verifying the cavity remainders.  These common inputs are reused in the strong pair and singleton proofs.

\subsection{Exact marked identities}
\label{sec:exact-marked-identities}

This section isolates the exact finite-\(N\) identity that turns a
marginal estimate into a marked partition-function estimate.  The point is to
separate the \(k\) observed variables from the remaining \(N-k\) variables before
any approximation is made.

With \(h_\mu\), \(I_\mu\), and \(\tl W\) as in \eqref{eq:h-I-tildeg-statement}--\eqref{eq:tildeg-def-statement}, for
\(\Xk=(x_1,\ldots,x_k)\), set
\begin{align}
 \Phi_{\Xk}(y)&\coloneqq \sum_{a=1}^k\tl W(x_a,y),\label{eq:Phi-x-def}\\
 D_k(\Xk)&\coloneqq \frac12\sum_{1\le a\ne b\le k}\tl W(x_a,x_b)-\sum_{a=1}^k h_\mu(x_a)+kI_\mu.\label{eq:Dk-def}
\end{align}

\begin{lemma}[Exact marked splitting]\label{lem:exact-splitting}
Let \(1\le k<N\), \(n=N-k\), \(\Xk=(x_1,\ldots,x_k)\), and \(Y_n=(y_1,\ldots,y_n)\).  Then
\begin{align}\label{eq:exact-splitting}
N\Fr_N((\Xk,Y_n),\mu)
=\frac{n}{N}\,n\Fr_n(Y_n,\mu)
+\frac1N\sum_{j=1}^n\Phi_{\Xk}(y_j)
+\frac1N D_k(\Xk).
\end{align}
\end{lemma}

\begin{proof}
By definition,
\begin{align}
N\Fr_N(z_1,\ldots,z_N;\mu)
=\frac1{2N}\sum_{1\le i\ne j\le N}\tl W(z_i,z_j)-\frac1N\sum_{i=1}^N h_\mu(z_i)+I_\mu.
\end{align}
Taking \((z_1,\ldots,z_N)=(\Xk,Y_n)\), split the off-diagonal sum into bulk--bulk, mark--bulk, and mark--mark pieces.  This gives
\begin{align}
N\Fr_N((\Xk,Y_n),\mu)
&=\frac1{2N}\sum_{1\le i\ne j\le n}\tl W(y_i,y_j)-\frac1N\sum_{j=1}^n h_\mu(y_j)+\frac{n}{N}I_\mu \notag\\
&\quad +\frac1N\sum_{a=1}^k\sum_{j=1}^n\tl W(x_a,y_j)
+\frac1{2N}\sum_{1\le a\ne b\le k}\tl W(x_a,x_b) \notag\\
&\quad -\frac1N\sum_{a=1}^k h_\mu(x_a)+\frac{k}{N}I_\mu.
\end{align}
The first line is \((n/N)n\Fr_n(Y_n,\mu)\), and the remaining terms are the final two terms in \eqref{eq:exact-splitting}.
\end{proof}

We now package the three factors in the splitting: the bulk
partition function, the interaction of the marks with the bulk, and the purely
marked remainder.  This notation is designed so that the marginal density is
just a normalized version of one marked weight.

For \(n=N-k\), define
\begin{align}
 \gamma_{N,k}&\coloneqq \be\frac{n}{N},\label{eq:gamma-Nk-def}\\
 \mathsf M_{N,k}(\Xk)&\coloneqq
 \E_{Y_n\sim\Q_{n,\gamma_{N,k}}(\mu)}
 \left[\exp\left(-\frac{\be}{N}\sum_{j=1}^n\Phi_{\Xk}(y_j)\right)\right],\label{eq:M-Nk-def}\\
 \Psi_{N,k}(\Xk)&\coloneqq
 \exp\left(-\frac{\be}{N}D_k(\Xk)\right)\mathsf M_{N,k}(\Xk).
 \label{eq:Psi-Nk-def}
\end{align}
We use the marginal densities defined in \eqref{eq:gibbs-marginal-density-ratio}.  For $k<N$, the marginal density is a normalized marked partition function.

\begin{lemma}[Normalized conditional identity]\label{lem:normalized-conditional-identity}
For \(1\le k<N\),
\begin{align}\label{eq:f-Psi-normalized}
 f_{N,\be}^{(k)}(\Xk)=\frac{\Psi_{N,k}(\Xk)}{\int\Psi_{N,k}\,\md\mu^{\otimes k}}.
\end{align}
Equivalently,
\begin{align}\label{eq:Psi-mean-partition-ratio}
 \int\Psi_{N,k}\,\md\mu^{\otimes k}=\frac{\K_{N,\be}(\mu)}{\K_{n,\gamma_{N,k}}(\mu)}.
\end{align}
\end{lemma}

\begin{proof}
By \Cref{lem:exact-splitting},
\begin{align}
f_{N,\be}^{(k)}(\Xk)
&=\frac1{\K_{N,\be}(\mu)}
\int \exp\{-\be N\Fr_N((\Xk,Y_n),\mu)\}\,\md\mu^{\otimes n}(Y_n)\notag\\
&=\frac{\K_{n,\gamma_{N,k}}(\mu)}{\K_{N,\be}(\mu)}
\exp\left(-\frac{\be}{N}D_k(\Xk)\right)
\E_{\Q_{n,\gamma_{N,k}}(\mu)}
\left[\exp\left(-\frac{\be}{N}\sum_{j=1}^n\Phi_{\Xk}(y_j)\right)\right]\notag\\
&=\frac{\K_{n,\gamma_{N,k}}(\mu)}{\K_{N,\be}(\mu)}\Psi_{N,k}(\Xk).
\end{align}
Integrating in \(\Xk\) gives \eqref{eq:Psi-mean-partition-ratio}; substituting the resulting constant gives \eqref{eq:f-Psi-normalized}.
\end{proof}

In particular,
\begin{align}\label{eq:chi2-Psi}
 \Chi\pa{\Q_{N,\be}^{(k)}(\mu)\mid\mu^{\otimes k}}
 =\left\|\frac{\Psi_{N,k}}{\int\Psi_{N,k}\,\md\mu^{\otimes k}}-1\right\|_{L^2(\mu^{\otimes k})}^2.
\end{align}
For the upper bound, only the case \(k=1\) will be estimated directly.

% STAGE2 MOVE: A2-onepoint-opening
\subsection{Orthogonal marginal projection}
\label{subsec:Hilbertian-all-marginal-reduction}

The following exact decomposition separates the first Hoeffding level from all higher levels.  It explains which estimates the Gibbs argument must supply: a sharp one-point bound and a bounded full-law chi-square divergence.

\begin{lemma}[Hoeffding projection reduction]\label{lem:hoeffding-reduction}
Let \(g_N\) be an exchangeable probability density on \(E^N\) with respect to \(\mu^{\otimes N}\), and let \(g_N^{(k)}\) be its labeled \(k\)-point marginal density.  Then
\begin{align}\label{eq:hoeffding-chi2-bound}
 \Chi(g_N^{(k)}\mu^{\otimes k}\mid\mu^{\otimes k})
 \le k\,\Chi(g_N^{(1)}\mu\mid\mu)
 +\left(\frac{k}{N}\right)^2\Chi(g_N\mu^{\otimes N}\mid\mu^{\otimes N}).
\end{align}
\end{lemma}

\begin{proof}
Let
\begin{align}\label{eq:hoeffding-decomposition-general}
 g_N=1+\sum_{\emptyset\ne A\subset\{1,\ldots,N\}}h_A
\end{align}
be the orthogonal Hoeffding decomposition of \(g_N\) in \(L^2(\mu^{\otimes N})\).  Thus, each \(h_A\) depends only on \((x_i)_{i\in A}\) and is degenerate in each of its variables.  By exchangeability, if \([r]:=\{1,\ldots,r\}\), then every \(h_A\) with \(|A|=r\) has the same \(L^2\)-norm as \(h_{[r]}\).  Conditional expectation onto the first \(k\) variables gives
\begin{align}\label{eq:hoeffding-marginal-decomposition}
 g_N^{(k)}-1=\sum_{\emptyset\ne A\subset\{1,\ldots,k\}}h_A.
\end{align}
For \(1\le k\le N\), orthogonality gives
\begin{align}\label{eq:hoeffding-marginal-chi2}
 \Chi(g_N^{(k)}\mu^{\otimes k}\mid\mu^{\otimes k})
 =\sum_{r=1}^k \binom{k}{r}\norm{h_{[r]}}_{L^2(\mu^{\otimes r})}^2.
\end{align}
The \(r=1\) contribution to \eqref{eq:hoeffding-marginal-chi2} is \(k\norm{h_{[1]}}_{L^2(\mu)}^2=k\Chi(g_N^{(1)}\mu\mid\mu)\).  For \(r\ge2\),
\begin{align}\label{eq:hoeffding-combinatorial-ratio}
 \frac{\binom{k}{r}}{\binom{N}{r}}
 \le \left(\frac{k}{N}\right)^r
 \le \left(\frac{k}{N}\right)^2.
\end{align}
This proves \eqref{eq:hoeffding-chi2-bound}.
\end{proof}

\subsection{Density moments and spectral stability}
\label{subsec:density-moments-spectral-stability}

The Gibbs density supplies the full-law moment bound.  A partition-function moment also gives spectral stability of the screened operator.  These two statements do not require the quantitative marginal hypotheses or a sign condition on the interaction.

% STAGE2 MOVE: A5-onepoint-input-reference
We use the operator $T_\mu$ defined in \eqref{eq:T-def}.  The following two consequences of the partition-function bound apply without a sign assumption.

For later use, introduce the centered energy
\begin{align}\label{eq:spectral-centered-energy}
\mathcal X_N\coloneqq N\Fr_N(\XN,\mu)+I_\mu
=\frac1N\sum_{1\le i<j\le N}\tl W(x_i,x_j)
-\frac1N\sum_{i=1}^N\pa{h_\mu(x_i)-2I_\mu}.
\end{align}
Its expectation under $\mu^{\otimes N}$ is zero.

% STAGE2 MOVE: A5-single-density-moment-lemma
\begin{lemma}[Uniform \(L^r\) bounds for the full density]\label{lem:full-density-Lr}
Recall from \eqref{eq:gibbs-marginal-density-ratio} that \(f_{N,\be}=f_{N,\be}^{(N)}\), and recall the normalized partition-function constant \(\mathfrak P_\gamma(\mu)\) from \eqref{eq:normalized-pf-constant}.  For every \(r\in[1,\infty)\), the following bound holds for the full density and, with the same constant, for every marginal density $f_{N,\be}^{(j)}$, $1\leq j\leq N$:
\begin{align}\label{eq:f-Lr-bound}
 \sup_{N\ge1}\norm{f_{N,\be}}_{L^r(\mu^{\otimes N})}\le \mathfrak P_{r\be}(\mu)^{1/r}.
\end{align}
\end{lemma}

\begin{proof}
For \(r=1\) there is nothing to prove.  For \(r>1\),
\begin{align}
 \norm{f_{N,\be}}_{L^r(\mu^{\otimes N})}^r
 =\frac{\K_{N,r\be}(\mu)}{\K_{N,\be}(\mu)^r}.
\end{align}
The numerator is bounded by \(e^{r\be I_\mu}\mathfrak P_{r\be}(\mu)\), while Jensen's inequality gives
\begin{align}\label{eq:jensen-lower-again}
\K_{N,\be}(\mu)\ge \exp\left(\be I_\mu\right).
\end{align}
Thus, \(\norm{f_{N,\be}}_{L^r}^r\le \mathfrak P_{r\be}(\mu)\), which proves \eqref{eq:f-Lr-bound}.  Marginalization is conditional expectation under the product law and is contractive in $L^r$, giving the marginal assertion.
\end{proof}

% STAGE2 MOVE: A5-classical-quadratic-limit-attribution
The quadratic-statistic limit used in the next lemma is classical; see Hoeffding~\cite{Hoeffding1948} and Dynkin and Mandelbaum~\cite{DynkinMandelbaum1983}.  We retain its spectral truncation proof to track the partition-function obstruction.

% STAGE2 MOVE: A2-onepoint-spectral-and-cavity-proof
\begin{lemma}[Spectral stability from a partition-function moment]
\label{lem:partition-spectral-stability}
Fix \(\be>0\) and \(\theta>1\).  Assume
\begin{align}\label{eq:partition-spectral-stability-assumptions}
 \tl W\in L^2(\mu^{\otimes2}),
 \qquad
 h_\mu-2I_\mu\in L^2(\mu),
 \qquad
 \sup_{N\ge1}\K_{N,\theta\be}(\mu)<\infty.
\end{align}
Then every spectral value \(\lambda\) of \(\Tc\) satisfies
\begin{align}\label{eq:partition-spectral-strict-lower}
 1+\theta\be\lambda>0.
\end{align}
Consequently, for every \(N\ge2\),
\begin{align}\label{eq:partition-spectral-resolvent-bound}
 \norm*{\pa*{I+\be\frac{N-1}{N}\Tc}^{-1}}_{L^2(\mu)\to L^2(\mu)}
 \le \frac{\theta N}{(\theta-1)N+1}
 \le \frac{\theta}{\theta-1}.
\end{align}
\end{lemma}

\begin{proof}
Let \((\lambda_m,e_m)_{m\ge1}\) be an orthonormal spectral representation of the nonzero part of \(\Tc\).  Since \(\Tc\one=0\), each \(e_m\) is centered.  Use the centered energy in \eqref{eq:spectral-centered-energy}.  
For \(\tl W_M=\sum_{m=1}^M\lambda_m e_m\otimes e_m\),
\begin{align}\label{eq:spectral-finite-rank-identity}
 \frac1N\sum_{1\le i<j\le N}\tl W_M(\xi_i,\xi_j)
 =\frac12\sum_{m=1}^M\lambda_m
 \left[
 \pa*{N^{-1/2}\sum_{i=1}^Ne_m(\xi_i)}^2
 -N^{-1}\sum_{i=1}^Ne_m(\xi_i)^2
 \right].
\end{align}
The multivariate central limit theorem and the law of large numbers therefore give convergence in distribution, for each fixed \(M\), to
\begin{align}\label{eq:spectral-finite-rank-limit}
 \frac12\sum_{m=1}^M\lambda_m(Z_m^2-1),
\end{align}
where \((Z_m)_{m\ge1}\) are independent standard Gaussian random variables.  If \(R_M=\tl W-\tl W_M\), canonicity gives the exact variance identity
\begin{align}\label{eq:spectral-tail-variance}
 \E_{\mu^{\otimes N}}\left[
 \left|\frac1N\sum_{1\le i<j\le N}R_M(\xi_i,\xi_j)\right|^2
 \right]
 =\frac{N-1}{2N}\norm{R_M}_{L^2(\mu^{\otimes2})}^2.
\end{align}
For clarity, write
\begin{align}\label{eq:spectral-finite-rank-statistic-def}
 \mathcal X_{N,M}\coloneqq\frac1N\sum_{1\le i<j\le N}\tl W_M(\xi_i,\xi_j).
\end{align}
Since \((\lambda_m)_{m\ge1}\in\ell^2\), define
\begin{align}
 \mathcal X_\infty&\coloneqq\frac12\sum_{m\ge1}\lambda_m(Z_m^2-1),
 \label{eq:spectral-infinite-chaos-def}\\
 \mathcal X_{\infty,M}&\coloneqq\frac12\sum_{m=1}^M\lambda_m(Z_m^2-1).
 \label{eq:spectral-truncated-chaos-def}
\end{align}
The first series converges in \(L^2\).  If \(\Phi:\R\to\R\) is bounded and Lipschitz, then \eqref{eq:spectral-tail-variance} and the centeredness of \(h_\mu-2I_\mu\) give
\begin{align}
 &\left|\E_{\mu^{\otimes N}}\left[\Phi(\mathcal X_N)\right]
 -\E_{\mu^{\otimes N}}\left[\Phi(\mathcal X_{N,M})\right]\right|\notag\\
 &\qquad\le \operatorname{Lip}(\Phi)
 \left[
 \left(\frac{N-1}{2N}\right)^{1/2}
 \norm{R_M}_{L^2(\mu^{\otimes2})}
 +\frac1{\sqrt N}\norm{h_\mu-2I_\mu}_{L^2(\mu)}
 \right].
 \label{eq:spectral-finite-N-triangular-bound}
\end{align}
Moreover, independence of the Gaussian coordinates yields
\begin{align}
 \E\left[\left|\mathcal X_\infty-\mathcal X_{\infty,M}\right|^2\right]
 =\frac12\sum_{m>M}\lambda_m^2
 =\frac12\norm{R_M}_{L^2(\mu^{\otimes2})}^2
 \longrightarrow0.
 \label{eq:spectral-limit-tail-variance}
\end{align}
For fixed \(M\), \eqref{eq:spectral-finite-rank-identity} gives \(\mathcal X_{N,M}\Longrightarrow\mathcal X_{\infty,M}\).  Combining this convergence with \eqref{eq:spectral-finite-N-triangular-bound} and \eqref{eq:spectral-limit-tail-variance}, then letting first \(N\to\infty\) and afterward \(M\to\infty\), proves convergence of bounded Lipschitz expectations.  Hence,
\begin{align}\label{eq:spectral-quadratic-chaos-limit}
 \mathcal X_N\Longrightarrow\mathcal X_\infty.
\end{align}

The identity \(\K_{N,\theta\be}(\mu)=e^{\theta\be I_\mu}\E[e^{-\theta\be\mathcal X_N}]\) and the assumed partition-function bound imply
\begin{align}\label{eq:spectral-uniform-exponential-moment}
 \sup_{N\ge1}\E_{\mu^{\otimes N}}\left[e^{-\theta\be\mathcal X_N}\right]<\infty.
\end{align}
By the Portmanteau theorem,
\begin{align}\label{eq:spectral-limit-exponential-moment}
 \E\left[e^{-\theta\be\mathcal X_\infty}\right]<\infty.
\end{align}
Fix \(m\).  The complementary series
\begin{align}\label{eq:spectral-complementary-series}
 R_m\coloneqq\frac12\sum_{n\ne m}\lambda_n(Z_n^2-1)
\end{align}
converges in \(L^2\) and almost surely and is independent of \(Z_m\).  Tonelli's theorem and independence reduce the corresponding one-coordinate factor in \eqref{eq:spectral-limit-exponential-moment} to
\begin{align}\label{eq:spectral-one-coordinate-integral}
 e^{\theta\be\lambda_m/2}(2\pi)^{-1/2}
 \int_{\R}e^{-(1+\theta\be\lambda_m)z^2/2}\,\md z.
\end{align}
This integral is finite exactly when \(1+\theta\be\lambda_m>0\); equality gives a constant integrand after combination with the Gaussian density, and a negative coefficient gives exponential growth.  Since the independent complementary factor is strictly positive, \eqref{eq:spectral-limit-exponential-moment} proves \eqref{eq:partition-spectral-strict-lower} for every nonzero eigenvalue.  The zero spectral value satisfies the same inequality trivially.

For a negative spectral value, \eqref{eq:partition-spectral-strict-lower} gives
\begin{align}\label{eq:spectral-finite-N-lower}
 1+\be\frac{N-1}{N}\lambda
 >1-\frac{N-1}{\theta N}
 =\frac{(\theta-1)N+1}{\theta N},
\end{align}
whereas a nonnegative spectral value gives a lower bound of one.  The functional calculus proves \eqref{eq:partition-spectral-resolvent-bound}.
\end{proof}

Until \Cref{subsec:signed-first-cavity-estimate}, let \(\theta>1\) be such that the partition-function bound at \(\theta\be\) holds, in addition to the quantitative moment and cavity hypotheses of \Cref{subsec:quantitative-hypotheses}.
  In the all-marginal theorem we take \(\theta=2\), while in \Cref{thm:fixed-marginal-Ltau} we take \(\theta=\tau\).  The case \(\be=0\) is immediate, so the screened argument below is needed only for \(\be>0\).

\noindent As in \Cref{sec:analytic-hypotheses-constants}, write \(\xi_1,\ldots,\xi_N\) for the coordinate random variables on \((\R^\d)^N\), reserving \(X_m\) for the configuration tuple \((x_1,\ldots,x_m)\).

\begin{lemma}[A one-dimensional moment input]\label{lem:fixed-x-moment}
This lemma is used only under the bounded lower-tail condition.  Here, \(\mathsf R_p\) denotes any admissible constant in the scalar Rosenthal inequality.  Under \eqref{eq:cavity-tail-bounded}, with \(p\) as in \eqref{eq:p-choice-quantitative},
\begin{align}\label{eq:centered-sum-Q-second}
 \sup_{N\ge2}\sup_{x\in\R^\d}
 \E_{\Xi_N\sim \Q_{N,\be}(\mu)}\left[\left|\frac1N\sum_{j=2}^N \tl W(x,\xi_j)\right|^2\right]
 \le
 \frac{\mathfrak P_{r\be}(\mu)^{1/r}(2\mathsf R_p)^{2/p}
 \norm{\tl W}_{L_x^\infty L_y^p(\mu^{\otimes2})}^2}{N}.
\end{align}
\end{lemma}

\begin{proof}
The mixed-norm bound \(\tl W\in L_x^\infty L_y^p(\mu^{\otimes2})\) is part of the bounded lower-tail condition \eqref{eq:cavity-tail-bounded}.

The H\"older conjugate of \(r=p/(p-2)\) is \(r'=p/2\).  By \Cref{lem:full-density-Lr},
\(\norm{f_{N,\be}}_{L^r}\le \mathfrak P_{r\be}(\mu)^{1/r}\).  For fixed \(x\), the random variables
\(\tl W(x,\xi_j)\) are centered under \(\mu\), because \(\tl W\) is centered in each variable.
The scalar Rosenthal inequality \cite{Rosenthal1970} and \eqref{eq:cavity-tail-bounded} imply
\begin{align}
 \E_{\mu^{\otimes N}}\left[\left|\frac1N\sum_{j=2}^N \tl W(x,\xi_j)\right|^p\right]
 \le 2\mathsf R_p
 \norm{\tl W}_{L_x^\infty L_y^p(\mu^{\otimes2})}^p
 N^{-p/2}.
\end{align}
Therefore, by H\"older's inequality,
\begin{align}
\E_{\Q_{N,\be}(\mu)}\left[\left|\frac1N\sum_{j=2}^N \tl W(x,\xi_j)\right|^2\right]
&\le \norm{f_{N,\be}}_{L^r(\mu^{\otimes N})}
\left(\E_{\mu^{\otimes N}}\left[\left|\frac1N\sum_{j=2}^N \tl W(x,\xi_j)\right|^p\right]\right)^{2/p}\notag\\
&\le
\frac{\mathfrak P_{r\be}(\mu)^{1/r}(2\mathsf R_p)^{2/p}
\norm{\tl W}_{L_x^\infty L_y^p(\mu^{\otimes2})}^2}{N},
\end{align}
which proves \eqref{eq:centered-sum-Q-second}.
\end{proof}

The joint \(L^p(\mu^{\otimes2})\) assumption in \eqref{eq:mu-controls} and the \(L_x^\infty L_y^p(\mu^{\otimes2})\) bound in \eqref{eq:cavity-tail-bounded} serve different purposes as assumptions.  The present fixed-\(x\) Rosenthal estimate uses the latter, whereas the partition-function estimate controls the \(L^r\)-norm of the full Gibbs density appearing in H\"older's inequality.

\subsection{Conditional density and cancellation}
\label{subsec:conditional-density-cancellation}

Recall that \(f_{N,\be}^{(1)}\) denotes the density of \(\Q_{N,\be}^{(1)}(\mu)\) with respect to \(\mu\).  Initially, we use only that it belongs to \(L^1(\mu)\).  Set
\begin{align}\label{eq:uN-def}
 u_N\coloneqq f_{N,\be}^{(1)}-1.
\end{align}
For \(N\ge2\), set \(\hat\Xi_1\coloneqq(\xi_2,\ldots,\xi_N)\) and condition on \(\sigma(\hat\Xi_1)\).  Up to a multiplicative factor depending only on \(\hat\Xi_1\), the \(x\)-dependent part of the Gibbs weight is
\begin{align}\label{eq:conditional-weight-x}
 \exp\left\{\frac{\be}{N}h_\mu(x)-\frac{\be}{N}\sum_{j=2}^N \tl W(x,\xi_j)\right\}.
\end{align}
Subtracting the \(\mu\)-mean of the exponent, which is \(2\be I_\mu/N\), gives exactly the cavity field \(B_N\) defined in \eqref{eq:cavity-B-def}.  In particular,
\begin{align}\label{eq:B-mean-zero}
 \int B_N(x;\hat\Xi_1)\,\md\mu(x)=0,
 \qquad \forall \hat\Xi_1,
\end{align}
and the conditional density of \(\xi_1\) given \(\hat\Xi_1\) is
\begin{align}\label{eq:conditional-density}
 \pi_N(x\mid \hat\Xi_1)=\frac{e^{B_N(x;\hat\Xi_1)}}{\int e^{B_N(z;\hat\Xi_1)}\,\md\mu(z)}.
\end{align}

We now prove the nonlinear cavity estimates under the two conditions stated in \Cref{subsec:quantitative-hypotheses}.

\parheading{Bounded lower-tail condition.}
Under \eqref{eq:cavity-tail-bounded},
\begin{align}\label{eq:g-tildeg-lb}
 \tl W(x,y)
 \ge \operatorname*{ess\,inf}_{(z,w)\in(\R^\d)^2}\tl W(z,w),
\end{align}
and therefore
\begin{align}\label{eq:B-upper-bound-expanded}
 B_N(x;\hat\Xi_1)
 \le \frac{\be}{N}\norm{h_\mu-2I_\mu}_{L^\infty}
 -\be\frac{N-1}{N}
 \operatorname*{ess\,inf}_{(z,w)\in(\R^\d)^2}\tl W(z,w).
\end{align}
Thus, with
\begin{align}\label{eq:b-cavity-bounded}
 \mathsf b_{\be,\mu}
 \coloneqq
 \be\pa{\norm{h_\mu-2I_\mu}_{L^\infty}
 -\operatorname*{ess\,inf}_{(z,w)\in(\R^\d)^2}\tl W(z,w)},
\end{align}
we have the deterministic upper bound
\begin{align}\label{eq:B-upper-bound}
 B_N(x;\hat\Xi_1)\le \mathsf b_{\be,\mu}
\end{align}
for all \(N\ge2\), for \(\mu\)-a.e. \(x\), and for \(\mu^{\otimes(N-1)}\)-a.e. \(\hat\Xi_1\).  Combining \eqref{eq:cavity-B-def} with \Cref{lem:fixed-x-moment}, define
\begin{align}\label{eq:B2-cavity-bounded-constant}
 \mathsf B_{2,\be,\mu}
 \coloneqq
 2\be^2\norm{h_\mu-2I_\mu}_{L^\infty}^2
 +2\be^2\mathfrak P_{r\be}(\mu)^{1/r}
       (2\mathsf R_p)^{2/p}
       \norm{\tl W}_{L_x^\infty L_y^p(\mu^{\otimes2})}^2.
\end{align}
Then
\begin{align}\label{eq:B-second-pointwise}
 \sup_{x\in\R^\d}\E_{\Q_{N,\be}(\mu)}\left[|B_N(x;\hat\Xi_1)|^2\right]
 \le \frac{\mathsf B_{2,\be,\mu}}{N}.
\end{align}

\parheading{Logarithmic lower-tail condition.}
Under \eqref{eq:cavity-tail-log}, the positive part of the cavity field satisfies \eqref{eq:B-log-upper}, and the nonlinear estimates below are controlled by the weighted constants \eqref{eq:log-weighted-cavity-constants}.  The sufficient condition stated in \Cref{lem:weighted-logarithmic-cavity} is proved in \Cref{app:weighted-log-cavity}.

The next lemma is the only nonlinear estimate in the one-point argument.  After centering, it identifies the normalized exponential conditional density with its exponent up to an \(O(N^{-1})\) error in \(L^2(\mu)\).  For \(b\ge0\), set
\begin{align}
 \Lambda_2(b)&\coloneqq\sup_{t\le b,\ t\ne0}\frac{e^t-1-t}{t^2},\label{eq:Lambda2-def}\\
 \Lambda_1(b)&\coloneqq\sup_{t\le b,\ t\ne0}\frac{e^t-1-t}{\abs{t}}.\label{eq:Lambda1-def}
\end{align}
We refer to these two definitions collectively as \eqref{eq:Lambda2-def}--\eqref{eq:Lambda1-def}.

\begin{lemma}[Abstract nonlinear cavity estimate]\label{lem:abstract-cavity-linearization}
For notational simplicity, write \(B(x)=B_N(x;\hat\Xi_1)\), and set
\begin{align}
 Z_B&\coloneqq \int e^{B(z)}\,\md\mu(z),\label{eq:ZB-def}\\
 \epsilon(b)&\coloneqq e^b-1-b,\label{eq:epsilon-def}\\
 \delta_B&\coloneqq \int \epsilon(B(z))\,\md\mu(z).\label{eq:deltaB-def}
\end{align}
Suppose that, for finite constants \(\mathsf E_2,\mathsf D_1,\mathsf D_2,\mathsf B_2\),
\begin{subequations}\label{eq:abstract-cavity-inputs}
\begin{align}
 \norm{\E\left[\epsilon(B(\cdot))\right]}_{L^2(\mu)}
 &\le \frac{\mathsf E_2}{N},\label{eq:abstract-cavity-input-E}\\
 \E\left[\delta_B\right]
 &\le \frac{\mathsf D_1}{N},\label{eq:abstract-cavity-input-D1}\\
 \E\left[\delta_B^2\right]
 &\le \frac{\mathsf D_2}{N},\label{eq:abstract-cavity-input-D2}\\
 \E\left[\int B(x)^2\,\md\mu(x)\right]
 &\le \frac{\mathsf B_2}{N}.\label{eq:abstract-cavity-input-B}
\end{align}
\end{subequations}
Then \(r_N\) defined by \eqref{eq:cavity-error-def} satisfies
\begin{align}\label{eq:abstract-cavity-output}
 \norm{r_N}_{L^2(\mu)}
 \le \frac{\mathsf E_2+\mathsf D_1+\sqrt{\mathsf B_2\mathsf D_2}}{N}.
\end{align}
\end{lemma}

\begin{proof}
Since \(\epsilon\ge0\), \eqref{eq:B-mean-zero} gives
\begin{align}
 Z_B&=1+\delta_B,\label{eq:ZB-deltaB}\\
 \delta_B&\ge0.\label{eq:deltaB-nonnegative}
\end{align}
Moreover,
\begin{align}
 \frac{e^{B(x)}}{Z_B}-1-B(x)
 &=\frac{1+B(x)+\epsilon(B(x))}{1+\delta_B}-1-B(x)\notag\\
 &=\frac{\epsilon(B(x))-\delta_B-B(x)\delta_B}{1+\delta_B}.\label{eq:normalized-exponential-remainder}
\end{align}
Thus, \(1+\delta_B\ge1\) and Cauchy's inequality yield
\begin{align}
 \norm{r_N}_{L^2(\mu)}
 &\le \norm{\E[\epsilon(B(\cdot))]}_{L^2(\mu)}
 +\E[\delta_B]+\norm{\E[|B(\cdot)|\delta_B]}_{L^2(\mu)}\notag\\
 &\le \frac{\mathsf E_2+\mathsf D_1}{N}
 +\E\left[\delta_B\norm{B}_{L^2(\mu)}\right]\notag\\
 &\le \frac{\mathsf E_2+\mathsf D_1}{N}
 +\pa{\E[\delta_B^2]}^{1/2}
  \pa{\E\left[\int B(x)^2\,\md\mu(x)\right]}^{1/2}.\label{eq:abstract-cavity-proof-bound}
\end{align}
This proves \eqref{eq:abstract-cavity-output}.
\end{proof}

\begin{lemma}[Cavity linearization]\label{lem:cavity-linearization}
Define
\begin{align}\label{eq:cavity-error-def}
 r_N(x)\coloneqq
 \E_{\Q_{N,\be}(\mu)}\left[
 \frac{e^{B_N(x;\hat\Xi_1)}}{\int e^{B_N(z;\hat\Xi_1)}\,\md\mu(z)}
 -1-B_N(x;\hat\Xi_1)\right].
\end{align}
Under the bounded lower-tail condition, set
\begin{align}\label{eq:R-cav-bounded-def}
 \mathsf R_{\mathrm{cav},\be,\mu}^{\mathrm{bd}}
 \coloneqq
 \mathsf B_{2,\be,\mu}
 \left(2\Lambda_2(\mathsf b_{\be,\mu})
 +\sqrt{2\Lambda_1(\mathsf b_{\be,\mu})\mathsf b_{\be,\mu}
          \Lambda_2(\mathsf b_{\be,\mu})}\right).
\end{align}
Under the logarithmic lower-tail condition, set
\begin{align}\label{eq:R-cav-log-def}
 \mathsf R_{\mathrm{cav},\be,\mu}^{\log}
 \coloneqq
 \mathsf E_{\log,\be,\mu}+\mathsf D_{\log,\be,\mu}
 +\sqrt{\mathsf B_{\log,\be,\mu}\mathsf V_{\log,\be,\mu}}.
\end{align}
Let \(\mathsf R_{\mathrm{cav},\be,\mu}\) denote the corresponding one of these two constants.  Then, for every \(N\ge2\),
\begin{align}\label{eq:cavity-error-bound}
 \norm{r_N}_{L^2(\mu)}\le \frac{\mathsf R_{\mathrm{cav},\be,\mu}}{N}.
\end{align}
\end{lemma}

\begin{proof}
It suffices to verify the hypotheses of \Cref{lem:abstract-cavity-linearization}.

\step{Step 1: Bounded lower-tail condition.}
Under this condition, \eqref{eq:B-upper-bound} and \eqref{eq:Lambda2-def}--\eqref{eq:Lambda1-def} give
\begin{subequations}\label{eq:epsilon-bounds}
\begin{align}
 0&\le \epsilon(B(x))\le \Lambda_2(\mathsf b_{\be,\mu})B(x)^2,\label{eq:epsilon-quadratic-bound}\\
 0&\le \epsilon(B(x))\le \Lambda_1(\mathsf b_{\be,\mu})|B(x)|.\label{eq:epsilon-linear-bound}
\end{align}
\end{subequations}
Together with \eqref{eq:B-second-pointwise}, this gives
\begin{subequations}\label{eq:bounded-cavity-abstract-inputs}
\begin{align}
 \norm{\E[\epsilon(B(\cdot))]}_{L^2(\mu)}
 &\le \frac{\Lambda_2(\mathsf b_{\be,\mu})\mathsf B_{2,\be,\mu}}{N},\label{eq:bounded-cavity-E}\\
 \E[\delta_B]
 &\le \frac{\Lambda_2(\mathsf b_{\be,\mu})\mathsf B_{2,\be,\mu}}{N},\label{eq:bounded-cavity-D1}\\
 \E\left[\int B(x)^2\,\md\mu(x)\right]
 &\le \frac{\mathsf B_{2,\be,\mu}}{N}.\label{eq:bounded-cavity-B}
\end{align}
\end{subequations}
Furthermore, \eqref{eq:B-mean-zero} implies \(\int B_{+}\,\md\mu=\int B_{-}\,\md\mu\), and therefore
\begin{align}\label{eq:B-L1-bounded-condition}
 \int |B|\,\md\mu=2\int B_+\,\md\mu\le 2\mathsf b_{\be,\mu}.
\end{align}
Using \eqref{eq:epsilon-linear-bound}, we get \(\delta_B\le 2\Lambda_1(\mathsf b_{\be,\mu})\mathsf b_{\be,\mu}\), whence
\begin{align}\label{eq:bounded-cavity-D2}
 \E[\delta_B^2]
 \le
 \frac{2\Lambda_1(\mathsf b_{\be,\mu})\mathsf b_{\be,\mu}
       \Lambda_2(\mathsf b_{\be,\mu})\mathsf B_{2,\be,\mu}}{N}.
\end{align}
\Cref{lem:abstract-cavity-linearization} gives \eqref{eq:cavity-error-bound} with \(\mathsf R_{\mathrm{cav},\be,\mu}=\mathsf R_{\mathrm{cav},\be,\mu}^{\mathrm{bd}}\).

\step{Step 2: Logarithmic lower-tail condition.}
Under this condition, the elementary inequalities
\begin{subequations}\label{eq:log-taylor-bounds}
\begin{align}
 0&\le e^t-1-t\le \frac12 e^{t_+}t^2,\label{eq:log-taylor-quadratic}\\
 0&\le e^t-1-t\le e^{t_+}|t|\label{eq:log-taylor-linear}
\end{align}
\end{subequations}
and the bound \eqref{eq:B-log-upper} imply
\begin{subequations}\label{eq:log-epsilon-bounds}
\begin{align}
 \epsilon(B(x))&\le \frac12 e^{\Phi_N(x;\hat\Xi_1)}B(x)^2,\label{eq:log-epsilon-quadratic}\\
 \epsilon(B(x))&\le e^{\Phi_N(x;\hat\Xi_1)}|B(x)|.\label{eq:log-epsilon-linear}
\end{align}
\end{subequations}
The weighted estimates \eqref{eq:log-weighted-cavity-constants} therefore give
\begin{subequations}\label{eq:log-cavity-abstract-inputs}
\begin{align}
 \norm{\E[\epsilon(B(\cdot))]}_{L^2(\mu)}
 &\le \frac{\mathsf E_{\log,\be,\mu}}{N},\label{eq:log-cavity-E}\\
 \E[\delta_B]
 &\le \frac{\mathsf D_{\log,\be,\mu}}{N},\label{eq:log-cavity-D1}\\
 \E[\delta_B^2]
 &\le \frac{\mathsf V_{\log,\be,\mu}}{N},\label{eq:log-cavity-D2}\\
 \E\left[\int B(x)^2\,\md\mu(x)\right]
 &\le \frac{\mathsf B_{\log,\be,\mu}}{N}.\label{eq:log-cavity-B}
\end{align}
\end{subequations}
Another application of \Cref{lem:abstract-cavity-linearization} gives \eqref{eq:cavity-error-bound} with \(\mathsf R_{\mathrm{cav},\be,\mu}=\mathsf R_{\mathrm{cav},\be,\mu}^{\log}\).
\end{proof}

We can now close the one-point estimate.  By \eqref{eq:conditional-density},
\begin{align}\label{eq:one-point-density-conditional-average}
 f_{N,\be}^{(1)}(x)=\E_{\Q_{N,\be}(\mu)}\left[\pi_N(x\mid\hat\Xi_1)\right].
\end{align}
Therefore, by \Cref{lem:cavity-linearization},
\begin{subequations}\label{eq:uN-linearized}
\begin{align}
 u_N(x)&=\E_{\Q_{N,\be}(\mu)}\left[B_N(x;\hat\Xi_1)\right]+r_N(x),\label{eq:uN-linearized-identity}\\
 \norm{r_N}_{L^2(\mu)}&\le \frac{\mathsf R_{\mathrm{cav},\be,\mu}}{N}.\label{eq:uN-linearized-error}
\end{align}
\end{subequations}
The bounded branch estimate \eqref{eq:B-second-pointwise} and the logarithmic estimate \eqref{eq:log-cavity-B} imply, respectively, that there is a finite constant \(\mathsf B_{\mathrm{cav},\be,\mu}\) such that
\begin{align}\label{eq:expected-cavity-square}
 \E_{\Q_{N,\be}(\mu)}\left[
 \norm{B_N(\cdot;\hat\Xi_1)}_{L^2(\mu)}^2
 \right]
 \le \frac{\mathsf B_{\mathrm{cav},\be,\mu}}{N}.
\end{align}
The section map is strongly measurable by approximation with rectangular simple functions.  It is therefore Bochner integrable in \(L^2(\mu)\), and Jensen's inequality gives
\begin{align}\label{eq:averaged-cavity-L2-preliminary}
 \norm*{\E_{\Q_{N,\be}(\mu)}\left[B_N(\cdot;\hat\Xi_1)\right]}_{L^2(\mu)}^2
 \le \frac{\mathsf B_{\mathrm{cav},\be,\mu}}{N}.
\end{align}
Together with \eqref{eq:uN-linearized-error}, this proves \(u_N\in L^2(\mu)\) before the screened equation is written.

Since \(f_{N,\be}^{(1)}=1+u_N\in L^2(\mu)\) and \(\tl W\in L^2(\mu^{\otimes2})\),
\begin{align}\label{eq:screened-Fubini-integrability}
 \int_{(\R^\d)^2}|\tl W(x,y)|f_{N,\be}^{(1)}(y)\,\md\mu(x)\md\mu(y)
 \le \norm{\tl W}_{L^2(\mu^{\otimes2})}
 \norm{f_{N,\be}^{(1)}}_{L^2(\mu)}<\infty.
\end{align}
Thus, Fubini's theorem and exchangeability justify
\begin{align}
\E_{\Q_{N,\be}(\mu)}\left[B_N(x;\hat\Xi_1)\right]
&=\frac{\be}{N}\pa{h_\mu(x)-2I_\mu}
-\frac{\be}{N}\sum_{j=2}^N \E_{\Q_{N,\be}(\mu)}\left[\tl W(x,\xi_j)\right]\notag\\
&=\frac{\be}{N}\pa{h_\mu(x)-2I_\mu}
-\be\frac{N-1}{N}\int \tl W(x,y)f_{N,\be}^{(1)}(y)\,\md\mu(y)\notag\\
&=\frac{\be}{N}\pa{h_\mu(x)-2I_\mu}
-\be\frac{N-1}{N}\Tc u_N(x).\label{eq:EB-computation}
\end{align}
Here, the last equality follows from \(\int \tl W(x,y)\,\md\mu(y)=0\).  Combining \eqref{eq:uN-linearized-identity} and \eqref{eq:EB-computation},
\begin{align}\label{eq:screened-one-point-equation}
 \pa*{I+\be\frac{N-1}{N}\Tc}u_N
 =\frac{\be}{N}\pa{h_\mu-2I_\mu}+r_N.
\end{align}
By \Cref{lem:partition-spectral-stability}, the screened operator is invertible and satisfies \eqref{eq:partition-spectral-resolvent-bound}.  Define
\begin{align}\label{eq:one-point-theta-constant}
 \mathsf U_{\theta,\be,\mu}
 \coloneqq
 \frac{\theta}{\theta-1}
 \left(
 \be\norm{h_\mu-2I_\mu}_{L^2(\mu)}
 +\mathsf R_{\mathrm{cav},\be,\mu}
 \right).
\end{align}
Then \eqref{eq:screened-one-point-equation} gives
\begin{align}\label{eq:one-point-L2-bound}
 \norm{u_N}_{L^2(\mu)}\le \frac{\mathsf U_{\theta,\be,\mu}}{N},
 \qquad N\ge2.
\end{align}
For the all-marginal theorem, set
\begin{align}\label{eq:one-point-constant}
 \mathsf U_{\be,\mu}
 \coloneqq \mathsf U_{2,\be,\mu}
 =2\left(
 \be\norm{h_\mu-2I_\mu}_{L^2(\mu)}
 +\mathsf R_{\mathrm{cav},\be,\mu}
 \right).
\end{align}
Thus, under the hypotheses of \Cref{thm:chi2-chaos-two-sided},
\begin{align}\label{eq:one-point-chi2}
 \Chi\pa{\Q_{N,\be}^{(1)}(\mu)\mid\mu}
 =\norm{f_{N,\be}^{(1)}-1}_{L^2(\mu)}^2
 \le \frac{\mathsf U_{\be,\mu}^2}{N^2}.
\end{align}
By \Cref{lem:normalized-conditional-identity} with \(k=1\), this is equivalently the one-mark estimate
\begin{align}\label{eq:one-mark-Psi-final}
 \left\|\frac{\Psi_{N,1}}{\int\Psi_{N,1}\,\md\mu}-1\right\|_{L^2(\mu)}
 \le \frac{\mathsf U_{\be,\mu}}{N}.
\end{align}

\begin{remark}
The seemingly dangerous term \(N^{-1}\sum_{j\ne1}\tl W(x,\xi_j)\) has a nonzero Gibbs average, and this average is exactly \(\frac{N-1}{N}\Tc u_N\).  Rather than estimating this contribution as an error, we absorb it into the left-hand side of \eqref{eq:screened-one-point-equation}.  The partition-function moment excludes spectrum at or beyond the singular threshold and gives the uniform inverse bound \eqref{eq:partition-spectral-resolvent-bound}, even when \(\Tc\) has both positive and negative spectral components.  Thus, no smallness condition on \(\be\) or on \(\norm{\Tc}\), and no positivity assumption, is needed.
\end{remark}

% STAGE2 MOVE: sec:hoeffding-projection-upper-bound
\subsection{Moment bounds for quantitative marginal approximation}
\label{subsec:signed-first-cavity-estimate}

In this subsection, assume the hypotheses of
\Cref{thm:abstract-quantitative-marginal}, independently of the order-bound
hypotheses used above.  We verify the same first-order cavity estimate
when the source is only bounded above and belongs to $L^6(\mu)$.

With $\mathcal X_N$ as in \eqref{eq:spectral-centered-energy}, set
\begin{align}
 \K^{\mathrm c}_{N,\gamma}
 &\coloneqq e^{-\gamma I_\mu}\K_{N,\gamma}(\mu)
 =\E_{\mu^{\otimes N}}\left[e^{-\gamma\mathcal X_N}\right],
 \label{eq:exact-centered-partition}\\
 f_{N,\be}&=\frac{e^{-\be\mathcal X_N}}{\K^{\mathrm c}_{N,\be}}.
 \label{eq:exact-centered-density}
\end{align}
The exponent is $-\be\mathsf W_N(\tl W)/2+
\be N^{-1}\sum_i\psi_\mu(x_i)$ and has product mean zero.  Hence,
\begin{equation}\label{eq:exact-Z-lower}
 \K^{\mathrm c}_{N,\gamma}\geq1,\qquad N\geq1,\quad 0\leq\gamma\leq\Gamma.
\end{equation}
H\"older extends \eqref{eq:signed-centered-pair-moment} to the interval
$[0,\Gamma]$.  The upper bound for $\psi_\mu$ and
\Cref{lem:full-density-Lr} therefore give
\begin{equation}\label{eq:exact-full-density-Lr}
 \sup_{N\geq2}\norm{f_{N,\be}}_{L^r(\mu^{\otimes N})}^r
 \leq e^{r\be C_+}\sup_{N\geq2}
     \E_{\mu^{\otimes N}}\left[e^{-r\be\mathsf W_N(\tl W)/2}\right]
 <\infty,\qquad 1\leq r\leq r_*.
\end{equation}
The same bound holds for every marginal by conditional expectation.  The
$N=1$ normalizer is finite because $\psi_\mu\leq C_+$, so
\Cref{lem:partition-spectral-stability} applies with $\theta=r_*$.  In
particular,
\begin{equation}\label{eq:signed-resolvent-bound}
 0<R_{\be,\mu}\leq\frac{r_*}{r_*-1}\mathrm{Id},\qquad
 \left\|\left(\mathrm{Id}+\be\frac{N-1}{N}T_\mu\right)^{-1}\right\|_{\mathrm{op}}
 \leq\frac{r_*}{r_*-1},\qquad N\geq2.
\end{equation}

\begin{lemma}[Cavity moments under the quantitative hypotheses]
\label{lem:signed-cavity-moments}
Use the field $B_N$ from \eqref{eq:cavity-B-def}.  It satisfies
$B_N\leq\be(C_++C_-)$ outside a product-null set.  Uniformly for $N\geq2$,
\begin{equation}\label{eq:signed-point-cavity-moment}
 \E_{\Q_{N,\be}(\mu)}\left[\abs{B_N(x)}^{2+\vartheta}\right]
 \leq C\left(N^{-1-\vartheta/2}
          +N^{-2-\vartheta}\abs{\psi_\mu(x)}^{2+\vartheta}\right)
 \quad\text{for $\mu$-almost every $x$}.
\end{equation}
Moreover, for every finite $q\geq2$,
\begin{equation}\label{eq:signed-Hilbert-cavity-moment}
 \E_{\Q_{N,\be}(\mu)}\left[\norm{B_N}_{L^2(\mu)}^q\right]
 \leq C_qN^{-q/2}.
\end{equation}
With $\delta_B$ defined in \eqref{eq:deltaB-def}, for every finite $q\geq1$,
\begin{equation}\label{eq:signed-delta-moment}
 0\leq\delta_B\leq C\norm{B_N}_2^2,\qquad
 \E_{\Q_{N,\be}(\mu)}[\delta_B^q]\leq C_qN^{-q}.
\end{equation}
\end{lemma}

\begin{proof}
The upper bound for the field follows from
\eqref{eq:signed-kernel-assumptions}--\eqref{eq:signed-source-assumptions}.
For fixed $x$, the variables $w_x(x_j)$ are independent and centered under
the product law.  Rosenthal's inequality~\cite{Rosenthal1970} and the row
bound give
\begin{equation}\label{eq:signed-product-row-moment}
 \E_{\mu^{\otimes N}}\left[\left|\frac1N\sum_{j=2}^Nw_x(x_j)\right|^p\right]
 \leq C_pM_p^pN^{-p/2}.
\end{equation}
Take $r_0=p/(p-2-\vartheta)$ and its conjugate $r_0'=p/(2+\vartheta)$.
The full-density bound \eqref{eq:exact-full-density-Lr} applies at $r_0$.
Separating the deterministic source in $B_N$ and applying H\"older proves
\eqref{eq:signed-point-cavity-moment}.  Lyapunov's inequality gives the
corresponding estimate at every exponent between $2$ and $2+\vartheta$.

By symmetry, the row map $x\mapsto w_x$ takes values in $L^2(\mu)$, has
norm at most $M_p$, and has Bochner mean zero.  Strong measurability follows
by approximation of the measurable square-integrable kernel with
rectangular simple functions.  For independent centered Hilbert-valued
variables $Y_j$ with $\norm{Y_j}\leq M_p$, symmetrization and the
Hilbert-space Khintchine--Kahane inequality~\cite{Kahane1985} give
\begin{equation}\label{eq:signed-Hilbert-product-sum}
 \E\left[\left\|\sum_{j=2}^NY_j\right\|^q\right]
 \leq C_q\E\left[\left(\sum_{j=2}^N\norm{Y_j}^2\right)^{q/2}\right]
 \leq C_qM_p^qN^{q/2}.
\end{equation}
Apply this at exponent $2q$ to $Y_j=w_{x_j}$.  H\"older with the
full-density $L^2$ bound transfers the estimate to Gibbs expectation.
The deterministic term requires only $\psi_\mu\in L^2(\mu)$.  This proves
\eqref{eq:signed-Hilbert-cavity-moment}.  Finally, on the upper-bounded
range of $B_N$, $0\leq e^{B_N}-1-B_N\leq CB_N^2$.  Integration and
\eqref{eq:signed-Hilbert-cavity-moment} prove
\eqref{eq:signed-delta-moment}.
\end{proof}

\begin{prop}[The first component under the quantitative hypotheses]
\label{prop:signed-onepoint-bound}
Under the hypotheses of \Cref{thm:abstract-quantitative-marginal},
\begin{equation}\label{eq:signed-onepoint-bound}
 \norm{r_N}_{L^2(\mu)}+\norm{u_N}_{L^2(\mu)}\leq\frac C N,
 \qquad N\geq2,
\end{equation}
where $u_N$ and $r_N$ are defined in \eqref{eq:uN-def} and
\eqref{eq:cavity-error-def}.
\end{prop}
\begin{proof}
We verify \Cref{lem:abstract-cavity-linearization}.  The exponent-$2$
consequence of \eqref{eq:signed-point-cavity-moment} and
$\psi_\mu\in L^4(\mu)$ give
\begin{equation}\label{eq:signed-expected-cavity-square}
 \left\|\E_{\Q_{N,\be}(\mu)}[B_N(\cdot)^2]\right\|_2\leq C/N.
\end{equation}
The upper field bound implies $\epsilon(B_N)\leq CB_N^2$.
The other three inputs follow from \eqref{eq:signed-delta-moment} and
\eqref{eq:signed-Hilbert-cavity-moment}; in particular,
$\E[\delta_B^2]\leq CN^{-2}\leq CN^{-1}$.
Thus, $\norm{r_N}_2\leq C/N$.  The full-density $L^2$ bound already gives
$u_N\in L^2(\mu)$.  The conditional-density and Fubini computation leading
to \eqref{eq:screened-one-point-equation} uses only this integrability,
canonicity, and the exact field, so the same equation holds here.  Applying
\eqref{eq:signed-resolvent-bound} proves the asserted estimate for $u_N$.
\end{proof}

\subsection{Proof of the marginal upper bounds}
\label{subsec:order-upper-conclusion}\label{sec:hoeffding-projection-upper-bound}

The full-space \(\chi^2\)-divergence is uniformly bounded by the partition-function estimate.  Indeed,
\begin{align}\label{eq:full-chi2-partition-ratio}
 1+\Chi\pa{\Q_{N,\be}(\mu)\mid\mu^{\otimes N}}
 =\frac{\K_{N,2\be}(\mu)}{\K_{N,\be}(\mu)^2}
 \le \mathfrak P_{2\be}(\mu),
\end{align}
using the definition of \(\mathfrak P_{2\be}(\mu)\) and Jensen's lower bound at \(\be\).

\begin{proof}[Proof of the upper bound in \Cref{thm:chi2-chaos-two-sided}]
The case \(\be=0\) is immediate.  If \(N=1\), then necessarily \(k=1\), and
\begin{align}\label{eq:N1-full-chi2-bound}
 1+\Chi\pa{\Q_{1,\be}(\mu)\mid\mu}
 =\frac{\K_{1,2\be}(\mu)}{\K_{1,\be}(\mu)^2}
 \le \mathfrak P_{2\be}(\mu).
\end{align}
so the asserted bound follows directly.  We may therefore assume \(N\ge2\).  Apply the one-point estimate with \(\theta=2\), and then apply \Cref{lem:hoeffding-reduction} with \(g_N=f_{N,\be}\).  Combining \eqref{eq:one-point-chi2} and \eqref{eq:full-chi2-partition-ratio}, we obtain
\begin{align}
\Chi\pa{\Q_{N,\be}^{(k)}(\mu)\mid\mu^{\otimes k}}
&\le k\frac{\mathsf U_{\be,\mu}^2}{N^2}
+\bigl(\mathfrak P_{2\be}(\mu)-1\bigr)\frac{k^2}{N^2}.
\end{align}
The \(\chi^2\)-upper constant used here is
\begin{align}\label{eq:chi-upper-constant}
 \mathsf C_{\chi,\be,\mu}^{\mathrm{up}}
 \coloneqq \mathsf U_{\be,\mu}^2+\bigl(\mathfrak P_{2\be}(\mu)-1\bigr).
\end{align}
Since \(k\le k^2\) for \(k\ge1\), this proves the \(\chi^2\)-line of \eqref{eq:unified-upper-final}
with the constant \(\mathsf C_{\chi,\be,\mu}^{\mathrm{up}}\).
\end{proof}

% STAGE2 MOVE: A4-subquadratic-extension-second
\section{Strong approximation of the pair kernel}
\label{sec:second-hoeffding-component}

The scalar Laplace comparison in \Cref{subsec:signed-scalar-Laplace}
applies to signed canonical kernels.  Under the quantitative Gibbs
hypotheses, it estimates the correlated-replica and screened-response
families.  These estimates identify the squared norm and cross term of
the pair kernel; polarization gives strong convergence.  No interacting fluctuation theorem or weak-covariance identification is used.

\subsection{Quantitative Laplace comparison without a sign condition}
\label{subsec:signed-scalar-Laplace}

We first prove a scalar estimate with assumptions independent of the Gibbs
hypotheses.  Let $(S,\rho)$ be a standard Borel probability space, and let
$c\in L^{2+\alpha}(\rho^{\otimes2})$, where $0<\alpha\leq1$, be real,
symmetric, and canonical.  Write $C$ for its integral operator and
$\mathsf W_N(c)=N^{-1}\sum_{i\ne j}c(x_i,x_j)$ for independent coordinates
with law $\rho$.  Let
\begin{equation}\label{eq:signed-general-Gaussian-form}
 \mathcal Q_C\coloneqq\sum_j\lambda_j(C)(Z_j^2-1)
 \quad\text{in }L^2,
\end{equation}
where the $Z_j$ are independent standard Gaussians.  The common construction
and Hilbert--Schmidt continuity of these forms are given in
\Cref{quant-tools:subsec-chaos-construction}.

\begin{lemma}[Smooth comparison for a canonical quadratic statistic]
\label{lem:quadratic-smooth-comparison}
If $h\in C^2(\R)$ has bounded first and second derivatives and globally
$\alpha$-H\"older second derivative, then
\begin{align}
 &\left|\E\left[h\left(\frac12\mathsf W_N(c)\right)\right]
          -\E\left[h\left(\frac12\mathcal Q_C\right)\right]\right|\notag\\
 &\qquad\leq C_\alpha[h'']_{C^{0,\alpha}}
         N^{-\alpha/2}\norm c_{2+\alpha}^{2+\alpha}
       +\frac{\norm{h'}_\infty\norm c_2}{\sqrt{2N}},\qquad N\geq2.
 \label{eq:quadratic-smooth-comparison}
\end{align}
The constant depends only on $\alpha$.
\end{lemma}

The proof in \Cref{quant-tools:subsec-smooth-proof} replaces one independent
coordinate at a time, following the Lindeberg strategy of
\cite[Section~2]{Chatterjee2006}.  Moment matching cancels the first two
Taylor terms.  Only polynomial moments are used under the intermediate
original/Gaussian product laws.  The same lemma will give the absolute-moment
estimate used for growing total variation.

\begin{thm}[Quantitative Laplace comparison]
\label{thm:quantitative-signed-Laplace}
Fix $\be>0$ and $q>1$, and suppose that
\begin{equation}\label{eq:signed-scalar-exponential-assumption}
 M\coloneqq\sup_{N\geq2}\E_{\rho^{\otimes N}}
                \left[e^{-q\be\mathsf W_N(c)/2}\right]<\infty.
\end{equation}
Then every spectral value $\lambda$ of $C$ satisfies $1+q\be\lambda>0$,
and, for every $N\geq2$,
\begin{align}
 &\left|\E_{\rho^{\otimes N}}\left[e^{-\be\mathsf W_N(c)/2}\right]
        -\detTwo(\mathrm{Id}+\be C)^{-1/2}\right|\notag\\
 &\qquad\leq C_{\alpha,\be,q}
       \left(M+\norm c_{2+\alpha}^{2+\alpha}+\norm c_2\right)
       N^{-\alpha(q-1)/(2q)}.
 \label{eq:signed-scalar-Laplace-bound}
\end{align}
The determinant is positive.  The constant is independent of the rank,
probability space, and $N$.
\end{thm}

\begin{proof}
\step{The limiting exponential and its determinant.}
Applying \Cref{lem:quadratic-smooth-comparison} to smooth bounded tests gives
$\mathsf W_N(c)/2\Longrightarrow\mathcal Q_C/2$.
The Portmanteau theorem and
\eqref{eq:signed-scalar-exponential-assumption} imply
\begin{equation}\label{eq:signed-Gaussian-exponential-moment}
 \E\left[e^{-q\be\mathcal Q_C/2}\right]\leq M.
\end{equation}
The one-mode argument in the proof of
\Cref{lem:partition-spectral-stability} now gives
$1+q\be\lambda>0$ for every eigenvalue, including strictness at the boundary:
the corresponding one-dimensional Gaussian integral would otherwise be
infinite, while the independent complementary exponential is positive.
The zero spectral value is harmless.  Hence,
\begin{equation}\label{eq:signed-scalar-spectral-gap}
 \mathrm{Id}+\be C\geq(1-q^{-1})\mathrm{Id}.
\end{equation}

Let $\mathcal Q_C^{(m)}$ be the spectral partial sum through the first $m$
Gaussian modes.  It is the conditional expectation of $\mathcal Q_C$ given
those modes.  Conditional Jensen gives
\begin{equation}\label{eq:signed-spectral-uniform-integrability}
 \E\left[\left(e^{-\be\mathcal Q_C^{(m)}/2}\right)^q\right]
 \leq\E\left[e^{-q\be\mathcal Q_C/2}\right]\leq M.
\end{equation}
The exponentials are therefore uniformly integrable.  Their convergence
in probability and finite-dimensional Gaussian integration yield
\begin{equation}\label{eq:signed-Gaussian-determinant}
 \E\left[e^{-\be\mathcal Q_C/2}\right]
 =\prod_j e^{\be\lambda_j(C)/2}(1+\be\lambda_j(C))^{-1/2}
 =\detTwo(\mathrm{Id}+\be C)^{-1/2}.
\end{equation}
The logarithms of the factors are $O(\lambda_j(C)^2)$ near zero;
\eqref{eq:signed-scalar-spectral-gap} controls the negative spectrum.
Thus, the product is finite and positive without a trace-class assumption.

\step{Truncating the scalar exponential.}
Choose a fixed smooth $\chi$ with $0\leq\chi\leq1$, equal to zero on
$(-\infty,-1]$ and to one on $[0,\infty)$.  For $L\geq0$, set
\begin{equation}\label{eq:signed-exponential-test-cutoff}
 h_L(x)\coloneqq e^{-\be x}\chi(x+L).
\end{equation}
Since $h_L(x)=e^{\be L}h_0(x+L)$ and $h_0$ has bounded derivatives through
order three,
\begin{equation}\label{eq:signed-cutoff-derivatives}
 \norm{h_L'}_\infty+[h_L'']_{C^{0,\alpha}}
 \leq C_{\alpha,\be}e^{\be L}.
\end{equation}
The error is bounded pointwise by
\begin{equation}\label{eq:signed-cutoff-tail}
 0\leq e^{-\be x}-h_L(x)
 \leq e^{-\be x}\one_{\{x<-L\}}
 \leq e^{-(q-1)\be L}e^{-q\be x}.
\end{equation}
Use \Cref{lem:quadratic-smooth-comparison} for $h_L$ and
\eqref{eq:signed-scalar-exponential-assumption},
\eqref{eq:signed-Gaussian-exponential-moment} for the two tails.  The total
error is at most
\begin{equation}\label{eq:signed-Laplace-preoptimization}
 C_{\alpha,\be}e^{\be L}
 \left(N^{-\alpha/2}\norm c_{2+\alpha}^{2+\alpha}
                +N^{-1/2}\norm c_2\right)
 +2M e^{-(q-1)\be L}.
\end{equation}
Taking $L=\alpha\log N/(2q\be)$ proves
\eqref{eq:signed-scalar-Laplace-bound}; the finite-Gaussian error is no
larger since $\alpha\leq1$.  The exponential moment has been used only at
the original and limiting laws, not during coordinate replacement.
\end{proof}

\subsection{Common inputs for the quantitative proof}
\label{subsubsec:exact-profile-partition-inputs}

From this subsection through \Cref{subsec:strong-pair-estimate}, assume the
hypotheses of \Cref{thm:abstract-quantitative-marginal}, with its fixed
exponents $p,\vartheta,\alpha$.  
We use the centered normalization
\eqref{eq:exact-centered-partition}--\eqref{eq:exact-centered-density}, the
density moments \eqref{eq:exact-full-density-Lr}, and the resolvent bound
\eqref{eq:signed-resolvent-bound}.  Define
\begin{equation}\label{eq:weak-second-model-v-def}
 v_{N,\be}(x,y)\coloneqq f_{N,\be}^{(2)}(x,y)
       -f_{N,\be}^{(1)}(x)-f_{N,\be}^{(1)}(y)+1,
\end{equation}
and abbreviate it by $v_N$.  As before, $u_N=f_{N,\be}^{(1)}-1$.
The first-component estimate is \eqref{eq:signed-onepoint-bound}.
\Cref{lem:hoeffding-reduction} and the density moment at exponent two give
\begin{equation}\label{eq:exact-level-control}
 \sum_{r=1}^N\binom Nr\norm{h_{[r]}}_{L^2(\mu^{\otimes r})}^2
 =\norm{f_{N,\be}-1}_{L^2(\mu^{\otimes N})}^2\leq C,
\end{equation}
where $h_{[1]}=u_N$ and $h_{[2]}=v_N$.
The scalar estimates below have rate
\begin{equation}\label{eq:quantitative-partition-error}
 \mathfrak e_N\coloneqq N^{-3\alpha/8}.
\end{equation}
This rate is distinct from the marginal rate
\eqref{eq:signed-profile-rate}.

\begin{lemma}[Insertion of a centered one-body term]
\label{lem:signed-source-insertion}
Let $A,V$ be real random variables with $V\leq C_0$,
$\E[V^2]\leq C_1/N$, and $\E[e^{2A}]\leq C_2$.  Then
\begin{equation}\label{eq:signed-source-insertion}
 \left|\E[e^{A+V}]-\E[e^A]\right|\leq C N^{-1/2},
\end{equation}
where $C$ depends only on $C_0,C_1,C_2$.
\end{lemma}
\begin{proof}
On the half-line $v\leq C_0$, one has $\abs{e^v-1}\leq C_{C_0}\abs v$.
Multiply this inequality by $e^A$ and apply Cauchy--Schwarz.
\end{proof}

\subsection{Correlated replicas and noise stability}
\label{subsubsec:exact-profile-strong-pair}

% STAGE2 MOVE: A3-replica-setup-role
Under the resampling coupling below, the squared norms of the Hoeffding components are the coefficients of a power series with nonnegative coefficients.  The replica representation expresses this series through partition functions.  A uniform estimate for the corresponding determinant will allow us to extract its degree-two coefficient.

% STAGE2 MOVE: A3-replica-coupling
For $0\leq t<1$, define the one-site coupling
\begin{equation}\label{eq:exact-nu-t}
 \nu_t(\md x\,\md y)
 \coloneqq
 (1-t)\mu(\md x)\mu(\md y)
 +t\mu(\md x)\delta_x(\md y).
\end{equation}
Both marginals are $\mu$.  Conditional on $X=x$, the variable $Y$ is retained with probability $t$ and otherwise resampled independently from $\mu$.

On $E=\R^{\d}\times\R^{\d}$, define
\begin{align}
 K_t\big((x,y),(x',y')\big)
 &\coloneqq\tl W(x,x')+\tl W(y,y'),\label{eq:exact-paired-kernel}\\
 \psi_t(x,y)&\coloneqq\psi_\mu(x)+\psi_\mu(y).\label{eq:exact-paired-source}
\end{align}
The kernel $K_t$ is symmetric, canonical with respect to $\nu_t$, and belongs to $L^2(\nu_t^{\otimes2})$.
For any canonical $r\in L^2(\mu^{\otimes2})$ and independent sites $(X,Y),(X',Y')$ with law $\nu_t$,
\begin{align}
 \E[r(Y,Y')\mid X,X']&=t^2r(X,X'),\label{eq:quantitative-paired-conditional}\\
 \norm{r(x,x')+r(y,y')}_{L^2(\nu_t^{\otimes2})}^2
 &=2(1+t^2)\norm r_2^2.\label{eq:quantitative-paired-tail-norm}
\end{align}
Every term with a resampled coordinate vanishes by canonicity. Thus, the coupled Hilbert--Schmidt estimate uses independence between sites, without requiring independence within a site.

% STAGE2 MOVE: A3-paired-spectrum-overlap-and-proposition
Let $\mathcal T_t$ be the integral operator with kernel $K_t$ on $L^2(\nu_t)$.

\begin{lemma}[Paired spectrum]\label{lem:exact-paired-spectrum}
If $(\lambda_j,e_j)$ are the nonzero spectral data of $T_\mu$, then the nonzero eigenvalues of $\mathcal T_t$ are
\begin{equation}\label{eq:exact-paired-spectrum}
 (1+t)\lambda_j
 \quad\text{and}\quad
 (1-t)\lambda_j,
 \qquad j\geq1.
\end{equation}
\end{lemma}

\begin{proof}
For centered $f,g\in L^2(\mu)$,
\begin{equation}\label{eq:exact-nu-cross}
 \int f(x)g(y)\,\md\nu_t(x,y)
 =t\int fg\,\md\mu.
\end{equation}
A direct conditional-expectation calculation gives
\begin{align}
 \mathcal T_t[e_j(x)]
 &=\lambda_je_j(x)+t\lambda_je_j(y),\label{eq:exact-Tt-ejx}\\
 \mathcal T_t[e_j(y)]
 &=t\lambda_je_j(x)+\lambda_je_j(y).
 \label{eq:exact-Tt-ejy}
\end{align}
Thus, $e_j(x)+e_j(y)$ and $e_j(x)-e_j(y)$ have the eigenvalues in \eqref{eq:exact-paired-spectrum}.  Their squared norms are $2(1+t)$ and $2(1-t)$, respectively.

For centered $f,g$, the additive map $J_t(f,g)=f(x)+g(y)$ satisfies
\begin{equation}\label{eq:quantitative-additive-equivalence}
 (1-t)(\norm f_2^2+\norm g_2^2)\leq\norm{J_t(f,g)}_{L^2(\nu_t)}^2
 \leq(1+t)(\norm f_2^2+\norm g_2^2).
\end{equation}
Its range is therefore closed for $t<1$.  The range of $\mathcal T_t$ is contained in this additive centered subspace, so every nonzero eigenfunction is additive.  The orthogonal complement is annihilated by self-adjointness, and the remaining additive directions arise from the kernel of $T_\mu$ and are also annihilated.  On the additive centered subspace, $\mathcal T_t$ is represented by
\begin{equation}\label{eq:exact-additive-matrix}
 (f,g)\longmapsto
 \big(T_\mu(f+tg),T_\mu(tf+g)\big).
\end{equation}
Decomposing $f,g$ into the spectral subspaces of $T_\mu$ shows that the displayed symmetric and antisymmetric eigenfunctions exhaust the nonzero spectrum.
\end{proof}


Define
\begin{align}
 \K^{\mathrm c,[2]}_{N,\be}(t)
 &\coloneqq\E_{\nu_t^{\otimes N}}\left[
 \exp\left\{-\frac{\be}{2N}\sum_{i\ne j}K_t(Z_i,Z_j)
 +\frac{\be}{N}\sum_{i=1}^N\psi_t(Z_i)\right\}\right].
 \label{eq:exact-paired-partition}
\end{align}

Let $\mathsf T_t$ be the product-resampling operator on $L^2(\mu^{\otimes N})$.  On a Hoeffding component indexed by $A\subset[N]$,
\begin{equation}\label{eq:exact-noise-eigenvalue}
 \mathsf T_t h_A=t^{\abs A}h_A.
\end{equation}
Thus, the overlap and its coefficients satisfy
\begin{align}
 A_N(t)&\coloneqq\ipp*{f_{N,\be},\mathsf T_tf_{N,\be}}-1
 =\sum_{r=1}^N c_{N,r}t^r,
 \qquad c_{N,r}\coloneqq\binom Nr\norm{h_{[r]}}_2^2\geq0,
 \label{eq:exact-noise-series}\\
 1+A_N(t)&=\frac{\K^{\mathrm c,[2]}_{N,\be}(t)}{(\K^{\mathrm c}_{N,\be})^2}.
 \label{eq:exact-overlap-ratio}
\end{align}
The coefficient sum is uniformly bounded by \eqref{eq:exact-level-control}.
The eigenvalues of $B_{\be,\mu}$ are $-\be\lambda/(1+\be\lambda)$.
By \eqref{eq:signed-resolvent-bound}, a negative $\lambda$ gives a value
in $[0,1/11)$, while a nonnegative $\lambda$ gives a value in $(-1,0]$.
Since $T_\mu$ is bounded, these inequalities imply
$\norm{B_{\be,\mu}}_{\mathrm{op}}<1$.  Put

\begin{equation}\label{eq:exact-overlap-limit}
 D(t)\coloneqq\det(\mathrm{Id}-t^2B_{\be,\mu}^2)^{-1/2}.
\end{equation}
Only the square of the Hilbert--Schmidt operator occurs in this ordinary Fredholm determinant.

\begin{prop}[Quantitative correlated-replica estimate]\label{prop:exact-paired-determinant}
Fix $0<t_0<1$.  Uniformly for $N\geq2$,
\begin{align}
 &\sup_{0\leq t\leq t_0}\left|\K^{\mathrm c,[2]}_{N,\be}(t)
 -\left[\detTwo(\mathrm{Id}+\be(1+t)T_\mu)
 \detTwo(\mathrm{Id}+\be(1-t)T_\mu)\right]^{-1/2}\right|
 \notag\\
 &\qquad\leq C\mathfrak e_N,
 \label{eq:exact-paired-determinant-limit}\\
 &\sup_{0\leq t\leq t_0}|A_N(t)-(D(t)-1)|\leq C\mathfrak e_N.
 \label{eq:quantitative-replica-rate}
\end{align}
\end{prop}
\begin{proof}
\step{The paired Laplace transform.}
By the marginal laws of $\nu_t$ and the triangle inequality,
\begin{equation}\label{eq:signed-paired-Lp}
 \norm{K_t}_{L^{2+\alpha}(\nu_t^{\otimes2})}
 \leq2\norm{\tl W}_{L^{2+\alpha}(\mu^{\otimes2})}.
\end{equation}
For the fourth exponential moment, Cauchy--Schwarz under $\nu_t^{\otimes N}$
gives
\begin{align}
 \E_{\nu_t^{\otimes N}}\left[e^{-2\be\mathsf W_N(K_t)}\right]
 &\leq\E_{\mu^{\otimes N}}
       \left[e^{-4\be\mathsf W_N(\tl W)}\right]\leq C.
 \label{eq:signed-paired-fourth-moment}
\end{align}
The last estimate uses \eqref{eq:signed-centered-pair-moment} at inverse
temperature $8\be\leq\Gamma$.  It uses independence between paired sites,
not independence of the coordinates within a site.  Apply
\Cref{thm:quantitative-signed-Laplace} with $q=4$ on $(E,\nu_t)$.  Its
constant is uniform for $0\leq t\leq t_0$ by
\eqref{eq:signed-paired-Lp}--\eqref{eq:signed-paired-fourth-moment}.  The
resulting error is $C\mathfrak e_N$.

To restore the source in \eqref{eq:exact-paired-partition}, use
\Cref{lem:signed-source-insertion} with
$V=\be N^{-1}\sum_i\psi_t(Z_i)$.  This variable is bounded above by
$2\be C_+$, and \eqref{eq:exact-nu-cross} gives
\begin{equation}\label{eq:signed-paired-source-variance}
 \E_{\nu_t^{\otimes N}}[V^2]
 =\frac{2\be^2(1+t)}N\norm{\psi_\mu}_2^2\leq C/N.
\end{equation}
The required second exponential moment follows from
\eqref{eq:signed-paired-fourth-moment}.  Since $N^{-1/2}\leq\mathfrak e_N$,
source insertion does not change the rate.
\Cref{lem:exact-paired-spectrum} identifies the limiting determinant and
proves \eqref{eq:exact-paired-determinant-limit}.  Its additive-space
argument is uniform on the fixed interval because $1-t\geq1-t_0>0$.

\step{The overlap ratio.}
Denote the displayed paired determinant by
$\K^{\mathrm c,[2]}_{\infty,\be}(t)$ and put
$\K^{\mathrm c}_{\infty,\be}=\detTwo(\mathrm{Id}+\be T_\mu)^{-1/2}$.
At $t=0$ the two coordinates are independent, so
\begin{equation}\label{quant-rep:eq-independent-endpoint}
 \K^{\mathrm c,[2]}_{N,\be}(0)=(\K^{\mathrm c}_{N,\be})^2,
 \qquad
 \K^{\mathrm c,[2]}_{\infty,\be}(0)=(\K^{\mathrm c}_{\infty,\be})^2.
\end{equation}
Thus, the first estimate also controls the squared denominator with error
$C\mathfrak e_N$.  Both denominators are at least one, by Jensen for the
centered finite-particle and Gaussian exponents.  The paired limiting
numerator is uniformly bounded by the Gaussian fourth-moment transfer
in \Cref{thm:quantitative-signed-Laplace}.  Consequently,
\begin{equation}\label{quant-rep:eq-stable-ratio}
 \sup_{0\leq t\leq t_0}
 \left|1+A_N(t)-
 \frac{\K^{\mathrm c,[2]}_{\infty,\be}(t)}
      {(\K^{\mathrm c}_{\infty,\be})^2}\right|
 \leq C\mathfrak e_N.
\end{equation}
For every spectral value $\lambda$ of $T_\mu$, the regularization factors
cancel and the remaining ratio is
\begin{equation}\label{quant-rep:eq-mode-ratio}
 \frac{1+\be\lambda}
 {\sqrt{(1+\be(1+t)\lambda)(1+\be(1-t)\lambda)}}
 =\left(1-t^2\left(\frac{\be\lambda}{1+\be\lambda}\right)^2\right)^{-1/2}.
\end{equation}
All square roots are positive by spectral stability.  Since
$B_{\be,\mu}^2$ is trace class and $\norm{B_{\be,\mu}}_{\mathrm{op}}<1$,
the product converges uniformly on this interval.  It gives
\begin{equation}\label{quant-rep:eq-final-det-ratio}
 \frac{\K^{\mathrm c,[2]}_{\infty,\be}(t)}
      {(\K^{\mathrm c}_{\infty,\be})^2}
 =\det(\mathrm{Id}-t^2B_{\be,\mu}^2)^{-1/2}=D(t).
\end{equation}
Together with \eqref{quant-rep:eq-stable-ratio}, this proves
\eqref{eq:quantitative-replica-rate} for every fixed $t_0<1$.
\end{proof}

\subsection{Quantitative screened response}
\label{subsec:screened-response-proof}

To estimate the pairing with $b_{\be,\mu}$, perturb the canonical pair
interaction in that direction and leave the one-body source fixed.  Set
\begin{align}
 A_\theta&\coloneqq\be T_\mu+\theta B_{\be,\mu},\qquad
 K_\theta(x,y)\coloneqq\be\tl W(x,y)+\theta b_{\be,\mu}(x,y),
 \label{eq:quantitative-response-operator}\\
 \K^{\mathrm c}_{N,\be}(\theta)
 &\coloneqq\E_{\mu^{\otimes N}}\left[
 \exp\left\{-\frac1{2N}\sum_{i\ne j}K_\theta(x_i,x_j)
                  +\frac\be N\sum_{i=1}^N\psi_\mu(x_i)\right\}\right],
 \label{eq:quantitative-response-partition}\\
 \Lambda_N(\theta)&\coloneqq\log\K^{\mathrm c}_{N,\be}(\theta),\qquad
 \Lambda(\theta)\coloneqq-\frac12\log\detTwo(\mathrm{Id}+A_\theta).
 \label{eq:quantitative-response-logpartition}
\end{align}
In particular, $\K^{\mathrm c}_{N,\be}(0)=\K^{\mathrm c}_{N,\be}$.

\begin{prop}[Quantitative screened response]
\label{prop:quantitative-screened-response}
There is $\theta_0>0$, depending only on the fixed hypothesis data, such that
\begin{equation}\label{eq:quantitative-response-rate}
 \sup_{\abs\theta\leq\theta_0}
       \abs{\Lambda_N(\theta)-\Lambda(\theta)}
 \leq C\mathfrak e_N,\qquad N\geq2.
\end{equation}
Both functions have uniformly bounded third derivatives on
$[-\theta_0/2,\theta_0/2]$.
\end{prop}

\begin{proof}
\step{The bounded correction and exponential moments.}
Write $T=T_\mu$ and $R=R_{\be,\mu}$ inside operator identities.  By
\eqref{eq:signed-resolvent-bound}, $0<R\leq(12/11)\mathrm{Id}$.  Therefore,
\begin{equation}\label{quant-resp:eq-Ckernel}
 \mathcal C\coloneqq\be^2TRT\geq0,\qquad
 c_{\mathrm{corr}}(x,y)\coloneqq\be^2\ipp*{w_x,Rw_y}
\end{equation}
has the centered feature $\be R^{1/2}w_x$.  With
$M_2=\operatorname*{ess\,sup}_x\norm{w_x}_2\leq M_p$, its squared norm is
bounded by
\begin{equation}\label{eq:signed-correction-bound}
 K\coloneqq\be^2\norm R_{\mathrm{op}}M_2^2,
 \qquad \abs{c_{\mathrm{corr}}(x,y)}\leq K.
\end{equation}
Strong measurability and centering follow as in
\Cref{lem:signed-cavity-moments}.  The resolvent identity gives
\begin{equation}\label{quant-resp:eq-Bdecomposition}
 B_{\be,\mu}=-\be T+\mathcal C,
 \qquad K_\theta=\be(1-\theta)\tl W+\theta c_{\mathrm{corr}}.
\end{equation}
Choose
\begin{equation}\label{quant-resp:eq-theta0}
 \theta_0\coloneqq\frac1{64(1+K)}.
\end{equation}
Cauchy--Schwarz gives
\begin{align}
 \E\left[e^{-2\mathsf W_N(K_\theta)}\right]
 &\leq\E\left[e^{-4\be(1-\theta)\mathsf W_N(\tl W)}\right]^{1/2}
       \E\left[e^{-4\theta\mathsf W_N(c_{\mathrm{corr}})}\right]^{1/2}.
 \label{quant-resp:eq-responseholder}
\end{align}
The first factor is controlled at inverse temperature
$8\be(1-\theta)\leq12\be\leq\Gamma$.
For $\theta\geq0$, the feature identity gives
$\mathsf W_N(c_{\mathrm{corr}})\geq-K$, so the second factor is bounded.
For $\theta<0$, apply \Cref{quant-resp:lem-boundedfeatureexp} with
$a=8\abs\theta$:
\begin{equation}\label{quant-resp:eq-attractivecorrection}
 \E\left[e^{4\abs\theta\mathsf W_N(c_{\mathrm{corr}})}\right]
 \leq(1-32\abs\theta K)^{-1/2}\leq\sqrt2.
\end{equation}
It follows that
\begin{equation}\label{quant-resp:eq-responseexp}
 \sup_{N\geq2}\sup_{\abs\theta\leq\theta_0}
       \E\left[e^{-2\mathsf W_N(K_\theta)}\right]<\infty.
\end{equation}
The same conclusion at every exponent between zero and four follows by
H\"older; convex interpolations of the kernels also preserve the bound.
Positivity has been used only for $\mathcal C$, not for $T$.

\step{The signed Laplace estimate and source insertion.}
The kernels $K_\theta$ are canonical and uniformly bounded in
$L^{2+\alpha}(\mu^{\otimes2})$ by
\eqref{quant-resp:eq-Bdecomposition} and
\eqref{eq:signed-correction-bound}.  Apply
\Cref{thm:quantitative-signed-Laplace} with target inverse temperature one
and $q=4$.  Equations \eqref{quant-resp:eq-responseexp} and
\eqref{eq:signed-scalar-Laplace-bound} give a uniform error
$C\mathfrak e_N$, and also
\begin{equation}\label{eq:signed-response-resolvent}
 \mathrm{Id}+A_\theta\geq\tfrac34\mathrm{Id},
 \qquad\abs\theta\leq\theta_0.
\end{equation}
The one-body average $\be N^{-1}\sum_i\psi_\mu(x_i)$ is bounded above,
is centered, and has variance $\be^2\norm{\psi_\mu}_2^2/N$.
\Cref{lem:signed-source-insertion} restores it at cost $CN^{-1/2}$.
The finite-particle partition function and limiting Gaussian expectation
are at least one by Jensen.  Since the logarithm is $1$-Lipschitz on
$[1,\infty)$, this proves \eqref{eq:quantitative-response-rate}.

\step{Third derivatives.}
Write the exponent as $H_N+\theta Y$, where
$Y=-\mathsf W_N(b_{\be,\mu})/2$ and $H_N$ is the exponent at $\theta=0$.
For every fixed integer $j\geq0$,
\begin{equation}\label{eq:signed-response-derivative-domination}
 \abs y^j e^{\theta y}
 \leq C_{j,\theta_0}(e^{\theta_0y}+e^{-\theta_0y}),
 \qquad \abs\theta\leq\theta_0/2.
\end{equation}
After multiplication by $e^{H_N}$, the right-hand side has uniformly
bounded integral by the endpoint estimates above and the source upper
bound.  Dominated convergence justifies differentiation.  If
$\E_{N,\theta}$ denotes expectation under this normalized tilt, then
\begin{equation}\label{eq:signed-response-third-derivative}
 \Lambda_N'''(\theta)
 =\E_{N,\theta}\left[(Y-\E_{N,\theta}[Y])^3\right],
 \qquad\abs{\Lambda_N'''(\theta)}\leq8\E_{N,\theta}[\abs Y^3]\leq C.
\end{equation}
Here, the denominator is at least one.  This argument does not require
$b_{\be,\mu}\in L^3(\mu^{\otimes2})$.  Spectral calculus, commutation of
$A_\theta$ and $B_{\be,\mu}$, and
\eqref{eq:signed-response-resolvent} give
\begin{equation}\label{eq:signed-limit-third-derivative}
 \Lambda'''(\theta)
 =-\Tr\left[((\mathrm{Id}+A_\theta)^{-1}B_{\be,\mu})^3\right],
 \quad
 \abs{\Lambda'''(\theta)}
 \leq(4/3)^3\norm{B_{\be,\mu}}_{\mathrm{op}}
                \norm{B_{\be,\mu}}_{\mathrm{HS}}^2.
\end{equation}
The trace is absolutely convergent.  This proves the derivative assertions.
\end{proof}

\subsection{Strong convergence of the pair kernel}
\label{subsec:strong-pair-estimate}

\begin{prop}[Quantitative strong pair limit]
\label{prop:exact-strong-second-hoeffding}
Uniformly for $N\geq2$,
\begin{align}
 \norm{Nv_N-b_{\be,\mu}}_2&\leq C N^{-\alpha/8},
 \label{eq:exact-strong-pair-limit}\\
 \left|\norm{Nv_N}_2^2-\norm{b_{\be,\mu}}_2^2\right|
 &\leq C N^{-\alpha/4},
 \label{eq:quantitative-pair-squared-norm}\\
 \left|\ipp*{Nv_N,b_{\be,\mu}}-\norm{b_{\be,\mu}}_2^2\right|
 &\leq C N^{-\alpha/4}.
 \label{eq:quantitative-pair-cross-term}
\end{align}
\end{prop}
\begin{proof}
\step{Squared norm from correlated replicas.}
The trace expansion near zero gives
\begin{equation}\label{eq:exact-small-t-logdet}
 \log D(t)=\frac12\sum_{m\geq1}\frac{t^{2m}}m\Tr(B_{\be,\mu}^{2m}),
 \qquad
 D(t)=1+\frac{t^2}{2}\norm{B_{\be,\mu}}_{\mathrm{HS}}^2+O(t^4).
\end{equation}
It converges uniformly near $[0,1/2]$ because $B_{\be,\mu}$ is
Hilbert--Schmidt and $\norm{B_{\be,\mu}}_{\mathrm{op}}<1$.
In particular, the sixth derivative of $D$ is bounded on that interval.
The nonnegative coefficients in \eqref{eq:exact-noise-series} and
\eqref{eq:exact-level-control} also give
\begin{equation}\label{eq:signed-replica-sixth-derivative}
 \sup_{0\leq t\leq1/2}\abs{A_N^{(6)}(t)}
 \leq C\sup_{r\geq6}r(r-1)\cdots(r-5)2^{6-r}<\infty,
\end{equation}
uniformly in $N$.  For $5h\leq1/2$, set
\begin{equation}\label{eq:quantitative-replica-stencil}
 \mathcal D_hF\coloneqq
 \frac{45F(0)-154F(h)+214F(2h)-156F(3h)+61F(4h)-10F(5h)}{24h^2}.
\end{equation}
Direct substitution shows that this expression equals one on $t^2$ and
zero on $1,t,t^3,t^4,t^5$.  Taylor's formula through degree five therefore
gives an error $Ch^4$ in extracting the quadratic coefficient.  Applying
this to $A_N$ and $D-1$, and using
\eqref{eq:quantitative-replica-rate} at $t_0=1/2$, gives
\begin{equation}\label{eq:quantitative-stencil-error}
 \left|c_{N,2}-\tfrac12\norm{B_{\be,\mu}}_{\mathrm{HS}}^2\right|
 \leq C\left(h^4+\mathfrak e_N/h^2\right).
\end{equation}
Choose $h=\mathfrak e_N^{1/6}$ for all sufficiently large $N$.  The exact
normalization is
\begin{equation}\label{eq:quantitative-pair-coefficient}
 \norm{Nv_N}_2^2=\frac{2N}{N-1}c_{N,2}.
\end{equation}
Thus, its error is $C(\mathfrak e_N^{2/3}+N^{-1})\leq CN^{-\alpha/4}$,
proving \eqref{eq:quantitative-pair-squared-norm}.

\step{Cross term from the screened response.}
Differentiation is justified by
\eqref{eq:signed-response-derivative-domination}.  Canonicity gives
\begin{equation}\label{eq:quantitative-response-finite-derivative}
 \Lambda_N'(0)
 =-\frac1{2N}\E_{\Q_{N,\be}(\mu)}
       \left[\sum_{i\ne j}b_{\be,\mu}(x_i,x_j)\right]
 =-\frac{N-1}{2}\ipp*{v_N,b_{\be,\mu}}.
\end{equation}
The product and singleton terms vanish.  Spectral differentiation gives
\begin{equation}\label{eq:quantitative-response-limit-derivative}
 \Lambda'(0)=-\frac12\Tr[(R_{\be,\mu}-\mathrm{Id})B_{\be,\mu}]
 =-\frac12\norm{b_{\be,\mu}}_2^2.
\end{equation}
The third-derivative bounds in \Cref{prop:quantitative-screened-response}
imply, for $F=\Lambda_N$ and $F=\Lambda$,
\begin{equation}\label{eq:signed-centered-response-difference}
 \left|\frac{F(h)-F(-h)}{2h}-F'(0)\right|\leq Ch^2,
 \qquad 0<h\leq\theta_0/2.
\end{equation}
Together with \eqref{eq:quantitative-response-rate}, this bounds the
difference of derivatives by $C(h^2+\mathfrak e_N/h)$.  Choose
$h=\mathfrak e_N^{1/3}$ for large $N$.  Converting from $(N-1)v_N$ to
$Nv_N$ in \eqref{eq:quantitative-response-finite-derivative} costs
$O(N^{-1})$, so \eqref{eq:quantitative-pair-cross-term} follows.
The finite initial range of $N$ in both arguments is covered by
\eqref{eq:exact-level-control}, with an enlarged constant.

Finally, polarization gives
\begin{equation}\label{eq:quantitative-pair-polarization}
 \norm{Nv_N-b_{\be,\mu}}_2^2
 =\norm{Nv_N}_2^2+\norm{b_{\be,\mu}}_2^2
             -2\ipp*{Nv_N,b_{\be,\mu}}\leq C N^{-\alpha/4}.
\end{equation}
Taking the square root proves \eqref{eq:exact-strong-pair-limit}.
\end{proof}

\section{Singleton profiles and statistical divergences}
\label{sec:singleton-quantitative-limits}

% STAGE2 MOVE: A2-singleton-and-assembly-roadmap
Assume the hypotheses of \Cref{thm:abstract-quantitative-marginal} and
use $u_N,v_N$ from \Cref{subsubsec:exact-profile-partition-inputs}.
The common second-order identity and cavity variance expansion below,
together with \Cref{lem:signed-cavity-moments}, identify the singleton
profile.  We then assemble the two kernels and prove the marginal and
divergence conclusions of the abstract theorem.  The model hypotheses are verified afterward in \Cref{subsubsec:signed-model-verification}.
% STAGE2 MOVE: A2-common-normalized-exponential-identity
\subsection{The second-order normalized-exponential identity}
\label{subsubsec:exact-profile-normalized-exponential}
To identify the singleton profile, we expand the normalized one-particle conditional density to second order.  Throughout this cavity argument, expectations without a subscript are under $\Q_{N,\be}(\mu)$.  We use $B_N$, $\pi_N$, $r_N$, $\epsilon$, and $\delta_B$ from \eqref{eq:cavity-B-def}, \eqref{eq:conditional-density}, \eqref{eq:cavity-error-def}, \eqref{eq:epsilon-def}, and \eqref{eq:deltaB-def}.  Each cavity field is centered, so $\delta_B\geq0$.

For a fixed reservoir, abbreviate $B=B_N(\cdot;\hat\Xi_1)$ and set
\begin{align}
 m_2(B)&\coloneqq\int B^2\,\md\mu,\label{eq:exact-m2-def}\\
 \rho(t)&\coloneqq e^t-1-t-\frac12t^2,\label{eq:exact-rho-def}\\
 \bar\rho(B)&\coloneqq\int\rho(B)\,\md\mu.\label{eq:exact-rho-bar-def}
\end{align}
Then
\begin{equation}\label{eq:exact-delta-second}
 \delta_B=\frac12m_2(B)+\bar\rho(B).
\end{equation}
Substituting this identity into \eqref{eq:normalized-exponential-remainder} gives
\begin{multline}\label{eq:exact-second-exact-identity}
 \pi_N(x\mid\hat\Xi_1)-1-B(x)
 -\frac12\left(B(x)^2-m_2(B)\right)\\
 =\frac{
 \rho(B(x))-\bar\rho(B)
 -B(x)\delta_B
 -\frac12\delta_B\left(B(x)^2-m_2(B)\right)
 }{1+\delta_B}.
\end{multline}
No asymptotic approximation has yet been made.

Define
\begin{equation}\label{eq:exact-qN-def}
 q_N(x)\coloneqq\E\left[B_N(x;\hat\Xi_1)^2\right],
 \qquad
 \bar q_N\coloneqq\int q_N\,\md\mu.
\end{equation}

% Exact-profile integration module: Riesz scalar and Hilbert cavity moments

% STAGE2 MOVE: subsec:common-cavity-variance
\subsection{The common cavity variance identity}
\label{subsec:common-cavity-variance}

% STAGE2 MOVE: A2-common-variance-integrability
The quantitative row assumption gives $M_2=\operatorname*{ess\,sup}_x\norm{w_x}_2\leq M_p<\infty$.  The density moments in \eqref{eq:exact-full-density-Lr} and \Cref{lem:signed-cavity-moments} justify the expectations below for $\mu$-almost every $x$.  In particular, the row-square term is integrable by H\"older at the available exponent $p/(p-2)$, and the distinct-pair term by the $L^2$ row and density bounds.


% STAGE2 MOVE: A2-common-variance-map
Define the bounded linear map
\begin{equation}\label{eq:exact-M-map}
 (\mathcal M g)(x)
 \coloneqq
 \iint w_x(y)w_x(z)g(y,z)\,\md\mu(y)\md\mu(z),
 \qquad g\in L^2(\mu^{\otimes2}).
\end{equation}
By the uniform row bound,
\begin{equation}\label{eq:exact-riesz-M-bounded}
 \norm{\mathcal M g}_{L^2(\mu)}
 \leq M_2^2\norm{g}_{L^2(\mu^{\otimes2})}.
\end{equation}

% STAGE2 MOVE: A2-common-variance-expansion
Let $S_x=\sum_{j=2}^Nw_x(\xi_j)$.  Expanding \eqref{eq:cavity-B-def},
\begin{equation}\label{eq:exact-qN-first-expansion}
 Nq_N(x)
 =\frac{\be^2}{N}
 \left[
 \psi_\mu(x)^2
 -2\psi_\mu(x)\E[S_x]
 +\E[S_x^2]
 \right].
\end{equation}
Exchangeability and canonicity give
\begin{equation}\label{eq:exact-S-first-moment}
 \E[S_x]=(N-1)T_\mu u_N(x).
\end{equation}
For the second moment,
\begin{align}
 \E[S_x^2]
 &=(N-1)\int w_x(y)^2\big(1+u_N(y)\big)\,\md\mu(y)
 \notag\\
 &\quad+(N-1)(N-2)
 \iint w_x(y)w_x(z)v_N(y,z)\,\md\mu(y)\md\mu(z).
 \label{eq:exact-S-second-moment}
\end{align}
Indeed, in the distinct-label contribution, the product term and the two first-Hoeffding terms vanish because $w_x$ is centered.

Substituting \eqref{eq:exact-S-first-moment}--\eqref{eq:exact-S-second-moment} into \eqref{eq:exact-qN-first-expansion},
\begin{align}
 Nq_N
 &=\be^2\Bigg[
 \frac{\psi_\mu^2}{N}
 -2\frac{N-1}{N}\psi_\mu T_\mu u_N
 +\frac{N-1}{N}\int w_{\cdot}(y)^2\,\md\mu(y)
 \notag\\
 &\qquad
 +\frac{N-1}{N}\int w_{\cdot}(y)^2u_N(y)\,\md\mu(y)
 +\frac{(N-1)(N-2)}{N^2}\mathcal M(Nv_N)
 \Bigg].
 \label{eq:exact-qN-complete-expansion}
\end{align}

\subsection{The second-order remainder and singleton profile}
\label{subsec:signed-second-cavity}

\begin{lemma}[Second-order cavity remainder]
\label{lem:signed-second-cavity-reduction}
Under the hypotheses of \Cref{thm:abstract-quantitative-marginal},
\begin{equation}\label{eq:signed-second-cavity-reduction}
 N\norm{r_N-\tfrac12(q_N-\bar q_N)}_{L^2(\mu)}
 \leq C N^{-\vartheta/2}.
\end{equation}
\end{lemma}
\begin{proof}
Use the exact identity \eqref{eq:exact-second-exact-identity}.  On the
upper-bounded range of $B_N$, $\abs{\rho(B_N)}\leq C\abs{B_N}^{2+\vartheta}$.
The pointwise moment estimate \eqref{eq:signed-point-cavity-moment} and
$\psi_\mu\in L^{4+2\vartheta}(\mu)$ give
\begin{equation}\label{eq:signed-rho-bound}
 \norm{\E[\abs{\rho(B_N(\cdot))}]}_2
       +\E[\abs{\bar\rho(B_N)}]\leq C N^{-1-\vartheta/2}.
\end{equation}
The source assumption supplies this moment because $\vartheta<1$.
By Minkowski and \eqref{eq:signed-delta-moment},
\begin{equation}\label{eq:signed-Bdelta-bound}
 \norm{\E[\abs{B_N}\delta_B]}_2
 \leq C\E[\norm{B_N}_2^3]\leq CN^{-3/2}.
\end{equation}
For the remaining mixed term, take the conjugate exponents
$q=(2+\vartheta)/\vartheta$ and $q'=(2+\vartheta)/2$.  Then
\begin{equation}\label{eq:signed-delta-square-bound}
 \E[\delta_B B_N(x)^2]
 \leq\E[\delta_B^q]^{1/q}
          \E[\abs{B_N(x)}^{2+\vartheta}]^{1/q'}
 \leq C\left(N^{-2}+N^{-3}\abs{\psi_\mu(x)}^2\right).
\end{equation}
This is $O(N^{-2})$ in $L^2(\mu)$, and
$\E[\delta_Bm_2(B_N)]\leq C\E[\norm{B_N}_2^4]\leq CN^{-2}$.
Taking absolute values in \eqref{eq:exact-second-exact-identity}, then
expectation and the spatial norm, and using $1+\delta_B\geq1$, proves
\eqref{eq:signed-second-cavity-reduction}.  No fourth pointwise moment
of the interaction is used.
\end{proof}

\begin{prop}[Quantitative singleton profile]
\label{prop:signed-strong-singleton}
Under the hypotheses of \Cref{thm:abstract-quantitative-marginal},
\begin{equation}\label{eq:signed-strong-singleton}
 \norm{Nu_N-a_{\be,\mu}}_{L^2(\mu)}\leq C\mathfrak r_N,
 \qquad N\geq2,
\end{equation}
with $\mathfrak r_N$ defined in \eqref{eq:signed-profile-rate}.
\end{prop}
\begin{proof}
We evaluate the common variance identity
\eqref{eq:exact-qN-complete-expansion}.  The term $\psi_\mu^2/N$ is
$O(N^{-1})$ in $L^2$ because $\psi_\mu\in L^4$.
The next term has the same bound, since
$\norm{T_\mu u_N}_\infty\leq M_2\norm{u_N}_2\leq C/N$ and
$\psi_\mu\in L^2$.
For the square-row term, take the same $q,q'$ as in
\eqref{eq:signed-delta-square-bound}.  Then $2q'=2+\vartheta<p$ and
$2q\leq r_*$.  H\"older and interpolation yield
\begin{align}
 \operatorname*{ess\,sup}_x
 \left|\int w_x(y)^2u_N(y)\,\md\mu(y)\right|
 &\leq M_p^2\norm{u_N}_q\notag\\
 &\leq M_p^2\norm{u_N}_2^{1/(q-1)}
             \norm{u_N}_{2q}^{1-1/(q-1)}
 \leq C N^{-\vartheta/2}.
 \label{eq:signed-square-row-error}
\end{align}
Here, \eqref{eq:exact-full-density-Lr} bounds the higher marginal norm,
\eqref{eq:signed-onepoint-bound} gives the $L^2$ bound, and
$1/(q-1)=\vartheta/2$.
The pair term is controlled by \eqref{eq:exact-riesz-M-bounded} and
\eqref{eq:exact-strong-pair-limit}.  Its limit satisfies
\begin{equation}\label{eq:signed-variance-identification}
 \norm{w_x}_2^2+(\mathcal M b_{\be,\mu})(x)
 =\ipp*{w_x,(\mathrm{Id}+B_{\be,\mu})w_x}
 =\ipp*{w_x,R_{\be,\mu}w_x}.
\end{equation}
All finite-particle coefficients differ from their limits by $O(N^{-1})$,
including the coefficient zero at $N=2$ for distinct reservoir labels.
Thus,
\begin{equation}\label{eq:signed-cavity-variance-limit}
 \norm{Nq_N-\omega_{\be,\mu}}_2
 \leq C\left(N^{-\alpha/8}+N^{-\vartheta/2}\right)=C\mathfrak r_N.
\end{equation}
Centering is an $L^2$ contraction, so
\Cref{lem:signed-second-cavity-reduction} identifies $Nr_N$ with
$\omega_{\be,\mu}^\circ/2$ with this error.
The resolvent identity and \eqref{eq:signed-resolvent-bound} give
\begin{equation}\label{eq:signed-finite-resolvent-difference}
 \left\|\left(\mathrm{Id}+\be\frac{N-1}{N}T_\mu\right)^{-1}
                  -R_{\be,\mu}\right\|_{\mathrm{op}}
 \leq\frac\be N\left(\frac{r_*}{r_*-1}\right)^2
                       \norm{T_\mu}_{\mathrm{op}}.
\end{equation}
Multiply \eqref{eq:screened-one-point-equation} by $N$, apply the bounded
inverse, and compare with \eqref{eq:exact-singleton-profile-def}.
Equations \eqref{eq:signed-cavity-variance-limit} and
\eqref{eq:signed-finite-resolvent-difference} prove
\eqref{eq:signed-strong-singleton}.
\end{proof}
\subsection{Uniform marginal assembly and divergence estimates}
\label{subsubsec:exact-profile-assembly}

% STAGE2 MOVE: A2-common-Hoeffding-assembly
\Cref{prop:signed-strong-singleton,prop:exact-strong-second-hoeffding}
give
\begin{equation}\label{eq:quantitative-kernel-rates}
 \norm{Nu_N-a_{\be,\mu}}_{L^2(\mu)}
 +\norm{Nv_N-b_{\be,\mu}}_{L^2(\mu^{\otimes2})}
 \leq C\mathfrak r_N.
\end{equation}
We also retain the sharper squared-norm estimate
\eqref{eq:quantitative-pair-squared-norm}.
For every $1\leq k\leq N$, marginalization of \eqref{eq:hoeffding-decomposition-general} gives
\begin{equation}\label{eq:exact-marginal-Hoeffding}
 f_{N,\be}^{(k)}-1
 =\sum_{i=1}^ku_N(x_i)+\sum_{1\leq i<j\leq k}v_N(x_i,x_j)
 +\mathcal R_{N,k}^{(\geq3)},
\end{equation}
where
\begin{equation}\label{eq:exact-higher-remainder}
 \mathcal R_{N,k}^{(\geq3)}\coloneqq
 \sum_{r=3}^k\sum_{\substack{A\subset[k]\\|A|=r}}h_A.
\end{equation}
An empty sum is zero.  For $3\leq r\leq k\leq N$,
\begin{equation}\label{eq:quantitative-binomial-ratio}
 \frac{\binom kr}{\binom Nr}
 =\prod_{j=0}^{r-1}\frac{k-j}{N-j}
 \leq\left(\frac kN\right)^r\leq\left(\frac kN\right)^3.
\end{equation}
Consequently, orthogonality and the total mass in \eqref{eq:exact-level-control} give
\begin{align}
 \norm{\mathcal R_{N,k}^{(\geq3)}}_2^2
 &=\sum_{r=3}^k\binom kr\norm{h_{[r]}}_2^2
 \leq\left(\frac kN\right)^3\sum_{r=3}^N\binom Nr\norm{h_{[r]}}_2^2
 \leq C\left(\frac kN\right)^3,
 \label{eq:exact-higher-tail-fixed}\\
 N^2\norm{\mathcal R_{N,k}^{(\geq3)}}_2^2&\leq C\frac{k^3}{N}.
 \label{eq:exact-higher-tail-scaled}
\end{align}
This estimate uses the total higher-order Hoeffding mass, so it remains valid for $k=N$; summing a separate bound for each level is unnecessary.

The profile norm is given by \eqref{eq:exact-Gk-L2}.
% STAGE2 MOVE: subsubsec:profile-moments
\subsubsection{Moments of the comparison profile}
\label{subsubsec:profile-moments}

% STAGE2 MOVE: A3-profile-moment-role
To pass from the density profile to relative entropy, we need a higher moment of $G_{k,\be,\mu}/k$ that is uniform in $k$.  The canonical pair sum is a martingale sum under the product law.  Applying \Cref{quant-tools:lem-martingale} to it gives the required moment bound without the extra factor of $k$ from the triangle inequality.

% STAGE2 MOVE: A3-profile-moments-and-gas-verification
\begin{lemma}[Canonical profile moments]\label{quant-tools:lem-profilemoments}
Let $a\in L^p(\mu)$ have mean zero and let $b\in L^p(\mu^{\otimes2})$ be symmetric and canonical, where $2\leq p<\infty$. For independent $X_1,\ldots,X_k$ with law $\mu$,
\begin{align}
 \left\|\sum_{i=1}^ka(X_i)\right\|_p&\leq\sqrt{(p-1)k}\norm a_p,\label{quant-tools:eq-singletonmoment}\\
 \left\|\sum_{i<j}b(X_i,X_j)\right\|_p&\leq(p-1)\sqrt{\binom k2}\norm b_p.\label{quant-tools:eq-pairmoment}
\end{align}
Consequently, for $G_k=\sum_i a(X_i)+\sum_{i<j}b(X_i,X_j)$,
\begin{equation}\label{quant-tools:eq-uniformprofilemoment}
 \left\|\frac{G_k}k\right\|_p
 \leq\sqrt{\frac{p-1}{k}}\norm a_p
 +(p-1)\sqrt{\frac{k-1}{2k}}\norm b_p.
\end{equation}
For the profiles in \Cref{thm:abstract-quantitative-marginal}, the normalized moment is uniformly finite at every exponent $2\leq q\leq\min\{p,6\}$, where $p$ is the row exponent of that theorem.

\end{lemma}
\begin{proof}
The first estimate follows immediately from \Cref{quant-tools:lem-martingale}. For the second, put $D_j=\sum_{i<j}b(X_i,X_j)$. The sequence is a martingale-difference sequence for $\sigma(X_1,\ldots,X_j)$ by canonicity in the second variable. To bound each $D_j$, reveal $X_j$ first and then $X_1,\ldots,X_{j-1}$. The summands $b(X_i,X_j)$ are martingale differences for this filtration by canonicity in the first variable. Thus, two applications of \eqref{quant-tools:eq-martingale} give
\begin{equation}\label{quant-tools:eq-twomartingales}
 \left\|\sum_{j=2}^kD_j\right\|_p^2
 \leq(p-1)\sum_{j=2}^k\norm{D_j}_p^2
 \leq(p-1)^2\sum_{j=2}^k(j-1)\norm b_p^2.
\end{equation}
This proves \eqref{quant-tools:eq-pairmoment}, and Minkowski gives \eqref{quant-tools:eq-uniformprofilemoment}.

For the abstract profiles, write $R=R_{\be,\mu}$ and $T=T_\mu$.
The resolvent identity and the bounded correction from
\eqref{quant-resp:eq-Ckernel} give
\begin{equation}\label{quant-tools:eq-bdecomposition}
 b(x,y)=-\be\tl W(x,y)+\be^2\ipp*{w_x,Rw_y}.
\end{equation}
Thus, $b\in L^p(\mu^{\otimes2})$ and $\omega_{\be,\mu}\in L^\infty(\mu)$.
For $F=\be\psi_\mu+\omega_{\be,\mu}^{\circ}/2$, one has $F\in L^6(\mu)$
and
\begin{equation}\label{quant-tools:eq-amoments}
 a=RF=F-\be TRF,\qquad
 \norm{TRF}_\infty\leq M_p\norm R_{\mathrm{op}}\norm F_2.
\end{equation}
Hence, $a\in L^6(\mu)$, proving the abstract profile assertion.
For the exact models, the row and source bounds in
\Cref{lem:exact-riesz-riesz-row-bounds,quant-resp:lem-logrows} give the
stronger conclusions: $a\in L^\infty$ and $b\in L^q$ for every
$2\leq q<\d/\s$ in the Riesz case, and $a,b\in L^q$ for every finite
$q\geq2$ in the logarithmic case.  Thus, the corresponding normalized
profile moments remain uniform in $k$ throughout these model ranges.
\end{proof}

% STAGE2 MOVE: A3-entropy-comparison-in-body
\subsubsection{Uniform relative-entropy comparison}\label{quant-tools:sec-entropy}

The next comparison is applied to the marginal density at scale $k/N$.  Its error is quadratic in the density-profile approximation, and its moment above order two is required only for the explicit comparison profile.

\begin{lemma}[Entropy comparison at an arbitrary scale]\label{quant-tools:lem-entropy}
Let $\nu$ be a probability measure, $t>0$, $H\in L^2(\nu)$ have mean zero, and $f=1+tH\geq0$. Let $P\in L^{2+\alpha}(\nu)$ for some $0<\alpha\leq1$. Then
\begin{equation}\label{quant-tools:eq-entropycomparison}
 \left|t^{-2}\Hr(f\nu\mid\nu)-\frac12\norm H_2^2\right|
 \leq4\norm{H-P}_2^2+2^{2+\alpha}t^\alpha\norm P_{2+\alpha}^{2+\alpha}.
\end{equation}
No $L^{2+\alpha}$ moment of $H$ is required.
\end{lemma}
\begin{proof}
For $z\geq-1$, use
\begin{equation}\label{quant-tools:eq-psientropy}
 \Psi(z)=\begin{cases}((1+z)\log(1+z)-z)/z^2,&z\ne0,\\1/2,&z=0,\end{cases}
\end{equation}
with $0\log0=0$. Twice integrating the second derivative of the numerator gives
\begin{equation}\label{quant-tools:eq-psirepresentation}
 \Psi(z)=\int_0^1\frac{1-r}{1+rz}\,\md r\qquad(z>-1).
\end{equation}
For $z\geq0$, this is between zero and $1/2$. For $-1<z<0$, the denominator is at least $1-r$, so it is between zero and one; the same holds at $z=-1$ by continuity. For $\abs z\leq1/2$, subtracting $1/2$ in \eqref{quant-tools:eq-psirepresentation} gives
$\abs{\Psi(z)-1/2}\leq\abs z/3$. For $\abs z\geq1/2$, the bound $0\leq\Psi\leq1$ gives
\begin{equation}\label{quant-tools:eq-psibound}
 \abs{\Psi(z)-1/2}\leq\min\{\abs z,1\}\qquad(z\geq-1).
\end{equation}
Since $\int H\,\md\nu=0$, and the entropy integrand after subtracting the linear term is at most $t^2H^2$,
\begin{equation}\label{quant-tools:eq-entropyexact}
 t^{-2}\Hr(f\nu\mid\nu)=\int H^2\Psi(tH)\,\md\nu.
\end{equation}
All quantities are finite, including the value at $f=0$.

On the set $\abs H>2\abs P$, one has $H^2\leq4\abs{H-P}^2$. On the complementary set, $H^2\leq4P^2$ and
$\min\{t\abs H,1\}\leq\min\{2t\abs P,1\}\leq(2t\abs P)^\alpha$.
Therefore, pointwise,
\begin{equation}\label{quant-tools:eq-entropypointwise}
 H^2\min\{t\abs H,1\}
 \leq4\abs{H-P}^2+2^{2+\alpha}t^\alpha\abs P^{2+\alpha}.
\end{equation}
Integrating \eqref{quant-tools:eq-entropypointwise} and using \eqref{quant-tools:eq-psibound}--\eqref{quant-tools:eq-entropyexact} proves \eqref{quant-tools:eq-entropycomparison}.
\end{proof}

% STAGE2 MOVE: subsubsec:quantitative-profile-proof
\subsubsection{Proof of the quantitative profile theorem}
\label{subsubsec:quantitative-profile-proof}

% STAGE2 MOVE: A2-profile-theorem-proof
\begin{proof}[Proof of \Cref{thm:abstract-quantitative-marginal}: marginal estimates]

\step{The density profile.}
Subtract \eqref{eq:exact-fixed-profile-def} from $N$ times \eqref{eq:exact-marginal-Hoeffding}.  Hoeffding orthogonality gives the exact identity
\begin{align}
 &\left\|N(f_{N,\be}^{(k)}-1)-G_{k,\be,\mu}\right\|_2^2\notag\\
 &\quad=k\norm{Nu_N-a_{\be,\mu}}_2^2
 +\binom{k}{2}\norm{Nv_N-b_{\be,\mu}}_2^2
 +N^2\norm{\mathcal R_{N,k}^{(\geq3)}}_2^2.
 \label{eq:exact-profile-orthogonal-error}
\end{align}
Equations~\eqref{eq:quantitative-kernel-rates} and
\eqref{eq:exact-higher-tail-scaled} bound this by
$C(k^2\mathfrak r_N^2+k^3/N)$.  Consequently,
\begin{equation}\label{eq:quantitative-uniform-profile}
 \norm{N(f_{N,\be}^{(k)}-1)-G_{k,\be,\mu}}_2
 \leq C\left(k\mathfrak r_N+\frac{k^{3/2}}{\sqrt N}\right).
\end{equation}
Division by $k$ proves \eqref{eq:quantitative-natural-profile}.


\step{Chi-square divergence.}
The exact squared norm is
\begin{equation}\label{eq:exact-chi-as-L2}
 N^2\Chi\!\left(\Q_{N,\be}^{(k)}(\mu)\mid\mu^{\otimes k}\right)
 =k\norm{Nu_N}_2^2+\binom{k}{2}\norm{Nv_N}_2^2
 +N^2\norm{\mathcal R_{N,k}^{(\geq3)}}_2^2.
\end{equation}
The singleton norms are uniformly bounded, so their difference of squares is $O(\mathfrak r_N)$ by \eqref{eq:quantitative-kernel-rates}.  For the pair term, use the sharper bound \eqref{eq:quantitative-pair-squared-norm}, not only the strong pair estimate.  Subtracting \eqref{eq:exact-Gk-L2} from \eqref{eq:exact-chi-as-L2} gives
\begin{equation}\label{eq:quantitative-chi-assembly}
 \left|N^2\Chi-\norm{G_{k,\be,\mu}}_2^2\right|
 \leq C\left(k\mathfrak r_N+k^2\mathfrak r_N^2+\frac{k^2}{N}+\frac{k^3}{N}\right),
\end{equation}
which implies \eqref{eq:quantitative-uniform-chi}.

\step{Total variation.}
The normalization of total variation and the reverse triangle inequality give
\begin{align}
 &\left|N\TV\!\left(\Q_{N,\be}^{(k)}(\mu),\mu^{\otimes k}\right)-\frac12\norm{G_{k,\be,\mu}}_1\right|\notag\\
 &\qquad\leq\frac12\left\|N(f_{N,\be}^{(k)}-1)-G_{k,\be,\mu}\right\|_1
 \leq\frac12\left\|N(f_{N,\be}^{(k)}-1)-G_{k,\be,\mu}\right\|_2.
 \label{eq:exact-TV-proof}
\end{align}
This proves \eqref{eq:quantitative-uniform-TV}.

\step{Relative entropy.}
Set $t=k/N$, $H=(N/k)(f_{N,\be}^{(k)}-1)$, and $P=G_{k,\be,\mu}/k$.  Then $\int H\,\md\mu^{\otimes k}=0$, $1+tH\geq0$, and \Cref{quant-tools:lem-profilemoments} gives a uniform $L^{2+\alpha}$ bound for $P$.  Applying \Cref{quant-tools:lem-entropy} and \eqref{eq:quantitative-natural-profile},
\begin{align}
 &\left|\frac{N^2}{k^2}\Hr\!\left(\Q_{N,\be}^{(k)}(\mu)\mid\mu^{\otimes k}\right)
 -\frac12\frac{N^2}{k^2}\Chi\!\left(\Q_{N,\be}^{(k)}(\mu)\mid\mu^{\otimes k}\right)\right|\notag\\
 &\qquad\leq C_\alpha\left(\mathfrak r_N^2+\frac{k}{N}+(k/N)^\alpha\right)
 \leq C_\alpha\left(\mathfrak r_N^2+(k/N)^\alpha\right).
 \label{eq:quantitative-entropy-chi-comparison}
\end{align}
Here, $k/N\leq(k/N)^\alpha$.  Multiply by $k^2$ and combine with \eqref{eq:quantitative-uniform-chi}; the term $k^3/N$ is absorbed by $k^2(k/N)^\alpha$.  This proves \eqref{eq:quantitative-uniform-entropy}.  The comparison uses a higher moment only for the explicit profile, not for $H$.

The growing-marginal conclusions are proved below.
\end{proof}

% STAGE2 MOVE: A5-natural-scale-consequence
For any integer sequence $1\leq k_N\leq N$ with $k_N/N\to0$, \eqref{eq:quantitative-natural-profile} gives
\begin{align}
 \sup_{1\leq k\leq k_N}
 \left\|\frac Nk(f_{N,\be}^{(k)}-1)-\frac1kG_{k,\be,\mu}\right\|_{L^2(\mu^{\otimes k})}
 &\leq C\left(\mathfrak r_N+\sqrt{\frac{k_N}{N}}\right)
 \longrightarrow0\qquad(N\to\infty).
 \label{eq:profile-pass-natural-uniform-limit}
\end{align}
Thus, the normalized error vanishes throughout every sublinear marginal range in either model, without a condition on $k_N\mathfrak r_N$.

% STAGE2 MOVE: subsec:growing-divergence-proof
\subsection{Divergence limits for growing marginals}
\label{subsec:growing-divergence-proof}

% STAGE2 MOVE: A3-absolute-moment-statement-interface
The chi-square and relative-entropy limits use the exact squared norm of the profile and the estimates just proved.  The total-variation limit also requires the following approximation of the absolute first moment of the canonical pair statistic.  For a self-adjoint Hilbert--Schmidt operator $C$, let $\mathcal Q_C$ be the centered Gaussian quadratic form defined in \eqref{eq:quantitative-centered-Gaussian-quadratic-form}, with $C$ in place of $B_{\be,\mu}$.  \Cref{quant-tools:sec-gaussian} gives its joint construction and proves the approximation.

% STAGE2 MOVE: A3-absolute-moment-proposition-in-body
\begin{prop}[Absolute-moment rate]\label{quant-tools:prop-absolute}
Let $0<\alpha\leq1$ and let $c\in L^{2+\alpha}(\mu^{\otimes2})$ be real, symmetric, and canonical. For independent $X_1,\ldots,X_k$ with law $\mu$, and the associated Gaussian quadratic form $\mathcal Q_C$,
\begin{equation}\label{quant-tools:eq-absolute}
 \left|\E\left[\left|\frac1k\sum_{i<j}c(X_i,X_j)\right|\right]
 -\frac12\E[\abs{\mathcal Q_C}]\right|
 \leq C_\alpha\norm c_{2+\alpha}\,k^{-\alpha/(2(2+\alpha))}
 \qquad(k\geq1).
\end{equation}
The constant depends only on $\alpha$, and in particular is independent of the rank of the kernel.
\end{prop}

% STAGE2 MOVE: A2-growing-divergence-theorem-proof
\begin{proof}[Completion of the proof of \Cref{thm:abstract-quantitative-marginal}]

Divide \eqref{eq:exact-Gk-L2} by $k^2$:
\begin{equation}\label{eq:quantitative-normalized-profile-variance}
 \left\|\frac1kG_{k,\be,\mu}\right\|_2^2
 =\frac1k\norm{a_{\be,\mu}}_2^2
 +\frac{k-1}{2k}\norm{B_{\be,\mu}}_{\mathrm{HS}}^2.
\end{equation}
Equations~\eqref{eq:quantitative-uniform-chi} and \eqref{eq:quantitative-normalized-profile-variance} give \eqref{eq:quantitative-growing-chi}, since $\mathfrak r_N/k\leq2/k$.
  The entropy estimate \eqref{eq:quantitative-growing-entropy} then follows from \eqref{eq:quantitative-entropy-chi-comparison}.

For total variation, \eqref{eq:quantitative-uniform-TV} gives
\begin{equation}\label{eq:quantitative-growing-TV-profile-error}
 \left|\frac Nk\TV-\frac12\left\|\frac1kG_{k,\be,\mu}\right\|_1\right|
 \leq C\left(\mathfrak r_N+\sqrt{k/N}\right).
\end{equation}
The singleton part satisfies
\begin{equation}\label{eq:quantitative-growing-singleton-L1}
 \E_{\mu^{\otimes k}}\left[\left|\frac1k\sum_{i=1}^ka_{\be,\mu}(X_i)\right|\right]
 \leq k^{-1/2}\norm{a_{\be,\mu}}_2.
\end{equation}
By \Cref{quant-tools:lem-profilemoments}, $b_{\be,\mu}\in L^{2+\alpha}(\mu^{\otimes2})$.  \Cref{quant-tools:prop-absolute} gives
\begin{align}
 &\left|\E_{\mu^{\otimes k}}\left[\left|\frac1k\sum_{i<j}b_{\be,\mu}(X_i,X_j)\right|\right]
 -\frac12\E[|\mathcal Q_{\be,\mu}|]\right|\notag\\
 &\qquad\leq C_\alpha k^{-\alpha/(2(2+\alpha))}.
 \label{eq:quantitative-growing-pair-absolute-moment}
\end{align}
Combining \eqref{eq:quantitative-growing-TV-profile-error}--\eqref{eq:quantitative-growing-pair-absolute-moment} proves \eqref{eq:quantitative-growing-TV}.  The factor $1/4$ is the product of the $1/2$ in total variation and the $1/2$ in the canonical quadratic $U$-statistic limit.  Only the absolute first moment of that statistic is compared.  If $k_N\to\infty$ and $k_N/N\to0$, then
\begin{align}
 \mathfrak r_N^2+\frac1{k_N}+\left(\frac{k_N}{N}\right)^\alpha
 &\longrightarrow0,\label{eq:profile-pass-growing-information-errors}\\
 \mathfrak r_N+\sqrt{\frac{k_N}{N}}+k_N^{-\alpha/(2(2+\alpha))}
 &\longrightarrow0\qquad(N\to\infty).
 \label{eq:profile-pass-growing-TV-errors}
\end{align}
Since $k_N/N\leq(k_N/N)^\alpha$, the three estimates therefore imply the asserted divergence limits.  In the logarithmic case, $\alpha=1$, giving the stated simplification.
\end{proof}

% STAGE2 MOVE: subsubsec:exact-profile-growing-errors
\subsection{Sufficient windows for undivided errors}
\label{subsubsec:exact-profile-growing-errors}

% STAGE2 MOVE: A5-smaller-window-corollary-preserved
\begin{cor}[Sufficient windows for the undivided errors]
\label{cor:quantitative-literal-growing-window}
Under the model hypotheses of \Cref{thm:quantitative-exact-marginal-profile},
the undivided density error in \eqref{eq:quantitative-uniform-profile} and
the three divergence errors in
\eqref{eq:quantitative-uniform-chi}--\eqref{eq:quantitative-uniform-TV}
tend to zero along every integer sequence $1\leq k_N\leq N$ satisfying
\begin{equation}\label{eq:quantitative-literal-growing-window}
 k_N=\begin{cases}
 o(N^{\alpha/8}),&0<\s<\d/2,\\
 o(N^{1/8}),&\s=0.
 \end{cases}
\end{equation}
Here, $\alpha$ satisfies \eqref{eq:quantitative-alpha-riesz} in the Riesz
case.  These windows are sufficient and are not claimed optimal.
\end{cor}
\begin{proof}
Take $\alpha=1$ for the logarithmic model.  In both cases,
$k_N\mathfrak r_N\to0$ and $k_N^2\mathfrak r_N^2\to0$, and, as $N\to\infty$,
\begin{align}
 k_N^3/N&=o(N^{3\alpha/8-1})\longrightarrow0,
 \label{eq:signed-window-cubic}\\
 k_N^{2+\alpha}/N^\alpha
 &=o(N^{\alpha(\alpha-6)/8})\longrightarrow0.
 \label{eq:signed-window-entropy}
\end{align}
Both exponents are negative because $0<\alpha\leq1$.
The density and total-variation finite-size error is the square root of
$k_N^3/N$.  Apply \Cref{thm:quantitative-exact-marginal-profile} and
\eqref{eq:quantitative-uniform-profile}.
\end{proof}
% STAGE2 MOVE: A2-ultraviolet-section-preserved
\section{Verification for logarithmic and Riesz interactions}
\label{sec:verification-log-riesz}

We collect the potential and row estimates once, then apply \Cref{thm:abstract-quantitative-marginal} to prove the model approximation and divergence theorems.  The final subsection verifies the separate order criterion, retaining its weaker logarithmic hypothesis and weighted alternative.  Assumptions made for the quantitative logarithmic inputs are confined to that subsection and its application; they do not govern the order verification.

\subsection{Riesz row and potential estimates}
\label{subsubsec:riesz-model-verification}

% STAGE2 MOVE: A5-Riesz-row-lemma-single-home
\begin{lemma}[Uniform row integrability]\label{lem:exact-riesz-riesz-row-bounds}
Under the Riesz hypotheses in \Cref{thm:quantitative-exact-marginal-profile}, for every $1\leq q<\d/\s$, there is a finite constant $C_q$, depending only on $\d,\s,q$, and $\norm{\mu}_{L^\infty}$, such that
\begin{equation}\label{eq:exact-riesz-row-Lq}
 \sup_{x\in\R^{\d}}
 \norm{\tl W(x,\cdot)}_{L^q(\mu)}\leq C_q.
\end{equation}
Moreover,
\begin{equation}\label{eq:exact-riesz-bounded-h-psi}
 \norm{h_\mu}_{L^\infty}
 +\norm{\psi_\mu}_{L^\infty}<\infty,
 \qquad
 \operatorname*{ess\,inf}_{x,y}\tl W(x,y)>-\infty.
\end{equation}
In particular,
\begin{equation}\label{eq:exact-riesz-uniform-row-L2}
 M_2\coloneqq
 \sup_x\norm{w_x}_{L^2(\mu)}<\infty.
\end{equation}
\end{lemma}

\begin{proof}
Split the Riesz potential at unit distance.  Uniformly in $x$,
\begin{align}
 \int W(x,y)\,\md\mu(y)
 &\leq
 \frac{\norm{\mu}_{L^\infty}}{\s}
 \int_{\abs z\leq1}\abs z^{-\s}\,\md z
 +\frac1\s,
 \label{eq:exact-riesz-h-bound-proof}
\end{align}
which is finite because $\s<\d$.  This proves the first bound in \eqref{eq:exact-riesz-bounded-h-psi}; the remaining bounded one-body terms follow from the definitions.

For $q<\d/\s$, the same decomposition gives
\begin{equation}\label{eq:exact-riesz-raw-row-Lq}
 \sup_x\int W(x,y)^q\,\md\mu(y)
 \leq
 \frac{\norm{\mu}_{L^\infty}}{\s^q}
 \int_{\abs z\leq1}\abs z^{-q\s}\,\md z
 +\frac1{\s^q}<\infty.
\end{equation}
Adding the bounded centering terms in \eqref{eq:centered-W-explicit} proves \eqref{eq:exact-riesz-row-Lq}.  Since $W\geq0$ and $h_\mu,I_\mu$ are bounded, the lower bound in \eqref{eq:exact-riesz-bounded-h-psi} follows.  Finally, $2<\d/\s$ by hypothesis, so \eqref{eq:exact-riesz-uniform-row-L2} is the case $q=2$.
\end{proof}

% STAGE2 MOVE: A5-Riesz-order-verification
\subsection{Logarithmic potential estimates}
\label{subsubsec:log-order-verification}

% STAGE2 MOVE: A5-common-logarithmic-estimates-assumptions
For the elementary estimates below, assume only that $\mu$ has bounded density and a finite polynomial moment of some positive order.  Set

% STAGE2 MOVE: A5-elementary-log-estimates-single-home
\begin{align}\label{eq:log-local-singularity-weight}
 \ell(z)\coloneqq\pa{-\log|z|}_+,
 \qquad
 L(x)\coloneqq\log(e+|x|).
\end{align}
The elementary bounds
\begin{align}\label{eq:log-kernel-elementary-bounds}
 \pa{\log|x-y|}_+\le L(x)+L(y),
 \qquad
 |\log|x-y||\le \ell(x-y)+L(x)+L(y)
\end{align}
hold for all \(x,y\).  Since \(\mu\in L^\infty\), for every finite \(m\ge1\),
\begin{align}\label{eq:uniform-local-log-moments}
 \sup_{x\in\R^\d}\int_{\R^\d}\ell(x-y)^m\,\md\mu(y)
 \le \norm{\mu}_{L^\infty}
 \int_{|z|<1}\pa{-\log|z|}^m\,\md z<\infty.
\end{align}
Moreover, the positive polynomial moment implies that \(L\in L^m(\mu)\) for every finite \(m\).  It follows from \eqref{eq:log-kernel-elementary-bounds}--\eqref{eq:uniform-local-log-moments} that the absolute logarithmic energy is finite.

The logarithmic potential satisfies
\begin{align}\label{eq:log-potential-growth}
 -L(x)-\int L\,\md\mu
 \le h_\mu(x)
 \le \sup_{z\in\R^\d}\int \ell(z-y)\,\md\mu(y)<\infty.
\end{align}
Hence, \(|h_\mu(x)-2I_\mu|\le C_\mu(1+L(x))\).  Also,
\begin{align}\label{eq:centered-log-growth}
 |\tl W(x,y)|
 \le C_\mu+\ell(x-y)+L(x)+L(y).
\end{align}
In particular, the finiteness assertions in \eqref{eq:mu-controls} follow.  The same estimates give
\begin{align}
 (\tl W)_-(x,y)&\le C_\mu+L(x)+L(y),\label{eq:centered-log-negative-part}\\
 (h_\mu-2I_\mu)_+(x)&\le C_\mu,\label{eq:log-potential-positive-part}
\end{align}
which verifies \eqref{eq:cavity-tail-log} for suitable finite constants.

\subsection{Logarithmic inputs for marginal approximation}
\label{subsubsec:log-quantitative-model-estimates}

% STAGE2 MOVE: A5-stronger-log-hypotheses-boundary
We now assume the logarithmic hypotheses of \Cref{thm:quantitative-exact-marginal-profile}.  The following potential and centered-row estimates verify the source and kernel assumptions of the quantitative theorem.  The weighted hypotheses of the order theorem above are unchanged.


% STAGE2 MOVE: A5-log-analytic-statement
\begin{lemma}[Logarithmic analytic input]\label{lem:exact-log-analytic}
Under the logarithmic hypotheses in \Cref{thm:quantitative-exact-marginal-profile}, including both \eqref{eq:exact-log-polynomial-moment} and \eqref{eq:exact-log-moment-threshold}, the following assertions hold.

\begin{enumerate}[label=\textup{(\roman*)}]
\item The absolute logarithmic energy is finite, $h_\mu$ is finite $\mu$-almost everywhere, and
\begin{equation}\label{eq:exact-log-h-two-sided}
 -L(x)-\int L\,\md\mu
 \leq h_\mu(x)\leq C_{\mathrm{loc}}
 \quad\text{for every }x\text{ for which }h_\mu(x)\text{ is defined}.
\end{equation}
In particular, for a finite constant $C$,
\begin{equation}\label{eq:exact-log-psi-growth}
 \abs{\psi_\mu(x)}\leq C\bigl(1+L(x)\bigr).
\end{equation}
\item For every finite $p\geq1$,
\begin{equation}\label{eq:exact-log-w-Lp}
 \tl W\in L^p(\mu^{\otimes2}).
\end{equation}
Moreover, with
\begin{align}
 a_0&\coloneqq2C_{\mathrm{loc}}+2\abs{I_\mu},
 &a_1&\coloneqq1,
 &a_h&\coloneqq C_{\mathrm{loc}}+2\abs{I_\mu},
 \label{eq:exact-log-envelope-constants}
\end{align}
one has
\begin{align}
 (\tl W)_-(x,y)
 &\leq a_0+a_1\bigl(L(x)+L(y)\bigr),\label{eq:exact-log-w-negative-envelope}\\
 (\psi_\mu)_+(x)
 &\leq a_h\bigl(1+L(x)\bigr).
 \label{eq:exact-log-psi-positive-envelope}
\end{align}
Thus, $\mathsf c_{\log,\be,\mu}=\be(a_0+a_1+a_h)$.
\item The logarithmic determinant and fluctuation inputs of~\cite{DGR26} satisfy
\begin{equation}\label{eq:exact-log-DGR-log-inputs}
 \int e^{-4\be h_\mu}\,\md\mu<\infty,
 \qquad
 \tl W\in L^2(\mu^{\otimes2}).
\end{equation}
\item The operator $T_\mu$ is self-adjoint, Hilbert--Schmidt, and nonnegative on $L^2(\mu)$.
\end{enumerate}
\end{lemma}

% STAGE2 MOVE: A5-log-analytic-proof-open
\begin{proof}

% STAGE2 MOVE: A5-log-analytic-proof-reuse
The estimates \eqref{eq:log-kernel-elementary-bounds}--\eqref{eq:centered-log-growth} prove finite absolute logarithmic energy, \eqref{eq:exact-log-h-two-sided}, and \eqref{eq:exact-log-psi-growth}.  They also give \eqref{eq:exact-log-w-Lp}, since every finite power of $L$ is integrable under a positive polynomial moment.  Using $h_\mu\leq C_{\mathrm{loc}}$ in the centering formula gives the precise envelopes \eqref{eq:exact-log-w-negative-envelope}--\eqref{eq:exact-log-psi-positive-envelope} with the constants in \eqref{eq:exact-log-envelope-constants}.

% STAGE2 MOVE: A5-log-analytic-exponential-and-positivity-proof
The lower bound in \eqref{eq:exact-log-h-two-sided} yields
\begin{equation}\label{eq:exact-log-exponential-h-bound}
 e^{-4\be h_\mu(x)}
 \leq
 e^{4\be\int L\,\md\mu}(e+\abs x)^{4\be}.
\end{equation}
The first part of \eqref{eq:exact-log-moment-threshold} proves the first assertion in \eqref{eq:exact-log-DGR-log-inputs}; the second follows from \eqref{eq:exact-log-w-Lp}.

Only nonnegativity of $T_\mu$ remains.  Let $f\in L^2(\mu)$ satisfy $\int f\,\md\mu=0$.  Since $\tl W\in L^2(\mu^{\otimes2})$, the quadratic form
\begin{equation}\label{eq:exact-log-quadratic-form-continuity}
 f\longmapsto\iint\tl W(x,y)f(x)f(y)\,\md\mu(x)\md\mu(y)
\end{equation}
is continuous on $L^2(\mu)$.  Approximate $f$ in $L^2(\mu)$ by bounded, compactly supported, centered functions $f_n$.  Indeed, truncate $f$ in value and space and subtract its mean times a fixed bounded compactly supported function having $\mu$-mean one.  Put $\nu_n=f_n\mu$.

For $0<s<\d$, the Riesz Fourier representation gives
\begin{equation}\label{eq:exact-log-riesz-positive-approximation}
 \frac1s\iint\abs{x-y}^{-s}\,\md\nu_n(x)\md\nu_n(y)
 =c_{\d,s}\int_{\R^{\d}}
 \abs{\widehat{\nu_n}(\xi)}^2\abs{\xi}^{s-\d}\,\md\xi
 \geq0,
\end{equation}
where $c_{\d,s}>0$.  Since $\nu_n(\R^{\d})=0$, the left-hand side is unchanged if $\abs{x-y}^{-s}$ is replaced by $\abs{x-y}^{-s}-1$.  As $s\downarrow0$,
\begin{equation}\label{eq:exact-log-riesz-to-log}
 \frac{\abs{x-y}^{-s}-1}{s}
 \longrightarrow-\log\abs{x-y}.
\end{equation}
The measures $\nu_n$ have bounded compactly supported densities.  Taking $0<s\leq s_0<\d$, the integrands in \eqref{eq:exact-log-riesz-to-log} are dominated near the diagonal by an integrable multiple of $\abs{x-y}^{-s_0}(1+\abs{\log\abs{x-y}})$, and they are uniformly bounded away from the diagonal on the common compact support.  Dominated convergence in \eqref{eq:exact-log-riesz-positive-approximation} therefore gives
\begin{equation}\label{eq:exact-log-log-positive-compact}
 \iint-\log\abs{x-y}\,\md\nu_n(x)\md\nu_n(y)\geq0.
\end{equation}
The one-variable centering terms vanish against $\nu_n$, so the left-hand side equals the quadratic form of $\tl W$.  Letting $n\to\infty$ in \eqref{eq:exact-log-quadratic-form-continuity} proves nonnegativity on the centered subspace.  Since $T_\mu\one=0$, decomposing an arbitrary $L^2(\mu)$ function into its mean and centered part proves nonnegativity on all of $L^2(\mu)$.  Self-adjointness and the Hilbert--Schmidt property follow from symmetry and \eqref{eq:exact-log-w-Lp} with $p=2$.

% STAGE2 MOVE: A5-log-analytic-proof-close
\end{proof}

% STAGE2 MOVE: A5-uniform-log-rows
\begin{lemma}[Uniform logarithmic rows]\label{quant-resp:lem-logrows}
Assume the logarithmic hypotheses in \Cref{thm:quantitative-exact-marginal-profile}. Then $h_\mu(x)+L(x)$ is uniformly bounded, $\tl W$ has a uniform lower bound, and, for every finite $p\geq1$,
\begin{equation}\label{quant-resp:eq-logrows}
 \sup_x\norm{w_x}_{L^p(\mu)}<\infty.
\end{equation}
All statements are understood outside null sets where a chosen kernel representative is undefined.
\end{lemma}
\begin{proof}
Let $\ell(z)=(-\log\abs z)_+$ and $C_0=\log(1+e)$. The triangle inequality implies
\begin{equation}\label{quant-resp:eq-logdistance}
 -L(y)\leq-\log\abs{x-y}+L(x)\leq L(y)+C_0+\ell(x-y).
\end{equation}
For the lower bound, use $\abs{x-y}\leq(e+\abs x)(e+\abs y)$. For the upper bound, use $e+\abs x\leq(e+\abs{x-y})(e+\abs y)$ and
$\log(1+e/\abs z)\leq C_0+\ell(z)$.

The local logarithmic moments are bounded uniformly by \eqref{eq:uniform-local-log-moments}.
Moreover, \eqref{eq:exact-log-moment-threshold} implies $L\in L^p(\mu)$ for every finite $p$. Integrating \eqref{quant-resp:eq-logdistance}, with $m_L=\int L\,\md\mu$, gives
\begin{equation}\label{quant-resp:eq-hshift}
 -m_L\leq h_\mu(x)+L(x)\leq m_L+C_0+C_{\mathrm{loc}}.
\end{equation}
Rewrite the centered kernel as
\begin{align}
 \tl W(x,y)
 &=\bigl[-\log\abs{x-y}+L(x)+L(y)\bigr]\notag\\
 &\quad-\bigl[h_\mu(x)+L(x)\bigr]-\bigl[h_\mu(y)+L(y)\bigr]+2I_\mu.
 \label{quant-resp:eq-logcenteredcancel}
\end{align}
The first bracket is nonnegative, and \eqref{quant-resp:eq-hshift} proves a uniform lower bound for $\tl W$. Its upper bound, by \eqref{quant-resp:eq-logdistance} and \eqref{quant-resp:eq-hshift}, is $C+2L(y)+\ell(x-y)$. Therefore,
\begin{equation}\label{quant-resp:eq-logonesidedrow}
 \abs{\tl W(x,y)}\leq C+2L(y)+\ell(x-y).
\end{equation}
Combining \eqref{quant-resp:eq-logonesidedrow} with \eqref{eq:uniform-local-log-moments} proves \eqref{quant-resp:eq-logrows}. The displayed one-sided envelope does not assert that the kernel itself is bounded; its collision singularity remains.
\end{proof}

\subsection{Verification of the quantitative hypotheses}
\label{subsubsec:signed-model-verification}

\begin{proof}[Proof of \Cref{thm:quantitative-exact-marginal-profile,thm:quantitative-growing-marginal-constants}]
It suffices to verify the hypotheses of
\Cref{thm:abstract-quantitative-marginal}.  The estimates used here concern
the interaction and the reference measure; no marginal estimate enters
this verification.

For Riesz interactions, fix $\alpha$ as in
\eqref{eq:quantitative-alpha-riesz}, choose $2+\alpha<p<\d/\s$, and take
$\vartheta=\alpha/2$.  \Cref{lem:exact-riesz-riesz-row-bounds} gives
\eqref{eq:signed-kernel-assumptions}--\eqref{eq:signed-source-assumptions}.
Proposition~5.1, equation~(5.4), of~\cite{DGR26}, applied at zero cutoff
and at the fixed inverse temperature $\Gamma$ in
\eqref{eq:signed-temperature-amplification}, gives
\eqref{eq:signed-centered-pair-moment}.  Its ordered statistic has the
normalization \eqref{eq:signed-ordered-statistic}.  There is no additional
spatial-moment assumption in this case, and
\begin{equation}\label{eq:signed-Riesz-model-rate}
 N^{-\alpha/8}+N^{-\vartheta/2}
 =N^{-\alpha/8}+N^{-\alpha/4}\leq2N^{-\alpha/8}.
\end{equation}

For the logarithmic model, take $p=4$, $\vartheta=1/2$, and $\alpha=1$.
Then $r_*=12$ and $\Gamma=12\be$.  The potential bounds in
\Cref{lem:exact-log-analytic} and the centered row estimates in
\Cref{quant-resp:lem-logrows} give
\eqref{eq:signed-kernel-assumptions}--\eqref{eq:signed-source-assumptions}.
In particular, $\psi_\mu$ is bounded above, and belongs to every finite
$L^q(\mu)$, although it need not be bounded below.  The logarithmic
centered-family bound in \cite[Proposition~5.1]{DGR26} requires finite
absolute logarithmic energy and $\int e^{-2\Gamma h_\mu}\,\md\mu<\infty$.
The former was proved above.  Since $\mathsf c_{\log,\be,\mu}\geq\be$,
\eqref{eq:exact-log-moment-threshold} implies $q>24\be$, and
\eqref{eq:log-potential-growth} gives
\begin{equation}\label{eq:signed-log-centered-moment-check}
 \int e^{-24\be h_\mu(x)}\,\md\mu(x)
 \leq e^{24\be\int L\,\md\mu}
       \int(e+\abs x)^{24\be}\,\md\mu(x)<\infty.
\end{equation}
Thus, the zero-cutoff conclusion of that proposition supplies
\eqref{eq:signed-centered-pair-moment}, and
\begin{equation}\label{eq:signed-log-model-rate}
 N^{-\alpha/8}+N^{-\vartheta/2}
 =N^{-1/8}+N^{-1/4}\leq2N^{-1/8}.
\end{equation}
Applying \Cref{thm:abstract-quantitative-marginal} proves all the asserted
estimates and limits.  The product-law case $\be=0$ is immediate.
\end{proof}
% Exact-profile integration module: shared partition normalization and exact external inputs

% STAGE2 MOVE: A2-exact-marked-identities
\subsection{Verification of the order criterion}
\label{subsec:model-order-verification}

The following applications of \Cref{thm:chi2-chaos-two-sided} give the sharp scales recorded in the introduction.  They do not use the quantitative marginal theorem.  In particular, the logarithmic criterion below retains its original, weaker polynomial-moment requirement and the more general weighted formulation.

\subsubsection{Riesz order bounds}
For every fixed $\be\geq0$, $0<\s<\d/2$, and bounded probability density $\mu$, the hypotheses of \Cref{thm:chi2-chaos-two-sided} hold for $W(x,y)=\g_\s(x-y)$.  Thus, its upper bounds hold for every $1\leq k\leq N$; for $\be>0$, its lower bounds hold for every $2\leq k\leq N$ once $N\geq N_0$, with $N_0$ independent of $k$.

\begin{proof}[Riesz order bounds]
Choose $2<p<\d/\s$.  \Cref{lem:exact-riesz-riesz-row-bounds} supplies \eqref{eq:mu-controls} and the bounded lower-tail condition \eqref{eq:cavity-tail-bounded}.  To verify the separate nontriviality condition \eqref{eq:interaction-nontriviality}, let \(\nu=f\mu\), where \(f\) is a nonzero bounded centered function.  The standard Fourier representation gives, whenever the energy is finite,
\begin{align}\label{eq:riesz-zero-mass-fourier}
 \int_{(\R^\d)^2}W(x,y)\,\md\nu(x)\md\nu(y)
 =c_{\d,\s}\int_{\R^\d}\abs{\widehat\nu(\xi)}^2\abs{\xi}^{\s-\d}\,\md\xi,
\end{align}
with \(c_{\d,\s}>0\); see, for example, \cite[Chapter~4, Section~4.1]{SerfatyLN}.  The right-hand side is strictly positive.  Since the one-variable centering terms vanish against a zero-mass measure, the same quadratic form is generated by \(\tl W\); in particular, \(\tl W\) cannot vanish in \(L^2(\mu^{\otimes2})\).  Applying the companion uniform-bound theorem~\cite[Theorem~2.2\textup{(i)}]{DGR26} at the two fixed inverse temperatures \(r\be\) and \(2\be\) supplies \eqref{eq:pf-input}.  The claim follows from \Cref{thm:chi2-chaos-two-sided}.
\end{proof}

% STAGE2 MOVE: subsubsec:log-order-verification
\subsubsection{Logarithmic order bounds}
\label{subsubsec:logarithmic-order-application}

Use the exponents $p,r$ from \eqref{eq:p-choice-quantitative} and the logarithmic lower-tail data of \eqref{eq:cavity-tail-log}.  These are the order-theory parameters, independent of the choices in the preceding quantitative verification.

% STAGE2 MOVE: A5-weighted-criterion-and-log-moment-statement
The following lemma gives an explicit, nonoptimal sufficient condition for the weighted estimates \eqref{eq:log-weighted-cavity-constants}.  Its auxiliary notation and the stronger pointwise estimate used in the proof are deferred to \Cref{app:weighted-log-cavity}.

\begin{lemma}[Weighted logarithmic cavity estimates]\label{lem:weighted-logarithmic-cavity}
Assume the logarithmic lower-tail condition \eqref{eq:cavity-tail-log} and \(\mathfrak P_{r\be}(\mu)<\infty\), and write \(c\coloneqq\mathsf c_{\log,\be,\mu}\).  Suppose that
\begin{subequations}\label{eq:weighted-logarithmic-inputs}
\begin{align}
 &\int_{\R^\d}(e+|x|)^{p^2c}\,\md\mu(x)<\infty,\label{eq:weighted-logarithmic-moment}\\
 &\left\|(e+|\cdot|)^{2c}
 \left(
 |h_\mu-2I_\mu|^2
 +\left[\int_{\R^\d}|\tl W(\cdot,y)|^p(e+|y|)^{p^2c}\,\md\mu(y)\right]^{2/p}
 \right)\right\|_{L^2(\mu)}<\infty.\label{eq:weighted-logarithmic-section}
\end{align}
\end{subequations}
Then the weighted cavity estimates \eqref{eq:log-weighted-cavity-constants} hold.
\end{lemma}

The proof of \Cref{lem:weighted-logarithmic-cavity} is given in \Cref{app:weighted-log-cavity}.

\begin{lemma}[Polynomial moments imply the logarithmic cavity hypotheses]
\label{lem:log-polynomial-moment-criterion}
Fix \(\be\ge0\) and \(p>2\), let \(W(x,y)=-\log\abs{x-y}\), and let
\(\mu\in L^\infty\cap\Pp(\R^\d)\).  Suppose that
\begin{align}\label{eq:log-polynomial-moment-condition}
 \int_{\R^\d}(e+|x|)^q\,\md\mu(x)<\infty
\end{align}
for some \(q>0\).  Then the logarithmic lower-tail condition \eqref{eq:cavity-tail-log} holds for suitable finite constants \(a_0,a_1,a_h\).  Fix one such choice and let \(c=\mathsf c_{\log,\be,\mu}\) be as in \eqref{eq:log-cavity-weight-constant}.  If
\begin{align}\label{eq:log-polynomial-moment-threshold}
 q>\max\left\{p^2c,\,4(p+1)c\right\},
\end{align}
then the finiteness requirements in \eqref{eq:mu-controls}, finite absolute logarithmic energy, and the assumptions \eqref{eq:weighted-logarithmic-inputs} all hold.  Consequently, if \(\mathfrak P_{r\be}(\mu)<\infty\), then the weighted cavity estimates \eqref{eq:log-weighted-cavity-constants} hold.
\end{lemma}

% STAGE2 MOVE: A5-log-moment-proof-open
\begin{proof}

% STAGE2 MOVE: A5-log-moment-proof-transition
The elementary estimates above give finite absolute logarithmic energy, \eqref{eq:mu-controls}, and the logarithmic lower-tail condition.  We verify the remaining weighted requirements.

% STAGE2 MOVE: A5-log-moment-proof-weighted-part
The first condition follows directly from \eqref{eq:log-polynomial-moment-condition} and \(q>p^2c\).  By \eqref{eq:centered-log-growth} and the inequality \((e+|y|)^{p^2c}\le C_{p,c}(e+|x|)^{p^2c}\) when \(|x-y|<1\),
\begin{align}\label{eq:weighted-logarithmic-section-bound}
 &\left[\int_{\R^\d}|\tl W(x,y)|^p(e+|y|)^{p^2c}\,\md\mu(y)\right]^{1/p}\notag\\
 &\qquad\le C_{\be,p,\mu}
 \left(1+L(x)+(e+|x|)^{pc}\right).
\end{align}
Together with \eqref{eq:log-potential-growth}, this gives
\begin{align}\label{eq:weighted-logarithmic-input-bound}
 &(e+|x|)^{2c}
 \left(
 |h_\mu(x)-2I_\mu|^2
 +\left[\int_{\R^\d}|\tl W(x,y)|^p(e+|y|)^{p^2c}\,\md\mu(y)\right]^{2/p}
 \right)\notag\\
 &\qquad\le C_{\be,p,\mu}
 \left[(e+|x|)^{2c}(1+L(x))^2+(e+|x|)^{2(p+1)c}\right].
\end{align}
After squaring, the right-hand side is integrable under \(\mu\) whenever \(q>4(p+1)c\); logarithmic factors are absorbed by the strict inequality.  This proves the second condition in \eqref{eq:weighted-logarithmic-inputs}.  The final assertion follows from \Cref{lem:weighted-logarithmic-cavity}.

% STAGE2 MOVE: A5-log-moment-proof-close
\end{proof}

% STAGE2 MOVE: A5-log-order-verification-and-general-corollary
For $W(x,y)=-\log\abs{x-y}$ and a bounded probability density satisfying the criterion of \Cref{lem:log-polynomial-moment-criterion} for some $p>2$, the conclusions of \Cref{thm:chi2-chaos-two-sided} hold.  Its constants and onset are independent of the marginal size.  This statement uses the order criterion, not the stronger logarithmic hypothesis of \Cref{thm:quantitative-exact-marginal-profile}.

\begin{proof}[Logarithmic order bounds]
Choose \(p>2\) and \(q\) as in \Cref{lem:log-polynomial-moment-criterion}.  That lemma supplies the quantitative integrability, logarithmic lower-tail, and weighted cavity hypotheses.  It also gives finite absolute logarithmic energy, which is precisely the additional logarithmic hypothesis in the companion uniform-bound theorem~\cite[Theorem~2.2\textup{(i)}]{DGR26}.  Applying that theorem at \(r\be\) and \(2\be\) supplies \eqref{eq:pf-input}.  If \(\nu=f\mu\) has zero total mass, the logarithmic energy admits the Fourier representation
\begin{align}\label{eq:log-zero-mass-fourier}
 -\int_{(\R^\d)^2}\log\abs{x-y}\,\md\nu(x)\md\nu(y)
 =c_\d\int_{\R^\d}\abs{\widehat\nu(\xi)}^2\abs{\xi}^{-\d}\,\md\xi,
\end{align}
whenever finite, with \(c_\d>0\); see, for example, \cite[Chapter~4, Section~4.1]{SerfatyLN}.  Choose any nonzero bounded centered \(f\) and set \(\nu=f\mu\).  The right-hand side is strictly positive, and the same quadratic form is generated by the Hoeffding-centered kernel because the centering terms vanish against \(\nu\).  Thus, \eqref{eq:interaction-nontriviality} holds, and \Cref{thm:chi2-chaos-two-sided} applies.
\end{proof}

\begin{remark}[Compactly supported logarithmic backgrounds]
If \(\mu\) has compact support, then \eqref{eq:log-polynomial-moment-condition} holds for every \(q>0\), so one may choose \(q\) above the threshold \eqref{eq:log-polynomial-moment-threshold}.  Thus, compact support satisfies the logarithmic order criterion, but is not required.  In this special case, one may alternatively use the simpler bounded lower-tail condition \eqref{eq:cavity-tail-bounded}.
\end{remark}

\begin{cor}[General weighted logarithmic backgrounds]\label{cor:intro-log-weighted}
Let \(\be\ge0\), let \(W(x,y)=-\log\abs{x-y}\), and let \(\mu\in L^\infty\cap\Pp(\R^\d)\) satisfy
\begin{align}\label{eq:absolute-log-energy-corollary}
 \int_{(\R^\d)^2}\abs{\log\abs{x-y}}\,\md\mu(x)\md\mu(y)<\infty.
\end{align}
Suppose that, for some \(p>2\), the quantitative integrability assumptions \eqref{eq:mu-controls}, the logarithmic lower-tail condition \eqref{eq:cavity-tail-log}, and the weighted cavity estimates \eqref{eq:log-weighted-cavity-constants} hold.  Then the conclusions of \Cref{thm:chi2-chaos-two-sided} hold.  The conditions of \Cref{lem:weighted-logarithmic-cavity} are sufficient for the weighted cavity estimates.
\end{cor}

\begin{proof}
The zero-mass Fourier representation used above verifies the nontriviality condition \eqref{eq:interaction-nontriviality} under the stated integrability assumptions.  The remaining external input is \eqref{eq:pf-input}, which follows by applying the logarithmic uniform-bound theorem~\cite[Theorem~2.2\textup{(i)}]{DGR26} at \(r\be\) and \(2\be\); the finite absolute logarithmic energy required there is assumed in \eqref{eq:absolute-log-energy-corollary}.  The weighted estimates are assumed directly or supplied by \Cref{lem:weighted-logarithmic-cavity}; hence, \Cref{thm:chi2-chaos-two-sided} applies.
\end{proof}

% STAGE2 MOVE: subsubsec:log-quantitative-model-estimates
\section{Fixed marginals under a subquadratic moment}\label{sec:subquadratic-marginals}
\label{subsec:fixed-marginal-U-statistic-transfer}

We first estimate canonical $U$-statistics under the product law and then pair them with the full Gibbs density.  The degree-two estimate preserves the $N^{-1}$ scale; the higher-order terms are negligible at that scale for each fixed marginal.

Let \((E,\mathcal E,\mu)\) be a probability space.  For integers \(1\le j\le N\), write
\begin{align}\label{eq:falling-factorial-definition}
 (N)_j\coloneqq N(N-1)\cdots(N-j+1).
\end{align}
A measurable kernel \(h:E^j\to\R\) is \emph{canonical} if its integral in each variable vanishes for almost every choice of the remaining variables.  For independent \(X_1,\ldots,X_N\sim\mu\), define the ordered canonical \(U\)-statistic
\begin{align}\label{eq:ordered-canonical-U-statistic}
 \mathcal U_{N,j}h
 \coloneqq\frac1{(N)_j}
 \sum_{\substack{1\le i_1,\ldots,i_j\le N\\
 i_1,\ldots,i_j\ \mathrm{all\ distinct}}}
 h(X_{i_1},\ldots,X_{i_j}).
\end{align}

\begin{lemma}[Canonical \(U\)-statistics of arbitrary fixed order]
\label{lem:fixed-order-canonical-U-moment}
For every integer \(j\ge2\) and every \(q\ge2\), there is a finite constant \(\mathsf A_{j,q}\), depending only on \(j\) and \(q\), such that, for every \(N\ge j\) and every canonical \(h\in L^q(\mu^{\otimes j})\),
\begin{align}\label{eq:canonical-U-moment}
 \norm{\mathcal U_{N,j}h}_{L^q(\mu^{\otimes N})}
 \le \frac{\mathsf A_{j,q}}{\sqrt{(N)_j}}
 \norm{h}_{L^q(\mu^{\otimes j})}.
\end{align}
If \(\mathsf C_j\) is the constant in the upper-tail comparison of de la Pe\~na and Montgomery-Smith~\cite[Theorem~1]{dPmS1995}, and \(\mathsf K_q\) is an admissible scalar Khintchine constant, one may take
\begin{align}\label{eq:canonical-U-moment-constants}
 \mathsf D_{j,q}\coloneqq \mathsf C_j^{1+1/q},
 \qquad
 \mathsf A_{j,q}\coloneqq \mathsf D_{j,q}2^j\mathsf K_q^j.
\end{align}
\end{lemma}

\begin{proof}
Let \(S^{\mathrm{cpl}}\) denote the unnormalized sum in \eqref{eq:ordered-canonical-U-statistic}, and let
\begin{align}\label{eq:decoupled-canonical-U-statistic}
 S^{\mathrm{dec}}
 \coloneqq
 \sum_{\substack{1\le i_1,\ldots,i_j\le N\\
 i_1,\ldots,i_j\ \mathrm{all\ distinct}}}
 h\pa{X_{i_1}^{(1)},\ldots,X_{i_j}^{(j)}},
\end{align}
where the \(j\) coordinate families are independent copies.  The upper direction of the cited tail-decoupling theorem requires no symmetry of \(h\) and gives
\begin{align}\label{eq:canonical-U-tail-decoupling}
 \Pb\pa{|S^{\mathrm{cpl}}|\ge t}
 \le \mathsf C_j\Pb\pa{\mathsf C_j|S^{\mathrm{dec}}|\ge t}.
\end{align}
Layer-cake integration yields
\begin{align}\label{eq:canonical-U-layer-cake}
 \E\left[|S^{\mathrm{cpl}}|^q\right]
 &\le \mathsf C_jq\int_0^\infty t^{q-1}
 \Pb\pa{|S^{\mathrm{dec}}|\ge t/\mathsf C_j}\,\md t\notag\\
 &=\mathsf C_j^{q+1}\E\left[|S^{\mathrm{dec}}|^q\right].
\end{align}
Thus, \(\norm{S^{\mathrm{cpl}}}_q\le \mathsf D_{j,q}\norm{S^{\mathrm{dec}}}_q\), which is the tail-to-moment conversion in \eqref{eq:canonical-U-moment-constants}.

Complete degeneracy permits successive symmetrization in the \(j\) independent coordinate families:
\begin{align}\label{eq:canonical-U-successive-symmetrization}
 \norm{S^{\mathrm{dec}}}_{L^q}
 \le 2^j
 \norm*{
 \sum_{\substack{1\le i_1,\ldots,i_j\le N\\
 i_1,\ldots,i_j\ \mathrm{all\ distinct}}}
 \varepsilon_{i_1}^{(1)}\cdots\varepsilon_{i_j}^{(j)}
 h\pa{X_{i_1}^{(1)},\ldots,X_{i_j}^{(j)}}}_{L^q}.
\end{align}
Conditionally on the \(X\)-variables, repeated Khintchine inequality gives
\begin{align}\label{eq:canonical-U-rademacher-bound}
 &\norm*{
 \sum_{\substack{1\le i_1,\ldots,i_j\le N\\
 i_1,\ldots,i_j\ \mathrm{all\ distinct}}}
 \varepsilon_{i_1}^{(1)}\cdots\varepsilon_{i_j}^{(j)}
 h\pa{X_{i_1}^{(1)},\ldots,X_{i_j}^{(j)}}}_{L^q(\varepsilon)}\notag\\
 &\qquad\le \mathsf K_q^j
 \left(
 \sum_{\substack{1\le i_1,\ldots,i_j\le N\\
 i_1,\ldots,i_j\ \mathrm{all\ distinct}}}
 \left|h\pa{X_{i_1}^{(1)},\ldots,X_{i_j}^{(j)}}\right|^2
 \right)^{1/2}.
\end{align}
Since \(q/2\ge1\), Minkowski's inequality implies
\begin{align}\label{eq:canonical-U-square-function}
 &\norm*{
 \left(
 \sum_{\substack{1\le i_1,\ldots,i_j\le N\\
 i_1,\ldots,i_j\ \mathrm{all\ distinct}}}
 \left|h\pa{X_{i_1}^{(1)},\ldots,X_{i_j}^{(j)}}\right|^2
 \right)^{1/2}}_{L^q(X)}\notag\\
 &\qquad\le \sqrt{(N)_j}\norm{h}_{L^q(\mu^{\otimes j})}.
\end{align}
Combining these estimates and dividing by \((N)_j\) proves \eqref{eq:canonical-U-moment}.
\end{proof}

\begin{lemma}[Refined canonical pair moment]
\label{lem:refined-canonical-pair-moment}
For every \(q>2\), there is a finite constant \(\mathsf B_q\), depending only on \(q\), such that, for every \(N\ge2\) and every canonical \(h\in L^q(\mu^{\otimes2})\),
\begin{align}\label{eq:refined-canonical-pair-moment}
 \norm{\mathcal U_{N,2}h}_{L^q(\mu^{\otimes N})}
 \le \mathsf B_q
 \left(
 \frac{\norm{h}_{L^2(\mu^{\otimes2})}}{N-1}
 +\frac{N^{1/q-1/2}}{N-1}
 \norm{h}_{L^q(\mu^{\otimes2})}
 \right).
\end{align}
\end{lemma}

\begin{proof}
Let \((X_i)_{i\le N}\) and \((Y_j)_{j\le N}\) be independent samples with law \(\mu\).  The order-two case of the moment decoupling obtained from~\cite[Theorem~1]{dPmS1995} reduces the coupled ordered sum, up to a constant depending only on \(q\), to
\begin{align}\label{eq:refined-pair-decoupled-sum}
 S\coloneqq\sum_{i\ne j}h(X_i,Y_j).
\end{align}
Condition on \(Y=(Y_1,\ldots,Y_N)\) and set
\begin{align}\label{eq:refined-pair-Zi}
 Z_i\coloneqq\sum_{j\ne i}h(X_i,Y_j).
\end{align}
The variables \(Z_i\) are conditionally independent and centered.  Conditional scalar Rosenthal inequality~\cite{Rosenthal1970} gives
\begin{align}\label{eq:refined-pair-conditional-Rosenthal}
 \norm{S}_{L^q}
 \le C_q\left\{
 \norm*{\pa*{\sum_i\E[Z_i^2\mid Y]}^{1/2}}_{L^q(Y)}
 +\pa*{\sum_i\E\left[|Z_i|^q\right]}^{1/q}
 \right\}.
\end{align}
For the first term, set
\begin{align}\label{eq:refined-pair-Ai}
 A_i(Y)\coloneqq
 \int_E\left|\sum_{j\ne i}h(x,Y_j)\right|^2\,\md\mu(x).
\end{align}
The \(L^2(\mu)\)-valued variables \(h(\cdot,Y_j)\) are independent and centered.  A Hilbert-space-valued Rosenthal inequality, for example the two-smooth Banach-space estimate in~\cite{Pinelis2015Rosenthal}, yields
\begin{align}\label{eq:refined-pair-Hilbert-Rosenthal}
 \norm{A_i^{1/2}}_{L^q(Y)}
 \le C_q\left(
 \sqrt N\norm{h}_{L^2(\mu^{\otimes2})}
 +N^{1/q}\norm{h}_{L^q(\mu^{\otimes2})}
 \right).
\end{align}
Here, the \(L^q\)-norm of \(y\mapsto\norm{h(\cdot,y)}_{L^2(\mu)}\) is bounded by \(\norm h_{L^q(\mu^{\otimes2})}\).  Minkowski's inequality in \(L^{q/2}\) therefore gives
\begin{align}\label{eq:refined-pair-first-term}
 \norm*{\pa*{\sum_iA_i}^{1/2}}_{L^q(Y)}
 \le C_q\left(
 N\norm{h}_{L^2(\mu^{\otimes2})}
 +N^{1/2+1/q}\norm{h}_{L^q(\mu^{\otimes2})}
 \right).
\end{align}
For the second term in \eqref{eq:refined-pair-conditional-Rosenthal}, condition on \(X_i\) and apply scalar Rosenthal inequality to the centered sum over \(j\).  Since
\begin{align}\label{eq:refined-pair-mixed-bound}
 \int_E\left(
 \int_E|h(x,y)|^2\,\md\mu(y)
 \right)^{q/2}\md\mu(x)
 \le \norm{h}_{L^q(\mu^{\otimes2})}^q,
\end{align}
we obtain
\begin{align}\label{eq:refined-pair-second-term}
 \pa*{\sum_i\E\left[|Z_i|^q\right]}^{1/q}
 \le C_qN^{1/2+1/q}\norm{h}_{L^q(\mu^{\otimes2})}.
\end{align}
Combining the preceding estimates, restoring the decoupling constant, and dividing by \(N(N-1)\) proves \eqref{eq:refined-canonical-pair-moment}.
\end{proof}

\begin{prop}[Exchangeable fixed-marginal \(L^\tau\) estimate]
\label{prop:abstract-fixed-marginal-Ltau}
Fix \(1<\tau<2\) and set \(q=\tau/(\tau-1)\).  For each \(N\), let \(F_N\) be an exchangeable probability density with respect to \(\mu^{\otimes N}\), with labeled marginal density \(F_N^{(k)}\).  Suppose that
\begin{align}\label{eq:abstract-fixed-marginal-assumptions}
 \sup_{N\ge1}\norm{F_N}_{L^\tau(\mu^{\otimes N})}\le M_\tau,
 \qquad
 \norm{F_N^{(1)}-1}_{L^2(\mu)}\le\frac{U}{N},
 \quad N\ge2.
\end{align}
Then, for every \(N\ge2\) and \(2\le k\le N\),
\begin{align}\label{eq:abstract-fixed-marginal-explicit}
 \norm{F_N^{(k)}-1}_{L^\tau(\mu^{\otimes k})}
 &\le \frac{U\sqrt{k}}{N}
 +(M_\tau+1)\mathsf B_q
 \left[
 \frac{\sqrt{\binom{k}{2}}}{N-1}
 +\frac{4\binom{k}{2}N^{1/q-1/2}}{N-1}
 \right]\notag\\
 &\quad +(M_\tau+1)\sum_{j=3}^k\binom{k}{j}
 \frac{2^j\mathsf A_{j,q}}{\sqrt{(N)_j}}.
\end{align}
The empty sum is zero.  For \(k=1\),
\begin{align}\label{eq:abstract-fixed-marginal-one-point}
 \norm{F_N^{(1)}-1}_{L^\tau(\mu)}\le\frac{U}{N}.
\end{align}
For every fixed \(k\ge2\),
\begin{align}\label{eq:abstract-fixed-marginal-limsup}
 \limsup_{N\to\infty}
 N\norm{F_N^{(k)}-1}_{L^\tau(\mu^{\otimes k})}
 \le U\sqrt{k}+(M_\tau+1)\mathsf B_q\sqrt{\binom{k}{2}}.
\end{align}
In particular, there is a finite constant \(C=C(\tau,M_\tau,U)\), independent of \(k\), such that, for every fixed \(k\), there is an onset \(N_0=N_0(k,\tau,M_\tau,U)\) for which
\begin{align}\label{eq:abstract-fixed-marginal-linear-k}
 \norm{F_N^{(k)}-1}_{L^\tau(\mu^{\otimes k})}
 \le C\frac{k}{N},
 \qquad N\ge N_0.
\end{align}
\end{prop}

\begin{proof}
By \(L^\tau\)--\(L^q\) duality, pair \(F_N^{(k)}-1\) against \(\varphi\in L^q(\mu^{\otimes k})\) with \(\norm\varphi_q\le1\).  Permutation averaging allows us to assume that \(\varphi\) is symmetric.  Its Hoeffding decomposition has the form
\begin{align}\label{eq:test-Hoeffding-decomposition}
 \varphi
 =\E_{\mu^{\otimes k}}\left[\varphi\right]
 +\sum_{j=1}^k\sum_{\substack{A\subset\{1,\ldots,k\}\\|A|=j}}
 \varphi_j(x_A),
\end{align}
with symmetric canonical kernels \(\varphi_j\).  Inclusion--exclusion and contractivity of conditional expectation give
\begin{align}\label{eq:test-Hoeffding-Lq-bounds}
 \norm{\varphi_j}_{L^q(\mu^{\otimes j})}\le2^j.
\end{align}
Since \(\norm\varphi_{L^2}\le1\), orthogonality also gives
\begin{align}\label{eq:test-Hoeffding-L2-bounds}
 \norm{\varphi_1}_{L^2(\mu)}\le k^{-1/2},
 \qquad
 \norm{\varphi_2}_{L^2(\mu^{\otimes2})}
 \le \binom{k}{2}^{-1/2}.
\end{align}
Marginal consistency and exchangeability yield
\begin{align}\label{eq:fixed-marginal-pairing-levels}
 \int_{E^k}\varphi(F_N^{(k)}-1)\,\md\mu^{\otimes k}
 =\sum_{j=1}^k\binom{k}{j}
 \E_{F_N\mu^{\otimes N}}\left[\varphi_j(X_1,\ldots,X_j)\right].
\end{align}
The first level is bounded by \(U\sqrt{k}/N\).  For \(j\ge2\), exchangeability and canonicity give
\begin{align}\label{eq:fixed-marginal-U-duality}
 \E_{F_N\mu^{\otimes N}}\left[\varphi_j(X_1,\ldots,X_j)\right]
 =\int_{E^N}(F_N-1)\mathcal U_{N,j}\varphi_j\,\md\mu^{\otimes N},
\end{align}
and hence
\begin{align}\label{eq:fixed-marginal-Holder-U}
 \left|\E_{F_N\mu^{\otimes N}}\left[\varphi_j(X_1,\ldots,X_j)\right]\right|
 \le(M_\tau+1)\norm{\mathcal U_{N,j}\varphi_j}_{L^q(\mu^{\otimes N})}.
\end{align}
For \(j=2\), apply \Cref{lem:refined-canonical-pair-moment} together with \eqref{eq:test-Hoeffding-Lq-bounds}--\eqref{eq:test-Hoeffding-L2-bounds}.  For \(j\ge3\), apply \Cref{lem:fixed-order-canonical-U-moment} and \eqref{eq:test-Hoeffding-Lq-bounds}.  Taking the supremum over \(\varphi\) proves \eqref{eq:abstract-fixed-marginal-explicit}.

For fixed \(k\), the second pair term is \(o(N^{-1})\) because \(1/q-1/2<0\), and every term of degree \(j\ge3\) is \(O_{k,\tau}(N^{-j/2})\).  This proves \eqref{eq:abstract-fixed-marginal-limsup}.  The last assertion follows from \(\sqrt{k}\le k\) and \(\sqrt{\binom{k}{2}}\le k/\sqrt2\); the onset may depend on \(k\).
\end{proof}

\begin{proof}[Proof of \Cref{thm:fixed-marginal-Ltau}]
By \Cref{lem:full-density-Lr} with exponent \(\tau\),
\begin{align}\label{eq:fixed-marginal-full-Gibbs-Ltau}
 \sup_{N\ge1}\norm{f_{N,\be}}_{L^\tau(\mu^{\otimes N})}
 \le \mathsf M_{\tau,\be,\mu}.
\end{align}
Moreover, \Cref{lem:normalized-partition-interpolation} with \(s=r\) and \(t=\tau\) gives
\begin{align}\label{eq:fixed-marginal-rbeta-input}
 \mathfrak P_{r\be}(\mu)^{1/r}
 \le \mathfrak P_{\tau\be}(\mu)^{1/\tau}.
\end{align}
Thus, the bounded cavity proof and the sufficient weighted logarithmic proof have exactly the required full-density \(L^r\) input, and the literal hypothesis \(\mathfrak P_{r\be}(\mu)<\infty\) in \Cref{lem:weighted-logarithmic-cavity} is satisfied.  No partition-function moment at \(2\be\) is used.  Applying the one-point argument in \Cref{sec:one-point-upper-estimate} with \(\theta=\tau\) gives
\begin{align}\label{eq:fixed-marginal-one-point-L2}
 \norm{f_{N,\be}^{(1)}-1}_{L^2(\mu)}
 \le\frac{\mathsf U_{\tau,\be,\mu}}{N}.
\end{align}
Now apply \Cref{prop:abstract-fixed-marginal-Ltau} with \(M_\tau=\mathsf M_{\tau,\be,\mu}\) and \(U=\mathsf U_{\tau,\be,\mu}\).  This gives \eqref{eq:fixed-marginal-Gibbs-explicit}--\eqref{eq:fixed-marginal-Gibbs-limsup} and the \(L^\tau\) line of \eqref{eq:fixed-marginal-Gibbs-rates}.  Since \(\mu^{\otimes k}\) is a probability measure,
\begin{align}\label{eq:fixed-marginal-TV-from-Ltau}
 2\TV\pa{\Q_{N,\be}^{(k)}(\mu),\mu^{\otimes k}}
 =\norm{f_{N,\be}^{(k)}-1}_{L^1(\mu^{\otimes k})}
 \le\norm{f_{N,\be}^{(k)}-1}_{L^\tau(\mu^{\otimes k})},
\end{align}
which proves the total-variation line.

The first line of \eqref{eq:fixed-marginal-Gibbs-divergences} is the \(\tau\)-th power of the \(L^\tau\) estimate.  Finally, there is a finite constant \(C_\tau\) such that
\begin{align}\label{eq:entropy-Ltau-global-comparison}
 t\log t-t+1\le C_\tau|t-1|^\tau,
 \qquad t\ge0.
\end{align}
Indeed, the quotient of the left-hand side by \(|t-1|^\tau\) extends continuously at \(t=1\) with value zero, is finite at \(t=0\), and tends to zero as \(t\to\infty\).  Integrating \eqref{eq:entropy-Ltau-global-comparison} proves the second line of \eqref{eq:fixed-marginal-Gibbs-divergences}.
\end{proof}

% STAGE2 MOVE: sec:chi2-tv-lower-bound-proof
\section{Matching lower bounds}
\label{sec:chi2-tv-lower-bound-proof}

% STAGE2 MOVE: A2-lower-bound-roadmap
We prove the lower estimates in \Cref{thm:chi2-chaos-two-sided} under that theorem's own hypotheses, without using the quantitative marginal theorem.  On the common hypotheses with a nontrivial interaction and $\be>0$, the quantitative estimates give another deduction: the two-particle limit covers finitely many small marginal sizes, direct subtraction gives a lower bound up to a sufficiently small fixed fraction of $N$, and contraction under marginalization covers every larger size.  The resulting lower bound holds for every $2\leq k\leq N$ once $N\geq N_0$, with $N_0$ independent of $k$.  This deduction would, however, impose the quantitative theorem's row, source, and amplified centered-pair moment assumptions.  The direct argument below preserves the distinct order criterion, including its weighted logarithmic lower-tail alternative.  In \Cref{sec:strict-product-law-jensen-gap}, a strict Jensen gap and convexity give a lower bound for the derivative of the product-law log-Laplace transform.  \Cref{sec:pair-statistic-lower-bound} uses this derivative to estimate the pair response and tests with the pair sum.  \Cref{sec:bounded-two-body-test-tv-lower-bound} obtains total variation by a bounded truncation, and then relative entropy by Pinsker's inequality.

% STAGE2 MOVE: A2-lower-proof-unchanged
\subsection{Strict product-law Jensen gap}
\label{sec:strict-product-law-jensen-gap}

Recall the centered energy $\mathcal X_N$ from \eqref{eq:spectral-centered-energy}.
For \(\XN\sim\mu^{\otimes N}\), one has \(\E\left[\mathcal X_N\right]=0\), and
\begin{align}\label{eq:KN-GN}
 \K_{N,\be}(\mu)
 =e^{\be I_\mu}\E_{\mu^{\otimes N}}\left[e^{-\be \mathcal X_N}\right].
\end{align}
Denote the log-Laplace transform of \(\mathcal X_N\), equivalently the cumulant generating function of \(-\mathcal X_N\), by
\begin{align}\label{eq:GN-def}
 G_N(\alpha)\coloneqq \log \E_{\mu^{\otimes N}}\left[e^{-\alpha \mathcal X_N}\right],
 \qquad \alpha\ge0.
\end{align}


\begin{lemma}[Finiteness and differentiability of the log partition function]\label{lem:GN-regularity}
Fix \(\be>0\) and assume \(\K_{N,2\be}(\mu)<\infty\).  Then \(G_N\) is finite on \([0,2\be]\), is continuously differentiable on \((0,2\be)\), and
\begin{align}\label{eq:GN-derivative-justified}
 G_N'(\alpha)
 =-\frac{\E_{\mu^{\otimes N}}\left[\mathcal X_Ne^{-\alpha\mathcal X_N}\right]}
 {\E_{\mu^{\otimes N}}\left[e^{-\alpha\mathcal X_N}\right]}
 =-\E_{\Q_{N,\alpha}(\mu)}[\mathcal X_N],
 \qquad 0<\alpha<2\be.
\end{align}
\end{lemma}

\begin{proof}
The identity \eqref{eq:KN-GN} at \(2\be\) gives exponential integrability at the right endpoint.  H\"older's inequality and \(G_N(0)=0\) then imply finiteness throughout \([0,2\be]\).

Write \(J=[a,b]\Subset(0,2\be)\), and choose \(\epsilon>0\) so that
\(a-\epsilon>0\) and \(b+\epsilon<2\be\).  For every \(\alpha\in J\),
\begin{align}\label{eq:GN-differentiation-dominator}
 |x|e^{-\alpha x}
 \le C_\epsilon\left(e^{-(a-\epsilon)x}+e^{-(b+\epsilon)x}\right),
 \qquad x\in\R.
\end{align}
After increasing \(C_\epsilon\) if necessary, the right-hand side of
\eqref{eq:GN-differentiation-dominator} is integrable by finiteness of \(G_N\)
on \([0,2\be]\), and it is independent of \(\alpha\in J\).  Differentiation
under the integral sign, as well as continuity of the derivative on \(J\),
now follows by dominated convergence.

The preceding gives the first identity.  Since
\begin{align}\label{eq:GN-modulated-exponent-identity}
 e^{-\alpha N\Fr_N}=e^{\alpha I_\mu}e^{-\alpha\mathcal X_N},
\end{align}
the second identity follows from the definition of \(\Q_{N,\alpha}(\mu)\).
\end{proof}


We use \Cref{lem:fixed-order-canonical-U-moment} with \(j=2\) to control the pair part of \(\mathcal X_N\).  The resulting moment estimate serves only to make the Jensen gap uniform in \(N\): the variance gives a fixed nonzero \(L^2\) mass, while the \(L^p\) bound prevents that mass from escaping to rare large values.

\begin{lemma}[Uniform product-law nondegeneracy]\label{lem:product-nondegenerate}
Let \(p>2\) be fixed by \eqref{eq:p-choice-quantitative} under the integrability assumptions \eqref{eq:mu-controls}, and set
\begin{align}\label{eq:Mp-mu-def}
 \mathsf M_{p,\mu}
 \coloneqq
 \pa{\frac{\mathsf A_{2,p}}2\norm{\tl W}_{L^p(\mu^{\otimes2})}
 +\norm{h_\mu-2I_\mu}_{L^p(\mu)}}^p .
\end{align}
For all \(N\ge2\),
\begin{align}\label{eq:calX-var-lower}
 \E_{X_N\sim\mu^{\otimes N}}\left[\mathcal X_N^2\right]
 \ge \frac14\norm{\tl W}_{L^2(\mu^{\otimes2})}^2,
\end{align}
and
\begin{align}\label{eq:calX-Lp-upper}
 \E_{X_N\sim\mu^{\otimes N}}\left[|\mathcal X_N|^p\right]\le \mathsf M_{p,\mu}.
\end{align}
\end{lemma}

\begin{proof}
Write
\begin{align}
 A_N&\coloneqq \frac1N\sum_{1\le i<j\le N}\tl W(x_i,x_j),
 \label{eq:AN-product-def}\\
 B_N&\coloneqq \frac1N\sum_{i=1}^N\pa{h_\mu(x_i)-2I_\mu}.
 \label{eq:BN-product-def}
\end{align}
so that \(\mathcal X_N=A_N-B_N\).  Since \(\tl W\) is centered in each variable,
\begin{align}\label{eq:AN-product-moments}
 \E\left[A_N\right]=0,
 \qquad
 \E\left[A_N^2\right]=\frac{\binom{N}{2}}{N^2}\norm{\tl W}_{L^2(\mu^{\otimes2})}^2
 =\frac{N-1}{2N}\norm{\tl W}_{L^2(\mu^{\otimes2})}^2.
\end{align}
Moreover, \(\E\left[A_NB_N\right]=0\), again by the centering of \(\tl W\).  Therefore,
\begin{align}\label{eq:calX-second-lower}
 \E\left[\mathcal X_N^2\right]=\E\left[A_N^2\right]+\E\left[B_N^2\right]
 \ge \frac14\norm{\tl W}_{L^2(\mu^{\otimes2})}^2>0,
 \qquad \forall N\ge2,
\end{align}
where positivity is exactly the standing nondegeneracy input
\(\norm{\tl W}_{L^2(\mu^{\otimes2})}>0\).

Recall that \(p>2\) was fixed in \eqref{eq:p-choice-quantitative} so that the integrability assumptions in \eqref{eq:mu-controls} hold.  Since \(\tl W\) is symmetric and centered in each variable,
\begin{align}\label{eq:AN-as-ordered-U-statistic}
 A_N=\frac{N-1}{2}\mathcal U_{N,2}\tl W.
\end{align}
Therefore, \Cref{lem:fixed-order-canonical-U-moment} and Minkowski's inequality give
\begin{align}
 \norm{A_N}_{L^p(\mu^{\otimes N})}
 &\le \frac{\mathsf A_{2,p}}2\sqrt{\frac{N-1}{N}}
 \norm{\tl W}_{L^p(\mu^{\otimes2})}
 \le \frac{\mathsf A_{2,p}}2\norm{\tl W}_{L^p(\mu^{\otimes2})},
 \label{eq:AN-Lp-fixed-order}\\
 \norm{B_N}_{L^p(\mu^{\otimes N})}
 &\le \norm{h_\mu-2I_\mu}_{L^p(\mu)}.
 \label{eq:BN-Lp-bound}
\end{align}
Thus, by the triangle inequality,
\begin{align}\label{eq:calX-Lp-from-components}
 \norm{\mathcal X_N}_{L^p(\mu^{\otimes N})}
 \le \frac{\mathsf A_{2,p}}2\norm{\tl W}_{L^p(\mu^{\otimes2})}
 +\norm{h_\mu-2I_\mu}_{L^p(\mu)},
\end{align}
which is \eqref{eq:calX-Lp-upper}.
\end{proof}

The next lemma quantifies the gap in Jensen's inequality for the product law
\(\mu^{\otimes N}\) in terms of the second moment of \(\mathcal X_N\) and the
local curvature of the exponential function.  For \(\be>0\), define
\begin{align}
 \mathsf R_{p,\mu}
 &\coloneqq
 \max\left\{1,
 \left(\frac{8\mathsf M_{p,\mu}}
 {\norm{\tl W}_{L^2(\mu^{\otimes2})}^2}\right)^{1/(p-2)}\right\},
 \label{eq:jensen-cutoff-def}\\
 \mathsf a_{\be,\mu}
 &\coloneqq
 \inf_{0<|t|\le\mathsf R_{p,\mu}}
 \frac{e^{-\be t}-1+\be t}{t^2},
 \label{eq:jensen-curvature-def}\\
 \mathsf J_{\be,\mu}
 &\coloneqq
 \log\left(1+\frac18\mathsf a_{\be,\mu}
 \norm{\tl W}_{L^2(\mu^{\otimes2})}^2\right).
 \label{eq:jensen-gap-constant-def}
\end{align}

\begin{lemma}[Strict Jensen gap]\label{lem:jensen-gap}
For every \(\be>0\),
\begin{align}\label{eq:jensen-gap}
 G_N(\be)\ge \mathsf J_{\be,\mu},
 \qquad N\ge2.
\end{align}
Consequently,
\begin{align}\label{eq:GN-derivative-lower}
 G_N'(\be)\ge \frac{\mathsf J_{\be,\mu}}{\be},
 \qquad N\ge2.
\end{align}
\end{lemma}

\begin{proof}
By the \(L^p\) estimate \eqref{eq:calX-Lp-upper} in
\Cref{lem:product-nondegenerate} and the definition
\eqref{eq:jensen-cutoff-def} of \(\mathsf R_{p,\mu}\),
\begin{align}\label{eq:calX-tail-second-moment}
 \E\left[\mathcal X_N^2\one_{\{|\mathcal X_N|>\mathsf R_{p,\mu}\}}\right]
 \le \mathsf R_{p,\mu}^{2-p}\E\left[|\mathcal X_N|^p\right]
 \le \frac18\norm{\tl W}_{L^2(\mu^{\otimes2})}^2,
 \qquad N\ge2.
\end{align}
Combining this with the variance lower bound \eqref{eq:calX-var-lower} from
\Cref{lem:product-nondegenerate},
\begin{align}
 \E\left[\mathcal X_N^2
 \one_{\{|\mathcal X_N|\le\mathsf R_{p,\mu}\}}\right]
 \ge \frac18\norm{\tl W}_{L^2(\mu^{\otimes2})}^2.
 \label{eq:calX-truncated-var}
\end{align}
Since
\begin{align}\label{eq:jensen-phi-def}
 \varphi_\be(t)\coloneqq e^{-\be t}-1+\be t \ge 0,
 \qquad \forall t\in\R,
\end{align}
the definition \eqref{eq:jensen-curvature-def} of \(\mathsf a_{\be,\mu}\),
the truncated variance bound \eqref{eq:calX-truncated-var}, and
\(\E\left[\mathcal X_N\right]=0\) give
\begin{align}\label{eq:jensen-exponential-lower}
 \E\left[e^{-\be \mathcal X_N}\right]-1
 =\E\left[\varphi_\be(\mathcal X_N)\right]
 \ge \frac18\mathsf a_{\be,\mu}\norm{\tl W}_{L^2(\mu^{\otimes2})}^2.
\end{align}
Therefore,
\begin{align}\label{eq:GN-gap-proof}
 G_N(\be)=\log\E\left[e^{-\be \mathcal X_N}\right]
 \ge \log\pa{1+\frac18\mathsf a_{\be,\mu}
 \norm{\tl W}_{L^2(\mu^{\otimes2})}^2}=\mathsf J_{\be,\mu}.
\end{align}
By \Cref{lem:GN-regularity}, \(G_N\) is differentiable at \(\be\).  Since it
is convex and \(G_N(0)=0\),
\begin{align}\label{eq:GN-convex-derivative-bound}
 G_N'(\be)\ge \frac{G_N(\be)-G_N(0)}{\be}
 \ge \frac{\mathsf J_{\be,\mu}}{\be},
\end{align}
which proves \eqref{eq:GN-derivative-lower}.
\end{proof}


\subsection{Pair statistic and the chi-square lower bound}
\label{sec:pair-statistic-lower-bound}

We now convert the Jensen gap into a marginal lower bound.  Set
\begin{align}\label{eq:mN-def}
 m_N\coloneqq \E_{\Xi_N\sim\Q_{N,\be}(\mu)}
 \left[\tl W(\xi_1,\xi_2)\right],
 \qquad \forall N\ge2.
\end{align}
Thus, \(m_N\) records the response of one Hoeffding-centered pair interaction to the Gibbs tilt.  To detect this response at the level of $k$ observed labels, set
\begin{align}\label{eq:intro-centered-pair-statistic}
 A_k(x_1,\ldots,x_k)
 \coloneqq \sum_{1\le i<j\le k}\tl W(x_i,x_j).
\end{align}
Because \(\tl W\) is centered in each variable, distinct pair summands are orthogonal under \(\mu^{\otimes k}\), including when the two pairs share one index.  Consequently,
\begin{align}\label{eq:intro-pair-orthogonality}
 \norm{A_k}_{L^2(\mu^{\otimes k})}^2
 =\binom{k}{2}\norm{\tl W}_{L^2(\mu^{\otimes2})}^2.
\end{align}

By exchangeability,

\begin{align}\label{eq:Ak-expectation}
 \int A_k(f_{N,\be}^{(k)}-1)\,\md\mu^{\otimes k}
 =\E_{\Q_{N,\be}^{(k)}(\mu)}\left[A_k\right]
 =\binom{k}{2}m_N.
\end{align}
Cauchy--Schwarz inequality gives
\begin{align}
 \Chi\pa{\Q_{N,\be}^{(k)}(\mu)\mid\mu^{\otimes k}}
 &=\norm{f_{N,\be}^{(k)}-1}_{L^2(\mu^{\otimes k})}^2\notag\\
 &\ge
 \frac{\pa*{\int A_k(f_{N,\be}^{(k)}-1)\,\md\mu^{\otimes k}}^2}
 {\norm{A_k}_{L^2(\mu^{\otimes k})}^2}.
 \label{eq:chi2-duality-test}
\end{align}
Substituting \eqref{eq:intro-pair-orthogonality} and \eqref{eq:Ak-expectation}, we obtain
\begin{align}\label{eq:chi-lower-from-mN}
 \Chi\pa{\Q_{N,\be}^{(k)}(\mu)\mid\mu^{\otimes k}}
 \ge \binom{k}{2}
 \frac{m_N^2}{\norm{\tl W}_{L^2(\mu^{\otimes2})}^2}.
\end{align}
Thus, the chi-square lower bound of order \(k^2/N^2\) follows once we prove
that the Gibbs expectation \(m_N\) is negative and has magnitude of order
\(N^{-1}\).  The next lemma derives exactly this estimate from the strict
product-law Jensen gap.  Define the threshold
\begin{align}\label{eq:N-beta-mu-threshold}
 \mathsf N_{\be,\mu}
 \coloneqq
 \left\lceil \max\left\{2,
 \frac{2\be\norm{h_\mu-2I_\mu}_{L^2(\mu)}
       \mathsf U_{\be,\mu}}{\mathsf J_{\be,\mu}}\right\}\right\rceil .
\end{align}

\begin{lemma}[Negative pair response]\label{lem:pair-response-lower}
For every \(\be>0\),
\begin{align}\label{eq:mN-negative}
 m_N\le -\frac{\mathsf J_{\be,\mu}}{\be N},
 \qquad \forall N\ge \mathsf N_{\be,\mu}.
\end{align}
\end{lemma}

\begin{proof}
By \eqref{eq:KN-GN},
\begin{align}\label{eq:log-partition-GN}
 \log\K_{N,\be}(\mu)=\be I_\mu+G_N(\be).
\end{align}
By \Cref{lem:GN-regularity}, differentiating in \(\be\) is justified and gives
\begin{align}\label{eq:GN-prime-energy}
 G_N'(\be)
 =-\E_{\Q_{N,\be}(\mu)}\left[N\Fr_N(\XN,\mu)+I_\mu\right].
\end{align}
Using the definition of \(\Fr_N\) and exchangeability,
\begin{align}\label{eq:energy-expectation-mN}
 \E_{\Q_{N,\be}(\mu)}\left[N\Fr_N(\XN,\mu)\right]
 =\frac{N-1}{2}m_N-
 \E_{\Q_{N,\be}^{(1)}(\mu)}\left[h_\mu\right]+I_\mu.
\end{align}
Combining \eqref{eq:GN-prime-energy} and \eqref{eq:energy-expectation-mN},
\begin{align}\label{eq:GN-prime-mN}
 G_N'(\be)
 =-\frac{N-1}{2}m_N+
 \pa{\E_{\Q_{N,\be}^{(1)}(\mu)}\left[h_\mu\right]-2I_\mu}.
\end{align}
Cauchy--Schwarz inequality and the one-point upper bound
\eqref{eq:one-point-L2-bound} give
\begin{align}
 \left|\E_{\Q_{N,\be}^{(1)}(\mu)}\left[h_\mu\right]-2I_\mu\right|
 &=\left|\int (h_\mu-2I_\mu)(f_{N,\be}^{(1)}-1)\,\md\mu\right|\notag\\
 &\le \norm{h_\mu-2I_\mu}_{L^2(\mu)}
      \frac{\mathsf U_{\be,\mu}}{N}.
 \label{eq:h-one-point-error}
\end{align}
By \Cref{lem:jensen-gap},
\begin{align}
 -\frac{N-1}{2}m_N
 &=G_N'(\be)-\pa{\E_{\Q_{N,\be}^{(1)}(\mu)}\left[h_\mu\right]-2I_\mu}\notag\\
 &\ge \frac{\mathsf J_{\be,\mu}}{\be}
      -\frac{\norm{h_\mu-2I_\mu}_{L^2(\mu)}
              \mathsf U_{\be,\mu}}{N}.
 \label{eq:mN-preliminary-lower}
\end{align}
If \(N\ge\mathsf N_{\be,\mu}\), the right-hand side is at least
\(\mathsf J_{\be,\mu}/(2\be)\).  Therefore,
\begin{align}\label{eq:mN-negative-conclusion}
 m_N\le -\frac{\mathsf J_{\be,\mu}}{\be(N-1)}
 \le -\frac{\mathsf J_{\be,\mu}}{\be N},
\end{align}
which proves \eqref{eq:mN-negative}.
\end{proof}

Combining \eqref{eq:chi-lower-from-mN} with
\Cref{lem:pair-response-lower}, we obtain
\begin{align}\label{eq:direct-chi-lower-quantitative}
 \Chi\pa{\Q_{N,\be}^{(k)}(\mu)\mid\mu^{\otimes k}}
 \ge
 \frac{\mathsf J_{\be,\mu}^2}
 {\be^2\norm{\tl W}_{L^2(\mu^{\otimes2})}^2}
 \frac{\binom{k}{2}}{N^2},
\end{align}
for every \(N\ge\mathsf N_{\be,\mu}\) and \(2\le k\le N\).
The total-variation estimate proved below would also recover a chi-square
lower bound of the same order.  We present the direct argument because it
identifies the \(L^2\) pair-response mechanism without the additional losses
introduced by truncation and by the chi-square upper bound used to control
the discarded tail.


\subsection{Bounded two-body test and the total-variation lower bound}
\label{sec:bounded-two-body-test-tv-lower-bound}

The preceding lower bound tested the marginal against \(A_k\), which need not be
bounded.  This is harmless for \(\chi^2\), where the Cauchy--Schwarz inequality uses
the \(L^2\)-norm of \(A_k\).  For total variation, however, the test function
must be bounded.  Although \(A_k\) contains \(\binom{k}{2}\) summands, Hoeffding
centering yields an \(L^p(\mu^{\otimes k})\) norm of order \(k\), rather than
the order \(k^2\) given by the triangle inequality.  Hence, truncation at
height \(Rk\) discards only a controlled tail.

\begin{lemma}[Moment bound for the pair statistic]\label{lem:Ak-Lp-bound}
With \(p\) as in \eqref{eq:p-choice-quantitative},
\begin{align}\label{eq:Ak-Lp-bound}
 \norm{A_k}_{L^p(\mu^{\otimes k})}
 \le \frac{\mathsf A_{2,p}}2\norm{\tl W}_{L^p(\mu^{\otimes2})} k,
 \qquad \forall k\ge2.
\end{align}
Consequently, for every \(R\ge1\),
\begin{align}\label{eq:Ak-tail-L2}
 \norm{A_k\one_{\{|A_k|>Rk\}}}_{L^2(\mu^{\otimes k})}
 \le \pa{\frac{\mathsf A_{2,p}}2\norm{\tl W}_{L^p(\mu^{\otimes2})}}^{p/2}
 R^{1-p/2}k.
\end{align}
\end{lemma}

\begin{proof}
By the standing integrability assumptions in \eqref{eq:mu-controls}, \(\tl W\in L^p(\mu^{\otimes2})\).  By \eqref{eq:intro-centered-pair-statistic} and the symmetry of \(\tl W\),
\begin{align}\label{eq:Ak-as-ordered-U-statistic}
 A_k=\frac{k(k-1)}2\mathcal U_{k,2}\tl W.
\end{align}
Applying \Cref{lem:fixed-order-canonical-U-moment} gives \eqref{eq:Ak-Lp-bound}.  Finally,
\begin{align}\label{eq:Ak-tail-L2-squared}
 \norm{A_k\one_{\{|A_k|>Rk\}}}_{L^2}^2
 \le (Rk)^{2-p}\norm{A_k}_{L^p}^p
 \le \pa{\frac{\mathsf A_{2,p}}2\norm{\tl W}_{L^p(\mu^{\otimes2})}}^pR^{2-p}k^2,
\end{align}
and taking square roots gives \eqref{eq:Ak-tail-L2}.
\end{proof}

To convert this moment estimate into a total-variation lower bound, we truncate
\(A_k\) at height \(Rk\).  The preceding lemma controls the discarded tail,
while the truncated statistic is bounded and may therefore be used in the dual
characterization of total variation.  We now fix the constants used in this
argument:
\begin{align}
 \mathsf S_{\be,\mu}
 &\coloneqq \frac{\mathsf J_{\be,\mu}}{4\be},
 \label{eq:S-beta-mu-def}\\
 \mathsf R_{\TV,\be,\mu}
 &\coloneqq
 \max\left\{1,
 \left(\frac{
        2\pa{\frac{\mathsf A_{2,p}}2\norm{\tl W}_{L^p(\mu^{\otimes2})}}^{p/2}
        \sqrt{\mathsf C_{\chi,\be,\mu}^{\mathrm{up}}}}
       {\mathsf S_{\be,\mu}}
 \right)^{2/(p-2)}\right\},
 \label{eq:R-TV-def}\\
 \mathsf c_{\TV,\be,\mu}^{\mathrm{low}}
 &\coloneqq \frac{\mathsf S_{\be,\mu}}{4\mathsf R_{\TV,\be,\mu}} .
 \label{eq:c-TV-low-def}
\end{align}

\begin{prop}[Total-variation lower bound]\label{prop:TV-lower}
Let \(\be>0\).  For all \(N\ge \mathsf N_{\be,\mu}\) and all \(2\le k\le N\),
\begin{align}\label{eq:TV-lower-final}
 \TV\pa{\Q_{N,\be}^{(k)}(\mu),\mu^{\otimes k}}
 \ge \mathsf c_{\TV,\be,\mu}^{\mathrm{low}}\frac{k}{N}.
\end{align}
\end{prop}

\begin{proof}
Recall that
\(f_{N,\be}^{(k)}=\md\Q_{N,\be}^{(k)}(\mu)/\md\mu^{\otimes k}\).
By the chi-square upper bound in \eqref{eq:unified-upper-final}, proved in
\Cref{sec:hoeffding-projection-upper-bound}, with the admissible constant
\(\mathsf C_{\chi,\be,\mu}^{\mathrm{up}}\) fixed in
\eqref{eq:chi-upper-constant},
\begin{align}\label{eq:f-L2-upper-for-TV}
 \norm{f_{N,\be}^{(k)}-1}_{L^2(\mu^{\otimes k})}
 \le \sqrt{\mathsf C_{\chi,\be,\mu}^{\mathrm{up}}}\frac{k}{N}.
\end{align}
By \eqref{eq:Ak-expectation} and \Cref{lem:pair-response-lower},
\begin{align}\label{eq:Ak-pair-response-TV}
 \left|\int A_k(f_{N,\be}^{(k)}-1)\,\md\mu^{\otimes k}\right|
 =\binom{k}{2}|m_N|
 \ge \mathsf S_{\be,\mu}\frac{k^2}{N},
 \qquad 2\le k\le N,
\end{align}
where we used \(\binom{k}{2}\ge k^2/4\).  By \Cref{lem:Ak-Lp-bound} and
\eqref{eq:f-L2-upper-for-TV},
\begin{align}
 \left|\int A_k\one_{\{|A_k|>Rk\}}(f_{N,\be}^{(k)}-1)\,\md\mu^{\otimes k}\right|
 &\le \norm{A_k\one_{\{|A_k|>Rk\}}}_{L^2}
      \norm{f_{N,\be}^{(k)}-1}_{L^2} \notag\\
 &\le
      \pa{\frac{\mathsf A_{2,p}}2\norm{\tl W}_{L^p(\mu^{\otimes2})}}^{p/2}
      \sqrt{\mathsf C_{\chi,\be,\mu}^{\mathrm{up}}}
      R^{1-p/2}\frac{k^2}{N}.\label{eq:tail-pair-response-error-TV}
\end{align}
With \(R=\mathsf R_{\TV,\be,\mu}\), the right side of \eqref{eq:tail-pair-response-error-TV} is at most
one half of the right side of \eqref{eq:Ak-pair-response-TV}.
Set
\begin{align}\label{eq:AkR-def}
 A_{k,R}\coloneqq A_k\one_{\{|A_k|\le Rk\}}.
\end{align}
Then
\begin{align}\label{eq:truncated-pair-response-TV}
 \left|\int A_{k,R}(f_{N,\be}^{(k)}-1)\,\md\mu^{\otimes k}\right|
 \ge \frac12\mathsf S_{\be,\mu}\frac{k^2}{N}.
\end{align}
On the other hand, \(\norm{A_{k,R}}_{L^\infty}\le Rk\) and
\(\norm{f_{N,\be}^{(k)}-1}_{L^1(\mu^{\otimes k})}=2\TV\pa{\Q_{N,\be}^{(k)}(\mu),\mu^{\otimes k}}\).  Therefore, the dual formulation of total variation gives
\begin{align}\label{eq:truncated-test-TV-duality}
 \left|\int A_{k,R}(f_{N,\be}^{(k)}-1)\,\md\mu^{\otimes k}\right|
 \le 2\norm{A_{k,R}}_{L^\infty}\TV\pa{\Q_{N,\be}^{(k)}(\mu),\mu^{\otimes k}}
 \le 2Rk\,\TV\pa{\Q_{N,\be}^{(k)}(\mu),\mu^{\otimes k}}.
\end{align}
Combining this with \eqref{eq:truncated-pair-response-TV} and the definition of
\(\mathsf c_{\TV,\be,\mu}^{\mathrm{low}}\) proves \eqref{eq:TV-lower-final}.
\end{proof}

We now finish the proof of \Cref{thm:chi2-chaos-two-sided}.

\begin{proof}[Proof of \Cref{thm:chi2-chaos-two-sided}]
The chi-square upper bound was proved in
\Cref{sec:hoeffding-projection-upper-bound}.  The remaining upper bounds follow
from the standard inequalities
\begin{align}
 \Hr(\nu\mid\pi)&\le \Chi(\nu\mid\pi),
 \label{eq:entropy-upper-by-chi}\\
 \TV(\nu,\pi)&\le \frac12\Chi(\nu\mid\pi)^{1/2}.
 \label{eq:TV-upper-by-chi}
\end{align}
Thus, admissible upper constants are
\begin{align}
 \mathsf C_{\TV,\be,\mu}^{\mathrm{up}}
 &\coloneqq \frac12\sqrt{\mathsf C_{\chi,\be,\mu}^{\mathrm{up}}},
 \label{eq:C-TV-upper-def}\\
 \mathsf C_{\Hr,\be,\mu}^{\mathrm{up}}
 &\coloneqq \mathsf C_{\chi,\be,\mu}^{\mathrm{up}}.
 \label{eq:C-H-upper-def}
\end{align}
Assume now \(\be>0\).  For the lower bounds,
\eqref{eq:direct-chi-lower-quantitative} and
\(\binom{k}{2}\ge k^2/4\), for \(k\ge2\), give the chi-square estimate in
\eqref{eq:unified-lower-chi} with
\begin{align}\label{eq:c-chi-low-direct-def}
 \mathsf c_{\chi,\be,\mu}^{\mathrm{low}}
 \coloneqq
 \frac{\mathsf J_{\be,\mu}^2}
 {4\be^2\norm{\tl W}_{L^2(\mu^{\otimes2})}^2}.
\end{align}
The total-variation lower bound is \Cref{prop:TV-lower}.  Finally, Pinsker's
inequality,
\begin{align}\label{eq:pinsker-lower-entropy}
 \Hr(\nu\mid\pi)\ge 2\TV(\nu,\pi)^2,
\end{align}
gives the entropy lower bound with
\begin{align}\label{eq:c-H-low-def}
 \mathsf c_{\Hr,\be,\mu}^{\mathrm{low}}
 \coloneqq 2\pa{\mathsf c_{\TV,\be,\mu}^{\mathrm{low}}}^2.
\end{align}
This proves all assertions of the theorem.
\end{proof}

% STAGE2 MOVE: sec:second-hoeffding-component
\section{Ultraviolet obstructions}
\label{sec:uv-sharpness-counterexamples}

We now record whole-space Riesz examples showing that \(\tl W\in L^2(\mu^{\otimes2})\) is the exact ultraviolet threshold for the classical fixed-particle chi-square and entropy scales, and that the uniform \(k/N\) total-variation bound fails throughout \(\d/2\le\s<\d\).  After introducing the common Riesz estimates, the two-point density representation, and the ultraviolet scale in \Cref{subsec:uv-common-setup}, we prove the all-temperature two-point chi-square lower bounds in \Cref{subsec:uv-fixed-particle-chi2}.  \Cref{subsec:uv-fixed-particle-consequences} transfers these estimates to every fixed number of labels and to relative entropy and total variation, thereby proving \Cref{thm:intro-ultraviolet-sharpness}\textup{(i)}.  Finally, \Cref{subsec:uv-mesoscopic} proves the two mesoscopic assertions in \Cref{thm:intro-ultraviolet-sharpness}\textup{(ii)}: the strictly supercritical argument uses a union of depleted close-pair events, whereas the critical argument uses sequential Hellinger accumulation across transition annuli.

\subsection{Common setup and the two-point density}
\label{subsec:uv-common-setup}

Throughout this section,
\begin{align}\label{eq:uv-standing-riesz-range}
 \frac{\d}{2}\le \s<\d,
 \qquad
 W(x,y)=\g_\s(x-y)=\frac1\s\abs{x-y}^{-\s},
\end{align}
and \(\mu\) is a probability measure on \(\R^\d\) with bounded Lebesgue density.

The Riesz potential is bounded without any compact-support assumption.  Indeed, splitting at an arbitrary radius and optimizing that radius gives
\begin{align}\label{eq:uv-Hmu-bound}
 H_\mu
 \coloneqq \norm{h_\mu}_{L^\infty}
 &\le \frac1\s\inf_{R>0}
 \left\{
 \frac{\abs{\mathbb S^{\d-1}}}{\d-\s}
 \norm{\mu}_{L^\infty}R^{\d-\s}+R^{-\s}
 \right\}
 \le C_{\d,\s}\norm{\mu}_{L^\infty}^{\s/\d}.
\end{align}
Since \(W\ge0\), it follows that
\begin{align}
 0\le 2I_\mu=\int h_\mu\,\md\mu&\le H_\mu,
 \label{eq:uv-Imu-bound}\\
 \norm{h_\mu-2I_\mu}_{L^\infty}&\le H_\mu,
 \label{eq:uv-one-body-bound}\\
 \tl W(x,y)&\ge-2H_\mu
 \quad\text{for \(\mu^{\otimes2}\)-a.e. \((x,y)\)}.
 \label{eq:uv-centered-lower-bound}
\end{align}

The same bounds make the failure of the Hilbert--Schmidt condition immediate.  Near the diagonal,
\begin{align}\label{eq:uv-centered-diagonal-asymptotic}
 \tl W(x,y)=\frac1\s\abs{x-y}^{-\s}+O_\mu(1).
\end{align}
Choose a bounded Borel set \(A\subset\R^\d\) of positive Lebesgue measure and \(a>0\) such that \(\mu\ge a\) almost everywhere on \(A\).  Since \(\abs{A\cap(A-z)}\to\abs A\) as \(z\to0\), for a sufficiently small \(r_0>0\),
\begin{align}\label{eq:uv-non-HS-integral}
 \int_{(\R^\d)^2}\abs{\tl W(x,y)}^2\,\md\mu(x)\md\mu(y)
 &\ge c_{\mu,\s}
 \int_{\abs z<r_0}\abs z^{-2\s}\abs{A\cap(A-z)}\,\md z
 =+\infty,
\end{align}
where the last equality holds exactly because \(2\s\ge\d\).

We next specialize the marked representation of \Cref{sec:exact-marked-identities} to two observed particles.  For \(N\ge3\), define the residual two-point prefactor by
\begin{align}\label{eq:uv-residual-prefactor-def}
 \mathcal R_{N,\be}(x,y)
 \coloneqq
 \frac{
 \exp\left\{\frac\be N\bigl(h_\mu(x)+h_\mu(y)-2I_\mu\bigr)\right\}
 \mathsf M_{N,2}(x,y)}
 {\displaystyle\int_{(\R^\d)^2}\Psi_{N,2}\,\md\mu^{\otimes2}}.
\end{align}
Then \Cref{lem:normalized-conditional-identity} gives the exact factorization
\begin{align}\label{eq:uv-two-point-factorization}
 f_{N,\be}^{(2)}(x,y)
 =\mathcal R_{N,\be}(x,y)
 \exp\left\{-\frac\be N\tl W(x,y)\right\}.
\end{align}
For every fixed configuration of the remaining \(N-2\) particles, Hoeffding centering and Jensen's inequality give
\begin{align}\label{eq:uv-two-point-normalization-lower}
 \int_{(\R^\d)^2}\Psi_{N,2}\,\md\mu^{\otimes2}
 \ge \exp\left(\frac{2\be I_\mu}{N}\right)\ge1.
\end{align}
Moreover, \eqref{eq:uv-centered-lower-bound} and the definition of \(\mathsf M_{N,2}\) imply
\begin{align}\label{eq:uv-two-point-reservoir-bound}
 0\le \mathsf M_{N,2}(x,y)\le e^{4\be H_\mu}.
\end{align}
Combining \eqref{eq:uv-Hmu-bound}--\eqref{eq:uv-two-point-reservoir-bound}, we obtain
\begin{align}\label{eq:uv-residual-prefactor-bound}
 0\le \mathcal R_{N,\be}(x,y)\le e^{5\be H_\mu}
\end{align}
for \(\mu^{\otimes2}\)-a.e. \((x,y)\), and consequently
\begin{align}\label{eq:uv-two-point-density-bound}
 0\le f_{N,\be}^{(2)}(x,y)\le e^{6\be H_\mu}.
\end{align}
Recall the ultraviolet scale \(\mathfrak a_N\) from \eqref{eq:intro-uv-scale}.  All constants below may depend on \((\be,\mu,\d,\s)\), but never on \(N\) or on the number of observed particles.

\subsection{All-temperature fixed-particle \texorpdfstring{\(\chi^2\)}{chi-square} failure}
\label{subsec:uv-fixed-particle-chi2}

\subsubsection{The strictly supercritical case}
\label{subsubsec:uv-supercritical-two-point}

The mutual Boltzmann factor in \eqref{eq:uv-two-point-factorization} suppresses pairs separated by a distance of order \(N^{-1/\s}\).  Comparing the probability of such pairs under the Gibbs and product laws gives the supercritical lower bound.

\begin{prop}[Strictly supercritical two-point obstruction]
\label{prop:uv-two-point-supercritical}
Let \(\d/2<\s<\d\) and \(\be>0\).  Then there are constants \(c>0\) and \(N_0<\infty\) such that, for every \(N\ge N_0\),
\begin{align}\label{eq:uv-two-point-supercritical}
 \Chi\pa{\Q_{N,\be}^{(2)}(\mu)\mid\mu^{\otimes2}}
 \ge cN^{-\d/\s}.
\end{align}
In particular,
\begin{align}\label{eq:uv-two-point-supercritical-divergence}
 N^2\Chi\pa{\Q_{N,\be}^{(2)}(\mu)\mid\mu^{\otimes2}}
 \longrightarrow+\infty.
\end{align}
\end{prop}

\begin{proof}
Retain the bounded set \(A\), lower-density constant \(a\), and radius \(r_0\) chosen above.  After decreasing \(r_0\), if necessary, \eqref{eq:uv-centered-diagonal-asymptotic} gives a constant \(c_0>0\) such that
\begin{align}\label{eq:uv-supercritical-local-lower}
 \tl W(x,y)\ge c_0\abs{x-y}^{-\s}
 \qquad\text{for \((x,y)\in A\times A\) with \(\abs{x-y}\le r_0\)}.
\end{align}
Choose \(\rho_*>0\) sufficiently small that
\begin{align}\label{eq:uv-close-pair-theta}
 \theta
 \coloneqq \exp\left\{5\be H_\mu-\be c_0\rho_*^{-\s}\right\}
 <\frac12.
\end{align}
For all sufficiently large \(N\), set
\begin{align}
 r_N&\coloneqq\rho_*N^{-1/\s},
 \label{eq:uv-close-pair-radius}\\
 E_N&\coloneqq
 \bigl\{(x,y)\in A\times A:\abs{x-y}\le r_N\bigr\}.
 \label{eq:uv-close-pair-set}
\end{align}
By \eqref{eq:uv-two-point-factorization}, \eqref{eq:uv-residual-prefactor-bound}, and \eqref{eq:uv-supercritical-local-lower},
\begin{align}\label{eq:uv-close-pair-pointwise-depletion}
 f_{N,\be}^{(2)}(x,y)\le\theta
 \qquad\text{for \((x,y)\in E_N\)}.
\end{align}
Hence,
\begin{align}\label{eq:uv-close-pair-gibbs-depletion}
 \Q_{N,\be}^{(2)}(\mu)(E_N)
 \le\theta\mu^{\otimes2}(E_N).
\end{align}

Boundedness of the density gives \(\mu^{\otimes2}(E_N)\le Cr_N^\d\).  For the reverse bound, use \(\mu\ge a\) on \(A\), the change of variables \(y=x+z\), and continuity of translations of \(\one_A\) in \(L^1(\R^\d)\).  We obtain
\begin{align}\label{eq:uv-close-pair-probability-scale}
 cN^{-\d/\s}
 \le a^2\int_{\abs z\le r_N}\abs{A\cap(A-z)}\,\md z
 \le \mu^{\otimes2}(E_N)
 \le CN^{-\d/\s}.
\end{align}
Therefore,
\begin{align}\label{eq:uv-supercritical-chi-lower-from-event}
 \Chi\pa{\Q_{N,\be}^{(2)}(\mu)\mid\mu^{\otimes2}}
 \ge(1-\theta)^2\mu^{\otimes2}(E_N)
 \ge cN^{-\d/\s}.
\end{align}
Since \(\d/\s<2\), \eqref{eq:uv-two-point-supercritical-divergence} follows.
\end{proof}

We retain the close-pair data \(A,a,\rho_*,\theta,r_N,E_N\) from the preceding proof.  The estimates \eqref{eq:uv-close-pair-gibbs-depletion} and \eqref{eq:uv-close-pair-probability-scale} will be used again in \Cref{subsubsec:uv-supercritical-mesoscopic}.

\subsubsection{The critical case}
\label{subsubsec:uv-critical-two-point}

At the critical exponent \(\s=\d/2\), the set \(\abs{x-y}\lesssim N^{-1/\s}\) has product probability only of order \(N^{-2}\).  The additional logarithm comes from the range of transition annuli on which the mutual Boltzmann factor differs perturbatively from one.  The factorization \eqref{eq:uv-two-point-factorization} gives the pointwise inequality
\begin{align}\label{eq:uv-critical-basic-decomposition}
 \bigl(1-f_{N,\be}^{(2)}\bigr)^2
 \ge\frac12\left(1-e^{-\be\tl W/N}\right)^2
 -\bigl(\mathcal R_{N,\be}-1\bigr)^2e^{-2\be\tl W/N}.
\end{align}
The first term has integral of order \(N^{-2}\log N\) across the transition annuli.  Thus, an error of order \(o(N^{-2}\log N)\) in the second term could not cancel the leading contribution.  For the actual residual prefactor, however, the natural estimate approximates \(\mathcal R_{N,\be}(x,y)\) by a function of one variable, rather than by the constant one.  We first isolate the scalar transition-annulus calculation used both here and in the mesoscopic critical argument.

\begin{lemma}[Scalar transition-annulus estimate]
\label{lem:uv-scalar-transition-annulus}
For every fixed \(\lambda>0\),
\begin{align}\label{eq:uv-scalar-transition-annulus-estimate}
 \inf_{v\in\R}\int_\epsilon^1
 \left(1-ve^{-\lambda t}\right)^2t^{-3}\,\md t
 =\lambda^2\log\frac1\epsilon+O_\lambda(1)
\end{align}
as \(\epsilon\downarrow0\).
\end{lemma}

\begin{proof}
Set
\begin{align}
 A_\epsilon&\coloneqq\int_\epsilon^1t^{-3}\,\md t,
 \label{eq:uv-scalar-A-def}\\
 B_{\epsilon,\lambda}&\coloneqq\int_\epsilon^1e^{-\lambda t}t^{-3}\,\md t,
 \label{eq:uv-scalar-B-def}\\
 C_{\epsilon,\lambda}&\coloneqq\int_\epsilon^1e^{-2\lambda t}t^{-3}\,\md t.
 \label{eq:uv-scalar-C-def}
\end{align}
The minimum in \eqref{eq:uv-scalar-transition-annulus-estimate} equals \(A_\epsilon-B_{\epsilon,\lambda}^2/C_{\epsilon,\lambda}\).  Taylor expansion at zero gives
\begin{align}
 A_\epsilon&=\frac1{2\epsilon^2}+O(1),
 \label{eq:uv-scalar-A-expansion}\\
 B_{\epsilon,\lambda}&=\frac1{2\epsilon^2}-\frac\lambda\epsilon
 +\frac{\lambda^2}{2}\log\frac1\epsilon+O_\lambda(1),
 \label{eq:uv-scalar-B-expansion}\\
 C_{\epsilon,\lambda}&=\frac1{2\epsilon^2}-\frac{2\lambda}\epsilon
 +2\lambda^2\log\frac1\epsilon+O_\lambda(1).
 \label{eq:uv-scalar-C-expansion}
\end{align}
Dividing these expansions yields \eqref{eq:uv-scalar-transition-annulus-estimate}.
\end{proof}

The following lemma records the corresponding one-variable prefactor criterion.

\begin{lemma}[One-variable prefactor criterion]
\label{lem:uv-critical-one-variable-criterion}
Let \(\s=\d/2\), let \(\be>0\), and suppose that \(\mu\ge m_B>0\) on a ball \(B\subset\R^\d\).  Let \(B'\Subset B\), fix \(0<\alpha<1/\s=2/\d\), and let \(0<c_-<c_+<\infty\).  Define
\begin{align}\label{eq:uv-critical-annulus-def}
 \mathcal T_N
 \coloneqq
 \left\{(x,y)\in B'\times B:
 c_-N^{-1/\s}\le\abs{x-y}\le c_+N^{-\alpha}\right\}.
\end{align}
Assume that
\begin{align}\label{eq:uv-one-variable-prefactor-assumption}
 \inf_{v\in L^2(B',\mu)}
 \int_{\mathcal T_N}
 \abs{\mathcal R_{N,\be}(x,y)-v(x)}^2
 e^{-2\be\tl W(x,y)/N}\,\md\mu(x)\md\mu(y)
 =o\left(N^{-2}\log N\right).
\end{align}
Then there are \(c>0\) and \(N_0<\infty\) such that
\begin{align}\label{eq:uv-critical-criterion-conclusion}
 \Chi\pa{\Q_{N,\be}^{(2)}(\mu)\mid\mu^{\otimes2}}
 \ge cN^{-2}\log N
\end{align}
for every \(N\ge N_0\).
\end{lemma}

\begin{proof}
For each \(N\), choose \(v_N\in L^2(B',\mu)\) such that the integral in \eqref{eq:uv-one-variable-prefactor-assumption} is at most the displayed infimum plus \(N^{-3}\).  By the hypothesis, this integral is \(o(N^{-2}\log N)\).  The pointwise identity
\begin{align}\label{eq:uv-critical-one-variable-decomposition}
 1-f_{N,\be}^{(2)}(x,y)
 =\left(1-v_N(x)e^{-\be\tl W(x,y)/N}\right)
 -\left(\mathcal R_{N,\be}(x,y)-v_N(x)\right)
 e^{-\be\tl W(x,y)/N}
\end{align}
holds independently of the domain of integration.  Therefore,
\begin{align}\label{eq:uv-critical-reduction-to-main-integral}
 \Chi\pa{\Q_{N,\be}^{(2)}(\mu)\mid\mu^{\otimes2}}
 &\ge\frac12
 \int_{\mathcal T_N}
 \left(1-v_N(x)e^{-\be\tl W(x,y)/N}\right)^2
 \md\mu(x)\md\mu(y)
 -o\left(N^{-2}\log N\right).
\end{align}
It remains to lower bound the displayed integral.

To this end, choose
\begin{align}\label{eq:uv-critical-alpha-bar-choice}
 \bar\alpha\in
 \left(\max\left\{\alpha,\frac1{2\s}\right\},\frac1\s\right).
\end{align}
Fix \(c_{\mathrm{in}}>c_-\) and then choose \(c_{\mathrm{out}}>0\) so small that \(x+r\omega\in B\) whenever \(x\in B'\), \(\omega\in\mathbb S^{\d-1}\), and \(0\le r\le c_{\mathrm{out}}\).  For all sufficiently large \(N\),
\begin{align}\label{eq:uv-critical-subannulus-inclusion}
 \overline{\mathcal T}_N
 \coloneqq
 \left\{(x,y)\in B'\times B:
 c_{\mathrm{in}}N^{-1/\s}\le\abs{x-y}
 \le c_{\mathrm{out}}N^{-\bar\alpha}\right\}
 \subset\mathcal T_N.
\end{align}
By \eqref{eq:uv-centered-diagonal-asymptotic}, write
\begin{align}\label{eq:uv-critical-singular-remainder}
 \tl W(x,y)=\frac1\s\abs{x-y}^{-\s}+R_\mu(x,y),
 \qquad
 \abs{R_\mu(x,y)}\le2H_\mu.
\end{align}
Set
\begin{align}
 \lambda&\coloneqq\frac{\be}{\s c_{\mathrm{in}}^\s},
 \label{eq:uv-critical-lambda-def}\\
 \epsilon_N&\coloneqq
 \left(\frac{c_{\mathrm{in}}}{c_{\mathrm{out}}}\right)^\s
 N^{\bar\alpha\s-1}.
 \label{eq:uv-critical-epsilon-def}
\end{align}
Under the change of variables
\begin{align}\label{eq:uv-critical-radial-change}
 t=\left(\frac{c_{\mathrm{in}}N^{-1/\s}}{r}\right)^\s,
 \qquad
 r=c_{\mathrm{in}}N^{-1/\s}t^{-1/\s},
\end{align}
the radial interval in \(\overline{\mathcal T}_N\) becomes \(\epsilon_N\le t\le1\), and criticality \(\d/\s=2\) gives
\begin{align}\label{eq:uv-critical-jacobian}
 r^{\d-1}\,\md r
 =\frac{c_{\mathrm{in}}^\d}{\s}N^{-2}t^{-3}\,\md t.
\end{align}
Define
\begin{align}\label{eq:uv-critical-vartheta-def}
 \vartheta_N(x,\omega,t)
 \coloneqq
 \exp\left\{-\frac\be N
 R_\mu\left(x,x+c_{\mathrm{in}}N^{-1/\s}t^{-1/\s}\omega\right)\right\}.
\end{align}
Then
\begin{align}\label{eq:uv-critical-vartheta-close}
 \norm{\vartheta_N-1}_{L^\infty}\le\frac{C_{\be,\mu}}{N},
\end{align}
and the mutual Boltzmann factor is \(e^{-\lambda t}\vartheta_N(x,\omega,t)\).

For \(0<\epsilon<1\), define the perturbed coefficients
\begin{align}
 B_{\epsilon,\lambda}^{\vartheta_N}(x)
 &\coloneqq\frac1{\abs{\mathbb S^{\d-1}}}
 \int_{\mathbb S^{\d-1}}\int_\epsilon^1
 e^{-\lambda t}\vartheta_N(x,\omega,t)t^{-3}\,\md t\,\md\omega,
 \label{eq:uv-critical-B-vartheta-def}\\
 C_{\epsilon,\lambda}^{\vartheta_N}(x)
 &\coloneqq\frac1{\abs{\mathbb S^{\d-1}}}
 \int_{\mathbb S^{\d-1}}\int_\epsilon^1
 e^{-2\lambda t}\vartheta_N(x,\omega,t)^2t^{-3}\,\md t\,\md\omega.
 \label{eq:uv-critical-C-vartheta-def}
\end{align}
Since \(\mu\ge m_B\) on \(B\), polar coordinates and \eqref{eq:uv-critical-jacobian} yield
\begin{align}\label{eq:uv-critical-quadratic-reduction}
 &\int_{\overline{\mathcal T}_N}
 \left(1-v_N(x)e^{-\be\tl W(x,y)/N}\right)^2
 \md\mu(x)\md\mu(y)\notag\\
 &\qquad\ge cN^{-2}
 \int_{B'}
 \left[
 A_{\epsilon_N}
 -2v_N(x)B_{\epsilon_N,\lambda}^{\vartheta_N}(x)
 +v_N(x)^2C_{\epsilon_N,\lambda}^{\vartheta_N}(x)
 \right]\,\md x.
\end{align}
For fixed \(x\), the bracket is a quadratic polynomial in the scalar \(v_N(x)\).  Completing the square shows that its minimum over \(v\in\R\) is
\begin{align}\label{eq:uv-critical-quadratic-minimum}
 A_{\epsilon_N}
 -\frac{B_{\epsilon_N,\lambda}^{\vartheta_N}(x)^2}
 {C_{\epsilon_N,\lambda}^{\vartheta_N}(x)}.
\end{align}

For the unperturbed factor \(\vartheta_N\equiv1\), use \(A_\epsilon,B_{\epsilon,\lambda},C_{\epsilon,\lambda}\) from \eqref{eq:uv-scalar-A-def}--\eqref{eq:uv-scalar-C-def}.  By \Cref{lem:uv-scalar-transition-annulus},
\begin{align}\label{eq:uv-critical-unperturbed-minimum}
 A_\epsilon-\frac{B_{\epsilon,\lambda}^2}{C_{\epsilon,\lambda}}
 =\lambda^2\log\frac1\epsilon+O_\lambda(1).
\end{align}
The bound \eqref{eq:uv-critical-vartheta-close} implies, uniformly in \(x\in B'\),
\begin{align}\label{eq:uv-critical-coefficient-perturbations}
 \abs{B_{\epsilon_N,\lambda}^{\vartheta_N}(x)-B_{\epsilon_N,\lambda}}
 +\abs{C_{\epsilon_N,\lambda}^{\vartheta_N}(x)-C_{\epsilon_N,\lambda}}
 \le C_{\be,\mu}N^{-1}\epsilon_N^{-2}.
\end{align}
Moreover,
\begin{align}\label{eq:uv-critical-coefficient-sizes}
 C_{\epsilon_N,\lambda}^{\vartheta_N}(x)\asymp C_{\epsilon_N,\lambda}\asymp\epsilon_N^{-2},
 \qquad
 \abs{B_{\epsilon_N,\lambda}^{\vartheta_N}(x)}+B_{\epsilon_N,\lambda}\lesssim\epsilon_N^{-2}.
\end{align}
Consequently,
\begin{align}\label{eq:uv-critical-perturbation-ratio}
 \left|
 \frac{B_{\epsilon_N,\lambda}^{\vartheta_N}(x)^2}
 {C_{\epsilon_N,\lambda}^{\vartheta_N}(x)}
 -\frac{B_{\epsilon_N,\lambda}^2}{C_{\epsilon_N,\lambda}}
 \right|
 \le C_{\be,\mu}N^{-1}\epsilon_N^{-2}.
\end{align}
Because \(\bar\alpha\s>1/2\),
\begin{align}\label{eq:uv-critical-perturbation-negligible}
 N^{-1}\epsilon_N^{-2}
 =C N^{1-2\bar\alpha\s}
 =o\left(\log\frac1{\epsilon_N}\right).
\end{align}
Also, \(\log(1/\epsilon_N)=(1-\bar\alpha\s)\log N+O(1)\).  Therefore, \eqref{eq:uv-critical-quadratic-reduction}--\eqref{eq:uv-critical-perturbation-negligible} give
\begin{align}\label{eq:uv-critical-main-integral-lower}
 \int_{\mathcal T_N}
 \left(1-v_N(x)e^{-\be\tl W(x,y)/N}\right)^2
 \md\mu(x)\md\mu(y)
 \ge cN^{-2}\log N.
\end{align}
Combining this lower bound with the negligible residual-prefactor error in \eqref{eq:uv-critical-reduction-to-main-integral} proves \eqref{eq:uv-critical-criterion-conclusion}.
\end{proof}

The next two lemmas provide the inputs needed to verify the one-variable prefactor criterion for \(\mathcal R_{N,\be}\).

\begin{lemma}[Uniform one-point density bound]
\label{lem:uv-uniform-one-point-bound}
For every \(\be_*<\infty\),
\begin{align}\label{eq:uv-uniform-one-point-bound}
 \sup_{M\ge1}\sup_{0\le\be'\le\be_*}
 \norm{f_{M,\be'}^{(1)}}_{L^\infty(\mu)}
 \le e^{3\be_*H_\mu}.
\end{align}
\end{lemma}

\begin{proof}
Assume first that \(M\ge2\).  The one-mark identity of \Cref{lem:normalized-conditional-identity} gives
\begin{align}\label{eq:uv-one-point-marked-representation}
 f_{M,\be'}^{(1)}(x)
 =\frac{
 e^{\be'(h_\mu(x)-I_\mu)/M}\mathsf M_{M,1}(x)}
 {\int\Psi_{M,1}\,\md\mu}.
\end{align}
By \eqref{eq:uv-centered-lower-bound},
\begin{align}\label{eq:uv-one-point-marked-numerator-bound}
 e^{\be'(h_\mu(x)-I_\mu)/M}\mathsf M_{M,1}(x)
 \le e^{3\be_*H_\mu}.
\end{align}
For each fixed configuration of the remaining particles, Jensen's inequality and Hoeffding centering imply
\begin{align}\label{eq:uv-one-point-marked-denominator-bound}
 \int\Psi_{M,1}\,\md\mu\ge1.
\end{align}
This proves the estimate for \(M\ge2\).  For \(M=1\), the same argument without the remaining particles gives
\begin{align}\label{eq:uv-one-point-M1-bound}
 f_{1,\be'}^{(1)}(x)
 =\frac{e^{\be'(h_\mu(x)-I_\mu)}}
 {\int e^{\be'(h_\mu-I_\mu)}\,\md\mu}
 \le e^{\be_*H_\mu},
\end{align}
which completes the proof.
\end{proof}

\begin{lemma}[Critical Riesz translation modulus]
\label{lem:uv-riesz-translation-modulus}
Let \(K\subset\R^\d\) be compact.  For \(0<r<1/2\), define the raw Riesz translation modulus
\begin{align}\label{eq:uv-translation-modulus-def}
 \omega_K(r)
 \coloneqq
 \sup_{\substack{x,y\in K\\\abs{x-y}\le r}}
 \int_{\R^\d}\abs{W(x,z)-W(y,z)}\,\md\mu(z).
\end{align}
At the critical exponent \(\s=\d/2\),
\begin{align}\label{eq:uv-translation-modulus-bound}
 \omega_K(r)\le C_K
 \begin{cases}
 r^{1/2},&\d=1,\\
 r\log(e/r),&\d=2,\\
 r,&\d\ge3.
 \end{cases}
\end{align}
Moreover, the centering identity gives
\begin{align}\label{eq:uv-centered-translation-modulus-bound}
 \sup_{\substack{x,y\in K\\\abs{x-y}\le r}}
 \int_{\R^\d}\abs{\tl W(x,z)-\tl W(y,z)}\,\md\mu(z)
 \le2\omega_K(r).
\end{align}
\end{lemma}

\begin{proof}
Set \(\delta=\abs{x-y}<1/2\).  On
\(\{\abs{x-z}\le2\delta\}\cup\{\abs{y-z}\le2\delta\}\), boundedness of the density gives a contribution of order \(\delta^{\d-\s}\).  On the region where both distances exceed \(2\delta\) and \(\abs{x-z}\le1\), the mean-value theorem gives
\begin{align}\label{eq:uv-riesz-mean-value-bound}
 \abs{W(x,z)-W(y,z)}
 \le C\delta\abs{x-z}^{-\s-1}.
\end{align}
Its integral is bounded by
\begin{align}\label{eq:uv-riesz-intermediate-integral}
 C\delta\int_{2\delta}^1r^{\d-\s-2}\,\md r.
\end{align}
On the remaining region, the same mean-value estimate is bounded by \(C\delta\), and integration uses only that \(\mu\) has mass one.  Thus, the raw translation modulus is of order \(\delta^{\d-\s}\) if \(\d-\s<1\), \(\delta\log(e/\delta)\) if \(\d-\s=1\), and \(\delta\) if \(\d-\s>1\).  Substituting \(\s=\d/2\) proves \eqref{eq:uv-translation-modulus-bound}.  Finally,
\begin{align}\label{eq:uv-centered-translation-identity}
 \tl W(x,z)-\tl W(y,z)
 =W(x,z)-W(y,z)-h_\mu(x)+h_\mu(y),
\end{align}
and
\begin{align}\label{eq:uv-potential-translation-bound}
 \abs{h_\mu(x)-h_\mu(y)}
 \le\int\abs{W(x,z)-W(y,z)}\,\md\mu(z),
\end{align}
which gives \eqref{eq:uv-centered-translation-modulus-bound}.
\end{proof}

\begin{lemma}[One-variable approximation of the residual prefactor]
\label{lem:uv-critical-residual-prefactor}
Let \(\s=\d/2\), let \(\be>0\), and suppose that \(\mu\ge m_B>0\) on a ball \(B\).  Let \(B'\Subset B\), and choose
\begin{align}\label{eq:uv-prefactor-alpha-range}
 \alpha\in
 \begin{cases}
 (1,2),&\d=1,\\
 (1/2,1),&\d=2,\\
 (2/(\d+2),2/\d),&\d\ge3.
 \end{cases}
\end{align}
Then there are constants \(0<c_-<c_+<\infty\) such that
\begin{align}\label{eq:uv-prefactor-one-variable-approximation}
 \inf_{v\in L^2(B',\mu)}
 \int_{\substack{(x,y)\in B'\times B\\
 c_-N^{-1/\s}\le\abs{x-y}\le c_+N^{-\alpha}}}
 \abs{\mathcal R_{N,\be}(x,y)-v(x)}^2
 e^{-2\be\tl W(x,y)/N}\,\md\mu(x)\md\mu(y)
 =o\left(N^{-2}\log N\right).
\end{align}
\end{lemma}

\begin{proof}
Let \(n=N-2\), \(\gamma=\be n/N\).  Define
\begin{align}\label{eq:uv-diagonal-reservoir-factor}
 \mathsf M_N^\Delta(x)
 \coloneqq
 \E_{\Q_{n,\gamma}(\mu)}\left[
 \exp\left\{-\frac{2\be}{N}\sum_{j=1}^n\tl W(x,Y_j)\right\}\right]
\end{align}
and
\begin{align}\label{eq:uv-one-variable-prefactor-def}
 v_N(x)
 \coloneqq
 \left(\int_{(\R^\d)^2}\Psi_{N,2}\,\md\mu^{\otimes2}\right)^{-1}
 \exp\left\{\frac{2\be}{N}\bigl(h_\mu(x)-I_\mu\bigr)\right\}
 \mathsf M_N^\Delta(x).
\end{align}
The one-body factors in \eqref{eq:uv-residual-prefactor-def} and \eqref{eq:uv-one-variable-prefactor-def} differ by at most \(C_{\be,\mu}/N\), uniformly for \((x,y)\in B\times B\).  For the reservoir factors, the mean-value theorem and \eqref{eq:uv-centered-lower-bound} give
\begin{align}\label{eq:uv-reservoir-factor-difference}
 \abs{\mathsf M_{N,2}(x,y)-\mathsf M_N^\Delta(x)}
 &\le\frac{C_{\be,\mu}}{N}
 \E_{\Q_{n,\gamma}(\mu)}\left[
 \sum_{j=1}^n
 \abs{\tl W(y,Y_j)-\tl W(x,Y_j)}\right].
\end{align}
By exchangeability, \Cref{lem:uv-uniform-one-point-bound} with \(\be_*=\be\), and \Cref{lem:uv-riesz-translation-modulus},
\begin{align}\label{eq:uv-reservoir-factor-modulus}
 \abs{\mathsf M_{N,2}(x,y)-\mathsf M_N^\Delta(x)}
 \le C_{\be,\mu}\omega_B(\abs{x-y}).
\end{align}
Since \(\int\Psi_{N,2}\,\md\mu^{\otimes2}\ge1\) by \eqref{eq:uv-two-point-normalization-lower} and all remaining factors are uniformly bounded, it follows that
\begin{align}\label{eq:uv-residual-prefactor-modulus}
 \abs{\mathcal R_{N,\be}(x,y)-v_N(x)}
 \le C_{\be,\mu}\left(\omega_B(\abs{x-y})+N^{-1}\right)
 \qquad\text{for \((x,y)\in B\times B\)}.
\end{align}
Also, \eqref{eq:uv-centered-lower-bound} gives
\begin{align}\label{eq:uv-critical-weight-upper}
 e^{-2\be\tl W(x,y)/N}\le e^{4\be H_\mu/N}.
\end{align}
Choose any fixed \(0<c_-<c_+<\infty\).  By \eqref{eq:uv-residual-prefactor-modulus}--\eqref{eq:uv-critical-weight-upper}, the integral in \eqref{eq:uv-prefactor-one-variable-approximation}, evaluated at \(v_N\), is bounded by
\begin{align}\label{eq:uv-prefactor-error-integral}
 C_{\be,\mu}
 \int_{\substack{(x,y)\in B'\times B\\
 \abs{x-y}\le c_+N^{-\alpha}}}
 \left(\omega_B(\abs{x-y})^2+N^{-2}\right)\,\md x\,\md y.
\end{align}
Using \eqref{eq:uv-translation-modulus-bound}, the contribution of \(\omega_B^2\) is
\begin{align}\label{eq:uv-prefactor-error-rates}
 \begin{cases}
 O(N^{-2\alpha}),&\d=1,\\
 O(N^{-4\alpha}(\log N)^2),&\d=2,\\
 O(N^{-\alpha(\d+2)}),&\d\ge3,
 \end{cases}
\end{align}
whereas the constant term is \(O(N^{-2-\alpha\d})\).  The range \eqref{eq:uv-prefactor-alpha-range} makes every contribution \(o(N^{-2}\log N)\), which proves \eqref{eq:uv-prefactor-one-variable-approximation}.
\end{proof}

\begin{prop}[Critical two-point obstruction]
\label{prop:uv-two-point-critical}
Let \(\s=\d/2\), let \(\be>0\), and suppose that \(\mu\ge m_B>0\) on a ball.  Then there are constants \(c>0\) and \(N_0<\infty\) such that, for every \(N\ge N_0\),
\begin{align}\label{eq:uv-two-point-critical}
 \Chi\pa{\Q_{N,\be}^{(2)}(\mu)\mid\mu^{\otimes2}}
 \ge cN^{-2}\log N.
\end{align}
In particular,
\begin{align}\label{eq:uv-two-point-critical-divergence}
 N^2\Chi\pa{\Q_{N,\be}^{(2)}(\mu)\mid\mu^{\otimes2}}
 \longrightarrow+\infty.
\end{align}
\end{prop}

\begin{proof}
Let \(B\) be a ball on which \(\mu\ge m_B>0\), and choose \(B'\Subset B\).  Select \(\alpha\) in the range \eqref{eq:uv-prefactor-alpha-range}.  \Cref{lem:uv-critical-residual-prefactor} verifies the hypothesis of \Cref{lem:uv-critical-one-variable-criterion}, which gives \eqref{eq:uv-two-point-critical}.  The divergence in \eqref{eq:uv-two-point-critical-divergence} is immediate.
\end{proof}

\subsection{Fixed-particle consequences for all three divergences}
\label{subsec:uv-fixed-particle-consequences}

The preceding propositions give two-point chi-square estimates.  Monotonicity under coordinate projection transfers them to every fixed number of observed particles, while the uniform upper bound \eqref{eq:uv-two-point-density-bound} converts the two-point chi-square estimate into relative-entropy and total-variation lower bounds.

\begin{cor}[Fixed-particle ultraviolet lower bounds]
\label{cor:uv-fixed-particle-all-divergences}
Let \(\be>0\) and fix an integer \(k\ge2\).  If \(\d/2<\s<\d\), assume the standing Riesz hypotheses.  If \(\s=\d/2\), assume in addition that \(\mu\ge m_B>0\) on some ball.  Then there are positive constants \(c_\chi,c_H,c_{\mathrm{TV}}\) such that, for all sufficiently large \(N\),
\begin{align}
 \Chi\pa{\Q_{N,\be}^{(k)}(\mu)\mid\mu^{\otimes k}}
 &\ge c_\chi\mathfrak a_N,
 \label{eq:uv-fixed-particle-chi}\\
 \Hr\pa{\Q_{N,\be}^{(k)}(\mu)\mid\mu^{\otimes k}}
 &\ge c_H\mathfrak a_N,
 \label{eq:uv-fixed-particle-H}\\
 \TV\pa{\Q_{N,\be}^{(k)}(\mu),\mu^{\otimes k}}
 &\ge c_{\mathrm{TV}}\mathfrak a_N.
 \label{eq:uv-fixed-particle-TV}
\end{align}
\end{cor}

\begin{proof}
Coordinate projection from \(k\) labels to two and monotonicity of \(f\)-divergences give
\begin{align}\label{eq:uv-chi-projection-monotonicity}
 \Chi\pa{\Q_{N,\be}^{(k)}(\mu)\mid\mu^{\otimes k}}
 \ge
 \Chi\pa{\Q_{N,\be}^{(2)}(\mu)\mid\mu^{\otimes2}}.
\end{align}
Thus, \eqref{eq:uv-fixed-particle-chi} follows from \Cref{prop:uv-two-point-supercritical,prop:uv-two-point-critical}.

By \eqref{eq:uv-two-point-density-bound}, \(0\le f_{N,\be}^{(2)}\le e^{6\be H_\mu}\).  With the convention \(0\log0=0\), compactness of \([0,e^{6\be H_\mu}]\) and the quadratic expansion at one give a constant \(c_{\be,\mu}>0\) such that
\begin{align}\label{eq:uv-entropy-chi-comparison}
 t\log t-t+1\ge c_{\be,\mu}(t-1)^2
 \qquad\text{for every \(0\le t\le e^{6\be H_\mu}\)},
\end{align}
and
\begin{align}\label{eq:uv-TV-chi-comparison}
 (t-1)^2\le(e^{6\be H_\mu}+1)\abs{t-1}
 \qquad\text{for every \(0\le t\le e^{6\be H_\mu}\)}.
\end{align}
Since \(\int(f_{N,\be}^{(2)}-1)\,\md\mu^{\otimes2}=0\),
\begin{align}
 \Hr\pa{\Q_{N,\be}^{(2)}(\mu)\mid\mu^{\otimes2}}
 &\ge c_{\be,\mu}
 \Chi\pa{\Q_{N,\be}^{(2)}(\mu)\mid\mu^{\otimes2}},
 \label{eq:uv-two-point-H-from-chi}\\
 \TV\pa{\Q_{N,\be}^{(2)}(\mu),\mu^{\otimes2}}
 &\ge\frac1{2(e^{6\be H_\mu}+1)}
 \Chi\pa{\Q_{N,\be}^{(2)}(\mu)\mid\mu^{\otimes2}}.
 \label{eq:uv-two-point-TV-from-chi}
\end{align}
Relative entropy and total variation are also monotone under coordinate projection.  Applying \Cref{prop:uv-two-point-supercritical,prop:uv-two-point-critical} in \eqref{eq:uv-two-point-H-from-chi}--\eqref{eq:uv-two-point-TV-from-chi}, and then projecting from \(k\) labels to two, proves \eqref{eq:uv-fixed-particle-H}--\eqref{eq:uv-fixed-particle-TV}.
\end{proof}

\subsection{Mesoscopic total-variation obstructions}
\label{subsec:uv-mesoscopic}

Recall the common scale \(\mathfrak a_N\) from \eqref{eq:intro-uv-scale}.  In both the critical and strictly supercritical regimes, we prove
\begin{align}\label{eq:uv-mesoscopic-common-form}
 \TV\pa{\Q_{N,\be}^{(k)}(\mu),\mu^{\otimes k}}
 \ge c k^2\mathfrak a_N
 \qquad\text{for \(2\le k\le c_0\mathfrak a_N^{-1/2}\)}.
\end{align}
The scale is the same, but the mechanisms are different.  In the strictly supercritical case, we take a union of the depleted close-pair events constructed in \Cref{subsubsec:uv-supercritical-two-point}.  At criticality, we accumulate one-step conditional Hellinger deficits along the chain of labeled marginals.

\subsubsection{The strictly supercritical case: close-pair accumulation}
\label{subsubsec:uv-supercritical-mesoscopic}

The strictly supercritical two-point argument becomes stronger when the number of observed labels grows.  At the birthday-paradox scale, the product law sees at least one close pair with probability of order one, whereas the Gibbs law suppresses each such pair by the fixed factor in \eqref{eq:uv-close-pair-gibbs-depletion}.

\begin{prop}[Strictly supercritical mesoscopic obstruction]
\label{prop:uv-mesoscopic-close-pair}
Let \(\d/2<\s<\d\) and \(\be>0\).  There are constants \(c_0,c_1>0\) and \(N_0<\infty\) such that, for every \(N\ge N_0\) and every integer
\begin{align}\label{eq:uv-mesoscopic-supercritical-k-range}
 2\le k\le c_0\mathfrak a_N^{-1/2},
\end{align}
one has
\begin{align}\label{eq:uv-mesoscopic-supercritical-TV-lower}
 \TV\pa{\Q_{N,\be}^{(k)}(\mu),\mu^{\otimes k}}
 \ge c_1k^2\mathfrak a_N.
\end{align}
\end{prop}

\begin{proof}
Fix the set \(A\), the radii \(r_N\), and the events \(E_N\) from the proof of \Cref{prop:uv-two-point-supercritical}.  The estimates \eqref{eq:uv-close-pair-gibbs-depletion} and \eqref{eq:uv-close-pair-probability-scale} give
\begin{align}\label{eq:uv-recalled-close-pair-data}
 \Q_{N,\be}^{(2)}(\mu)(E_N)&\le\theta\mu^{\otimes2}(E_N),
 \\
 cN^{-\d/\s}&\le\mu^{\otimes2}(E_N)\le CN^{-\d/\s},
 \qquad 0<\theta<1.
\end{align}
For \(1\le i<j\le k\), define
\begin{align}
 E_{ij}&\coloneqq
 \bigl\{(x_1,\ldots,x_k):(x_i,x_j)\in E_N\bigr\},
 \label{eq:uv-observed-pair-event}\\
 E_{N,k}&\coloneqq\bigcup_{1\le i<j\le k}E_{ij}.
 \label{eq:uv-observed-close-pair-union}
\end{align}
For each pair \(i<j\),
\begin{align}\label{eq:uv-single-pair-product-probability}
 \mu^{\otimes k}(E_{ij})=\mu^{\otimes2}(E_N).
\end{align}
If the index pairs \(\{i,j\}\) and \(\{\ell,m\}\) are disjoint, then the corresponding events are independent.  If they share one vertex, then, for example,
\begin{align}\label{eq:uv-overlapping-pair-intersection}
 \mu^{\otimes k}(E_{12}\cap E_{13})
 &=\int_A\mu\bigl(A\cap B(x,r_N)\bigr)^2\,\md\mu(x)\notag\\
 &\le Cr_N^\d\mu^{\otimes2}(E_N)
 \le C\mu^{\otimes2}(E_N)^2,
\end{align}
where the last inequality uses the lower bound in \eqref{eq:uv-close-pair-probability-scale}.  The same upper bound holds for disjoint pairs.  Bonferroni's inequality therefore gives
\begin{align}\label{eq:uv-bonferroni}
 \mu^{\otimes k}(E_{N,k})
 \ge\binom{k}{2}\mu^{\otimes2}(E_N)
 -C\binom{k}{2}^{\!2}\mu^{\otimes2}(E_N)^2.
\end{align}
On the other hand, the union bound, exchangeability, and \eqref{eq:uv-close-pair-gibbs-depletion} give
\begin{align}\label{eq:uv-gibbs-union}
 \Q_{N,\be}^{(k)}(\mu)(E_{N,k})
 \le\theta\binom{k}{2}\mu^{\otimes2}(E_N).
\end{align}
Hence,
\begin{align}\label{eq:uv-event-probability-gap}
 \mu^{\otimes k}(E_{N,k})-\Q_{N,\be}^{(k)}(\mu)(E_{N,k})
 \ge\binom{k}{2}\mu^{\otimes2}(E_N)
 \left(1-\theta-C\binom{k}{2}\mu^{\otimes2}(E_N)\right).
\end{align}
Choose \(c_0>0\) sufficiently small.  For \(k\le c_0\mathfrak a_N^{-1/2}\), the upper bound in \eqref{eq:uv-close-pair-probability-scale} makes the final error at most \((1-\theta)/2\).  The lower bound in the same display then yields
\begin{align}\label{eq:uv-event-gap-lower}
 \mu^{\otimes k}(E_{N,k})-\Q_{N,\be}^{(k)}(\mu)(E_{N,k})
 \ge ck^2\mathfrak a_N.
\end{align}
The event characterization of total variation proves \eqref{eq:uv-mesoscopic-supercritical-TV-lower}.
\end{proof}

\subsubsection{The critical case: sequential Hellinger accumulation}
\label{subsubsec:uv-critical-mesoscopic}

\begin{prop}[Critical mesoscopic total-variation obstruction]
\label{prop:uv-critical-mesoscopic-TV}
Let \(\s=\d/2\), let \(\be>0\), and let \(\mu\) be a probability measure on \(\R^\d\) with bounded Lebesgue density.  Assume that \(\mu\ge m_B>0\) on some ball.  Then there are constants \(c_0,c_1>0\) and \(N_0\ge3\), depending only on \((\be,\mu,\d)\), such that, for every \(N\ge N_0\) and every integer
\begin{align}\label{eq:uv-critical-k-window}
 2\le k\le c_0\mathfrak a_N^{-1/2},
\end{align}
one has
\begin{align}\label{eq:uv-critical-mesoscopic-TV-lower}
 \TV\pa{\Q_{N,\be}^{(k)}(\mu),\mu^{\otimes k}}
 \ge c_1k^2\mathfrak a_N.
\end{align}
\end{prop}

Fix a ball \(B\subset\R^\d\) and \(m_B>0\) such that
\begin{align}\label{eq:uv-critical-local-density}
 \mu(x)\ge m_B
 \qquad\text{for Lebesgue-a.e. \(x\in B\)}.
\end{align}
Since \(W\ge0\), the common bound \eqref{eq:uv-Hmu-bound} also gives \(0\le h_\mu\le H_\mu\).

The proof has three parts.  First, an abstract Hellinger recursion converts one-step conditional deficits into a lower bound for the total variation of the \(k\)-particle marginal.  Second, the raw-kernel marked-reservoir factorization shows that each isolated observed particle with controlled far-field interaction contributes a conditional deficit of order \(\mathfrak a_N\).  Third, a product-law estimate shows that a positive fraction of the observed particles satisfy these conditions in expectation.

\parheading{Hellinger affinity and accumulation of conditional deficits.}

We use Hellinger affinity at the iterative stage because the square root is compatible with disintegration.  Indeed, the conditional factorization in \eqref{eq:uv-critical-conditional-density-chain} below yields an exact multiplicative recursion for the affinity.  Total variation has no comparable multiplicative conditional identity.  After accumulating the resulting affinity deficits, we return to total variation through \eqref{eq:uv-critical-TV-affinity-lower}.

For \(0\le j\le N\), recall the marginal density ratio \(f_{N,\be}^{(j)}\) from \eqref{eq:gibbs-marginal-density-ratio} and define its Hellinger affinity with the product law, also called the Bhattacharyya coefficient, by
\begin{align}\label{eq:uv-critical-affinity-def}
 \mathcal A_{N,j}
 \coloneqq
 \int\sqrt{f_{N,\be}^{(j)}}\,\md\mu^{\otimes j}.
\end{align}
With the normalization in which the squared Hellinger distance is \(\frac12\int\left(\sqrt{\frac{\md P}{\md\lambda}}-\sqrt{\frac{\md Q}{\md\lambda}}\right)^2\,\md\lambda\), it equals \(1-\mathcal A(P,Q)\); see \cite[Section~2.4, pp.~83--86]{Tsybakov2009}.  We use the conventions \(\mathcal A_{N,0}=1\) and \(f_{N,\be}^{(0)}=1\).

For \(0\le j<N\), let \(q_{N,j}(x\mid X_j)\) denote the conditional density, relative to \(\mu\), of the \((j+1)\)-st coordinate under \(\Q_{N,\be}^{(j+1)}(\mu)\) given \(X_j=(x_1,\ldots,x_j)\).\footnote{For each pair \((N,j)\), we take the explicit conditional-density version below on the full-measure set of pairwise-distinct marked configurations and extend it arbitrarily elsewhere.  Since the family of pairs \((N,j)\) is countable and every square-root bridge introduced below is absolutely continuous with respect to \(\mu^{\otimes j}\), these versions may be chosen on one common full-measure set.}  Thus,
\begin{align}\label{eq:uv-critical-conditional-density-chain}
 f_{N,\be}^{(j+1)}(X_j,x)
 =f_{N,\be}^{(j)}(X_j)q_{N,j}(x\mid X_j)
\end{align}
for \(\mu^{\otimes(j+1)}\)-a.e. \((X_j,x)\).  Define
\begin{align}\label{eq:uv-critical-conditional-affinity-deficit}
 \delta_{N,j}(X_j)
 \coloneqq
 1-\int\sqrt{q_{N,j}(x\mid X_j)}\,\md\mu(x)
 =\frac12\int\pa{\sqrt{q_{N,j}(x\mid X_j)}-1}^2\,\md\mu(x)
\end{align}
and the probability measure
\begin{align}\label{eq:uv-critical-bridge-measure}
 \md\nu_{N,j}(X_j)
 \coloneqq
 \frac{\sqrt{f_{N,\be}^{(j)}(X_j)}}{\mathcal A_{N,j}}\,
 \md\mu^{\otimes j}(X_j).
\end{align}
We call \(\nu_{N,j}\) the \emph{square-root bridge} because its density is the normalized geometric mean \(\sqrt{f_{N,\be}^{(j)}\cdot1}\) of the Gibbs marginal density and the product density.  It is precisely the probability measure under which the one-step conditional affinity deficits are averaged in the recursion below.

\begin{lemma}[Affinity recursion, bridge estimate, and accumulation criterion]
\label{lem:uv-critical-affinity-recursion}
For every \(0\le j<N\),
\begin{align}\label{eq:uv-critical-affinity-recursion}
 \mathcal A_{N,j+1}
 =\mathcal A_{N,j}\pa{1-\E_{\nu_{N,j}}\left[\delta_{N,j}\right]}.
\end{align}
Moreover,
\begin{align}\label{eq:uv-critical-TV-affinity-lower}
 \TV\pa{\Q_{N,\be}^{(j)}(\mu),\mu^{\otimes j}}
 \ge1-\mathcal A_{N,j},
\end{align}
and
\begin{align}\label{eq:uv-critical-bridge-TV-bound}
 \TV\pa{\nu_{N,j},\mu^{\otimes j}}
 \le\frac12\sqrt{\mathcal A_{N,j}^{-2}-1}.
\end{align}

Suppose, for some \(N\ge3\) and \(k\ge2\), that measurable functions \(G_{N,j}\), \(1\le j<k\), and constants \(c_\delta,c_G>0\) satisfy
\begin{align}
 0\le G_{N,j}&\le j,
 \label{eq:uv-critical-abstract-count-range}\\
 \delta_{N,j}&\ge c_\delta\mathfrak a_NG_{N,j},
 \label{eq:uv-critical-abstract-deficit}\\
 \E_{\mu^{\otimes j}}\left[G_{N,j}\right]&\ge c_Gj.
 \label{eq:uv-critical-abstract-count}
\end{align}
Then there are constants \(c_0,c_1>0\), depending only on \(c_\delta\) and \(c_G\), such that
\begin{align}\label{eq:uv-critical-abstract-accumulation}
 \TV\pa{\Q_{N,\be}^{(k)}(\mu),\mu^{\otimes k}}
 \ge c_1k^2\mathfrak a_N
 \qquad\text{whenever \(2\le k\le c_0\mathfrak a_N^{-1/2}\)}.
\end{align}
\end{lemma}

\begin{proof}
Using \eqref{eq:uv-critical-conditional-density-chain} and then \eqref{eq:uv-critical-bridge-measure},
\begin{align}
 \mathcal A_{N,j+1}
 &=\int\sqrt{f_{N,\be}^{(j)}(X_j)}
 \left(\int\sqrt{q_{N,j}(x\mid X_j)}\,\md\mu(x)\right)
 \md\mu^{\otimes j}(X_j)\notag\\
 &=\mathcal A_{N,j}\pa{1-\E_{\nu_{N,j}}\left[\delta_{N,j}\right]}.
 \label{eq:uv-critical-affinity-recursion-proof}
\end{align}
Since \(\min\{t,1\}\le\sqrt t\) for \(t\ge0\),
\begin{align}\label{eq:uv-critical-TV-affinity-proof}
 \TV\pa{\Q_{N,\be}^{(j)}(\mu),\mu^{\otimes j}}
 =1-\int\min\{f_{N,\be}^{(j)},1\}\,\md\mu^{\otimes j}
 \ge1-\mathcal A_{N,j}.
\end{align}
The density of \(\nu_{N,j}\) relative to \(\mu^{\otimes j}\) is \(\sqrt{f_{N,\be}^{(j)}}/\mathcal A_{N,j}\), and hence
\begin{align}\label{eq:uv-critical-bridge-chi-square}
 \Chi\pa{\nu_{N,j}\mid\mu^{\otimes j}}
 =\mathcal A_{N,j}^{-2}-1.
\end{align}
The standard bound \(\TV\le\frac12\sqrt{\Chi}\) gives \eqref{eq:uv-critical-bridge-TV-bound}.

It remains to prove the accumulation assertion.  Choose \(a_*\in(0,1)\) sufficiently close to one that
\begin{align}\label{eq:uv-critical-a-star-choice}
 \frac12\sqrt{a_*^{-2}-1}\le\frac{c_G}{8}.
\end{align}
If \(\mathcal A_{N,j}<a_*\) for some \(1\le j<k\), then \eqref{eq:uv-critical-affinity-recursion} shows that the affinities are nonincreasing, and \eqref{eq:uv-critical-TV-affinity-lower} gives
\begin{align}\label{eq:uv-critical-affinity-dichotomy-large-TV}
 \TV\pa{\Q_{N,\be}^{(k)}(\mu),\mu^{\otimes k}}
 \ge1-a_*.
\end{align}

Suppose instead that \(\mathcal A_{N,j}\ge a_*\) for every \(1\le j<k\).  Since \(0\le G_{N,j}\le j\), \eqref{eq:uv-critical-bridge-TV-bound}, \eqref{eq:uv-critical-a-star-choice}, and \eqref{eq:uv-critical-abstract-count} give
\begin{align}\label{eq:uv-critical-bridge-good-count}
 \E_{\nu_{N,j}}\left[G_{N,j}\right]
 \ge\E_{\mu^{\otimes j}}\left[G_{N,j}\right]
 -2j\TV\pa{\nu_{N,j},\mu^{\otimes j}}
 \ge\frac{c_G}{2}j.
\end{align}
Combining this estimate with \eqref{eq:uv-critical-abstract-deficit},
\begin{align}\label{eq:uv-critical-bridge-deficit-lower}
 \E_{\nu_{N,j}}\left[\delta_{N,j}\right]
 \ge c j\mathfrak a_N,
 \qquad 1\le j<k.
\end{align}
Iteration of \eqref{eq:uv-critical-affinity-recursion} and the inequality \(1-u\le e^{-u}\) yield
\begin{align}\label{eq:uv-critical-affinity-exponential-decay}
 \mathcal A_{N,k}
 \le\exp\left\{-c\mathfrak a_N\sum_{j=1}^{k-1}j\right\}
 \le\exp\left\{-c'k^2\mathfrak a_N\right\}.
\end{align}
Consequently,
\begin{align}\label{eq:uv-critical-TV-nonlinear-lower}
 \TV\pa{\Q_{N,\be}^{(k)}(\mu),\mu^{\otimes k}}
 \ge1-e^{-c'k^2\mathfrak a_N}.
\end{align}
Choose \(c_0>0\) so that \(c'c_0^2\le1\), and then choose
\begin{align}\label{eq:uv-critical-final-c1-choice}
 c_1\le\min\left\{\frac{c'}2,\frac{1-a_*}{c_0^2}\right\}.
\end{align}
For \(k\le c_0\mathfrak a_N^{-1/2}\), the right-hand side of \eqref{eq:uv-critical-TV-nonlinear-lower} is at least \((c'/2)k^2\mathfrak a_N\), while \eqref{eq:uv-critical-affinity-dichotomy-large-TV} gives the same conclusion in the alternative case.  This proves \eqref{eq:uv-critical-abstract-accumulation}.
\end{proof}

\parheading{Marked-reservoir factorization and uniform bounds.}

The centered Gibbs exponent can be expanded back into the nonnegative raw Riesz interaction.  Up to a normalizing constant independent of the configuration,
\begin{align}\label{eq:uv-critical-raw-gibbs-density}
 f_{N,\be}(X_N)
 \propto
 \exp\left\{
 \be\sum_{i=1}^Nh_\mu(x_i)
 -\frac\be N\sum_{1\le i<\ell\le N}W(x_i,x_\ell)
 \right\}.
\end{align}
Indeed, this follows by inserting \(\tl W(x,y)=W(x,y)-h_\mu(x)-h_\mu(y)+2I_\mu\) in the modulated Gibbs exponent.  The raw form \eqref{eq:uv-critical-raw-gibbs-density} is advantageous because every interaction term in the exponential is nonpositive.

Fix \(1\le j<N\), \(X_j=(x_1,\ldots,x_j)\), and set \(n=N-j-1\).  When \(n\ge1\), let \(\mathsf R_{N,j}^{X_j}\) be the probability measure on \((\R^\d)^n\) whose density relative to \(\mu^{\otimes n}\) is proportional to
\begin{multline}\label{eq:uv-critical-marked-reservoir-law}
 \exp\Bigg\{
 \be\sum_{\ell=1}^nh_\mu(y_\ell)
 -\frac\be N\sum_{1\le\ell<m\le n}W(y_\ell,y_m)\\
 -\frac\be N\sum_{a=1}^j\sum_{\ell=1}^nW(x_a,y_\ell)
 \Bigg\}.
\end{multline}
For \(n=0\), this is the point mass on the empty configuration.  Set
\begin{align}\label{eq:uv-critical-marked-H-def}
 H_{N,j}(x;X_j)
 \coloneqq
 e^{\be h_\mu(x)}
 \E_{\mathsf R_{N,j}^{X_j}}\left[
 \exp\left\{-\frac\be N\sum_{\ell=1}^nW(x,Y_\ell)\right\}
 \right]
\end{align}
and
\begin{align}\label{eq:uv-critical-marked-normalization-def}
 \mathcal N_{N,j}(X_j)
 \coloneqq
 \int H_{N,j}(x;X_j)
 \exp\left\{-\frac\be N\sum_{a=1}^jW(x_a,x)\right\}\md\mu(x).
\end{align}
A direct splitting of \eqref{eq:uv-critical-raw-gibbs-density} gives the exact conditional density
\begin{align}\label{eq:uv-critical-marked-conditional-factorization}
 q_{N,j}(x\mid X_j)
 =\frac{H_{N,j}(x;X_j)}{\mathcal N_{N,j}(X_j)}
 \exp\left\{-\frac\be N\sum_{a=1}^jW(x_a,x)\right\}.
\end{align}

\begin{lemma}[Uniform marked-reservoir bounds]
\label{lem:uv-critical-marked-reservoir}
For every compact \(K\subset\R^\d\), there is \(C=C(\be,\mu,K)<\infty\) such that, for every \(N\ge2\), every \(1\le j<N\), and \(\mu^{\otimes j}\)-a.e. \(X_j\),
\begin{align}
 C^{-1}\le H_{N,j}(x;X_j)&\le C
 \quad\text{for \(\mu\)-a.e. \(x\)},
 \label{eq:uv-critical-H-uniform-bound}\\
 C^{-1}\le \mathcal N_{N,j}(X_j)&\le C.
 \label{eq:uv-critical-normalization-uniform-bound}
\end{align}
Moreover,
\begin{align}\label{eq:uv-critical-conditional-uniform-upper}
 0\le q_{N,j}(x\mid X_j)\le C
 \quad\text{for \(\mu\)-a.e. \(x\)},
\end{align}
and, for \(x,x'\in K\) with \(\abs{x-x'}\le r\),
\begin{align}\label{eq:uv-critical-H-modulus}
 \abs{H_{N,j}(x;X_j)-H_{N,j}(x';X_j)}
 \le C\omega_K(r).
\end{align}
\end{lemma}

\begin{proof}
Suppose first that \(n\ge1\).  Conditional on all reservoir coordinates except \(Y_\ell\), the density of \(Y_\ell\) relative to \(\mu\) is proportional to
\begin{align}\label{eq:uv-critical-reservoir-coordinate-weight}
 \exp\left\{
 \be h_\mu(y)-\frac\be N\left(
 \sum_{m\ne\ell}W(y,Y_m)+\sum_{a=1}^jW(x_a,y)
 \right)\right\}.
\end{align}
The numerator in \eqref{eq:uv-critical-reservoir-coordinate-weight} is at most \(e^{\be H_\mu}\).  By Jensen's inequality, its normalizing denominator is at least the exponential of
\begin{align}\label{eq:uv-critical-reservoir-coordinate-denominator}
 2\be I_\mu
 -\frac\be N\left(\sum_{m\ne\ell}h_\mu(Y_m)+\sum_{a=1}^jh_\mu(x_a)\right)
 \ge-\be H_\mu,
\end{align}
because the two sums contain \(n-1+j=N-2\) terms.  Thus, every one-coordinate conditional density, and hence every one-coordinate marginal density of \(\mathsf R_{N,j}^{X_j}\), is bounded above by
\begin{align}\label{eq:uv-critical-reservoir-onepoint-upper}
 e^{2\be H_\mu}.
\end{align}

The upper bound on \(H_{N,j}\) follows from \(W\ge0\) and \eqref{eq:uv-Hmu-bound}.  For the lower bound, Jensen's inequality and \eqref{eq:uv-critical-reservoir-onepoint-upper} give
\begin{align}\label{eq:uv-critical-H-lower}
 H_{N,j}(x;X_j)
 &\ge
 \exp\left\{
 \be h_\mu(x)-\frac{\be n}{N}
 \int W(x,y)\,\md(\mathsf R_{N,j}^{X_j})^{(1)}(y)
 \right\}\notag\\
 &\ge \exp\left\{-\be e^{2\be H_\mu}H_\mu\right\}.
\end{align}
The same bounds are immediate when \(n=0\).  The upper bound for \(\mathcal N_{N,j}\) follows from \(H_{N,j}\le C\) and \(W\ge0\).  For the lower bound, Jensen's inequality gives
\begin{align}\label{eq:uv-critical-normalization-lower}
 \mathcal N_{N,j}(X_j)
 &\ge C^{-1}\int
 \exp\left\{-\frac\be N\sum_{a=1}^jW(x_a,x)\right\}\md\mu(x)\notag\\
 &\ge C^{-1}\exp\left\{-\frac\be N\sum_{a=1}^jh_\mu(x_a)\right\}
 \ge C^{-1}e^{-\be H_\mu}.
\end{align}
Equation \eqref{eq:uv-critical-conditional-uniform-upper} follows from \eqref{eq:uv-critical-marked-conditional-factorization}, \eqref{eq:uv-critical-H-uniform-bound}--\eqref{eq:uv-critical-normalization-uniform-bound}, and \(W\ge0\).

To prove \eqref{eq:uv-critical-H-modulus}, note that the marked-reservoir law is independent of the new point \(x\), while the exponential in \eqref{eq:uv-critical-marked-H-def} belongs to \([0,1]\).  Therefore,
\begin{align}
 &\left|
 \E_{\mathsf R_{N,j}^{X_j}}\left[
 \exp\left\{-\frac\be N\sum_{\ell=1}^nW(x,Y_\ell)\right\}\right]
 -\E_{\mathsf R_{N,j}^{X_j}}\left[
 \exp\left\{-\frac\be N\sum_{\ell=1}^nW(x',Y_\ell)\right\}\right]
 \right|\notag\\
 &\qquad\le\frac\be N\sum_{\ell=1}^n
 \E_{\mathsf R_{N,j}^{X_j}}\left[
 \abs{W(x,Y_\ell)-W(x',Y_\ell)}
 \right]
 \le C\omega_K(r).
 \label{eq:uv-critical-reservoir-expectation-modulus}
\end{align}
Also, \(\abs{h_\mu(x)-h_\mu(x')}\le\omega_K(r)\), and hence \eqref{eq:uv-critical-H-modulus} follows from \eqref{eq:uv-critical-marked-H-def} and the boundedness of \(h_\mu\).
\end{proof}

\parheading{Conditional deficit from good anchors.}

Choose concentric balls \(B_0\Subset B_1\Subset B\), and fix \(\alpha\) in the range \eqref{eq:uv-prefactor-alpha-range}.  Set
\begin{align}\label{eq:uv-critical-rho-def}
 \rho_N\coloneqq N^{-\alpha},
\end{align}
and fix \(c_->0\) sufficiently small that \(\be(2\s)^{-1}c_-^{-\s}>1\).  The upper inequalities in \eqref{eq:uv-prefactor-alpha-range} ensure that \(N^{-1/\s}\ll\rho_N\).

For \(X_j=(x_1,\ldots,x_j)\) and \(1\le a\le j\), define the truncated gradient load
\begin{align}\label{eq:uv-critical-Lambda-def}
 \Lambda_{N,j,a}(X_j)
 \coloneqq
 \frac1N\sum_{\substack{1\le b\le j\\b\ne a}}
 \abs{x_b-x_a}^{-\s-1}
 \one_{\{\abs{x_b-x_a}\ge4\rho_N\}}.
\end{align}
A good anchor is an observed particle lying in the interior ball, isolated from the other observed particles at scale \(\rho_N\), and subject to a controlled far-field gradient load.  Precisely, call the index \(a\) a good anchor if
\begin{align}\label{eq:uv-critical-good-anchor-def}
 x_a\in B_0,
 \qquad
 \min_{b\ne a}\abs{x_b-x_a}\ge4\rho_N,
 \qquad
 \rho_N^{\d+2}\Lambda_{N,j,a}(X_j)^2
 \le\eta\mathfrak a_N,
\end{align}
where \(\eta>0\) will be fixed sufficiently small.  When \(j=1\), the other-particle sum is empty, \(\Lambda_{N,1,1}=0\), and the isolation condition is vacuous; equivalently, the minimum over the empty set is interpreted as \(+\infty\).  Let \(G_{N,j}(X_j)\) denote the number of good anchors.

\begin{lemma}[Good-anchor conditional deficit]
\label{lem:uv-critical-good-anchor-deficit}
After choosing \(\eta>0\) sufficiently small and then increasing \(N_0\), for every \(N\ge N_0\), every \(1\le j<N\), and \(\mu^{\otimes j}\)-a.e. \(X_j\),
\begin{align}\label{eq:uv-critical-conditional-deficit-good-count}
 \delta_{N,j}(X_j)\ge c\mathfrak a_NG_{N,j}(X_j).
\end{align}
\end{lemma}

\begin{proof}
Fix a good anchor \(a\).  For \(x\in B(x_a,\rho_N)\), factor \eqref{eq:uv-critical-marked-conditional-factorization} as
\begin{align}\label{eq:uv-critical-sqrt-q-anchor-factorization}
 \sqrt{q_{N,j}(x\mid X_j)}
 =U_{N,j,a}(x;X_j)
 \exp\left\{-\frac\be{2N}W(x_a,x)\right\},
\end{align}
where
\begin{align}\label{eq:uv-critical-U-anchor-def}
 U_{N,j,a}(x;X_j)
 \coloneqq
 \mathcal N_{N,j}(X_j)^{-1/2}H_{N,j}(x;X_j)^{1/2}
 \exp\left\{-\frac\be{2N}\sum_{b\ne a}W(x_b,x)\right\}.
\end{align}
By \Cref{lem:uv-critical-marked-reservoir}, \(U_{N,j,a}\) is uniformly bounded.  If \(r=\abs{x-x_a}\le\rho_N\), the isolation condition in \eqref{eq:uv-critical-good-anchor-def} and the mean-value theorem give
\begin{align}\label{eq:uv-critical-U-anchor-modulus}
 \abs{U_{N,j,a}(x;X_j)-U_{N,j,a}(x_a;X_j)}
 \le C\pa{\omega_{B_1}(r)+r\Lambda_{N,j,a}(X_j)}.
\end{align}
Indeed, the first term follows from \eqref{eq:uv-critical-H-modulus}; for the other marked particles, the segment joining \(x_a\) to \(x\) stays a fixed proportion of \(\abs{x_b-x_a}\) away from \(x_b\), and \(\abs{\nabla W(z)}\le C\abs z^{-\s-1}\).

Let
\begin{align}\label{eq:uv-critical-anchor-annulus}
 \mathcal T_{N,a}
 \coloneqq
 \left\{x:\ c_-N^{-1/\s}\le\abs{x-x_a}\le\rho_N\right\}.
\end{align}
For large \(N\), this annulus lies in \(B\).  Set
\(v_a=U_{N,j,a}(x_a;X_j)\) and \(t(r)=\be(2\s N)^{-1}r^{-\s}\).  For \(x\in\mathcal T_{N,a}\), \eqref{eq:uv-critical-sqrt-q-anchor-factorization} gives the exact decomposition
\begin{align}
 \sqrt{q_{N,j}(x\mid X_j)}-1
 &=\left(v_ae^{-t(r)}-1\right)
 +\left(U_{N,j,a}(x;X_j)-v_a\right)e^{-t(r)}.
 \label{eq:uv-critical-anchor-decomposition}
\end{align}
The local lower bound \eqref{eq:uv-critical-local-density}, polar coordinates around \(x_a\), and the critical identity \(\d=2\s\) imply, with
\(\epsilon_N=\be(2\s)^{-1}N^{\alpha\s-1}\),
\begin{align}
 &\inf_{v\ge0}\int_{\mathcal T_{N,a}}
 \left|ve^{-t(\abs{x-x_a})}-1\right|^2\,\md\mu(x)\notag\\
 &\qquad\ge cN^{-2}\inf_{v\in\R}
 \int_{\epsilon_N}^1\left|ve^{-t}-1\right|^2t^{-3}\,\md t
 \ge cN^{-2}\log N
 =c\mathfrak a_N.
 \label{eq:uv-critical-anchor-scalar-main}
\end{align}
Here, the upper endpoint in the \(t\)-variable exceeds one by the choice of \(c_-\), while \(\epsilon_N\downarrow0\) because \(\alpha\s<1\); the second inequality is \Cref{lem:uv-scalar-transition-annulus} with \(\lambda=1\).  On the other hand, \eqref{eq:uv-critical-U-anchor-modulus} and \(e^{-t(r)}\le1\) give
\begin{align}
 &\int_{\mathcal T_{N,a}}
 \left|U_{N,j,a}(x;X_j)-v_a\right|^2e^{-2t(r)}\,\md\mu(x)\notag\\
 &\qquad\le C\int_0^{\rho_N}\omega_{B_1}(r)^2r^{\d-1}\,\md r
 +C\rho_N^{\d+2}\Lambda_{N,j,a}(X_j)^2.
 \label{eq:uv-critical-anchor-separated-errors}
\end{align}
Applying \((u+w)^2\ge\frac12u^2-w^2\) to \eqref{eq:uv-critical-anchor-decomposition}, and then using \eqref{eq:uv-critical-anchor-scalar-main} and \eqref{eq:uv-critical-anchor-separated-errors}, gives
\begin{align}
 \int_{\mathcal T_{N,a}}\pa{\sqrt{q_{N,j}(x\mid X_j)}-1}^2\,\md\mu(x)
 &\ge c\mathfrak a_N\notag\\
 &\quad-C\int_0^{\rho_N}\omega_{B_1}(r)^2r^{\d-1}\,\md r
 -C\rho_N^{\d+2}\Lambda_{N,j,a}(X_j)^2.
 \label{eq:uv-critical-anchor-deficit-with-errors}
\end{align}
The first error is the translation-modulus contribution.  By \eqref{eq:uv-translation-modulus-bound},
\begin{align}\label{eq:uv-critical-deterministic-modulus-error}
 \int_0^{\rho_N}\omega_{B_1}(r)^2r^{\d-1}\,\md r
 \le C
 \begin{cases}
 \rho_N^2,&\d=1,\\
 \rho_N^4\log^2(e/\rho_N),&\d=2,\\
 \rho_N^{\d+2},&\d\ge3.
 \end{cases}
\end{align}
The lower inequalities in \eqref{eq:uv-prefactor-alpha-range} make \eqref{eq:uv-critical-deterministic-modulus-error} equal to \(o(\mathfrak a_N)\).  The last term in \eqref{eq:uv-critical-anchor-deficit-with-errors} is at most \(C\eta\mathfrak a_N\) by goodness.  Taking \(\eta\) small proves a lower bound \(c\mathfrak a_N\) on each good-anchor annulus.

The annuli associated with distinct good anchors are disjoint because their centers are separated by at least \(4\rho_N\).  Summing their contributions and using \eqref{eq:uv-critical-conditional-affinity-deficit} proves \eqref{eq:uv-critical-conditional-deficit-good-count}.
\end{proof}

\parheading{Product-law abundance of good anchors.}

\begin{prop}[Uniform good-anchor count]
\label{prop:uv-critical-good-anchor-count}
There are \(c_B>0\), \(c_0>0\), and \(N_0\) such that, for every \(N\ge N_0\) and every integer satisfying \eqref{eq:uv-critical-k-window},
\begin{align}\label{eq:uv-critical-expected-good-count}
 \E_{\mu^{\otimes j}}\left[G_{N,j}\right]\ge c_Bj,
 \qquad 1\le j\le k.
\end{align}
\end{prop}

\begin{proof}
For a fixed anchor \(a\), the probability that \(X_a\in B_0\) and another observed point lies within \(4\rho_N\) is at most
\begin{align}\label{eq:uv-critical-close-anchor-probability}
 Cj\rho_N^\d.
\end{align}
Conditional on \(X_a=x\in B_0\), the remaining variables are independent.  If \(Y\sim\mu\), then boundedness of the density and the critical identity \(2\s=\d\) give, uniformly in \(x\in B_0\),
\begin{align}
 \E_\mu\left[\abs{x-Y}^{-2\s-2}
 \one_{\{\abs{x-Y}\ge4\rho_N\}}\right]
 &\le C\rho_N^{-2},
 \label{eq:uv-critical-gradient-kernel-second-moment}\\
 \E_\mu\left[\abs{x-Y}^{-\s-1}
 \one_{\{\abs{x-Y}\ge4\rho_N\}}\right]
 &\le C
 \begin{cases}
 \rho_N^{-1/2},&\d=1,\\
 \log(e/\rho_N),&\d=2,\\
 1,&\d\ge3.
 \end{cases}
 \label{eq:uv-critical-gradient-kernel-first-moment}
\end{align}
The first estimate follows because the radial integrand is proportional to \(r^{-3}\); the second follows by integrating \(r^{\d/2-2}\).  Expanding the square in \eqref{eq:uv-critical-Lambda-def} and using conditional independence, we obtain
\begin{align}\label{eq:uv-critical-Lambda-second-moment}
 \E_{\mu^{\otimes j}}\left[
 \Lambda_{N,j,a}(X_j)^2\,\middle|\,X_a=x
 \right]
 \le\frac C{N^2}\left[
 j\rho_N^{-2}+j^2
 \begin{cases}
 \rho_N^{-1},&\d=1,\\
 \log^2(e/\rho_N),&\d=2,\\
 1,&\d\ge3.
 \end{cases}
 \right].
\end{align}
Markov's inequality and \(\mathfrak a_N=N^{-2}\log N\) therefore show that the probability that \(X_a\in B_0\) and the gradient-load condition in \eqref{eq:uv-critical-good-anchor-def} fails is at most
\begin{align}\label{eq:uv-critical-gradient-bad-probability}
 \frac C{\log N}\left[
 j\rho_N^\d+j^2\rho_N^{\d+2}
 \begin{cases}
 \rho_N^{-1},&\d=1,\\
 \log^2(e/\rho_N),&\d=2,\\
 1,&\d\ge3.
 \end{cases}
 \right].
\end{align}
For \(j\le c_0N/\sqrt{\log N}\), the expressions in \eqref{eq:uv-critical-close-anchor-probability} and \eqref{eq:uv-critical-gradient-bad-probability} converge to zero uniformly in \(j\).  More explicitly, the potentially largest terms are bounded, up to logarithmic factors, by
\begin{align}\label{eq:uv-critical-dimension-rate-table}
 \begin{cases}
 N^{1-\alpha}\quad\text{and}\quad N^{2-2\alpha},&\d=1,\\
 N^{1-2\alpha}\quad\text{and}\quad N^{2-4\alpha},&\d=2,\\
 N^{1-\alpha\d}\quad\text{and}\quad N^{2-\alpha(\d+2)},&\d\ge3.
 \end{cases}
\end{align}
Every exponent in \eqref{eq:uv-critical-dimension-rate-table} is negative by \eqref{eq:uv-prefactor-alpha-range}; for \(\d\ge3\), note that \(2/(\d+2)>1/\d\).  Therefore, for all sufficiently large \(N\), a fixed anchor is good with probability at least \(\mu(B_0)/2\), uniformly in \(j\) throughout the stated window.  Exchangeability yields \eqref{eq:uv-critical-expected-good-count} with \(c_B=\mu(B_0)/2\).
\end{proof}

\begin{proof}[Proof of \Cref{prop:uv-critical-mesoscopic-TV}]
By \Cref{lem:uv-critical-good-anchor-deficit,prop:uv-critical-good-anchor-count}, the functions \(G_{N,j}\) satisfy \eqref{eq:uv-critical-abstract-count-range}--\eqref{eq:uv-critical-abstract-count} throughout the window \eqref{eq:uv-critical-k-window}, after decreasing \(c_0\) and increasing \(N_0\), if necessary.  The accumulation conclusion \eqref{eq:uv-critical-abstract-accumulation} of \Cref{lem:uv-critical-affinity-recursion} proves \eqref{eq:uv-critical-mesoscopic-TV-lower}.
\end{proof}

\begin{cor}[Mesoscopic failure of all three classical local scales]
\label{cor:uv-mesoscopic-all-divergences}
Assume either that \(\s=\d/2\) and \(\mu\ge m_B>0\) on some ball, or that \(\d/2<\s<\d\).  For every sufficiently small fixed \(\kappa>0\), set
\begin{align}\label{eq:uv-mesoscopic-common-kN}
 k_N\coloneqq\left\lfloor\kappa\mathfrak a_N^{-1/2}\right\rfloor.
\end{align}
Then
\begin{align}
 \liminf_{N\to\infty}
 \TV\pa{\Q_{N,\be}^{(k_N)}(\mu),\mu^{\otimes k_N}}&>0,
 \label{eq:uv-mesoscopic-TV-liminf}\\
 \liminf_{N\to\infty}
 \Hr\pa{\Q_{N,\be}^{(k_N)}(\mu)\mid\mu^{\otimes k_N}}&>0,
 \label{eq:uv-mesoscopic-H-liminf}\\
 \liminf_{N\to\infty}
 \Chi\pa{\Q_{N,\be}^{(k_N)}(\mu)\mid\mu^{\otimes k_N}}&>0.
 \label{eq:uv-mesoscopic-chi-liminf}
\end{align}
Moreover,
\begin{align}
 \frac{k_N}{N}
 &\asymp
 \begin{cases}
 (\log N)^{-1/2},&\s=\d/2,\\
 N^{\d/(2\s)-1},&\d/2<\s<\d,
 \end{cases}
 \longrightarrow0,
 \label{eq:uv-mesoscopic-vanishing-fraction}\\
 \frac{k_N^2}{N^2}
 &\asymp
 \begin{cases}
 (\log N)^{-1},&\s=\d/2,\\
 N^{\d/\s-2},&\d/2<\s<\d,
 \end{cases}
 \longrightarrow0.
 \label{eq:uv-mesoscopic-subcritical-scale-vanishes}
\end{align}
\end{cor}

\begin{proof}
The total-variation lower bound follows from \Cref{prop:uv-mesoscopic-close-pair,prop:uv-critical-mesoscopic-TV}.  Pinsker's inequality and \(\Chi(\nu\mid\pi)\ge4\TV(\nu,\pi)^2\) give the relative-entropy and chi-square conclusions.  The final two limits follow from \eqref{eq:intro-uv-scale} and \eqref{eq:uv-mesoscopic-common-kN}.
\end{proof}

Consequently, no estimate of the form
\begin{align}\label{eq:uv-impossible-k-over-N-bound}
 \TV\pa{\Q_{N,\be}^{(k)}(\mu),\mu^{\otimes k}}
 \le C_{\be,\mu}\frac{k}{N}
\end{align}
can hold uniformly in \(1\le k\le N\) throughout \(\d/2\le\s<\d\).

\begin{proof}[Proof of \Cref{thm:intro-ultraviolet-sharpness}]
Part~\textup{(i)} is \Cref{cor:uv-fixed-particle-all-divergences}.  Part~\textup{(ii)} follows from \Cref{prop:uv-mesoscopic-close-pair,prop:uv-critical-mesoscopic-TV} and \Cref{cor:uv-mesoscopic-all-divergences}.
\end{proof}

% STAGE2 MOVE: sec:further-questions
% BEGIN WITHHELD SECTION 11 -- author instruction 2026-09-20; source-only, not for publication
% WITHHELD11 \section{Further questions}
% WITHHELD11 \label{sec:further-questions}
% WITHHELD11 
% WITHHELD11 % STAGE2 MOVE: A7-conclusion-synthesis
% WITHHELD11 Below the Hilbert--Schmidt threshold, the order theorem controls every marginal size, while the quantitative profile theorem identifies the leading correction to independence, the exact fixed-marginal limits, and the normalized divergence limits on growing sublinear marginals.  The ultraviolet theorem gives fixed-particle and mesoscopic lower obstructions, but not matching asymptotics.  The questions below concern the remaining scales, sharper rates, and extensions of the interaction class.
% WITHHELD11 
% WITHHELD11 % STAGE2 MOVE: A7-five-further-questions-preserved
% WITHHELD11 \parheading{Sharp ultraviolet asymptotics and mesoscopic crossover.}
% WITHHELD11 With \(\mathfrak a_N\) as in \eqref{eq:intro-uv-scale}, \Cref{thm:intro-ultraviolet-sharpness} identifies the fixed-particle lower scale and proves a total-variation lower bound of order \(k^2\mathfrak a_N\) up to the mesoscopic scale \(k\asymp\mathfrak a_N^{-1/2}\), but it does not provide matching upper bounds or limiting profiles.  Natural questions are to determine leading constants and limiting laws under local nondegeneracy assumptions on \(\mu\), both for fixed \(k\) and in triangular regimes in which \(k=k_N\) grows with \(N\).  In total variation, one would like to know whether there is first a crossover from an ordinary mean-field contribution of order \(k/N\) to the ultraviolet contribution of order \(k^2\mathfrak a_N\), before the later scale \(k\asymp\mathfrak a_N^{-1/2}\) on which chaos is lost altogether.
% WITHHELD11 
% WITHHELD11 The critical and strictly supercritical regimes should be treated separately.  At criticality, the lower bound is produced by multiscale accumulation of conditional Hellinger deficits, whereas, in the strictly supercritical regime, it is produced by close-pair depletion.  A sharper theory should also clarify the relation between the local marginal obstruction and the divergence of the modulated partition function observed by the authors in \cite{DGR26}, for example through the exact marked identities or differentiation with respect to the inverse temperature.
% WITHHELD11 
% WITHHELD11 \parheading{Optimal rates and finer growing-marginal asymptotics.}
% WITHHELD11 \Cref{thm:quantitative-exact-marginal-profile,thm:quantitative-growing-marginal-constants} give explicit rates for the complete singleton-plus-pair profile and identify all three normalized divergence limits for every growing sublinear marginal, in both the Riesz and logarithmic cases under the stated hypotheses.  The natural-scale profile requires only $k_N=o(N)$, whereas \Cref{cor:quantitative-literal-growing-window} gives narrower sufficient windows for the undivided errors.  The convergence exponents and these latter windows are not claimed optimal.  Determining optimal rates and the largest undivided-error windows remains open, as does identifying further correction terms.  A separate question is to describe the crossover when the observed marginal is a nonvanishing fraction of the full system; the present sublinear limits do not identify that regime.
% WITHHELD11 
% WITHHELD11 \parheading{Ultraviolet thresholds for higher-order statistical divergences.}
% WITHHELD11 For a fixed \(\tau>2\), the direct-pair calculation suggests that the relevant Riesz threshold is \(\s=\d/\tau\), rather than the Hilbert--Schmidt threshold \(\d/2\).  For \(c>0\), the model integral is
% WITHHELD11 \begin{align}\label{eq:further-Ltau-direct-pair-model}
% WITHHELD11  \int_0^1
% WITHHELD11  \left|1-\exp\left(-\frac{c}{Nr^\s}\right)\right|^\tau
% WITHHELD11  r^{\d-1}\,\md r,
% WITHHELD11 \end{align}
% WITHHELD11 which has the predicted order
% WITHHELD11 \begin{align}\label{eq:further-Ltau-direct-pair-orders}
% WITHHELD11  \begin{cases}
% WITHHELD11   N^{-\tau},&\tau\s<\d,\\[0.2em]
% WITHHELD11   N^{-\tau}\log N,&\tau\s=\d,\\[0.2em]
% WITHHELD11   N^{-\d/\s},&\tau\s>\d.
% WITHHELD11  \end{cases}
% WITHHELD11 \end{align}
% WITHHELD11 This suggests a family of statistical-divergence-dependent ultraviolet transitions.  Since \(\d/\tau<\d/2\), the classical fixed-particle estimate \(\mathscr D_\tau\asymp(k/N)^\tau\) cannot be expected throughout the Hilbert--Schmidt regime.  Proving the positive result below \(\d/\tau\), together with sharp critical and supercritical obstructions and the corresponding growing-marginal behavior, would require new arguments.  The present method is intrinsically \(L^2\)-based: the screened one-point equation is solved in \(L^2(\mu)\), while the uniform all-marginal transfer uses Hilbert-space orthogonality and the exact decomposition of chi-square divergence into squared \(L^2\)-norms.  An extension would therefore require new \(L^\tau\) cavity estimates and non-Hilbertian control of the Hoeffding levels, or a direct \(L^\tau\) analysis of the marked partition function.  For the logarithmic kernel, no finite-\(\tau\) diagonal threshold is suggested by local integrability, since every finite power of the logarithmic singularity is locally integrable.  These predictions remain conjectural.
% WITHHELD11 
% WITHHELD11 \parheading{Universality of the ultraviolet mechanism.}
% WITHHELD11 The ultraviolet results in \Cref{thm:intro-ultraviolet-sharpness} are proved for homogeneous Riesz kernels, whereas the subcritical criterion is formulated for a much broader class of interactions.  It is natural to ask which parts of the critical and supercritical conclusions depend only on the local behavior of the interaction near the collision diagonal.  Relevant examples include regularly varying radial interactions, Riesz singularities plus smoother perturbations, and anisotropic singular kernels.
% WITHHELD11 
% WITHHELD11 In the strictly supercritical regime, the leading scale may be governed largely by the product probability of the close-pair set and the amount by which the mutual Boltzmann factor depletes it.  At criticality, however, the proof also uses a family of transition annuli, translation-modulus estimates, and uniform marked-reservoir bounds.  Thus, a tail condition on \(\tl W(X,Y)\) alone is unlikely to characterize the critical universality class.  The appropriate problem is to identify local geometric and conditional regularity hypotheses under which the same scales and mechanisms persist.
% WITHHELD11 
% WITHHELD11 \parheading{Stable mixed-sign interactions.}
% WITHHELD11 The partition-function spectral-stability result in \Cref{lem:partition-spectral-stability} shows that positivity of \(T_\mu\) is not essential for the screened one-point estimate.  The remaining model problem is to verify the partition-function and cavity hypotheses for concrete stable interactions for which \(T_\mu\) has both positive and negative spectral components.  The attractive logarithmic interaction is a principal test case, especially throughout the full inverse-temperature range in which the relevant modulated partition functions remain uniformly controlled.
% WITHHELD11 
% WITHHELD11 The global \(L^2\) identity requires control of the partition function at inverse temperature \(2\be\), but this may be stronger than the integrability needed for a fixed marginal.  Indeed, squaring the marked representation of a \(k\)-particle marginal produces two independent inverse-temperature-\(\be\) reservoirs sharing the same \(k\) marked particles, rather than one reservoir at inverse temperature \(2\be\).  It is therefore natural to ask whether fixed, or quantitatively growing, marginals retain local \(L^2\) control throughout the full stable regime, even when the complete Gibbs density has only subquadratic integrability.  Such a result would recover the quadratic local chi-square and relative-entropy upper bounds and would raise the corresponding question of matching lower bounds.  The fixed-marginal estimate in \Cref{thm:fixed-marginal-Ltau} supplies order-\(N^{-1}\) total-variation control under a subquadratic partition-function moment, but it does not give a quantitative growing-marginal range or the sharp quadratic information scale.  Resolving these questions would clarify which parts of the sharp size-of-chaos theory are specifically Hilbertian and which persist throughout the full stable mean-field regime.
% WITHHELD11 
% END WITHHELD SECTION 11

% STAGE2 MOVE: A3-appendix-delimiter
\appendix
\crefalias{section}{appendix}

% STAGE2 MOVE: A3-weighted-log-order-appendix-preserved
\section{Weighted logarithmic cavity estimates}
\label{app:weighted-log-cavity}

The statement of \Cref{lem:weighted-logarithmic-cavity} is kept with the model verification in \Cref{sec:verification-log-riesz}.  We give its proof here so that the main one-particle argument is not interrupted by this secondary weighted extension.

\begin{proof}[Proof of \Cref{lem:weighted-logarithmic-cavity}]
We first introduce the notation used only in this appendix.  For \(\theta\ge0\), set
\begin{align}\label{eq:weighted-log-proof-weights}
 W_\theta(x)&\coloneqq e^{\theta L(x)}=(e+|x|)^\theta,
 &
 a(x)&\coloneqq h_\mu(x)-2I_\mu,\\
 c&\coloneqq \mathsf c_{\log,\be,\mu},
 &
 \Theta&\coloneqq pc,
 \qquad
 \Theta_*\coloneqq p\Theta=p^2c.
\end{align}
Define
\begin{align}\label{eq:weighted-log-proof-GA}
 G_*(x)&\coloneqq
 \left[\int_{\R^\d}|\tl W(x,y)|^p W_{\Theta_*}(y)\,\md\mu(y)\right]^{1/p},
 &
 A_*(x)&\coloneqq |a(x)|^2+G_*(x)^2,
\end{align}
and
\begin{align}\label{eq:weighted-log-proof-M}
 M_*\coloneqq\int_{\R^\d}W_{\Theta_*}\,\md\mu.
\end{align}
The hypotheses \eqref{eq:weighted-logarithmic-inputs} say precisely that
\begin{align}\label{eq:weighted-log-proof-inputs}
 M_*<\infty,
 \qquad
 \norm{W_{2c}A_*}_{L^2(\mu)}<\infty.
\end{align}
We shall prove the stronger pointwise estimate
\begin{align}\label{eq:weighted-log-pointwise-B2}
 \E_{\Q_{N,\be}(\mu)}\left[
 e^{\lambda\Phi_N(x;\hat\Xi_1)}B_N(x;\hat\Xi_1)^2\right]
 \le
 \frac{\mathsf C_{\mathrm{wt},\be,p,\mu}}{N}
 W_{\lambda c}(x)A_*(x)
\end{align}
for every \(N\ge2\), every \(\lambda\in\{0,1,2\}\), and \(\mu\)-a.e. \(x\), where the constant depends only on the quantities in \eqref{eq:weighted-log-proof-inputs}, on \(p\), \(\be\), and on \(\mathfrak P_{r\be}(\mu)\).

\step{Step 1: Product-space reduction.}
Fix \(N\ge2\), \(\lambda\in\{0,1,2\}\), and \(x\).  Set \(s\coloneqq \lambda p/2\) and \(\eta_N\coloneqq sc/N\).  By \Cref{lem:full-density-Lr} and H\"older's inequality with conjugate exponents \(r=p/(p-2)\) and \(p/2\),
\begin{align}\label{eq:weighted-log-holder-step}
&\E_{\Q_{N,\be}(\mu)}\left[
 e^{\lambda\Phi_N(x;\hat\Xi_1)}B_N(x;\hat\Xi_1)^2\right] \notag\\
&\qquad\le
 \mathfrak P_{r\be}(\mu)^{1/r}
 \left\{\E_{\mu^{\otimes(N-1)}}\left[
 e^{s\Phi_N(x;\hat\Xi_1)}|B_N(x;\hat\Xi_1)|^p\right]\right\}^{2/p}.
\end{align}
Since
\begin{align}\label{eq:weighted-log-product-factorization}
 e^{s\Phi_N(x;\hat\Xi_1)}
 =e^{sc}W_{sc}(x)\prod_{j=2}^N W_{\eta_N}(\xi_j),
\end{align}
let
\begin{align}\label{eq:weighted-log-tilted-law}
 Z_\eta\coloneqq \int W_\eta\,\md\mu,
 \qquad
 \md\mu_\eta\coloneqq Z_\eta^{-1}W_\eta\,\md\mu.
\end{align}
Log-convexity of \(\eta\mapsto Z_\eta\), together with \(Z_0=1\), \(0\le\eta_N\le \Theta/N\), and \eqref{eq:weighted-log-proof-inputs}, gives
\begin{align}\label{eq:weighted-log-product-weight-bound}
 Z_{\eta_N}^{2(N-1)/p}
 \le \max\{1,M_*\}^{2/p}.
\end{align}
Consequently, \eqref{eq:weighted-log-holder-step} yields
\begin{align}\label{eq:weighted-log-holder-reduced}
&\E_{\Q_{N,\be}(\mu)}\left[
 e^{\lambda\Phi_N(x;\hat\Xi_1)}B_N(x;\hat\Xi_1)^2\right] \notag\\
&\qquad\le
 C_{\be,p,\mu} W_{\lambda c}(x)
 \left\{\E_{\mu_{\eta_N}^{\otimes(N-1)}}\left[
 |B_N(x;\hat\Xi_1)|^p\right]\right\}^{2/p}.
\end{align}
Here and below, \(C_{\be,p,\mu}\) denotes a finite constant depending only on the quantities in \eqref{eq:weighted-log-proof-inputs}, on \(p\), \(\be\), and on \(\mathfrak P_{r\be}(\mu)\).

\step{Step 2: Moment estimate under the tilted product law.}
Write \(k_x(y)\coloneqq \tl W(x,y)\).  Since \(\int k_x\,\md\mu=0\), the mean under the tilted law satisfies
\begin{align}\label{eq:weighted-log-tilted-mean}
 m_{\eta_N}(x)
 \coloneqq \int k_x(y)\,\md\mu_{\eta_N}(y)
 =Z_{\eta_N}^{-1}\int k_x(y)\pa{W_{\eta_N}(y)-1}\,\md\mu(y).
\end{align}
Moreover, \(Z_{\eta_N}\ge1\) and
\begin{align}\label{eq:weighted-log-weight-increment}
 0\le W_{\eta_N}(y)-1
 \le \eta_N L(y)W_{\eta_N}(y)
 \le \eta_N L(y)W_\Theta(y).
\end{align}
If \(c=0\), then \(\eta_N=0\) and \(m_{\eta_N}=0\).  If \(c>0\), H\"older's inequality and \(\Theta_*=p\Theta\) give
\begin{align}\label{eq:weighted-log-tilted-mean-bound}
 |m_{\eta_N}(x)|
 &\le \eta_N
 \int |k_x(y)|L(y)W_\Theta(y)\,\md\mu(y)\notag\\
 &\le \eta_N G_*(x)\norm{L}_{L^{p/(p-1)}(\mu)}\notag\\
 &\le \frac{C_{\be,p,\mu}}{N}G_*(x).
\end{align}
The logarithmic moment is finite by \eqref{eq:weighted-log-proof-inputs}, and the same bound is trivial when \(c=0\).

Also,
\begin{align}\label{eq:weighted-log-weighted-section}
 \left[\int |k_x(y)|^p\,\md\mu_{\eta_N}(y)\right]^{1/p}
 \le G_*(x),
\end{align}
because \(Z_{\eta_N}\ge1\) and \(W_{\eta_N}\le W_{\Theta_*}\).  Applying the scalar Rosenthal inequality under \(\mu_{\eta_N}^{\otimes(N-1)}\) to the centered variables \(k_x(\xi_j)-m_{\eta_N}(x)\), we have
\begin{align}\label{eq:weighted-log-centered-sum-bound}
 \left\{\E_{\mu_{\eta_N}^{\otimes(N-1)}}\left[
 \left|\frac1N\sum_{j=2}^N\pa{k_x(\xi_j)-m_{\eta_N}(x)}\right|^p\right]\right\}^{2/p}
 \le \frac{C_p}{N}G_*(x)^2.
\end{align}
Since
\begin{align}\label{eq:weighted-log-cavity-decomposition}
 B_N(x;\hat\Xi_1)
 =\frac{\be}{N}a(x)
 -\frac{\be}{N}\sum_{j=2}^N\pa{k_x(\xi_j)-m_{\eta_N}(x)}
 -\be\frac{N-1}{N}m_{\eta_N}(x),
\end{align}
\eqref{eq:weighted-log-tilted-mean-bound} and \eqref{eq:weighted-log-centered-sum-bound} give
\begin{align}\label{eq:weighted-log-Bp-bound}
 \left\{\E_{\mu_{\eta_N}^{\otimes(N-1)}}\left[
 |B_N(x;\hat\Xi_1)|^p\right]\right\}^{2/p}
 \le
 \frac{C_{\be,p,\mu}}{N}A_*(x).
\end{align}
Combining \eqref{eq:weighted-log-holder-reduced} and \eqref{eq:weighted-log-Bp-bound} proves \eqref{eq:weighted-log-pointwise-B2}.

\step{Step 3: Verification of the weighted cavity constants.}
It remains to check that \eqref{eq:weighted-log-pointwise-B2} implies \eqref{eq:log-weighted-cavity-constants}.  Taking \(\lambda=1\) gives
\begin{align}\label{eq:weighted-log-E-bound}
 N\norm{\E_{\Q_{N,\be}(\mu)}\left[
 e^{\Phi_N(\cdot;\hat\Xi_1)}B_N(\cdot;\hat\Xi_1)^2\right]}_{L^2(\mu)}
 \le
 \mathsf C_{\mathrm{wt},\be,p,\mu}
 \norm{W_cA_*}_{L^2(\mu)}<\infty,
\end{align}
and similarly
\begin{align}\label{eq:weighted-log-D-bound}
 N\E_{\Q_{N,\be}(\mu)}\left[
 \int e^{\Phi_N(x;\hat\Xi_1)}
 B_N(x;\hat\Xi_1)^2\,\md\mu(x)\right]
 \le
 \mathsf C_{\mathrm{wt},\be,p,\mu}
 \int W_cA_*\,\md\mu<\infty.
\end{align}
Taking \(\lambda=0\) gives the unweighted \(B_N^2\) estimate.  Finally, since \(\mu\) is a probability measure,
\begin{align}\label{eq:weighted-log-V-cauchy}
 \left(\int e^{\Phi_N(x;\hat\Xi_1)}|B_N(x;\hat\Xi_1)|\,\md\mu(x)\right)^2
 \le
 \int e^{2\Phi_N(x;\hat\Xi_1)}B_N(x;\hat\Xi_1)^2\,\md\mu(x),
\end{align}
and the case \(\lambda=2\) gives the estimate defining \(\mathsf V_{\log,\be,\mu}\).  The assumption \(\norm{W_{2c}A_*}_{L^2(\mu)}<\infty\) implies all the displayed integrability bounds above, and hence \eqref{eq:log-weighted-cavity-constants}.
\end{proof}

% STAGE2 MOVE: A3-constant-ledger-preserved
\section{Dependence of constants}\label{app:dependence-constants}

This appendix records one admissible choice of the constants in \Cref{thm:chi2-chaos-two-sided}.  The numerical values are not optimized; the purpose is to make the dependence on the quantitative hypotheses explicit.  All essential suprema and infima are taken with respect to \(\mu\) or \(\mu^{\otimes2}\), according to the number of variables.

Let \(\mathsf R_p\) be an admissible constant in the scalar Rosenthal inequality, and let \(\mathsf A_{2,p}\) be the fixed-order canonical \(U\)-statistic constant from \Cref{lem:fixed-order-canonical-U-moment}.  We also use the functions \(\Lambda_1\) and \(\Lambda_2\) defined in \eqref{eq:Lambda2-def}--\eqref{eq:Lambda1-def}.

Let \(\mathsf R_{\mathrm{cav},\be,\mu}\) denote the cavity constant supplied by \Cref{lem:cavity-linearization}; explicit choices are given in \eqref{eq:R-cav-bounded-def} and \eqref{eq:R-cav-log-def}.  Define
\begin{align}\label{eq:upper-constant-defs}
 \mathsf U_{\be,\mu}
 &\coloneqq
 2\left(
 \be\norm{h_\mu-2I_\mu}_{L^2(\mu)}
 +\mathsf R_{\mathrm{cav},\be,\mu}
 \right),\\
 \mathsf C_{\chi,\be,\mu}^{\mathrm{up}}
 &\coloneqq\mathsf U_{\be,\mu}^2+\pa{\mathfrak P_{2\be}(\mu)-1},\\
 \mathsf C_{\TV,\be,\mu}^{\mathrm{up}}
 &\coloneqq\frac12\sqrt{\mathsf C_{\chi,\be,\mu}^{\mathrm{up}}},
 &
 \mathsf C_{\Hr,\be,\mu}^{\mathrm{up}}
 &\coloneqq\mathsf C_{\chi,\be,\mu}^{\mathrm{up}}.
\end{align}

For the fixed-marginal branch, let \(1<\tau<2\), \(q=\tau/(\tau-1)\), and \(p\ge2q\).  The constants \(\mathsf A_{j,q}\) and \(\mathsf B_q\) are those in \Cref{lem:fixed-order-canonical-U-moment,lem:refined-canonical-pair-moment}, with the admissible choice \eqref{eq:canonical-U-moment-constants}.  Under the hypotheses of \Cref{thm:fixed-marginal-Ltau}, use \(\mathsf M_{\tau,\be,\mu}\) and \(\mathsf U_{\tau,\be,\mu}\) from \eqref{eq:fixed-marginal-Gibbs-constants}.  Substitution into \eqref{eq:abstract-fixed-marginal-explicit} gives exactly \eqref{eq:fixed-marginal-Gibbs-explicit}; no optimization of \(\mathsf A_{j,q}\), \(\mathsf B_q\), or the fixed-\(k\) onset is intended.

For \(\be>0\), recall \(\mathsf M_{p,\mu}\) from
\eqref{eq:Mp-mu-def} and \(\mathsf R_{p,\mu}\),
\(\mathsf a_{\be,\mu}\), and \(\mathsf J_{\be,\mu}\) from
\eqref{eq:jensen-cutoff-def}--\eqref{eq:jensen-gap-constant-def}.  The latter
three constants are strictly positive.  The lower-bound constants may be
chosen as
\begin{align}
 \mathsf N_{\be,\mu}
 &\coloneqq
 \left\lceil\max\left\{2,
 \frac{2\be\norm{h_\mu-2I_\mu}_{L^2(\mu)}\mathsf U_{\be,\mu}}
 {\mathsf J_{\be,\mu}}\right\}\right\rceil,
 \label{eq:lower-constant-N}\\
 \mathsf S_{\be,\mu}
 &\coloneqq\frac{\mathsf J_{\be,\mu}}{4\be},
 \label{eq:lower-constant-S}\\
 \mathsf R_{\TV,\be,\mu}
 &\coloneqq
 \max\left\{1,
 \left(\frac{2\pa{\frac{\mathsf A_{2,p}}2\norm{\tl W}_{L^p(\mu^{\otimes2})}}^{p/2}
 \sqrt{\mathsf C_{\chi,\be,\mu}^{\mathrm{up}}}}
 {\mathsf S_{\be,\mu}}\right)^{2/(p-2)}\right\},
 \label{eq:lower-constant-R-TV}\\
 \mathsf c_{\TV,\be,\mu}^{\mathrm{low}}
 &\coloneqq\frac{\mathsf S_{\be,\mu}}{4\mathsf R_{\TV,\be,\mu}},
 \label{eq:lower-constant-c-TV}\\
 \mathsf c_{\Hr,\be,\mu}^{\mathrm{low}}
 &\coloneqq2\pa{\mathsf c_{\TV,\be,\mu}^{\mathrm{low}}}^2,
 \label{eq:lower-constant-c-H}\\
 \mathsf c_{\chi,\be,\mu}^{\mathrm{low}}
 &\coloneqq
 \frac{\mathsf J_{\be,\mu}^2}
 {4\be^2\norm{\tl W}_{L^2(\mu^{\otimes2})}^2}.
 \label{eq:lower-constant-c-chi}
\end{align}

% STAGE2 MOVE: app:quantitative-replica-response
\section{Probability estimates}
\label{app:probability-estimates}
\label{app:quantitative-growing-marginals}

The bounded-feature estimate controls the positive correction in the two-sided response.  The martingale estimate gives the canonical-profile moments and the mixed-coordinate bounds in the smooth Gaussian comparison.  That comparison is used both by the signed Laplace theorem and by the absolute-moment approximation for growing total variation.

\subsection{Positive exponentials of bounded-feature statistics}
\label{app:quantitative-replica-response}
\label{subsec:bounded-feature-exponentials}

The following estimate controls the bounded positive correction for
negative response parameters in \Cref{subsec:screened-response-proof}.
It is not applied to the signed core interaction.
\begin{lemma}[Small positive exponential of a bounded-feature statistic]\label{quant-resp:lem-boundedfeatureexp}
Let $Y_1,\ldots,Y_N$ be independent copies of a centered random variable in a separable real Hilbert space, with $\norm{Y_i}^2\leq K$. If $4aK<1$ and $a\geq0$, then
\begin{equation}\label{quant-resp:eq-boundedfeatureexp}
 \E\left[\exp\left\{\frac a{2N}\sum_{i\ne j}\ipp*{Y_i,Y_j}\right\}\right]
 \leq(1-4aK)^{-1/2}.
\end{equation}
\end{lemma}
\begin{proof}
For $u\geq0$, introduce independent copies $Y_i'$. Convexity of $z\mapsto e^{u\norm z^2/N}$ and conditional Jensen give
\begin{equation}\label{quant-resp:eq-symmetrize}
 \E\left[e^{u\norm{\sum_iY_i}^2/N}\right]
 \leq\E\left[e^{u\norm{\sum_i(Y_i-Y_i')}^2/N}\right].
\end{equation}
The independent differences are symmetric. Conditional on their values $d_i$, insert independent Rademacher signs $\epsilon_i$. On the finite-dimensional span of the $d_i$, let $G$ be a standard Gaussian vector. The Gaussian moment-generating identity and $\cosh z\leq e^{z^2/2}$ give
\begin{align}
 \E_\epsilon\left[e^{u\norm{\sum_i\epsilon_id_i}^2/N}\right]
 &=\E_G\left[\prod_i\cosh\left(\sqrt{2u/N}\ipp*{G,d_i}\right)\right]\notag\\
 &\leq\E_G\left[e^{(u/N)\sum_i\ipp*{G,d_i}^2}\right].\label{quant-resp:eq-Gaussianbound}
\end{align}
Let $\lambda_j\geq0$ be the eigenvalues of $\sum_i d_i\otimes d_i$. Their sum is at most $4NK$. If $8uK<1$, the last expectation is
\begin{equation}\label{quant-resp:eq-Gaussianproduct}
 \prod_j(1-2u\lambda_j/N)^{-1/2}
 \leq\left(1-\frac{2u}N\sum_j\lambda_j\right)^{-1/2}
 \leq(1-8uK)^{-1/2}.
\end{equation}
The first inequality follows from $\prod_j(1-t_j)\geq1-\sum_jt_j$ for nonnegative $t_j$ with sum less than one, proved by induction. All calculations take place in a space of dimension at most $N$, so no infinite-dimensional standard Gaussian measure is being assumed.

Finally, $\sum_{i\ne j}\ipp*{Y_i,Y_j}=\norm{\sum_iY_i}^2-\sum_i\norm{Y_i}^2\leq\norm{\sum_iY_i}^2$. Apply \eqref{quant-resp:eq-symmetrize}--\eqref{quant-resp:eq-Gaussianproduct} with $u=a/2$.
\end{proof}

\subsection{A martingale moment inequality}
\label{quant-tools:sec-moments}

% STAGE2 MOVE: A3-general-martingale-proof-preserved
\begin{lemma}[A scalar martingale moment bound]\label{quant-tools:lem-martingale}
Let $(D_j)_{j=1}^m$ be real martingale differences in $L^p$, where $2\leq p<\infty$. Then
\begin{equation}\label{quant-tools:eq-martingale}
 \left\|\sum_{j=1}^mD_j\right\|_p^2\leq(p-1)\sum_{j=1}^m\norm{D_j}_p^2.
\end{equation}
The initial sigma algebra may be nontrivial.
\end{lemma}
\begin{proof}
For $p=2$, the result is orthogonality. Suppose $p>2$. Let $S$ be measurable with respect to a sigma algebra $\mathcal F$, let $D\in L^p$ satisfy $\E[D\mid\mathcal F]=0$, and suppose $S\in L^p$. For $\epsilon>0$, put
\begin{equation}\label{quant-tools:eq-smoothnorm}
 A_\epsilon(t)=\E[\abs{S+tD}^p]+\epsilon,
 \qquad J_\epsilon(t)=A_\epsilon(t)^{2/p},\qquad 0\leq t\leq1.
\end{equation}
Differentiation twice is valid by H\"older and domination by powers of $\abs S+\abs D$. If $V=S+tD$, direct calculation gives
\begin{align}
 J_\epsilon''(t)
 &=2(2-p)A_\epsilon(t)^{2/p-2}
       \bigl(\E[\abs V^{p-2}VD]\bigr)^2\notag\\
 &\quad+2(p-1)A_\epsilon(t)^{2/p-1}
       \E[\abs V^{p-2}D^2]
 \leq2(p-1)\norm D_p^2.\label{quant-tools:eq-smoothnormsecond}
\end{align}
The first term is nonpositive. By H\"older, the remaining expectation is at most
$(\E[\abs V^p])^{(p-2)/p}\norm D_p^2$. This cancels the nonpositive power of $A_\epsilon$. Moreover,
\begin{equation}\label{quant-tools:eq-smoothnormfirst}
 J_\epsilon'(0)=2A_\epsilon(0)^{2/p-1}\E[\abs S^{p-2}SD]=0,
\end{equation}
by conditional centering. Integrating \eqref{quant-tools:eq-smoothnormsecond} twice and sending $\epsilon\downarrow0$ gives
$\norm{S+D}_p^2\leq\norm S_p^2+(p-1)\norm D_p^2$.
Apply this successively with $S=\sum_{i<j}D_i$ to obtain \eqref{quant-tools:eq-martingale}. This is the usual $L^p$ smoothness argument; the underlying differential inequality is also given in~\cite[Proposition 2.1]{Pinelis1994}.
\end{proof}

% STAGE2 MOVE: quant-tools:sec-gaussian
\subsection{Gaussian quadratic forms and smooth comparison}
\label{quant-tools:sec-gaussian}

\subsubsection{The centered Gaussian quadratic form}
\label{quant-tools:subsec-chaos-construction}

Let $(S,\rho)$ be a standard Borel probability space and let $c\in L^2(\rho^{\otimes2})$ be real, symmetric, and canonical, with Hilbert--Schmidt operator $C$ on $L^2_0(\rho)$. For a finite-rank self-adjoint operator with matrix $(C_{rs})$ in an orthonormal basis of a finite-dimensional containing subspace, define
\begin{equation}\label{quant-tools:eq-finitechaos}
 \mathcal Q_C=\sum_{r,s}C_{rs}(Z_rZ_s-\delta_{rs}),
\end{equation}
using the coordinates of a single isonormal Gaussian process on $L^2_0(\rho)$. The elementary fourth-moment identity for independent standard Gaussians gives
\begin{equation}\label{quant-tools:eq-chaosisometry}
 \E[(Z_rZ_s-\delta_{rs})(Z_uZ_v-\delta_{uv})]
 =\delta_{ru}\delta_{sv}+\delta_{rv}\delta_{su},
 \qquad \E[\mathcal Q_C\mathcal Q_D]=2\Tr(CD).
\end{equation}
The first identity follows by independence when an index occurs only once, by $\E[Z^2]=1$ when indices occur in two distinct pairs, and by $\E[Z^4]=3$ when all indices agree. The second follows by expanding the finite sums and using symmetry. Orthogonal changes of basis preserve \eqref{quant-tools:eq-finitechaos}. Finite-rank approximation in Hilbert--Schmidt norm thus extends the definition to all self-adjoint Hilbert--Schmidt $C$, and
\begin{equation}\label{quant-tools:eq-chaoscontinuity}
 \norm{\mathcal Q_C-\mathcal Q_D}_2=\sqrt2\norm{C-D}_{\mathrm{HS}}.
\end{equation}
Diagonalizing $C$ gives the law $\sum_j\lambda_j(C)(Z_j^2-1)$ in $L^2$. This construction also justifies \eqref{eq:quantitative-centered-Gaussian-quadratic-form}.
\subsubsection{Proof of the smooth comparison}
\label{quant-tools:subsec-smooth-proof}

\begin{proof}[Proof of \Cref{lem:quadratic-smooth-comparison}]
Set $p=2+\alpha$.  All kernel norms in this proof are relative to
$\rho^{\otimes2}$.

\step{Finite partitions and mixed moments.}
Choose increasing finite measurable partitions whose generated sigma
algebra is the Borel sigma algebra modulo $\rho$.  Let $c_m$ be the
conditional expectation of $c$ onto the corresponding product partition,
ignoring zero-mass cells.  Then $c_m$ is symmetric and canonical, and
\begin{equation}\label{quant-tools:eq-partitionbounds}
 \norm{c_m}_p\leq\norm c_p,\qquad
 \norm{c_m-c}_2\longrightarrow0.
\end{equation}
The first assertion is conditional Jensen; the second follows by density
of product-partition simple functions in $L^2(\rho^{\otimes2})$.
Integrating the conditional expectation in either variable proves
canonicity.  An orthonormal basis $\phi_1,\ldots,\phi_\ell$ of the
centered functions constant on the partition cells gives
\begin{equation}\label{quant-tools:eq-finitefeature}
 c_m(x,y)=\phi(x)^{\mathsf T}A\phi(y),\qquad
 \E[\phi(x)]=0,\qquad \E[\phi(x)\phi(x)^{\mathsf T}]=\mathrm{Id},
 \qquad\norm A_{\mathrm{HS}}=\norm{c_m}_2,
\end{equation}
where $A$ is real symmetric.  There is no restriction on its sign.
A zero-dimensional centered subspace gives the zero kernel.

Let $\xi=\phi(x)$, and let $g$ be an independent standard Gaussian in
$\R^\ell$.  If $V,W$ are independent and each has either law $\xi$ or
law $g$, then
\begin{equation}\label{quant-tools:eq-mixedbilinear}
 \norm{V^{\mathsf T}AW}_p\leq C_p\norm{c_m}_p.
\end{equation}
For two original vectors this is the kernel norm.  For one original and
one Gaussian vector, conditioning on $\xi$ gives, with $Z$ a scalar
standard Gaussian,
\begin{align}
 \E[\abs{\xi^{\mathsf T}Ag}^p]
 &=\E[\abs Z^p]\int
       \left(\int\abs{c_m(x,y)}^2\,\md\rho(y)\right)^{p/2}\md\rho(x)
 \notag\\
 &\leq\E[\abs Z^p]\norm{c_m}_p^p.
 \label{quant-tools:eq-mixedoriginal}
\end{align}
For two Gaussian vectors, condition on the first and diagonalize $A$.
Minkowski in $L^{p/2}$ bounds
$\norm{\abs{Ag}}_p$ by $\norm Z_p\norm A_{\mathrm{HS}}$.  Hence,
\begin{equation}\label{quant-tools:eq-mixedGaussian}
 \norm{g^{\mathsf T}Ag'}_p
 \leq\norm Z_p^2\norm A_{\mathrm{HS}}\leq\norm Z_p^2\norm{c_m}_p.
\end{equation}
This proves \eqref{quant-tools:eq-mixedbilinear} with a dimension-independent
constant.

\step{Replacement of one coordinate at a time.}
For independent vectors $V_1,\ldots,V_N$ with these two possible laws, set
\begin{equation}\label{quant-tools:eq-quadraticfinite}
 U(V_1,\ldots,V_N)\coloneqq\frac1N\sum_{i<j}V_i^{\mathsf T}AV_j.
\end{equation}
When replacing $V_i=\xi_i$ by $g_i$, condition on all the other vectors.
The sum is $z+y$, where $z$ does not involve $V_i$ and
$y=N^{-1}V_i^{\mathsf T}A\sum_{j\ne i}V_j$.  The two candidate laws have
the same conditional first and second moments:
\begin{equation}\label{quant-tools:eq-matching}
 \E[y\mid(V_j)_{j\ne i}]=0,\qquad
 \E[y^2\mid(V_j)_{j\ne i}]
 =N^{-2}\left|A\sum_{j\ne i}V_j\right|^2.
\end{equation}
To estimate its $p$-moment, reveal $V_i$ first and the remaining vectors
afterward.  The terms $V_i^{\mathsf T}AV_j$ are martingale differences.
\Cref{quant-tools:lem-martingale} and
\eqref{quant-tools:eq-mixedbilinear} imply
\begin{equation}\label{quant-tools:eq-incrementmoment}
 \norm y_p\leq C_pN^{-1/2}\norm{c_m}_p.
\end{equation}
The fractional Taylor remainder satisfies
\begin{equation}\label{eq:signed-smooth-Taylor-remainder}
 \left|h(z+y)-h(z)-h'(z)y-\tfrac12h''(z)y^2\right|
 \leq\frac{[h'']_{C^{0,\alpha}}}{(1+\alpha)(2+\alpha)}\abs y^{2+\alpha}.
\end{equation}
Indeed, subtracting the quadratic term in Taylor's integral formula gives
$y^2\int_0^1(1-r)(h''(z+ry)-h''(z))\,\md r$.
The linear and quadratic terms cancel after conditional expectation.
Summing the $N$ replacement errors yields
\begin{equation}\label{eq:signed-smooth-Lindeberg}
 \left|\E[h(U(\xi_1,\ldots,\xi_N))]
       -\E[h(U(g_1,\ldots,g_N))]\right|
 \leq C_\alpha[h'']_{C^{0,\alpha}}N^{-\alpha/2}\norm{c_m}_p^p.
\end{equation}
All expectations are finite: $h$ is Lipschitz, its first two derivatives
are bounded, and \eqref{quant-tools:eq-incrementmoment} supplies the
Taylor-remainder integrability.  This is the one-coordinate moment-matching
argument of \cite[Section~2]{Chatterjee2006}, with the fractional and
vector-valued estimates made explicit here.

\step{The Gaussian statistic and passage to the kernel.}
An orthogonal change in the coordinate index, with first vector
$N^{-1/2}(1,\ldots,1)$, gives
\begin{equation}\label{quant-tools:eq-GaussianU}
 U(g_1,\ldots,g_N)\overset{\mathrm{law}}=
 \frac{N-1}{2N}\mathcal Q_{C_m}^{(0)}
 -\frac1{2N}\sum_{r=1}^{N-1}\mathcal Q_{C_m}^{(r)},
\end{equation}
where $C_m$ is the operator of $c_m$ and the superscripts denote independent
copies.  Diagonalize $A$ and use
$\sum_{i<j}g_ig_j=((\sum_i g_i)^2-\sum_i g_i^2)/2$ in each mode.
The trace constants cancel, proving the identity for either sign of $A$.
Couple the right-hand side with $\mathcal Q_{C_m}^{(0)}/2$.
Their difference is $-(2N)^{-1}\sum_{r=0}^{N-1}\mathcal Q_{C_m}^{(r)}$,
whose $L^2$ norm is $\norm{c_m}_2/\sqrt{2N}$ by
\eqref{quant-tools:eq-chaosisometry}.  Thus,
\begin{equation}\label{eq:signed-smooth-Gaussian-error}
 \left|\E[h(U(g_1,\ldots,g_N))]-\E[h(\mathcal Q_{C_m}/2)]\right|
 \leq\frac{\norm{h'}_\infty\norm{c_m}_2}{\sqrt{2N}}.
\end{equation}
For each fixed $N$, canonicity and the common Gaussian construction give
\begin{align}
 \left\|\tfrac12\mathsf W_N(c-c_m)\right\|_2
 &=\sqrt{\frac{N-1}{2N}}\norm{c-c_m}_2,
 \label{eq:signed-kernel-statistic-continuity}\\
 \left\|\tfrac12(\mathcal Q_C-\mathcal Q_{C_m})\right\|_2
 &=\frac{\norm{c-c_m}_2}{\sqrt2}.
 \label{eq:signed-kernel-Gaussian-continuity}
\end{align}
These quantities tend to zero.  Since $h$ is Lipschitz, we may send
$m\to\infty$ in \eqref{eq:signed-smooth-Lindeberg} and
\eqref{eq:signed-smooth-Gaussian-error}, using
\eqref{quant-tools:eq-partitionbounds}.  This proves
\eqref{eq:quadratic-smooth-comparison}.  No quantitative rate of the
partition approximation is required.
\end{proof}

\subsubsection{The absolute first moment}
\label{quant-tools:subsec-absolute-proof}

\begin{proof}[Proof of \Cref{quant-tools:prop-absolute}]
Set $p=2+\alpha$.  The zero-kernel case is immediate; for $k=1$, use
$\E[\abs{\mathcal Q_C}]/2\leq\norm c_2/\sqrt2$.  Suppose $k\geq2$ and
$\norm c_p>0$.  For $\delta>0$, set $h_\delta(x)=\sqrt{x^2+\delta^2}$.
Direct differentiation gives
\begin{equation}\label{quant-tools:eq-smoothing}
 0\leq h_\delta(x)-\abs x\leq\delta,\qquad
 \norm{h_\delta'}_\infty\leq1,\qquad
 \norm{h_\delta''}_\infty\leq\delta^{-1},\qquad
 \norm{h_\delta'''}_\infty\leq3\delta^{-2}.
\end{equation}
Taking the minimum of the uniform and Lipschitz bounds for the difference
of $h_\delta''$ yields
$[h_\delta'']_{C^{0,\alpha}}\leq C_\alpha\delta^{-1-\alpha}$.
Apply \Cref{lem:quadratic-smooth-comparison} with $N=k$ and $\rho=\mu$:
\begin{align}
 &\left|\E_{\mu^{\otimes k}}\left[
    \left|\frac1k\sum_{i<j}c(x_i,x_j)\right|\right]
       -\tfrac12\E[\abs{\mathcal Q_C}]\right|\notag\\
 &\qquad\leq2\delta+C_\alpha\delta^{-1-\alpha}
             k^{-\alpha/2}\norm c_p^p+\frac{\norm c_2}{\sqrt{2k}}.
 \label{quant-tools:eq-presmoothing}
\end{align}
Taking $\delta=\norm c_p k^{-\alpha/(2(2+\alpha))}$ proves the claim,
since $\norm c_2\leq\norm c_p$ and
$\alpha/(2(2+\alpha))\leq1/6<1/2$.
\end{proof}

% STAGE2 MOVE: bibliography-commands
\bibliographystyle{amsalpha}
\bibliography{sharp_chaos_references}

% STAGE2 MOVE: end-document
\end{document}
